\documentclass[11pt,a4paper,final,oneside]{book}
% TODO: publish black/white version
\usepackage{luatextra}
\usepackage[utf8x]{luainputenc}

\usepackage{unicode-math}
\usepackage{amsmath} % dep: after unicode-math
\usepackage[toc,page]{appendix}
\usepackage[autostyle=true]{csquotes}
\usepackage{polyglossia} \setmainlanguage{english} \setdefaultlanguage{english} % before biblatex
\usepackage[backend=bibtex,natbib=true]{biblatex}
\usepackage{booktabs}
\usepackage[centerlast,small,sc]{caption}
\usepackage{fontspec}
\usepackage{graphicx}
\usepackage[makeindex]{imakeidx}  % dep: before hyperref
\usepackage{hyperref} % dep: after biblatex
\usepackage{import}
\usepackage{luatexja-fontspec}
\usepackage{listings}
\usepackage{longtable}
\usepackage{enumitem}
\usepackage{parcolumns}
\usepackage{microtype}
\usepackage{mdframed}
\usepackage{ntheorem}
%\usepackage[oldstyle]{libertine} % (or 'lining')  % alternatives: Garamond, Caladea
\usepackage{subfig}
\usepackage{thmtools}
\usepackage{titling}
\usepackage{newunicodechar}
\usepackage{pxfonts}
\usepackage{pdfpages}
\usepackage{pdflscape}
\usepackage{xcolor}
\usepackage{xunicode}

%\usepackage[lf]{Baskervaldx}
%\usepackage[bigdelims,vvarbb]{newtxmath}
%\usepackage[cal=boondoxo]{mathalfa} % mathcal from STIX, unslanted a bit
%\renewcommand*\oldstylenums[1]{\textosf{#1}}

\captionsetup{justification=raggedright}

\newmdtheoremenv[%
  backgroundcolor=white,
  linecolor=white!60!black,
  linewidth=3pt]{todo}{Todo}

\parskip3pt
\setlength\parskip{0.7em}

\setromanfont[Ligatures=TeX,Numbers=OldStyle,Variant=01,Fractions=On]{Linux Libertine O}
\setromanfont[Ligatures=TeX,Numbers=OldStyle,Variant=01,Fractions=On]{Linux Libertine O}
\setmonofont{Inconsolata}
% FAIL does not power \mathbb N
% http://tex.stackexchange.com/a/47300/18015
\setmathfont{Asana Math}

\makeatletter
\def\@makechapterhead#1{%
  % REMARK I removed parcolumns, because it collides with
  % lualatexja-fontspec. See issue 978
  % http://tracker.luatex.org/view.php?id=978
  \noindent
  \begin{minipage}{0.5\textwidth}
  %\begin{parcolumns}[nofirstindent,colwidths={1=0.5\textwidth, 2=0.4\textwidth}, distance={0.1\textwidth}]{2}
  %  \colchunk[1]{%
      \ifx\chappic\@empty%
        \vspace*{50\p@}%
      \else%
        \vspace*{5\p@}%
          \includegraphics[width=100pt]\chappic %\\[15pt]
      \fi%
  %  }%
    \vspace{20pt}
  \end{minipage}%
  \begin{minipage}{0.5\textwidth}
  %  \colchunk[2]{%
      \ifx\chapquote\@empty\else%
        \vspace*{5\p@}
          \begin{flushright}
            %\begin{minipage}{0.3\textwidth}
              \enquote{\textsc{\chapquote}}\\ ---\chapquotesrc%
            %\end{minipage}
          \end{flushright}
      \fi%
  %  }
  %  \colplacechunks
  %\end{parcolumns}
  \end{minipage}

  {\parindent \z@ \raggedright \normalfont
    \ifnum \c@secnumdepth >\m@ne
      \if@mainmatter
        \huge\bfseries \@chapapp\space \thechapter
        \par\nobreak
        \vskip 20\p@
      \fi
    \fi
    \interlinepenalty\@M
    \Huge \bfseries #1\par\nobreak
    \vskip 40\p@
  }
}
\makeatother
\newcommand*\chappic{}
\newcommand*\chapquote{}
\newcommand*\chapquotesrc{}

\makeindex

\captionsetup[figure]{labelfont=bf}

\definecolor{linkcolor}{cmyk}{0.0,0.0,0.0,0.65}
\definecolor{citecolor}{cmyk}{0.0,0.0,0.0,0.82}
\definecolor{urlcolor}{cmyk}{0.0,0.0,0.0,0.60}

\definecolor{gray}{rgb}{0.7,0.7,0.7}
\newcommand{\dnW}[1][A]{\texttt{#1}}
\newcommand{\dnI}[1]{\texttt{#1}}
\newcommand{\dnCx}{\texttt{\color{darkgray}{x}}}
\newcommand{\dnCu}{\texttt{\color{darkgray}{u}}}
\newcommand{\dnCn}{\texttt{\color{darkgray}{n}}}
\newcommand{\dnCz}{\texttt{\color{black}{0}}}
\newcommand{\dnCo}{\texttt{\color{black}{1}}}
\newcommand{\dnCq}{\texttt{\color{black}{?}}}
\newcommand{\dnCa}{\texttt{\color{black}{A}}}
\newcommand{\dnCb}{\texttt{\color{black}{B}}}
\newcommand{\dnCc}{\texttt{\color{black}{C}}}
\newcommand{\dnCd}{\texttt{\color{black}{D}}}
\newcommand{\dnCh}{\texttt{\color{black}{-}}}

\newcommand{\cNP}{$\mathcal{N\!P}$}
\newcommand{\cP}{$\mathcal{P}$}
\newcommand{\cNPC}{$\mathcal{N\!PC}$}
\newcommand{\cPneqNP}{$\mathcal{P} \neq \mathcal{N\!P}$}

\newcommand\concat{\,\|\,}
\newcommand\yes{$✓$}
\newcommand\no{$✗$}  % REMARK not supported by any loaded font, but missing character looks visually fine
\newcommand\set[1]{\left\{#1\right\}}
\newcommand\boolf[1]{\texttt{#1}}
\newcommand\satfeature[1]{#1}
\newcommand\timeout{$\top$}
\newcommand\unknown{---}

\hyphenation{DIMACS}
\hyphenation{CNF}

\DeclareMathOperator\IF{\boolf{IF}}
\DeclareMathOperator\MAJ{\boolf{MAJ}}
\DeclareMathOperator\XOR{\boolf{XOR}}
\DeclareMathOperator\SigO{\Sigma_0\,}
\DeclareMathOperator\SigT{\Sigma_1\,}
\DeclareMathOperator\sigO{\sigma_0\,}
\DeclareMathOperator\sigT{\sigma_1\,}

\hypersetup{
  pdfpagemode={UseOutlines},
  pdftitle={Differential cryptanalysis with SAT solvers},
  pdfauthor={Lukas Prokop},
  pdfkeywords={master thesis, satisfiability, differential cryptanalysis},
  bookmarksopen=true,
  bookmarksopenlevel=0,
  hypertexnames=false,
  colorlinks=true,
  citecolor=citecolor,
  linkcolor=linkcolor,
  urlcolor=urlcolor,
  pdfstartview={FitV},
  unicode,
  breaklinks=true
}

\addbibresource{thesis.bib}
\setcounter{tocdepth}{3}

\author{Lukas Prokop}
\title{Differential cryptanalysis with SAT solvers}

\declaretheorem[thmbox={style=L,bodystyle=\noindent\normalfont\setlength{\parskip}{0.3\baselineskip},cut=false},style=plain,name={Theorem}]{theorem}
\declaretheorem[thmbox={style=L,bodystyle=\noindent\normalfont\setlength{\parskip}{0.3\baselineskip},cut=false},style=plain,name={Definition}]{defi}
\declaretheorem[thmbox={style=L,bodystyle=\noindent\normalfont\setlength{\parskip}{0.3\baselineskip},cut=false},style=plain,name={Claim}]{prop}
\declaretheorem[thmbox={style=L,bodystyle=\noindent\normalfont\setlength{\parskip}{0.3\baselineskip},cut=false},style=plain,name={Properties}]{pro}
%\newtheorem{theorem}{Theorem}
%\newtheorem{defi}{Definition}
%\newtheorem{prop}{Proposition}
%\newtheorem{pro}{Properties}

\numberwithin{theorem}{chapter}
\numberwithin{defi}{chapter}
\numberwithin{prop}{chapter}

\lstset{
  basicstyle=\small\ttfamily,
  frame=single
}

%\newcommand\theauthor{Lukas Prokop}
\newcommand\authormail{admin@lukas-prokop.at}
\newcommand\institute{Institute of Applied Information Processing and Communications}
\newcommand\university{University of Technology, Graz}

\begin{document}

\frontmatter
\begin{titlepage}
  \begin{center}
    {\scshape\LARGE \university\par}\vspace{1.2cm}
    \textsc{\Large Master Thesis}\\[0.5cm]

    \hrulefill \\[0.4cm]
    {\huge \bfseries Differential cryptanalysis \\ with SAT solvers\\}\vspace{0.1cm}
    \hrulefill \\[1.5cm]

    \begin{minipage}[t]{0.4\textwidth}
      \begin{flushleft} \large
        \emph{Author:}\\[3pt]
        %\href{mailto:\authormail}
        {\theauthor}
      \end{flushleft}
    \end{minipage}
    \begin{minipage}[t]{0.4\textwidth}
      \begin{flushright} \large
        \emph{Supervisor:} \\[3pt]
        {Maria Eichlseder \\ Florian Mendel}
      \end{flushright}
    \end{minipage}\\[3cm]

    \large \textit{A thesis submitted in fulfillment of the requirements \\ for the master's degree in Computer Science}\\[0.3cm]
    \textit{at the}\\[0.4cm]
    \begin{minipage}{0.4\textwidth}
      \centering
      \institute
    \end{minipage}
    \vspace{3.3cm}

    {\large \today}\\[7pt]
    \begin{minipage}{0.4\textwidth}
      \centering
      \includegraphics[width=\textwidth]{img/iaik.pdf}
    \end{minipage}

    \vfill
  \end{center}
\end{titlepage}
\thispagestyle{empty}

\newcommand{\nextoddpage}{
  \if@openright\cleardoublepage\else\clearpage\fi
  \ifdef{\phantomsection}{\phantomsection}{}
}

\newenvironment{declaration}{
  {\nextoddpage}
  %{\space Affidavit}
  {\thispagestyle{plain}}
  {\null\vfil}
  {\noindent%\huge\bfseries{
    \huge{\textsc{Affidavit}}
  %}
  \par\vspace{10pt}}
}{}

\NewDocumentEnvironment{acknowledgements}{}{%
  {\nextoddpage}
  %\tttypeout{\acknowledgementname}
  \thispagestyle{plain}
  \begin{center}{
      %\huge\textit{
      \Large\textsc{Acknowledgements}
      %}
    \par}\end{center}
}{\vfil\vfil\null}

\NewDocumentEnvironment{abstract}{}{%
  {\nextoddpage}
  \thispagestyle{plain}
  \begin{center}{\huge\textsc{Abstract}\par}\end{center}
}{\vfil\vfil\null}

\clearpage
\ifodd\value{page}\else\hbox{}\newpage\fi
\includepdf[pages=-]{Vorlage_TU.pdf}

\begin{abstract}%
  \parskip5pt
  Hash functions are ubiquitous in the modern information age. They provide preimage, second preimage
  and collision resistance which are needed in a wide range of applications.

  In August 2006, Wang et al. showed efficient attacks against several hash function designs including
  MD4, MD5, HAVAL-128 and RIPEMD. With these results differential cryptanalysis has been shown useful
  to break collision resistance in hash functions. Over the years, advanced attacks based on those
  differential approaches have been developed.

  To find collisions like Wang et al., a cryptanalyst needs to specify a differential
  characteristic as starting point, whose differences cancel out in the output.
  This starting point has a huge impact on the runtime.
  Once such a differential characteristic was discovered,
  in a second step the actual values for those differences
  are found yielding an actual hash collision.

  Evaluating differential characteristics can be a cumbersome and tedious task.
  Dedicated tools can implement search heuristics. Performance shortcuts
  can be made by studying the hash algorithm's differential behavior, etc.

  SAT solvers inherently implement a search heuristic and should derive those shortcuts
  on their own. In this thesis, we use SAT solvers to solve differential cryptanalysis
  problems. We show that SAT solvers are, in general, not able to derive those shortcuts on their
  own, but on the other hand we discuss approaches which significantly improve the runtime.
  SAT encoding remains an important topic to improve SAT solving runtimes.

  We wrote a library generating CNFs for differential characteristics and
  optionally modify the CNF for certain SAT encodings. Those modifications
  allowed a significant runtime improvement helping us to solve full-rounds
  hash collisions in MD4 and 24-rounds hash collisions in SHA-256.
  Finally, we also provide a small CNF analysis library to compare encoded problems
  with each other.

  \vspace{3pt}
  \textbf{Keywords:}
    hash function, differential cryptanalysis, differential characteristic,
    MD4, SHA-256, collision resistance,
    satisfiability, SAT solver
\end{abstract}

%% Skipped because TU Graz form already included
%\begin{declaration}
%  I declare that I have authored this thesis independently, that I have not used other
%  than the declared sources/resources, and that I have explicitly indicated all material
%  which has been quoted either literally or by content from the sources used.
%  The text document uploaded to TUGRAZonline is identical to the present master‘s
%  thesis dissertation.

%  %\rule[0.5em]{5em}{0.5pt} \hfill{} \rule[0.5em]{5em}{0.5pt} \\
%  %Date \hfill{} Digital signature
%  \vspace{120pt}
%  \hfill{}\raggedright{(digital signature)}
%\end{declaration}

\begin{acknowledgements}%
  \parskip5pt
  First of all I would like to thank my academic advisor for his continuous support
  during this project. Many hours of debugging were involved in writing this
  master thesis project, but thanks to Florian Mendel, this project came to a release
  with nice results. Also thanks for continuously reviewing this document.

  I would also like to thank Maria Eichlseder for her great support. Her unique
  way to ask questions brought me back on track several times.
  Mate Soos supported me during my bachelor thesis with SAT related issues
  and his support continued with this master thesis in private conversations.

  Also thanks to Roderick Bloem and Armin Biere who organized a meeting
  one year before submitting this work defining the main approaches involved
  in this thesis. Armin Biere released custom lingeling versions for us,
  which allowed us to influence the guessing strategy in lingeling.
  He also shared his thoughts about our testcases with us.

  And finally I am grateful for the support by Martina,
  who also supported me during good and bad days with this thesis,
  and my parents who provided a prosperous environment to me
  to be able to stand where I am today.

  Thank you. \\
  \indent どうもありがとうございました。
\end{acknowledgements}

All source codes are available at \href{http://lukas-prokop.at/proj/megosat}{lukas-prokop.at/proj/megosat}
and published under terms and conditions of Free/Libre Open Source Software.
This document was printed with \LuaLaTeX{} in the Linux Libertine typeface.

\tableofcontents
\mainmatter

\renewcommand*\chappic{img/intro.pdf}
\renewcommand*\chapquote{}
\renewcommand*\chapquotesrc{}
\chapter{Introduction}
\label{ch:intro}
%
Hash functions are used as cryptographic primitives in many applications and protocols.
They take an arbitrary input message and provide a hash value. Input message and hash value
are considered as byte strings in a particular encoding.
The hash value is of fixed length and satisfies several properties which make it useful
in a variety of applications.

In this thesis we will consider the hash algorithms MD4 and SHA-256 and represent
differential characteristics of hash collisions as SAT problem. If and only if
satisfiability is given, the particular differential state is achievable
using two different inputs leading to the same output. As far as SAT solvers
return an actual model satisfying that state, we get a complete hash collision
which can be verified and visualized.
If the internal state of the hash algorithm is too large, the attack can be easily
computationally simplified by modelling only a subset of steps of the hash algorithm
or changing the modelled differential path.

Gaining experience with these kind of problems with previous non-SAT-based tools
we try to apply best practices to a satisfiability setting.
We will discuss which SAT techniques lead to best performance characteristics
for our MD4 and SHA-256 testcases.

\newpage  % FIX manual pagebreak
\section{Cryptanalysis preliminaries}
\label{sec:intro-crypt-prelim}
%
\index{Hash value}
\index{Preimage}
\index{Hash function}
\begin{defi}[Hash function]
  A \emph{hash function} is a mapping $h: X \to Y$ with $X = \left\{0,1\right\}^*$ and
  $Y = \left\{0,1\right\}^n$ for some fixed $n \in {\mathbb{N}}_{\geq 1}$.
  \begin{itemize}[noitemsep,topsep=0pt]
    \item Let $x \in X$, then $h(x)$ is called \emph{hash value of $x$}.
    \item Let $h(x) = y \in Y$, then $x$ is called \emph{preimage of $y$}.
  \end{itemize}
\end{defi}

One example showing the use of hash functions as primitives are JSON Web Tokens (JWT)
specified in RFC~7519~\cite{rfc7519}. Section~8 defines implementation requirements
and refers to RFC~7518~\cite{rfc7518}, which specifies cryptographic algorithms to
be implemented. \enquote{HMAC SHA-256} (besides \enquote{none}) is the only signature
and MAC algorithm required to be implemented.
SHA-256 as hash algorithm is used as cryptographic primitive in this configuration.

A hash function has to satisfy the following security requirements:

\index{Preimage resistance}
\begin{defi}[Preimage resistance]
  Given $y \in Y$,
  a hash function $h$ is \emph{preimage resistant} iff it is computationally infeasible
  to find $x \in X$ such that $h(x) = y$.
\end{defi}

\index{Second-preimage resistance}
\begin{defi}[Second-preimage resistance]
  Given $x \in X$,
  a hash function $h$ is \emph{second-preimage resistant} iff it is computationally infeasible
  to find $x_2 \in X$ with $x \neq x_2$ such that $h(x) = h(x_2)$.
  $x_2$ is called \emph{second preimage}.
\end{defi}

\index{Collision resistance}
\begin{defi}[Collision resistance]
  A hash function $h$ is \emph{collision resistant} iff it is computationally infeasible to
  find any two $x \in X$ and $x_2 \in X$ with $x \neq x_2$ such that $h(x) = h(x_2)$.
\end{defi}

As far as hash functions accept input strings of arbitrary length, but return a fixed
size output string, existence of collisions is unavoidable~\cite{schlaffer}.
However, good hash functions make it very difficult to determine collisions or preimages.

The considered hash functions apply padding to their input to normalize their
input size to a multiple of its block size. The round function follows afterwards.
In the following we always consider input of block size meaning the round function
is run only for one block and padding is always considered part of this input message.
Padding is negligible, because once we have two colliding blocks,
the collision will be reflected in the output in these single-pipe Merkle-Darmg\aa rd designs.
This results in a length extension attack, making input padding negligible for cryptanalysis.

\section{Cryptanalysis of Hash Functions}
\label{sec:intro-cryptanalysis}
%
In August 2004, Wang et al. published results at Crypto'04~\cite{wang2004} which revealed
that MD4, MD5, HAVAL-128 and RIPEMD can be broken practically using differential cryptanalysis.
Their work is based on preliminary work by Hans Dobbertin~\cite{Dobbertin1998}.
On an IBM~P690 machine, an MD5 collision can be computed in about one hour using this approach.
Collisions for HAVAL-128, MD4 and RIPEMD were found as well. Patrick Stach's \texttt{md4coll.c}
program~\cite{md4coll} implements Wang's approach and can find MD4 collisions in few seconds
on my Thinkpad~x220 setup specified in \hyperref[app:setup]{Appendix~\ref{app:setup}}.

Let $n$ denote the digest size, i.e. the size of the hash value $h(x)$ as number of bits.
Due to the birthday paradox, a collision attack has a generic complexity of $2^{n/2}$
whereas preimage and second preimage attacks have generic complexities of $2^n$.
So it is cryptanalytically easier to find any two colliding hash values than the
preimage or second preimage for a given hash value.

Following results by Wang et al., differential cryptanalysis was shown as
powerful tool for cryptanalysis of hash algorithms. This thesis applies those
ideas to satisfiability approaches.

\begin{table}[bt]
  \begin{center}
    \begin{tabular}{cccc}
      \hline \hline
      \multicolumn{4}{c}{Message 1} \\
      \hline
      4d7a9c83 & \textbf{d6cb927a} & \textbf{29d5a578} & 57a7a5ee \\
      de748a3c & dcc366b3 & b683a020 & 3b2a5d9f \\
      c69d71b3 & f9e99198 & d79f805e & a63bb2e8 \\
      \textbf{45dc8e31} & 97e31fe5 & 2794bf08 & b9e8c3e9 \\
      \hline
      \multicolumn{4}{c}{Message 2} \\
      \hline
      4d7a9c83 & \textbf{56cb927a} & \textbf{b9d5a578} & 57a7a5ee \\
      de748a3c & dcc366b3 & b683a020 & 3b2a5d9f \\
      c69d71b3 & f9e99198 & d79f805e & a63bb2e8 \\
      \textbf{45dd8e31} & 97e31fe5 & 2794bf08 & b9e8c3e9 \\
      \hline
      \multicolumn{4}{c}{Hash value of Message 1 and Message 2} \\
      \hline
      5f5c1a0d & 71b36046 & 1b5435da & 9b0d807a \\
      \hline \hline
    \end{tabular}
    \caption[Hexadecimal values of one MD4 collisions given in paper~\cite{wang2004}]{
      One of two MD4 hash collisions provided in~\cite{wang2004}.
      Values are given in hexadecimal, message words are enumerated from
      left to right, top to bottom. Differences are highlighted in
      bold for illustration purposes. For comparison the first bits
      of Message 1 are \texttt{11000001\dots} and the last bits are
      \texttt{\dots10011101}.
      A message represents one block of 512~bits.
      %The hash value consists of 128~bits.
    }
    \label{tab:wang-md4-collision1}
  \end{center}
\end{table}

\newpage
\section{Differential cryptanalysis}
\label{sec:intro-diff-cryptanalysis}
%
\index{Hash collision}
\begin{defi}[Hash collision]
  Given a hash function $h$,
  a hash collision is a pair $(x, x_2)$ with $x \neq x_2$ such that
  $h(x) = h(x_2)$.
\end{defi}

Differential cryptanalysis is based on the idea to consider two execution states
of hash algorithms for slightly different input messages. We trace those difference
to learn about the propagation of message differences.

MD4 has a block size of 512~bits. Take some message in a particular encoding
and provide this padded byte string as 512~bit input to the hash algorithm.
Then use a copy of this byte string as input to a second hash algorithm implementation
with slight modifications in few bits. MD4 has 48 round function applications
in 3 rounds. Differential cryptanalysis considers the difference in the evaluation
state between the two instances.

Visualizing those differences helps the cryptanalyst to find modifications yielding
a small number of differences in the evaluation state.
The cryptanalyst consecutively modifies the input values to eventually
receive a collision in the output value.
If the number of differences in the evaluation state is small,
this trail is expected to result in a hash collision with higher probability.

\begin{figure}[bt]
  \begin{center}
    \fbox{\includegraphics[width=0.55\textwidth]{img/diff_cryptanalysis.pdf}}
    \caption[Typical attack setting for a collision attack]{
      Typical attack setting for a collision attack:
      Hash function $f$ is applied to two inputs $M$ and $M^*$ which differ
      by some predefined bits. A hash collision is given if and only if output values
      $C$ and $C^*$ show the same value. In differential cryptanalysis we observe
      the differences between two instances applying function $f$
      to inputs $M$ and $M^*$.
    }
  \end{center}
\end{figure}

\section{Satisfiability}
\label{sec:intro-sat}
%
\index{Boolean function}
\begin{defi}
  A \emph{Boolean function} is a mapping $h: X \to Y$ with $X = \left\{0,1\right\}^n$
  for $n \in \mathbb N_{\geq 1}$ and $Y = \left\{0,1\right\}$.
\end{defi}

The following definition gives three basic Boolean functions:

\index{AND}
\begin{defi}
  \begin{itemize}[noitemsep,topsep=0pt]
    \item
      \emph{AND} maps $X = \left\{0,1\right\}^2$
      to $1$ if all values of $X$ are $1$.
    \item
      \emph{OR} maps $X = \left\{0,1\right\}^2$
      to $1$ if any value of $X$ is $1$.
    \item
      \emph{NOT} maps $X = \left\{0,1\right\}^1$
      to $1$ if the single value of $X$ is $0$.
  \end{itemize}
  All functions return $0$ in the other case.
\end{defi}

\index{Truth table}
\begin{defi}
  A \emph{truth table} unambiguously defines a Boolean function
  by enlisting the evaluated truth value for all possible sets of
  inputs.

  Table~\ref{tab:andornot-truthtables} shows truth tables for
  \boolf{AND}, \boolf{OR} and \boolf{NOT}.
\end{defi}

\begin{figure}[!t]
  \centering
  \subfloat[\boolf{AND}]{%
    \begin{tabular}{cc|c}
      $v_1$ & $v_2$ & $f(v_1, v_2)$ \\
     \hline
      $1$ & $1$ & $1$ \\
      $1$ & $0$ & $0$ \\
      $0$ & $1$ & $0$ \\
      $0$ & $0$ & $0$
    \end{tabular}
  }
  ~
  \subfloat[\boolf{OR}]{%
    \begin{tabular}{cc|c}
      $v_1$ & $v_2$ & $f(v_1, v_2)$ \\
     \hline
      $1$ & $1$ & $1$ \\
      $1$ & $0$ & $1$ \\
      $0$ & $1$ & $1$ \\
      $0$ & $0$ & $0$
    \end{tabular}
  }
  ~
  \raisebox{13.6pt}{%
  \subfloat[\boolf{NOT}]{%
    \begin{tabular}{c|c}
      $v$ & $f(v)$ \\
     \hline
      $1$ & $0$ \\
      $0$ & $1$
    \end{tabular}
  }%
  }%
  \caption{Truth tables for \boolf{AND}, \boolf{OR} and \boolf{NOT}}
  \label{tab:andornot-truthtables}
\end{figure}

In the following we discuss how Boolean functions are related
to computation in general and hash algorithms specifically.

\index{Algorithm}
\begin{defi}
  An \emph{algorithm} is a step-wise set of instructions to solve a problem.
\end{defi}

One specific set of algorithms transform a given input to some output according
to some rules. Such algorithms are said to satisfy the \emph{I/O property}.
Hash algorithms satisfy the I/O property.

\begin{theorem}
  Every algorithm can be represented as Boolean function.
\end{theorem}

This is trivial to see considering that computers are built from logic gates.
However, in practice physical properties are used to employ storage of values,
which does not exist in the model of Boolean algebra. This can be mapped to
our Boolean model by constants for stored values and assuming an infinite memory.

\index{Satisfiability}
\index{Assignment}
\index{Model}
\begin{defi}
  A Boolean function $f$ is \emph{satisfiable} iff there exists at least one
  input $x \in X$ such that $f(x) = 1$.
  Every input $x \in X$ satisfying this property is called \emph{model}.
  Every element of $X$ is called \emph{assignment}.
\end{defi}

Consider a certain state during the computation of a hash value.
Represent the computation (i.e. hash algorithm) as Boolean equation system.
Represent the state as assignment to a subset of variables used in the computation.
Now replace, omit or invert some of these variables.
This modified state can be reached with some input if and only if
the SAT problem is satisfiable.

The generic complexity of SAT determination is given by $2^n$ for $n$ Boolean variables.
%This is closely related to the \cP $\overset{?}{\neq}$ \cNP question.
The corresponding tool to determine satisfiability is defined in the following.

\index{SAT solver}
\begin{defi}
  A \emph{SAT solver} is a tool to determine satisfiability of a Boolean function.
  If satisfiability is given, it returns some model.
\end{defi}

\section{Thesis Outline}
\label{sec:intro-outline}
%
This thesis is organized as follows:

\begin{description}
\item[In Chapter~\ref{ch:intro}] we discussed the basic properties and fundamentals
of the tools in discussion including hash functions and SAT solvers.

\item[In Chapter~\ref{ch:dc}] we introduce the MD4 and SHA-256 hash functions and
discuss possible approaches in differential cryptanalysis.

\item[In Chapter~\ref{ch:sat}] we discuss SAT solving and potential
approaches to speed up SAT solvers for cryptographic problems.

\item[In Chapter~\ref{ch:results}] we show results of our work
and discuss its implications.

\item[In Chapter~\ref{ch:summary}] we suggest future work based on our results.
\end{description}

\renewcommand*\chappic{img/hashalgos.pdf}
\renewcommand*\chapquote{}
\renewcommand*\chapquotesrc{}
%
\chapter{Hash algorithms}
\label{ch:hash}

In this chapter, we will define hash functions and their desired security
properties. In the following, we look at SHA-256 and MD4 as established hash functions.
MD4 unlike SHA-256 is practically broken, but has a comparably small internal state.
It therefore allows a good starting point to devise our attacks.
In a next step, we scaled up to SHA-256 which has an internal state size
at least twice as large.
In Chapter~\ref{ch:enc}, we will represent them with Boolean algebra to allow
to reason about states in those hash functions using SAT solvers.

\section{Preliminaries Redux}
\label{sec:hash-prelim}
%
\index{Hash value}
\index{Preimage}
\index{Hash function}
\begin{defi}[Hash function]
  A \emph{hash function} is a mapping $h: X \to Y$ with $X = \left\{0,1\right\}^*$ and
  $Y = \left\{0,1\right\}^n$ for some fixed $n \in {\symbol{"2115}}_{\geq 1}$.  % FAIL
  \begin{itemize}[noitemsep,topsep=0pt]
    \item Let $x \in X$, then $h(x)$ is called \emph{hash value of $x$}.
    \item Let $h(x) = y \in Y$, then $x$ is called \emph{preimage of $y$}.
  \end{itemize}
\end{defi}

Hash functions are considered as cryptographic primitives
used as building blocks in cryptographic protocols.
A hash function has to satisfy the following security requirements:

\index{Preimage resistance}
\begin{defi}[Preimage resistance]
  Given $y \in Y$,
  a hash function $h$ is \emph{preimage resistant} iff it is computationally infeasible
  to find $x \in X$ such that $h(x) = y$.
\end{defi}

\clearpage
\index{Second-preimage resistance}
\begin{defi}[Second-preimage resistance]
  Given $x \in X$,
  a hash function $h$ is \emph{second-preimage resistant} iff it is computationally infeasible
  to find $x_2 \in X$ with $x \neq x_2$ such that $h(x) = h(x_2)$.
  $x_2$ is called \emph{second preimage}.
\end{defi}

\index{Collision resistance}
\index{Collision}
\begin{defi}[Collision resistance]
  A hash function $h$ is \emph{collision resistant} iff it is computationally infeasible to
  find any two $x \in X$ and $x_2 \in X$ with $x \neq x_2$ such that $h(x) = h(x_2)$.
  Tuple $(x, x_2)$ is called \emph{collision}.
\end{defi}

As far as hash functions accept input strings of arbitrary length, but return a fixed
size output string, existence of collisions is unavoidable~\cite{schlaffer}.
However, good hash functions make it very difficult to find collisions or preimages.

Any digital data can be hashed (i.e. used as input to a hash function) by considering
it in binary representation. The format or encoding is not part of the hash function's
specification.

\subsection{Merkle-Damg\aa{}rd designs}
\label{sec:hash-md}
%
The Merkle-Damg\aa{}rd design is a particular design of hash functions providing the
following security guarantee:

\begin{defi}[Collision resistance inheritance]
  Let $F_0$ be a collision resistant compression function.
  A hash function in Merkle-Damg\aa{}rd design is collision resistant if $F_0$ is collision resistant.
\end{defi}

This motivates thorough research of collisions in compression functions.
The design was found independently by Ralph C. Merkle and Ivan B. Damg\aa{}rd.
It was described by Merkle in his PhD thesis~\cite[p. 13--15]{merkle1979secrecy}
and used in popular hash functions such as MD4, MD5 and the SHA2 hash function family.

The single-pipe design works as follows:
\begin{enumerate}\itemsep0pt
\item Split the input into blocks of uniform block size.
  If necessary, apply padding to the last block to achieve full block size.
\item Compression function $F_0$ is applied iteratively using the output $y_{i-1}$ of
  the previous iteration and the next input block $x_i$, denoted $y_i = F_0(y_{i-1}, x_i)$.
\item An optional postprocessing function is applied.
\end{enumerate}

\subsection{Padding and length extension attacks}
\label{sec:hash-length-ext-attack}
%
Hash functions of single-piped Merkle-Damg\aa{}rd design inherently suffer
from length extension attacks. MD4 and SHA-256 apply padding to their input
to normalize their input size to a multiple of its block size.
The compression function is applied afterwards. This design is vulnerable
to length extensions.

Consider some collision $(x_0, x_1)$ with $F_0(x_0) = y = F_0(x_1)$ where $x_0$ and
$x_1$ have a size of one block. Let $p$ be a suffix with size of one block.
Then also $(x_0 \,\|\, p, x_1 \,\|\, p)$ (where $\|$ denotes concatenation)
represents a collision in single-piped Merkle-Damg\aa{}rd designs, because
it holds that:
\[ F_0(F_0(x_0), p) = F_0(F_0(x_1), p) \iff F_0(y, p) = F_0(y, p) \]
Hence $(x_0 \,\|\, p, x_1 \,\|\, p)$ is a collision as well.
As far as $F_0$ is applied recursively to every block, $p$ can be of arbitrary size
and $(x_0, x_1)$ can be of arbitrary uniform size.

Because of this vulnerability, cryptanalysts focus on finding a collision
in compression functions. In our tests will only consider input of one block
and padding will be neglected due to this vulnerability.

\subsection{Example usage}
\label{sec:hash-usage}
%
The following list describes application examples:
%
\begin{description}
  \item[Digital signatures.]
    Digital signatures are schemes to guarantee data and origin integrity.
    They are used to verify authenticity of PDF~documents,
    encrypted emails or software distributions. A digital signature
    can be verified, showing that a particular user was involved in the
    signature process. No third party can forge a digital signature.
    Because the document itself might be too large, only its hash value
    is attached to the hash signature.
  \item[Storing passwords.]
    User account data such as passwords are commonly stored in databases.
    Those databases are subject to theft if network breaches take place
    and the passwords are known to the attackers. Hashing and salting
    those passwords with a hash algorithm before storing them in the database
    is a working countermeasure.
  \item[Commitment schemes.]
    Commitment schemes like Zero-Knowledge Proofs use hash algorithms to ensure the
    original message is possessed by a certain party without revealing the document
    itself. Coin-flipping or secure computation also use hash algorithms as
    cryptographic building blocks.
\end{description}

% One example showing the use of hash functions as primitives are JSON Web Tokens (JWT)
% specified in RFC~7519~\cite{rfc7519}. Its application allows web developers to
% represent claims to be transferred between two parties.
%
% Section~8 defines implementation requirements and refers to RFC~7518~\cite{rfc7518},
% which specifies cryptographic algorithms such as \enquote{HMAC SHA-256} to be
% implemented. It is (besides \enquote{none}) the only required signature and MAC algorithm.

\section{MD4}
\label{sec:dc-md4}
%
\index{MD4}
MD4 is a cryptographic hash function originally described in RFC~1186~\cite{rfc1186},
updated in RFC~1320~\cite{rfc1320} and declared obsolete by RFC~6150~\cite{rfc6150}. It was
invented by Ronald Rivest in 1990 with properties given in Table~\ref{tab:md4}.
In 1995~\cite{Dobbertin1998}, successful full-round attacks have been found to break collision resistance.
Three years later, preimage and second-preimage resistance in MD4 got broken as well.
Some of those attacks are described in~\cite{md4-2007} and \cite{cryptoeprint:2005:151}.
We derived a Python~3 implementation based on a Python~2 implementation
and made it available on github~\cite{md4-py3k}.

\begin{table}[h]
  \begin{center}
    \begin{tabular}{lcl}
      block size           & 512 bits       & namely variable \texttt{block} in RFC~1320~\cite{rfc1320} \\
      digest size          & 128 bits       & as per Section~3.5 in RFC~1320~\cite{rfc1320} \\
      internal state size  & 128 bits       & namely variables $A$, $B$, $C$ and $D$ \\
      word size            & 32 bits        & as per Section~2 in RFC~1320~\cite{rfc1320} \\
    \end{tabular}
    \caption{MD4 hash algorithm properties}
    \label{tab:md4}
  \end{center}
\end{table}

MD4 uses three auxiliary Boolean functions:
\begin{defi}
  The Boolean \boolf{IF} function is defined as follows:
  If the first argument is true, the second argument is returned.
  Otherwise the third argument is returned.

  The Boolean \boolf{MAJ} function returns true if the number
  of arguments with Boolean value true is in the majority.
  The Boolean \boolf{XOR} function returns true if the number
  of arguments with Boolean value true is odd.
\end{defi}

Using the logical operators $\land$ (\boolf{AND}), $\lor$ (\boolf{OR})
and $\neg$ (\boolf{NEG}) we can define them as (see Section~\ref{sec:sat-intro}
for a thorough discussion of these operators):
\begin{align}
  \IF(X,Y,Z) &\coloneqq (X \land Y) \lor (\neg X \land Z) \\
  \MAJ(X,Y,Z) &\coloneqq (X \land Y) \lor (X \land Z) \lor (Y \land Z) \\
  \XOR(X,Y,Z) &\coloneqq (X \land \neg Y \land \neg Z) \lor (\neg X \land Y \land \neg Z) \nonumber\\
              &\lor (\neg X \land \neg Y \land Z) \lor (X \land Y \land Z) \nonumber\\
              &\coloneqq (X \oplus Y \oplus Z)
\end{align}

In the following, a brief overview of MD4 is given.

\begin{description}
  \item[Padding and length extension.]
    First of all, padding is applied. A single bit $1$ is appended to the
    input. As long as the input does not reach a length congruent 448 modulo 512,
    bit $0$ is appended.
    Afterwards, length appending takes place.
    Append the first 64 less significant bits of the input length (without the previous modifications)
    represented in binary.
  \item[Initialization.]
    The message is split into 512-bit blocks (i.e. 16 32-bit words).
    State variables $A_{i}$ with $-4 \leq i < 0$ are initialized with these
    hexadecimal values:
    % remark: A_{-4} := A, A_{-1} := B, A_{-2} := C, A_{-3} := D on Wikipedia
    \[
      \begin{array}{cccc}
        A_{-4} = \texttt{01234567} &
        A_{-1} = \texttt{89abcdef} &
        A_{-2} = \texttt{fedcba98} &
        A_{-3} = \texttt{76543210}
      \end{array}
    \]
  \item[Round function and state variable updates.]
    The round function
    is applied in 3 rounds with 16 iterations.
    In every iteration, values $A_{-1}$, $A_{-2}$ and $A_{-3}$ are taken
    as arguments to function $F$. Function $F$ is \boolf{IF} in round~1,
    followed by \boolf{MAJ} for round~2 and \boolf{XOR} for the final round~3.
    The resulting value is added to $A_{-1}$, current message block $M$
    and constant $X$. Finally, the 32-bit sum will be left-rotated by $p$
    positions. Left rotation is formally defined in Definition~\ref{def:shifts}.
    We define $X$ and $p$:

    Let $i$ be the iteration counter between $1$ and $16$ and
    $r$ the round between $1$ and $3$. Then $X$ takes the value
    of the $i$-th column and $r$-th row of matrix $C$.
    $p$ takes the value of row $r$ and column $i \bmod{4}$ of matrix $P$.
    \[
      C = \left(\begin{array}{cccccccccccccccc}
        0 & 1 & 2 & 3 & 4 & 5 & 6 & 7 & 8 & 9 & 10 & 11 & 12 & 13 & 14 & 15 \\
        0 & 4 & 8 & 12 & 1 & 5 & 9 & 13 & 2 & 6 & 10 & 14 & 3 & 7 & 11 & 15 \\
        0 & 8 & 4 & 12 & 2 & 10 & 6 & 14 & 1 & 9 & 5 & 13 & 3 & 11 & 7 & 15
      \end{array}\right)
    \] \[
      P = \begin{pmatrix}
        3 & 7 & 9 & 11 \\
        3 & 5 & 9 & 13 \\
        3 & 9 & 11 & 15
      \end{pmatrix}
    \]
    This round function design is visualized in Figure~\ref{fig:md4-round-function}.
\end{description}

\begin{figure}[!b]
  \begin{center}
    \includegraphics[scale=0.7]{img/md4.pdf}
    \caption{MD4 round function updating state variables}
    \label{fig:md4-round-function}
  \end{center}
\end{figure}










\section{SHA-256}
\label{sec:dc-sha-256}
%
\index{SHA-256}
SHA-256 is a hash function from the SHA-2 family designed by the National Security Agency (NSA)
and published as FIPS PUB 180-4 (originally 2001)~\cite{fips-pub-180-4}.
It uses a Merkle-Damg\aa{}rd construction with a Davies-Meyer compression function.
The best known preimage attack was found in 2011
and breaks preimage resistance for 52~rounds~\cite{bicliques}. The best known collision attack
breaks collision resistance for 31~rounds of SHA-256~\cite{improving} and pseudo-collision
resistance for 46~rounds~\cite{high2011}. The best practical attack is a pseudo-collision
for 38 steps~\cite{mendel2013improving}.

\begin{table}[!hb]
  \begin{center}
    \begin{tabular}{lcl}
      block size           & 512 bits   & as per Section~1 of the standard~\cite{fips-pub-180-4} \\
      digest size          & 256 bits   & mentioned as Message Digest size~\cite{fips-pub-180-4} \\
      internal state size  & 256 bits   & as per Section~1 of the standard~\cite{fips-pub-180-4} \\
      word size            & 32 bits    & as per Section~1 of the standard~\cite{fips-pub-180-4}
    \end{tabular}
    \caption{SHA-256 hash algorithm properties}
    \label{tab:sha256}
  \end{center}
\end{table}
%
\index{Left-shift}
\index{Right-shift}
\index{Left-rotation}
\index{Right-rotation}
\begin{defi}[Shifts, rotations and a notational remark]
  \label{def:shifts}
  Consider a 32-bit word $X$ with 32 binary values $b_i$ with $0 \leq i \leq 31$.
  $b_0$ refers to the least significant bit. Shifting ($≪$ and $≫$) and
  rotation ($⋘$ and $⋙$) creates a new 32-bit word $Y$ with 32 binary values $a_i$.
  We define the following notations:
  \begin{align*}
    Y \coloneqq X ≪ n  &\iff a_i \coloneqq b_{i-n} \text{ if } 0 \leq i-n < 32 \text{ and } 0 \text{ otherwise } \\
    Y \coloneqq X ≫ n  &\iff a_i \coloneqq b_{i+n} \text{ if } 0 \leq i+n < 32 \text{ and } 0 \text{ otherwise } \\
    Y \coloneqq X ⋘ n  &\iff a_i \coloneqq b_{i-n \bmod{32}} \text{ as used in MD4} \\
    Y \coloneqq X ⋙ n  &\iff a_i \coloneqq b_{i+n \bmod{32}}
  \end{align*}
\end{defi}
%
Besides MD4's \boolf{MAJ} and \boolf{IF},
another four auxiliary functions are defined. Recognize that $\oplus$ denotes
the \boolf{XOR} function whereas $\boxplus$ denotes 32-bit addition.
\[
  \begin{array}{rcccccl}
    \Sigma_0(X) &\coloneqq& (X ⋙ 2) &\oplus& (X ⋙ 13) &\oplus& (X ⋙ 22) \\
    \Sigma_1(X) &\coloneqq& (X ⋙ 6) &\oplus& (X ⋙ 11) &\oplus& (X ⋙ 25) \\
    \sigma_0(X) &\coloneqq& (X ⋙ 7) &\oplus& (X ⋙ 18) &\oplus& (X ≫ 3) \\
    \sigma_1(X) &\coloneqq& (X ⋙ 17) &\oplus& (X ⋙ 19) &\oplus& (X ≫ 10)
  \end{array}
\]
%
\begin{description}
  \item[Padding and length extension.]
    The padding and length extension scheme of MD4 is used also in SHA-256.
    Append bit $1$, followed by a sequence of bit $0$ until the message reaches a length
    of 448 modulo 512 bits. Afterwards the first 64 bits of the binary representation
    of the original input are appended.
  \item[Initialization.]
    In a similar manner to MD4, initialization of internal state variables
    (called \enquote{working variables} in~\cite[Section~6.2.2]{fips-pub-180-4})
    takes place before running the round function. The eight state variables
    are initialized with the following hexadecimal values:
    \[
      \begin{array}{llll}
        A_{-1} = \texttt{6a09e667} &
        A_{-2} = \texttt{bb67ae85} &
        A_{-3} = \texttt{3c6ef372} &
        A_{-4} = \texttt{a54ff53a} \\
        E_{-1} = \texttt{510e527f} &
        E_{-2} = \texttt{9b05688c} &
        E_{-3} = \texttt{1f83d9ab} &
        E_{-4} = \texttt{5be0cd19}
      \end{array}
    \]
    Furthermore SHA-256 uses 64 constant values in its round function.
    We initialize step constants $K_i$ for $0 \leq i < 64$ with the following hexadecimal
    values (which must be read left to right and top to bottom):
    \[
      \begin{array}{cccccc}
        \texttt{428a2f98} & \texttt{71374491} & \texttt{b5c0fbcf} & \texttt{e9b5dba5} & \texttt{3956c25b} & \texttt{59f111f1} \\
        \texttt{923f82a4} & \texttt{ab1c5ed5} & \texttt{d807aa98} & \texttt{12835b01} & \texttt{243185be} & \texttt{550c7dc3} \\
        \texttt{72be5d74} & \texttt{80deb1fe} & \texttt{9bdc06a7} & \texttt{c19bf174} & \texttt{e49b69c1} & \texttt{efbe4786} \\
        \texttt{0fc19dc6} & \texttt{240ca1cc} & \texttt{2de92c6f} & \texttt{4a7484aa} & \texttt{5cb0a9dc} & \texttt{76f988da} \\
        \texttt{983e5152} & \texttt{a831c66d} & \texttt{b00327c8} & \texttt{bf597fc7} & \texttt{c6e00bf3} & \texttt{d5a79147} \\
        \texttt{06ca6351} & \texttt{14292967} & \texttt{27b70a85} & \texttt{2e1b2138} & \texttt{4d2c6dfc} & \texttt{53380d13} \\
        \texttt{650a7354} & \texttt{766a0abb} & \texttt{81c2c92e} & \texttt{92722c85} & \texttt{a2bfe8a1} & \texttt{a81a664b} \\
        \texttt{c24b8b70} & \texttt{c76c51a3} & \texttt{d192e819} & \texttt{d6990624} & \texttt{f40e3585} & \texttt{106aa070} \\
        \texttt{19a4c116} & \texttt{1e376c08} & \texttt{2748774c} & \texttt{34b0bcb5} & \texttt{391c0cb3} & \texttt{4ed8aa4a} \\
        \texttt{5b9cca4f} & \texttt{682e6ff3} & \texttt{748f82ee} & \texttt{78a5636f} & \texttt{84c87814} & \texttt{8cc70208} \\
        \texttt{90befffa} & \texttt{a4506ceb} & \texttt{bef9a3f7} & \texttt{c67178f2}
      \end{array}
    \]
  \item[Precomputation of W.]
    Let $W_i$ for $0 \leq i < 16$ be the sixteen 32-bit words of the padded input message.
    Then compute $W_i$ for $16 \leq i < 64$ the following way:
    \[ W_i \coloneqq \sigT(W_{i-2}) + W_{i-7} + \sigO(W_{i-15}) + W_{i-16} \]
  \item[Round function.]
    For every block of 512~bits, the round function is applied.
    The eight state variables are updated iteratively for $i$ from 0 to 63.
    $W_i$ and $K_i$ refer to the previously initialized values.
    \begin{align*}
      E_i     &\coloneqq A_{i-4} + E_{i-4} + \SigT(E_{i-1}) + \IF{(E_{i-1}, E_{i-2}, E_{i-3})} + K_i + W_i \\
      A_i     &\coloneqq E_i - A_{i-4} + \SigO(A_{i-1}) + \MAJ{(A_{i-1}, A_{i-2}, A_{i-3})}
    \end{align*}
%   \item[Computation of intermediate hash values.]
%     Intermediate hash values for the Davies-Meyer construction are
%     initialized with the following values:
%     \[
%       \begin{array}{cp{10pt}cp{10pt}cp{10pt}c}
%         H_0^{(0)} \coloneqq A_{-1}  & &  H_1^{(0)} \coloneqq A_{-2}  & &  H_2^{(0)} \coloneqq A_{-3}  & &  H_3^{(0)} \coloneqq A_{-4} \\
%         H_4^{(i)} \coloneqq E_{-1}  & &  H_5^{(i)} \coloneqq E_{-2}  & &  H_6^{(i)} \coloneqq E_{-3}  & &  H_7^{(i)} \coloneqq E_{-4} \\
%       \end{array}
%     \]
%     Every block creates its own $E_{i}$ and $A_{i}$ values for $60 \leq i < 64$.
%     These are used to compute the next intermediate values:
%     \[
%       \begin{array}{lp{10pt}r}
%         H_0^{(j)} \coloneqq A_{63} + H_0^{(i-1)}  & &  H_4^{(j)} \coloneqq E_{63} + H_4^{(i-1)} \\
%         H_1^{(j)} \coloneqq A_{62} + H_1^{(i-1)}  & &  H_5^{(j)} \coloneqq E_{62} + H_5^{(i-1)} \\
%         H_2^{(j)} \coloneqq A_{61} + H_2^{(i-1)}  & &  H_6^{(j)} \coloneqq E_{61} + H_6^{(i-1)} \\
%         H_3^{(j)} \coloneqq A_{60} + H_3^{(i-1)}  & &  H_7^{(j)} \coloneqq E_{60} + H_7^{(i-1)}
%       \end{array}
%     \]
%   \item[Finalization.]
%     The final hash digest of size 256 bits is provided as
%     \[
%       H_0^{(N)} \concat
%       H_1^{(N)} \concat
%       H_2^{(N)} \concat
%       H_3^{(N)} \concat
%       H_4^{(N)} \concat
%       H_5^{(N)} \concat
%       H_6^{(N)} \concat
%       H_7^{(N)}
%     \]
%     where $N$ denotes the index of the last block and operator $\|$ denotes
%     concatenation. Hence $H_0^{(N)}$ are the four least significant bytes of the digest.
\end{description}
\begin{figure}[!ht]
  \begin{center}
    \includegraphics[scale=0.7]{img/sha256.pdf}
    \caption{SHA-256 round function~\cite{analysisSHA256}}
    \label{fig:sha256-round-function}
  \end{center}
\end{figure}

\renewcommand*\chappic{img/diff-crypt.pdf}
\renewcommand*\chapquote{Just because it's automatic doesn't mean it works.}
\renewcommand*\chapquotesrc{Daniel J. Bernstein}
%
\chapter{Differential cryptanalysis}
\label{ch:dc}
%
In Chapter~\ref{ch:hash} we defined two hash functions. In this chapter
we consider such functions from a differential perspective. Differential
considerations will turn out to yield successful collision attacks on hash
functions. We introduce a notation to easily represent differential
characteristics.

\section{Motivation}
\label{sec:dc-motivation}
%
In August 2004, Wang et al. published results at Crypto'04~\cite{wang2004} which revealed
that MD4, MD5, HAVAL-128 and RIPEMD can be broken practically using differential cryptanalysis.
Their work is based on preliminary work by Hans Dobbertin~\cite{Dobbertin1998}.
On an IBM~P690 machine, an MD5 collision can be computed in about one hour using this approach.
Collisions for HAVAL-128, MD4 and RIPEMD were found as well. Patrick Stach's \texttt{md4coll.c}
program~\cite{md4coll} implements Wang's approach and can find MD4 collisions in few seconds
on my Thinkpad~x220 setup specified in \hyperref[app:setup]{Appendix~\ref{app:setup}}.

Let $n$ denote the digest size, i.e. the size of the hash value $h(x)$ in bits.
Due to the birthday paradox, a collision attack has a generic complexity of $2^{n/2}$
whereas preimage and second preimage attacks have generic complexities of $2^n$.
In other words it is computationally easier to find any two colliding hash values
than the preimage or second preimage for a given hash value.

Following results by Wang et al., differential cryptanalysis was shown as
powerful tool for cryptanalysis of hash algorithms. This thesis applies those
ideas to satisfiability approaches.

\begin{table}[bt]
  \begin{center}
    \begin{tabular}{cccc}
      \hline \hline
      \multicolumn{4}{c}{Message 1} \\
      \hline
      4d7a9c83 & \textbf{d6cb927a} & \textbf{29d5a578} & 57a7a5ee \\
      de748a3c & dcc366b3 & b683a020 & 3b2a5d9f \\
      c69d71b3 & f9e99198 & d79f805e & a63bb2e8 \\
      \textbf{45dc8e31} & 97e31fe5 & 2794bf08 & b9e8c3e9 \\
      \hline
      \multicolumn{4}{c}{Message 2} \\
      \hline
      4d7a9c83 & \textbf{56cb927a} & \textbf{b9d5a578} & 57a7a5ee \\
      de748a3c & dcc366b3 & b683a020 & 3b2a5d9f \\
      c69d71b3 & f9e99198 & d79f805e & a63bb2e8 \\
      \textbf{45dd8e31} & 97e31fe5 & 2794bf08 & b9e8c3e9 \\
      \hline
      \multicolumn{4}{c}{Hash value of Message 1 and Message 2} \\
      \hline
      5f5c1a0d & 71b36046 & 1b5435da & 9b0d807a \\
      \hline \hline
    \end{tabular}
    \caption[Hexadecimal values of one MD4 collisions given in paper~\cite{wang2004}]{
      One of two MD4 hash collisions provided in~\cite{wang2004}.
      A message represents one block of 512~bits.
      %The hash value consists of 128~bits.
      Values are given in hexadecimal, message words are enumerated from
      left to right, top to bottom. Differences are highlighted in
      bold for illustration purposes. For comparison the first bits
      of Message 1 are \texttt{11000001\dots} and the last bits are
      \texttt{\dots10011101}.
    }
    \label{tab:wang-md4-collision1}
  \end{center}
\end{table}

\newpage
\section{Fundamentals}
\label{sec:dc-fundamentals}
%
\index{Hash collision}
\begin{defi}[Hash collision]
  Given a hash function $h$,
  a hash collision is a pair $(x_1, x_2)$ with $x_1 \neq x_2$ such that
  $h(x_1) = h(x_2)$.
\end{defi}

\index{Pseudo collision}
Pseudo-collisions are also often considered when attacking hash functions.
A \emph{pseudo collision} is given if a hash collision can be found for a given
hash function, but the initial vectors (IV) can be chosen arbitrarily.

Hash algorithms consume input values as blocks of bits.
As far as the length of the input must not conform to the block size,
padding is applied. Now consider such a block of input values
and another copy of it. We use those two blocks as inputs for two
hash algorithm instances, but provide slight modifications in few bits.
Differential cryptanalysis is based on the idea to consider those execution states
and tracing those difference to learn about the propagation of message differences.
Compare this setup with Figure~\ref{tab:collision-attack}.

\index{Avalanche effect}
At the very beginning only the few defined differences are given. But as the
hash algorithm progresses in computation, differences are propagated to more
and more bits. Most likely the final value will differ in many bits, because
of a desirable hash algorithm property called \emph{avalanche effect}.
A small difference in the input should lead to a significant difference
in the output (i.e. visually recognizable).

Visualizing those differences helps the cryptanalyst to find modifications
yielding a small number of differences in the evaluation state. The propagation
of differences in a particular hash algorithm is called \emph{differential
path}. Empirical results in differential cryptanalysis indicate that spare
paths are desirable, because it is easier to cancel out few differences in the
output compared to many differences.
The cryptanalyst consecutively modifies the input values to eventually
receive a collision in the output value.

\index{Differential path}
\index{Differential characteristic}
\begin{defi}
  The propagation of differences in a particular function is called
  \emph{differential path}. The complete differential state during
  a computation is called \emph{differential characteristic}.
\end{defi}

\begin{theorem}
  Assuming the number of differences in a differential path is small,
  this path is expected to result in a hash collision with higher probability.
\end{theorem}

\begin{figure}[pbt]
  \begin{center}
    \fbox{\includegraphics[width=0.55\textwidth]{img/diff_cryptanalysis.pdf}}
    \caption[Common attack setting for a collision attack]{
      Common attack setting for a collision attack:
      Hash function $f$ is applied to two inputs $M$ and $M^*$ which differ
      by some predefined bits. $\Delta M$ describes the difference between
      these values. A hash collision is given if and only if output values
      $C$ and $C^*$ show the same value. In differential cryptanalysis we observe
      the differences between two instances applying function $f$
      to inputs $M$ and $M^*$.
    }
    \label{tab:collision-attack}
  \end{center}
\end{figure}

\section{Differential notation}
\label{sec:dc-notation}
%
\index{Bit condition}
\index{Generalized bit condition}
\index{Differential notation}
Differential notation helps us to visualize differential characteristics
by defining so-called \emph{generalized bit conditions}.
It was introduced by Christian Rechberger and Christophe de Canni\`ere
in 2006~\cite[Section 3.2]{char-2006}, inspired by \emph{signed differences} by
Wang~et al. and is shown in Table~\ref{tab:diff-notation}.

Consider two hash algorithm instances. Let $x_i$ be some bit
from the first instance and let $x_i^*$ be the corresponding bit
from the second instance. Differences are computed using a \boolf{XOR}
and commonly denoted as $\Delta x = x_i \oplus x_i^*$.
Bit conditions allow us to encode possible relations between bits $x_i$ and $x_i^*$.

For example, let us take a look at the original Wang et al. hash collision
in MD4 provided in Table~\ref{tab:wang-md4-collision1}.
We extract all values with differences and represent them using differential notation.
This gives us Table~\ref{tab:differential-wang-values}.

The following properties hold for bit conditions:
\begin{itemize}[noitemsep,topsep=0pt]
  \item If $x_i = x_i^*$ holds and some value is known, $\set{0,1}$ contains its bit condition.
  \item If $x_i \neq x_i^*$ holds and some value is known, $\set{u,n}$ contains its bit condition.
  \item If $x_i = x_i^*$ holds and the values are unknown, its bit condition is $\dnI{-}$.
  \item If $x_i \neq x_i^*$ holds and the values are unknown, its bit condition is $\dnI{x}$.
\end{itemize}
Applying this notation to hash collisions means that arbitrary bit conditions
(expect for \dnI{\#}) can be specified for the input values. In one of the
intermediate iterations, we enforce a difference using one of the bit conditions
$\set{u,n,x}$. This excludes trivial solutions with no differences from the set of
possible solutions. And the final values need to lack differences thus are
represented using a dash \dnI{-}.

\begin{table}[bt]
  \begin{center}
    \small
    \begin{tabular}{lll}
      \hline \hline
      bit & binary & hexadecimal representation / differential notation \\
      \hline \hline
      $x_0$ & \dnI{d6cb927a} & \dnI{11010110110010111001001001111010} \\
      $x_1$ & \dnI{29d5a578} & \dnI{00101001110101011010010101111000} \\
      $x_2$ & \dnI{45dc8e31} & \dnI{01000101110111001000111000110001} \\
      \hline
      $x_0^*$ & \dnI{56cb927a} & \dnI{01010110110010111001001001111010} \\
      $x_1^*$ & \dnI{b9d5a578} & \dnI{10111001110101011010010101111000} \\
      $x_2^*$ & \dnI{45dd8e31} & \dnI{01000101110111011000111000110001} \\
      \hline
      $\Delta x$ & & \dnI{u1010110110010111001001001111010} \\
                 & & \dnI{n01n1001110101011010010101111000} \\
                 & & \dnI{010001011101110n1000111000110001} \\
      \hline \hline
    \end{tabular}
    \caption[Bit differences in the original Wang et al. hash collision]{
      The three words different between Message~1 and Message~2 of the original
      MD4 hash collision by Wang et al. The last three lines show
      how differences can be
      written down using bit conditions. As far as 4 symbols are not from the
      set $\set{0, 1}$ it holds that the messages differ by 4~bits.
    }
    \label{tab:differential-wang-values}
  \end{center}
\end{table}

\begin{table}[p]
  \small
  \begin{center}
    \begin{tabular}{|c|cccc|}
    \hline\hline
      $(x_i, x_i^*)$ & $(0,0)$ & $(1,0)$ & $(0,1)$ & $(1,1)$ \\
    \hline
      \dnI{?}        & \yes    & \yes    & \yes    & \yes    \\
      \dnI{-}        & \yes    & \no     & \no     & \yes    \\
      \dnI{x}        & \no     & \yes    & \yes    & \no     \\
      \dnI{0}        & \yes    & \no     & \no     & \no     \\
      \dnI{u}        & \no     & \yes    & \no     & \no     \\
      \dnI{n}        & \no     & \no     & \yes    & \no     \\
      \dnI{1}        & \no     & \no     & \no     & \yes    \\
      \dnI{\#}       & \no     & \no     & \no     & \no     \\
    \hline \hline
    \end{tabular}
    \begin{tabular}{|c|cccc|}
    \hline\hline
      $(x_i, x_i^*)$ & $(0,0)$ & $(1,0)$ & $(0,1)$ & $(1,1)$ \\
    \hline
      \dnI{3}        & \yes    & \yes    & \no     & \no     \\
      \dnI{5}        & \yes    & \no     & \yes    & \no     \\
      \dnI{7}        & \yes    & \yes    & \yes    & \no     \\
      \dnI{A}        & \no     & \yes    & \no     & \yes    \\
      \dnI{B}        & \yes    & \yes    & \no     & \yes    \\
      \dnI{C}        & \no     & \no     & \yes    & \yes    \\
      \dnI{D}        & \yes    & \no     & \yes    & \yes    \\
      \dnI{E}        & \no     & \yes    & \yes    & \yes    \\
    \hline \hline
    \end{tabular}
    \caption[Differential notation as introduced in~\cite{char-2006}]{
      Differential notation as introduced in~\cite{char-2006}.
      The left-most column specifies a symbol called \enquote{bit condition}
      and right-side columns indicate which bit configurations
      are possible for two given bits $x_i$ and $x_i^*$.
    }
    \label{tab:diff-notation}
  \end{center}
\end{table}

\begin{table}[p]
  \begin{center}
    \begin{tabular}{cp{5cm}cl}
      $\Delta x$      & conjunctive normal form &
      $\Delta x$      & conjunctive normal form \\
    \hline
      \dnI{\#}        & $(x) \land (\neg x)$ &
      \dnI{1}         & $(x) \land (x^*)$ \\

      \dnI{0}         & $(\neg x) \land (\neg x^*)$ &
      \dnI{-}         & $\neg (x \oplus x^*)$ \\

      \dnI{u}         & $(x) \land (\neg x^*)$ &
      \dnI{A}         & $(x)$ \\

      \dnI{3}         & $(\neg x^*)$ &
      \dnI{B}         & $(x \lor \neg x^*)$ \\

      \dnI{n}         & $(\neg x) \land (x^*)$ &
      \dnI{C}         & $(x^*)$ \\

      \dnI{5}         & $(\neg x)$ &
      \dnI{D}         & $(\neg x \lor x^*)$ \\

      \dnI{x}         & $(x \oplus x^*)$ &
      \dnI{E}         & $(x \lor x^*)$ \\

      \dnI{7}         & $(\neg x \lor \neg x^*)$ &
      \dnI{?}         &  \\
    \end{tabular}
    \caption[Representation of bit conditions as CNF]{%
      All bit conditions represented as CNF using
      two Boolean variables $x$ and $x^*$ to represent
      two bits.
    }
    \label{tab:simple-eval-clauses}
  \end{center}
\end{table}


\section{A simple addition example}
\label{sec:dc-example}
%
Using this notation, we can now reason about the behavior of functions on differential values.
We start with 1-bit addition as most basic exercise to the reader. Consider a matrix with
two input rows and one output row. The values of the first two rows are added such that
the bit difference at the third row is created.

Figure~\ref{tab:1-bit-example} illustrates this example. Remember that symbols such as
\dnI{-} and \dnI{0} underlie semantics defined in Table~\ref{tab:diff-notation}.
It is also interesting to see how propagation of values can work. In
Figure~\ref{tab:1-bit-deduction} we see how an underspecified value \dnI{?} can be
strengthened once we have checked which values can be taken. Recognize that the
system is constrained by the function in use and the definition of the differential symbols.

Finally we can extend our testcases to 4~bits and retrieve testcases such as
Table~\ref{tab:4-bit-addition} and \ref{tab:sigma}.

\begin{table}[p]
  \begin{center}
    \begin{minipage}{20pt}\begin{tabular}{c} \dnI{-} \\ \dnI{-} \\ \hline \dnI{-} \end{tabular}\end{minipage}
    \hspace{20pt}$\Rightarrow$\hspace{20pt}
    \begin{minipage}{20pt}\begin{tabular}{c} \dnI{0}\dnI{0} \\ \dnI{0}\dnI{0} \\ \hline \dnI{0}\dnI{0} \end{tabular}\end{minipage}
    \begin{minipage}{20pt}\begin{tabular}{c} \dnI{0}\dnI{0} \\ \dnI{0}\dnI{0} \\ \hline \dnI{1}\dnI{1} \end{tabular}\end{minipage}
    \begin{minipage}{20pt}\begin{tabular}{c} \dnI{0}\dnI{0} \\ \dnI{1}\dnI{1} \\ \hline \dnI{0}\dnI{0} \end{tabular}\end{minipage}
    \begin{minipage}{20pt}\begin{tabular}{c} \dnI{0}\dnI{0} \\ \dnI{1}\dnI{1} \\ \hline \dnI{1}\dnI{1} \end{tabular}\end{minipage}
    \begin{minipage}{20pt}\begin{tabular}{c} \dnI{1}\dnI{1} \\ \dnI{0}\dnI{0} \\ \hline \dnI{0}\dnI{0} \end{tabular}\end{minipage}
    \begin{minipage}{20pt}\begin{tabular}{c} \dnI{1}\dnI{1} \\ \dnI{0}\dnI{0} \\ \hline \dnI{1}\dnI{1} \end{tabular}\end{minipage}
    \begin{minipage}{20pt}\begin{tabular}{c} \dnI{1}\dnI{1} \\ \dnI{1}\dnI{1} \\ \hline \dnI{0}\dnI{0} \end{tabular}\end{minipage}
    \begin{minipage}{20pt}\begin{tabular}{c} \dnI{1}\dnI{1} \\ \dnI{1}\dnI{1} \\ \hline \dnI{1}\dnI{1} \end{tabular}\end{minipage}
    \caption[A simple 1-bit addition example]{%
      A simple 1-bit addition example:
      On the left the differential characteristic is given.
      Two dashes, by definition, denote a missing difference in both arguments.
      The result of the addition most neither show a difference.
      This yields eight possible bit configurations where two values close to each other denote $(M, M^*)$ of Figure~\ref{tab:collision-attack}.
      Due to the behavior of addition, we know that configurations 2, 3, 5 and 8 (from left to right) are invalid.
    }
    \label{tab:1-bit-example}
  \end{center}
\end{table}
\begin{table}[p]
  \begin{center}
    \begin{minipage}{20pt}\begin{tabular}{c} \dnI{-} \\ \dnI{-} \\ \hline \dnI{?} \end{tabular}\end{minipage}
    \hspace{20pt}$\Rightarrow$\hspace{20pt}
    \begin{minipage}{20pt}\begin{tabular}{c} \dnI{0}\dnI{0} \\ \dnI{0}\dnI{0} \\ \hline \dnI{0}\dnI{0} \end{tabular}\end{minipage}
    \begin{minipage}{20pt}\begin{tabular}{c} \dnI{0}\dnI{0} \\ \dnI{1}\dnI{1} \\ \hline \dnI{1}\dnI{1} \end{tabular}\end{minipage}
    \begin{minipage}{20pt}\begin{tabular}{c} \dnI{1}\dnI{1} \\ \dnI{0}\dnI{0} \\ \hline \dnI{1}\dnI{1} \end{tabular}\end{minipage}
    \begin{minipage}{20pt}\begin{tabular}{c} \dnI{1}\dnI{1} \\ \dnI{1}\dnI{1} \\ \hline \dnI{0}\dnI{0} \end{tabular}\end{minipage}
    \caption[Propagation of bit conditions in a differential characteristic]{%
      Like Figure~\ref{tab:1-bit-example}, but any difference value for the result bit is possible.
      As such we consider any possible bit configuration, but eventually recognize that only four bit configurations
      are consistent with the behavior of addition. Because all resulting configurations show no bit difference
      in the output bit, we can strengthen \dnI{?} by replacing it with \dnI{-}. This illustrates how
      knowledge about differential states can be propagated.
    }
    \label{tab:1-bit-deduction}
  \end{center}
\end{table}

\begin{table}[p]
  \begin{center}
    \begin{minipage}{0.23\textwidth}
      \begin{tabular}{rl}
        A: & \dnI{0}\dnI{0}\dnI{1}\dnI{1} \\
        B: & \dnI{0}\dnI{1}\dnI{0}\dnI{1} \\
        S: & \dnI{1}\dnI{0}\dnI{0}\dnI{0}
      \end{tabular}
    \end{minipage}
    \begin{minipage}{0.23\textwidth}
      \begin{tabular}{rl}
        A: & \dnI{-}\dnI{-}\dnI{-}\dnI{x} \\
        B: & \dnI{-}\dnI{-}\dnI{-}\dnI{x} \\
        S: & \dnI{?}\dnI{?}\dnI{?}\dnI{?}
      \end{tabular}
    \end{minipage}
    \begin{minipage}{0.23\textwidth}
      \begin{tabular}{rl}
        A: & \dnI{-}\dnI{-}\dnI{-}\dnI{x} \\
        B: & \dnI{-}\dnI{-}\dnI{-}\dnI{x} \\
        S: & \dnI{?}\dnI{?}\dnI{?}\dnI{-}
      \end{tabular}
    \end{minipage}
    \begin{minipage}{0.23\textwidth}
      \begin{tabular}{rl}
        A: & \dnI{-}\dnI{-}\dnI{-}\dnI{x} \\
        B: & \dnI{-}\dnI{-}\dnI{-}\dnI{x} \\
        S: & \dnI{x}\dnI{?}\dnI{?}\dnI{?}
      \end{tabular}
    \end{minipage}
  \end{center}%

  \begin{center}
    \begin{minipage}{0.23\textwidth}
      \begin{tabular}{rl}
        A: & \dnI{0}\dnI{0}\dnI{1}\dnI{1} \\
        B: & \dnI{0}\dnI{1}\dnI{0}\dnI{1} \\
        S: & \dnI{0}\dnI{0}\dnI{0}\dnI{0}
      \end{tabular}
    \end{minipage}
    \begin{minipage}{0.23\textwidth}
      \begin{tabular}{rl}
        A: & \dnI{-}\dnI{-}\dnI{-}\dnI{x} \\
        B: & \dnI{-}\dnI{-}\dnI{-}\dnI{x} \\
        S: & \dnI{?}\dnI{?}\dnI{?}\dnI{x}
      \end{tabular}
    \end{minipage}
    \begin{minipage}{0.23\textwidth}
      \begin{tabular}{rl}
        A: & \dnI{-}\dnI{-}\dnI{-}\dnI{-} \\
        B: & \dnI{-}\dnI{-}\dnI{-}\dnI{x} \\
        S: & \dnI{x}\dnI{-}\dnI{?}\dnI{?}
      \end{tabular}
    \end{minipage}
  \end{center}%
  \caption[
    Testcases for 4-bit addition
  ]{
    Testcases for 4-bit addition:
    The upper line shows valid differential characteristics for 4-bit addition
    whereas the lower line show invalid ones for 4-bit addition.
    The rows are conventionally named using capital letters.
  }
  \label{tab:4-bit-addition}
\end{table}

\begin{table}[p]
  \begin{center}
    \begin{minipage}{0.23\textwidth}
      \begin{tabular}{rl}
        A: & \dnI{-}\dnI{-}\dnI{-}\dnI{-} \\
        S: & \dnI{0}\dnI{0}\dnI{0}\dnI{0}
      \end{tabular}
    \end{minipage}
    \begin{minipage}{0.23\textwidth}
      \begin{tabular}{rl}
        A: & \dnI{7}\dnI{C}\dnI{-}\dnI{3} \\
        S: & \dnI{-}\dnI{3}\dnI{u}\dnI{?}
      \end{tabular}
    \end{minipage}
    \begin{minipage}{0.23\textwidth}
      \begin{tabular}{rl}
        A: & \dnI{0}\dnI{u}\dnI{C}\dnI{D} \\
        S: & \dnI{A}\dnI{D}\dnI{C}\dnI{7} \\
      \end{tabular}
    \end{minipage}
  \end{center}

  \begin{center}
    \begin{minipage}{0.23\textwidth}
      \begin{tabular}{rl}
        A: & \dnI{-}\dnI{-}\dnI{-}\dnI{x} \\
        S: & \dnI{0}\dnI{0}\dnI{0}\dnI{0}
      \end{tabular}
    \end{minipage}
    \begin{minipage}{0.23\textwidth}
      \begin{tabular}{rl}
        A: & \dnI{x}\dnI{x}\dnI{x}\dnI{x} \\
        S: & \dnI{0}\dnI{0}\dnI{0}\dnI{0} \\
      \end{tabular}
    \end{minipage}
  \end{center}
  \caption[Differential characteristics for the SHA-2 Sigma function]{
    Differential characteristics for the SHA-2 Sigma function.
    The upper line shows valid states. The lower line shows invalid ones.
  }
  \label{tab:sigma}
\end{table}

\section{Differential characteristics in action}
\label{sec:dc-actual-chars}
%
In the previous section we illustrated how propagation with differential values
works and how differential characteristics are written down. It is always
important to keep in mind which function the characteristic illustrates,
because this is not documented with the characteristic.

Now consider MD4 as defined in Section~\ref{sec:dc-md4}. MD4 takes some input
message (in our case limited to size of one block), the state variables
are initialized and iteratively new $A_i$ are computed.

Similarly, SHA-256 takes a message block $M$ and initializes eight
variables with an initial vector (IV). The remaining $W_i$ are computed
and iteratively, values $A_i$ and $E_i$ are computed.

Those values are structured in differential characteristics illustrated
in Figure~\ref{img:char-layouts}. Those layouts are used to specify our
hash collisions we want to evaluate. Table~\ref{tab:wang-collision-propagated}
is also gives an application of the layout.

\begin{figure}[p]
  \begin{center}
    \includegraphics[width=0.4\linewidth]{img/md4_layout.pdf}
    \includegraphics[width=0.5\linewidth]{img/sha256_layout.pdf}
    \caption{Layout of MD4 and SHA-256 differential characteristics}
    \label{img:char-layouts}
  \end{center}
\end{figure}


\texttt{
\begin{table}[p]
\begin{center}
\scalebox{.35}{
\begin{tabular}{|r|c|c|c|c|}
\hline
  $i$  &      &  $\nabla{S_{i,0}}$  &  $\nabla{S_{i,1}}$  &  $\nabla{S_{i,2}}$ \\
\hline
 \dnI{-4} & \dnW: & {{\dnCz}{\dnCo}{\dnCo}{\dnCz}{\dnCz}{\dnCo}{\dnCo}{\dnCo}{\dnCz}{\dnCo}{\dnCz}{\dnCz}{\dnCz}{\dnCo}{\dnCz}{\dnCo}{\dnCz}{\dnCz}{\dnCo}{\dnCz}{\dnCz}{\dnCz}{\dnCo}{\dnCo}{\dnCz}{\dnCz}{\dnCz}{\dnCz}{\dnCz}{\dnCz}{\dnCz}{\dnCo}} & & \\
 \dnI{-3} & \dnW: & {{\dnCz}{\dnCz}{\dnCz}{\dnCo}{\dnCz}{\dnCz}{\dnCz}{\dnCz}{\dnCz}{\dnCz}{\dnCo}{\dnCo}{\dnCz}{\dnCz}{\dnCo}{\dnCz}{\dnCz}{\dnCo}{\dnCz}{\dnCo}{\dnCz}{\dnCo}{\dnCz}{\dnCz}{\dnCz}{\dnCo}{\dnCo}{\dnCo}{\dnCz}{\dnCo}{\dnCo}{\dnCz}} & & \\
 \dnI{-2} & \dnW: & {{\dnCo}{\dnCz}{\dnCz}{\dnCo}{\dnCo}{\dnCz}{\dnCz}{\dnCz}{\dnCo}{\dnCz}{\dnCo}{\dnCo}{\dnCo}{\dnCz}{\dnCo}{\dnCz}{\dnCo}{\dnCo}{\dnCz}{\dnCo}{\dnCo}{\dnCo}{\dnCz}{\dnCz}{\dnCo}{\dnCo}{\dnCo}{\dnCo}{\dnCo}{\dnCo}{\dnCo}{\dnCz}} & & \\
 \dnI{-1} & \dnW: & {{\dnCo}{\dnCo}{\dnCo}{\dnCz}{\dnCo}{\dnCo}{\dnCo}{\dnCo}{\dnCo}{\dnCo}{\dnCz}{\dnCz}{\dnCo}{\dnCo}{\dnCz}{\dnCo}{\dnCo}{\dnCz}{\dnCo}{\dnCz}{\dnCo}{\dnCz}{\dnCo}{\dnCo}{\dnCo}{\dnCz}{\dnCz}{\dnCz}{\dnCo}{\dnCz}{\dnCz}{\dnCo}} & & \\
 \dnI{0} & \dnW: & {{\dnCz}{\dnCo}{\dnCo}{\dnCz}{\dnCo}{\dnCz}{\dnCo}{\dnCo}{\dnCo}{\dnCo}{\dnCz}{\dnCo}{\dnCz}{\dnCo}{\dnCz}{\dnCz}{\dnCo}{\dnCo}{\dnCo}{\dnCz}{\dnCz}{\dnCo}{\dnCz}{\dnCz}{\dnCz}{\dnCz}{\dnCz}{\dnCo}{\dnCz}{\dnCz}{\dnCo}{\dnCz}} & \dnW[W]{}: & {{\dnCz}{\dnCo}{\dnCz}{\dnCz}{\dnCo}{\dnCo}{\dnCz}{\dnCo}{\dnCz}{\dnCo}{\dnCo}{\dnCo}{\dnCo}{\dnCz}{\dnCo}{\dnCz}{\dnCo}{\dnCz}{\dnCz}{\dnCo}{\dnCo}{\dnCo}{\dnCz}{\dnCz}{\dnCo}{\dnCz}{\dnCz}{\dnCz}{\dnCz}{\dnCz}{\dnCo}{\dnCo}} \\
 \dnI{1} & \dnW: & {{\dnCz}{\dnCo}{\dnCo}{\dnCo}{\dnCz}{\dnCo}{\dnCo}{\dnCz}{\dnCz}{\dnCo}{\dnCz}{\dnCz}{\dnCo}{\dnCo}{\dnCo}{\dnCo}{\dnCo}{\dnCo}{\dnCo}{\dnCz}{\dnCo}{\dnCo}{\dnCo}{\dnCz}{\dnCz}{\dnCu}{\dnCo}{\dnCo}{\dnCz}{\dnCz}{\dnCz}{\dnCo}} & \dnW[W]{}: & {{\dnCu}{\dnCo}{\dnCz}{\dnCo}{\dnCz}{\dnCo}{\dnCo}{\dnCz}{\dnCo}{\dnCo}{\dnCz}{\dnCz}{\dnCo}{\dnCz}{\dnCo}{\dnCo}{\dnCo}{\dnCz}{\dnCz}{\dnCo}{\dnCz}{\dnCz}{\dnCo}{\dnCz}{\dnCz}{\dnCo}{\dnCo}{\dnCo}{\dnCo}{\dnCz}{\dnCo}{\dnCz}} \\
 \dnI{2} & \dnW: & {{\dnCo}{\dnCz}{\dnCo}{\dnCz}{\dnCo}{\dnCz}{\dnCo}{\dnCo}{\dnCz}{\dnCo}{\dnCz}{\dnCz}{\dnCz}{\dnCz}{\dnCz}{\dnCz}{\dnCz}{\dnCo}{\dnCo}{\dnCo}{\dnCz}{\dnCu}{\dnCz}{\dnCo}{\dnCn}{\dnCo}{\dnCo}{\dnCo}{\dnCz}{\dnCz}{\dnCo}{\dnCz}} & \dnW[W]{}: & {{\dnCn}{\dnCz}{\dnCo}{\dnCn}{\dnCo}{\dnCz}{\dnCz}{\dnCo}{\dnCo}{\dnCo}{\dnCz}{\dnCo}{\dnCz}{\dnCo}{\dnCz}{\dnCo}{\dnCo}{\dnCz}{\dnCo}{\dnCz}{\dnCz}{\dnCo}{\dnCz}{\dnCo}{\dnCz}{\dnCo}{\dnCo}{\dnCo}{\dnCo}{\dnCz}{\dnCz}{\dnCz}} \\
 \dnI{3} & \dnW: & {{\dnCo}{\dnCz}{\dnCo}{\dnCz}{\dnCo}{\dnCo}{\dnCu}{\dnCo}{\dnCz}{\dnCz}{\dnCo}{\dnCo}{\dnCo}{\dnCo}{\dnCz}{\dnCo}{\dnCz}{\dnCo}{\dnCz}{\dnCo}{\dnCz}{\dnCz}{\dnCo}{\dnCz}{\dnCz}{\dnCo}{\dnCz}{\dnCo}{\dnCz}{\dnCz}{\dnCz}{\dnCo}} & \dnW[W]{}: & {{\dnCz}{\dnCo}{\dnCz}{\dnCo}{\dnCz}{\dnCo}{\dnCo}{\dnCo}{\dnCo}{\dnCz}{\dnCo}{\dnCz}{\dnCz}{\dnCo}{\dnCo}{\dnCo}{\dnCo}{\dnCz}{\dnCo}{\dnCz}{\dnCz}{\dnCo}{\dnCz}{\dnCo}{\dnCo}{\dnCo}{\dnCo}{\dnCz}{\dnCo}{\dnCo}{\dnCo}{\dnCz}} \\
 \dnI{4} & \dnW: & {{\dnCz}{\dnCz}{\dnCo}{\dnCz}{\dnCo}{\dnCo}{\dnCz}{\dnCz}{\dnCz}{\dnCo}{\dnCo}{\dnCz}{\dnCz}{\dnCz}{\dnCo}{\dnCo}{\dnCz}{\dnCo}{\dnCz}{\dnCo}{\dnCz}{\dnCo}{\dnCz}{\dnCo}{\dnCo}{\dnCo}{\dnCo}{\dnCo}{\dnCz}{\dnCz}{\dnCo}{\dnCz}} & \dnW[W]{}: & {{\dnCo}{\dnCo}{\dnCz}{\dnCo}{\dnCo}{\dnCo}{\dnCo}{\dnCz}{\dnCz}{\dnCo}{\dnCo}{\dnCo}{\dnCz}{\dnCo}{\dnCz}{\dnCz}{\dnCo}{\dnCz}{\dnCz}{\dnCz}{\dnCo}{\dnCz}{\dnCo}{\dnCz}{\dnCz}{\dnCz}{\dnCo}{\dnCo}{\dnCo}{\dnCo}{\dnCz}{\dnCz}} \\
 \dnI{5} & \dnW: & {{\dnCz}{\dnCz}{\dnCz}{\dnCo}{\dnCo}{\dnCz}{\dnCo}{\dnCz}{\dnCz}{\dnCo}{\dnCo}{\dnCz}{\dnCz}{\dnCz}{\dnCo}{\dnCz}{\dnCo}{\dnCz}{\dnCu}{\dnCo}{\dnCo}{\dnCz}{\dnCo}{\dnCz}{\dnCz}{\dnCz}{\dnCz}{\dnCz}{\dnCz}{\dnCz}{\dnCz}{\dnCo}} & \dnW[W]{}: & {{\dnCo}{\dnCo}{\dnCz}{\dnCo}{\dnCo}{\dnCo}{\dnCz}{\dnCz}{\dnCo}{\dnCo}{\dnCz}{\dnCz}{\dnCz}{\dnCz}{\dnCo}{\dnCo}{\dnCz}{\dnCo}{\dnCo}{\dnCz}{\dnCz}{\dnCo}{\dnCo}{\dnCz}{\dnCo}{\dnCz}{\dnCo}{\dnCo}{\dnCz}{\dnCz}{\dnCo}{\dnCo}} \\
 \dnI{6} & \dnW: & {{\dnCz}{\dnCz}{\dnCz}{\dnCo}{\dnCo}{\dnCz}{\dnCo}{\dnCo}{\dnCz}{\dnCz}{\dnCu}{\dnCn}{\dnCu}{\dnCu}{\dnCo}{\dnCo}{\dnCz}{\dnCz}{\dnCz}{\dnCo}{\dnCz}{\dnCz}{\dnCz}{\dnCz}{\dnCz}{\dnCo}{\dnCo}{\dnCo}{\dnCo}{\dnCz}{\dnCo}{\dnCz}} & \dnW[W]{}: & {{\dnCq}{\dnCq}{\dnCq}{\dnCq}{\dnCq}{\dnCq}{\dnCq}{\dnCq}{\dnCq}{\dnCq}{\dnCq}{\dnCq}{\dnCq}{\dnCq}{\dnCq}{\dnCq}{\dnCq}{\dnCq}{\dnCq}{\dnCq}{\dnCq}{\dnCq}{\dnCq}{\dnCq}{\dnCq}{\dnCq}{\dnCq}{\dnCq}{\dnCq}{\dnCq}{\dnCq}{\dnCq}} \\
 \dnI{7} & \dnW: & {{\dnCz}{\dnCz}{\dnCo}{\dnCz}{\dnCo}{\dnCz}{\dnCo}{\dnCo}{\dnCo}{\dnCz}{\dnCz}{\dnCz}{\dnCz}{\dnCz}{\dnCz}{\dnCo}{\dnCz}{\dnCu}{\dnCn}{\dnCn}{\dnCz}{\dnCo}{\dnCo}{\dnCz}{\dnCz}{\dnCo}{\dnCz}{\dnCo}{\dnCz}{\dnCz}{\dnCz}{\dnCz}} & \dnW[W]{}: & {{\dnCz}{\dnCz}{\dnCo}{\dnCo}{\dnCo}{\dnCz}{\dnCo}{\dnCo}{\dnCz}{\dnCz}{\dnCo}{\dnCz}{\dnCo}{\dnCz}{\dnCo}{\dnCz}{\dnCz}{\dnCo}{\dnCz}{\dnCo}{\dnCo}{\dnCo}{\dnCz}{\dnCo}{\dnCo}{\dnCz}{\dnCz}{\dnCo}{\dnCo}{\dnCo}{\dnCo}{\dnCo}} \\
 \dnI{8} & \dnW: & {{\dnCq}{\dnCq}{\dnCq}{\dnCq}{\dnCq}{\dnCq}{\dnCq}{\dnCq}{\dnCq}{\dnCq}{\dnCq}{\dnCq}{\dnCq}{\dnCq}{\dnCq}{\dnCq}{\dnCq}{\dnCq}{\dnCq}{\dnCq}{\dnCq}{\dnCq}{\dnCq}{\dnCq}{\dnCq}{\dnCq}{\dnCq}{\dnCq}{\dnCq}{\dnCq}{\dnCq}{\dnCq}} & \dnW[W]{}: & {{\dnCo}{\dnCo}{\dnCz}{\dnCz}{\dnCz}{\dnCo}{\dnCo}{\dnCz}{\dnCo}{\dnCz}{\dnCz}{\dnCo}{\dnCo}{\dnCo}{\dnCz}{\dnCo}{\dnCz}{\dnCo}{\dnCo}{\dnCo}{\dnCz}{\dnCz}{\dnCz}{\dnCo}{\dnCo}{\dnCz}{\dnCo}{\dnCo}{\dnCz}{\dnCz}{\dnCo}{\dnCo}} \\
 \dnI{9} & \dnW: & {{\dnCq}{\dnCq}{\dnCq}{\dnCq}{\dnCq}{\dnCq}{\dnCq}{\dnCq}{\dnCq}{\dnCq}{\dnCq}{\dnCq}{\dnCq}{\dnCq}{\dnCq}{\dnCq}{\dnCq}{\dnCq}{\dnCq}{\dnCq}{\dnCq}{\dnCq}{\dnCq}{\dnCq}{\dnCq}{\dnCq}{\dnCq}{\dnCq}{\dnCq}{\dnCq}{\dnCq}{\dnCq}} & \dnW[W]{}: & {{\dnCo}{\dnCo}{\dnCo}{\dnCo}{\dnCo}{\dnCz}{\dnCz}{\dnCo}{\dnCo}{\dnCo}{\dnCo}{\dnCz}{\dnCo}{\dnCz}{\dnCz}{\dnCo}{\dnCo}{\dnCz}{\dnCz}{\dnCo}{\dnCz}{\dnCz}{\dnCz}{\dnCo}{\dnCo}{\dnCz}{\dnCz}{\dnCo}{\dnCo}{\dnCz}{\dnCz}{\dnCz}} \\
 \dnI{10} & \dnW: & {{\dnCo}{\dnCz}{\dnCn}{\dnCz}{\dnCz}{\dnCo}{\dnCz}{\dnCz}{\dnCo}{\dnCz}{\dnCz}{\dnCz}{\dnCz}{\dnCo}{\dnCz}{\dnCo}{\dnCq}{\dnCq}{\dnCq}{\dnCq}{\dnCq}{\dnCq}{\dnCq}{\dnCq}{\dnCq}{\dnCq}{\dnCq}{\dnCq}{\dnCq}{\dnCo}{\dnCo}{\dnCz}} & \dnW[W]{}: & {{\dnCo}{\dnCo}{\dnCz}{\dnCo}{\dnCz}{\dnCo}{\dnCo}{\dnCo}{\dnCo}{\dnCz}{\dnCz}{\dnCo}{\dnCo}{\dnCo}{\dnCo}{\dnCo}{\dnCo}{\dnCz}{\dnCz}{\dnCz}{\dnCz}{\dnCz}{\dnCz}{\dnCz}{\dnCz}{\dnCo}{\dnCz}{\dnCo}{\dnCo}{\dnCo}{\dnCo}{\dnCz}} \\
 \dnI{11} & \dnW: & {{\dnCu}{\dnCo}{\dnCz}{\dnCz}{\dnCz}{\dnCo}{\dnCo}{\dnCz}{\dnCo}{\dnCz}{\dnCo}{\dnCo}{\dnCz}{\dnCo}{\dnCo}{\dnCz}{\dnCq}{\dnCq}{\dnCq}{\dnCq}{\dnCq}{\dnCq}{\dnCq}{\dnCq}{\dnCq}{\dnCq}{\dnCq}{\dnCq}{\dnCq}{\dnCo}{\dnCo}{\dnCo}} & \dnW[W]{}: & {{\dnCo}{\dnCz}{\dnCo}{\dnCz}{\dnCz}{\dnCo}{\dnCo}{\dnCz}{\dnCz}{\dnCz}{\dnCo}{\dnCo}{\dnCo}{\dnCz}{\dnCo}{\dnCo}{\dnCo}{\dnCz}{\dnCo}{\dnCo}{\dnCz}{\dnCz}{\dnCo}{\dnCz}{\dnCo}{\dnCo}{\dnCo}{\dnCz}{\dnCo}{\dnCz}{\dnCz}{\dnCz}} \\
 \dnI{12} & \dnW: & {{\dnCz}{\dnCz}{\dnCo}{\dnCz}{\dnCo}{\dnCo}{\dnCu}{\dnCz}{\dnCz}{\dnCu}{\dnCo}{\dnCz}{\dnCo}{\dnCz}{\dnCo}{\dnCo}{\dnCq}{\dnCq}{\dnCq}{\dnCq}{\dnCq}{\dnCq}{\dnCq}{\dnCq}{\dnCq}{\dnCq}{\dnCq}{\dnCq}{\dnCq}{\dnCz}{\dnCo}{\dnCo}} & \dnW[W]{}: & {{\dnCz}{\dnCo}{\dnCz}{\dnCz}{\dnCz}{\dnCo}{\dnCz}{\dnCo}{\dnCo}{\dnCo}{\dnCz}{\dnCo}{\dnCo}{\dnCo}{\dnCz}{\dnCn}{\dnCo}{\dnCz}{\dnCz}{\dnCz}{\dnCo}{\dnCo}{\dnCo}{\dnCz}{\dnCz}{\dnCz}{\dnCo}{\dnCo}{\dnCz}{\dnCz}{\dnCz}{\dnCo}} \\
 \dnI{13} & \dnW: & {{\dnCo}{\dnCz}{\dnCu}{\dnCn}{\dnCo}{\dnCn}{\dnCz}{\dnCo}{\dnCz}{\dnCz}{\dnCo}{\dnCo}{\dnCz}{\dnCz}{\dnCz}{\dnCo}{\dnCq}{\dnCq}{\dnCq}{\dnCq}{\dnCq}{\dnCq}{\dnCq}{\dnCq}{\dnCq}{\dnCq}{\dnCq}{\dnCq}{\dnCq}{\dnCo}{\dnCz}{\dnCo}} & \dnW[W]{}: & {{\dnCo}{\dnCz}{\dnCz}{\dnCo}{\dnCz}{\dnCo}{\dnCo}{\dnCo}{\dnCo}{\dnCo}{\dnCo}{\dnCz}{\dnCz}{\dnCz}{\dnCo}{\dnCo}{\dnCz}{\dnCz}{\dnCz}{\dnCo}{\dnCo}{\dnCo}{\dnCo}{\dnCo}{\dnCo}{\dnCo}{\dnCo}{\dnCz}{\dnCz}{\dnCo}{\dnCz}{\dnCo}} \\
 \dnI{14} & \dnW: & {{\dnCz}{\dnCz}{\dnCz}{\dnCz}{\dnCo}{\dnCz}{\dnCo}{\dnCz}{\dnCz}{\dnCo}{\dnCz}{\dnCo}{\dnCz}{\dnCz}{\dnCz}{\dnCo}{\dnCh}{\dnCh}{\dnCh}{\dnCh}{\dnCh}{\dnCh}{\dnCh}{\dnCh}{\dnCh}{\dnCh}{\dnCh}{\dnCh}{\dnCh}{\dnCo}{\dnCo}{\dnCz}} & \dnW[W]{}: & {{\dnCz}{\dnCz}{\dnCo}{\dnCz}{\dnCz}{\dnCo}{\dnCo}{\dnCo}{\dnCo}{\dnCz}{\dnCz}{\dnCo}{\dnCz}{\dnCo}{\dnCz}{\dnCz}{\dnCo}{\dnCz}{\dnCo}{\dnCo}{\dnCo}{\dnCo}{\dnCo}{\dnCo}{\dnCz}{\dnCz}{\dnCz}{\dnCz}{\dnCo}{\dnCz}{\dnCz}{\dnCz}} \\
 \dnI{15} & \dnW: & {{\dnCz}{\dnCz}{\dnCz}{\dnCo}{\dnCo}{\dnCo}{\dnCo}{\dnCz}{\dnCo}{\dnCz}{\dnCo}{\dnCz}{\dnCo}{\dnCu}{\dnCz}{\dnCo}{\dnCh}{\dnCh}{\dnCh}{\dnCh}{\dnCh}{\dnCh}{\dnCh}{\dnCh}{\dnCh}{\dnCh}{\dnCh}{\dnCh}{\dnCh}{\dnCo}{\dnCz}{\dnCz}} & \dnW[W]{}: & {{\dnCo}{\dnCz}{\dnCo}{\dnCo}{\dnCo}{\dnCz}{\dnCz}{\dnCo}{\dnCo}{\dnCo}{\dnCo}{\dnCz}{\dnCo}{\dnCz}{\dnCz}{\dnCz}{\dnCo}{\dnCo}{\dnCz}{\dnCz}{\dnCz}{\dnCz}{\dnCo}{\dnCo}{\dnCo}{\dnCo}{\dnCo}{\dnCz}{\dnCo}{\dnCz}{\dnCz}{\dnCo}} \\
 \dnI{16} & \dnW: & {{\dnCn}{\dnCz}{\dnCz}{\dnCn}{\dnCz}{\dnCu}{\dnCn}{\dnCz}{\dnCo}{\dnCo}{\dnCz}{\dnCo}{\dnCz}{\dnCz}{\dnCo}{\dnCz}{\dnCh}{\dnCh}{\dnCh}{\dnCh}{\dnCh}{\dnCh}{\dnCh}{\dnCh}{\dnCh}{\dnCh}{\dnCh}{\dnCh}{\dnCh}{\dnCo}{\dnCo}{\dnCo}} & & \\
 \dnI{17} & \dnW: & {{\dnCz}{\dnCz}{\dnCz}{\dnCo}{\dnCo}{\dnCo}{\dnCo}{\dnCo}{\dnCz}{\dnCz}{\dnCo}{\dnCo}{\dnCo}{\dnCz}{\dnCo}{\dnCz}{\dnCh}{\dnCh}{\dnCh}{\dnCh}{\dnCh}{\dnCh}{\dnCh}{\dnCh}{\dnCh}{\dnCh}{\dnCh}{\dnCh}{\dnCh}{\dnCo}{\dnCo}{\dnCz}} & & \\
 \dnI{18} & \dnW: & {{\dnCz}{\dnCo}{\dnCz}{\dnCo}{\dnCz}{\dnCo}{\dnCo}{\dnCo}{\dnCz}{\dnCz}{\dnCz}{\dnCz}{\dnCo}{\dnCo}{\dnCz}{\dnCo}{\dnCh}{\dnCh}{\dnCh}{\dnCh}{\dnCh}{\dnCh}{\dnCh}{\dnCh}{\dnCh}{\dnCh}{\dnCh}{\dnCh}{\dnCh}{\dnCo}{\dnCz}{\dnCz}} & & \\
 \dnI{19} & \dnW: & {{\dnCu}{\dnCo}{\dnCn}{\dnCo}{\dnCz}{\dnCz}{\dnCz}{\dnCz}{\dnCz}{\dnCz}{\dnCz}{\dnCo}{\dnCz}{\dnCo}{\dnCo}{\dnCo}{\dnCh}{\dnCh}{\dnCh}{\dnCh}{\dnCh}{\dnCh}{\dnCh}{\dnCh}{\dnCh}{\dnCh}{\dnCh}{\dnCh}{\dnCh}{\dnCo}{\dnCz}{\dnCz}} & & \\
 \dnI{20} & \dnW: & {{\dnCn}{\dnCo}{\dnCu}{\dnCn}{\dnCo}{\dnCz}{\dnCz}{\dnCo}{\dnCo}{\dnCo}{\dnCo}{\dnCo}{\dnCo}{\dnCo}{\dnCz}{\dnCo}{\dnCo}{\dnCo}{\dnCz}{\dnCo}{\dnCz}{\dnCz}{\dnCz}{\dnCz}{\dnCz}{\dnCz}{\dnCo}{\dnCo}{\dnCz}{\dnCo}{\dnCz}{\dnCz}} & & \\
 \dnI{21} & \dnW: & {{\dnCo}{\dnCo}{\dnCo}{\dnCo}{\dnCz}{\dnCz}{\dnCo}{\dnCo}{\dnCo}{\dnCz}{\dnCo}{\dnCo}{\dnCz}{\dnCz}{\dnCz}{\dnCz}{\dnCz}{\dnCo}{\dnCz}{\dnCo}{\dnCo}{\dnCo}{\dnCo}{\dnCo}{\dnCo}{\dnCo}{\dnCz}{\dnCo}{\dnCz}{\dnCo}{\dnCz}{\dnCz}} & & \\
 \dnI{22} & \dnW: & {{\dnCz}{\dnCo}{\dnCz}{\dnCo}{\dnCo}{\dnCo}{\dnCz}{\dnCo}{\dnCo}{\dnCo}{\dnCz}{\dnCz}{\dnCo}{\dnCo}{\dnCz}{\dnCo}{\dnCz}{\dnCz}{\dnCo}{\dnCo}{\dnCz}{\dnCz}{\dnCo}{\dnCo}{\dnCz}{\dnCz}{\dnCo}{\dnCo}{\dnCo}{\dnCz}{\dnCo}{\dnCz}} & & \\
 \dnI{23} & \dnW: & {{\dnCz}{\dnCo}{\dnCz}{\dnCo}{\dnCz}{\dnCz}{\dnCz}{\dnCz}{\dnCo}{\dnCo}{\dnCo}{\dnCz}{\dnCo}{\dnCo}{\dnCo}{\dnCz}{\dnCo}{\dnCo}{\dnCz}{\dnCz}{\dnCz}{\dnCo}{\dnCo}{\dnCo}{\dnCo}{\dnCz}{\dnCz}{\dnCz}{\dnCo}{\dnCo}{\dnCo}{\dnCo}} & & \\
 \dnI{24} & \dnW: & {{\dnCz}{\dnCz}{\dnCz}{\dnCz}{\dnCz}{\dnCz}{\dnCo}{\dnCz}{\dnCz}{\dnCz}{\dnCz}{\dnCo}{\dnCz}{\dnCz}{\dnCo}{\dnCz}{\dnCz}{\dnCz}{\dnCo}{\dnCo}{\dnCz}{\dnCo}{\dnCo}{\dnCo}{\dnCz}{\dnCz}{\dnCz}{\dnCo}{\dnCo}{\dnCz}{\dnCo}{\dnCz}} & & \\
 \dnI{25} & \dnW: & {{\dnCo}{\dnCz}{\dnCo}{\dnCo}{\dnCz}{\dnCz}{\dnCz}{\dnCz}{\dnCo}{\dnCz}{\dnCz}{\dnCo}{\dnCz}{\dnCo}{\dnCo}{\dnCz}{\dnCz}{\dnCz}{\dnCz}{\dnCo}{\dnCz}{\dnCo}{\dnCz}{\dnCz}{\dnCo}{\dnCo}{\dnCo}{\dnCz}{\dnCo}{\dnCz}{\dnCo}{\dnCz}} & & \\
 \dnI{26} & \dnW: & {{\dnCz}{\dnCz}{\dnCz}{\dnCz}{\dnCo}{\dnCz}{\dnCo}{\dnCz}{\dnCo}{\dnCz}{\dnCz}{\dnCz}{\dnCo}{\dnCz}{\dnCz}{\dnCo}{\dnCz}{\dnCo}{\dnCo}{\dnCo}{\dnCz}{\dnCo}{\dnCo}{\dnCo}{\dnCz}{\dnCo}{\dnCz}{\dnCz}{\dnCz}{\dnCz}{\dnCz}{\dnCo}} & & \\
 \dnI{27} & \dnW: & {{\dnCz}{\dnCz}{\dnCz}{\dnCz}{\dnCz}{\dnCo}{\dnCo}{\dnCz}{\dnCo}{\dnCo}{\dnCo}{\dnCz}{\dnCo}{\dnCo}{\dnCo}{\dnCo}{\dnCz}{\dnCo}{\dnCz}{\dnCo}{\dnCo}{\dnCz}{\dnCo}{\dnCz}{\dnCo}{\dnCz}{\dnCo}{\dnCo}{\dnCz}{\dnCz}{\dnCo}{\dnCo}} & & \\
 \dnI{28} & \dnW: & {{\dnCo}{\dnCz}{\dnCo}{\dnCo}{\dnCz}{\dnCo}{\dnCo}{\dnCz}{\dnCz}{\dnCo}{\dnCz}{\dnCo}{\dnCo}{\dnCo}{\dnCz}{\dnCo}{\dnCz}{\dnCo}{\dnCo}{\dnCz}{\dnCo}{\dnCo}{\dnCz}{\dnCz}{\dnCz}{\dnCz}{\dnCo}{\dnCz}{\dnCz}{\dnCo}{\dnCz}{\dnCo}} & & \\
 \dnI{29} & \dnW: & {{\dnCo}{\dnCz}{\dnCo}{\dnCz}{\dnCz}{\dnCz}{\dnCo}{\dnCz}{\dnCz}{\dnCz}{\dnCz}{\dnCz}{\dnCo}{\dnCo}{\dnCz}{\dnCo}{\dnCz}{\dnCo}{\dnCz}{\dnCz}{\dnCo}{\dnCz}{\dnCz}{\dnCz}{\dnCz}{\dnCo}{\dnCo}{\dnCz}{\dnCo}{\dnCz}{\dnCz}{\dnCo}} & & \\
 \dnI{30} & \dnW: & {{\dnCz}{\dnCz}{\dnCo}{\dnCz}{\dnCo}{\dnCz}{\dnCz}{\dnCo}{\dnCo}{\dnCo}{\dnCz}{\dnCo}{\dnCz}{\dnCo}{\dnCo}{\dnCo}{\dnCo}{\dnCo}{\dnCz}{\dnCz}{\dnCz}{\dnCo}{\dnCo}{\dnCo}{\dnCz}{\dnCo}{\dnCo}{\dnCz}{\dnCz}{\dnCz}{\dnCo}{\dnCo}} & & \\
 \dnI{31} & \dnW: & {{\dnCo}{\dnCo}{\dnCo}{\dnCo}{\dnCo}{\dnCo}{\dnCz}{\dnCz}{\dnCo}{\dnCz}{\dnCz}{\dnCo}{\dnCz}{\dnCz}{\dnCo}{\dnCz}{\dnCo}{\dnCo}{\dnCz}{\dnCo}{\dnCz}{\dnCo}{\dnCo}{\dnCo}{\dnCo}{\dnCz}{\dnCo}{\dnCo}{\dnCz}{\dnCo}{\dnCo}{\dnCz}} & & \\
 \dnI{32} & \dnW: & {{\dnCz}{\dnCo}{\dnCz}{\dnCz}{\dnCo}{\dnCo}{\dnCo}{\dnCo}{\dnCo}{\dnCo}{\dnCz}{\dnCo}{\dnCz}{\dnCz}{\dnCo}{\dnCz}{\dnCz}{\dnCo}{\dnCo}{\dnCz}{\dnCo}{\dnCz}{\dnCz}{\dnCz}{\dnCz}{\dnCz}{\dnCo}{\dnCz}{\dnCo}{\dnCo}{\dnCo}{\dnCo}} & & \\
 \dnI{33} & \dnW: & {{\dnCz}{\dnCz}{\dnCo}{\dnCo}{\dnCo}{\dnCz}{\dnCz}{\dnCz}{\dnCz}{\dnCz}{\dnCo}{\dnCo}{\dnCo}{\dnCo}{\dnCz}{\dnCo}{\dnCz}{\dnCo}{\dnCo}{\dnCz}{\dnCo}{\dnCo}{\dnCo}{\dnCz}{\dnCo}{\dnCo}{\dnCo}{\dnCz}{\dnCz}{\dnCo}{\dnCz}{\dnCz}} & & \\
 \dnI{34} & \dnW: & {{\dnCz}{\dnCz}{\dnCo}{\dnCz}{\dnCz}{\dnCz}{\dnCz}{\dnCz}{\dnCz}{\dnCo}{\dnCo}{\dnCo}{\dnCz}{\dnCo}{\dnCz}{\dnCo}{\dnCo}{\dnCo}{\dnCo}{\dnCz}{\dnCo}{\dnCz}{\dnCz}{\dnCz}{\dnCz}{\dnCz}{\dnCz}{\dnCo}{\dnCz}{\dnCo}{\dnCz}{\dnCo}} & & \\
 \dnI{35} & \dnW: & {{\dnCn}{\dnCz}{\dnCo}{\dnCz}{\dnCz}{\dnCz}{\dnCz}{\dnCz}{\dnCz}{\dnCz}{\dnCo}{\dnCo}{\dnCz}{\dnCz}{\dnCo}{\dnCo}{\dnCz}{\dnCz}{\dnCz}{\dnCz}{\dnCz}{\dnCo}{\dnCz}{\dnCz}{\dnCz}{\dnCo}{\dnCo}{\dnCo}{\dnCz}{\dnCz}{\dnCo}{\dnCz}} & & \\
 \dnI{36} & \dnW: & {{\dnCn}{\dnCz}{\dnCz}{\dnCz}{\dnCz}{\dnCo}{\dnCo}{\dnCo}{\dnCo}{\dnCo}{\dnCo}{\dnCz}{\dnCo}{\dnCz}{\dnCo}{\dnCo}{\dnCo}{\dnCo}{\dnCz}{\dnCo}{\dnCo}{\dnCo}{\dnCo}{\dnCz}{\dnCz}{\dnCo}{\dnCz}{\dnCo}{\dnCo}{\dnCz}{\dnCz}{\dnCo}} & & \\
 \dnI{37} & \dnW: & {{\dnCo}{\dnCo}{\dnCz}{\dnCz}{\dnCo}{\dnCz}{\dnCz}{\dnCz}{\dnCz}{\dnCz}{\dnCz}{\dnCo}{\dnCo}{\dnCz}{\dnCo}{\dnCz}{\dnCz}{\dnCo}{\dnCz}{\dnCz}{\dnCz}{\dnCz}{\dnCo}{\dnCo}{\dnCz}{\dnCz}{\dnCz}{\dnCz}{\dnCo}{\dnCo}{\dnCz}{\dnCz}} & & \\
 \dnI{38} & \dnW: & {{\dnCo}{\dnCz}{\dnCo}{\dnCo}{\dnCz}{\dnCz}{\dnCz}{\dnCz}{\dnCz}{\dnCo}{\dnCo}{\dnCz}{\dnCz}{\dnCo}{\dnCo}{\dnCo}{\dnCo}{\dnCo}{\dnCo}{\dnCz}{\dnCo}{\dnCz}{\dnCz}{\dnCo}{\dnCo}{\dnCz}{\dnCo}{\dnCz}{\dnCo}{\dnCo}{\dnCz}{\dnCz}} & & \\
 \dnI{39} & \dnW: & {{\dnCz}{\dnCz}{\dnCz}{\dnCo}{\dnCz}{\dnCz}{\dnCo}{\dnCz}{\dnCz}{\dnCz}{\dnCz}{\dnCz}{\dnCo}{\dnCz}{\dnCo}{\dnCz}{\dnCz}{\dnCz}{\dnCz}{\dnCo}{\dnCo}{\dnCz}{\dnCo}{\dnCo}{\dnCz}{\dnCz}{\dnCz}{\dnCo}{\dnCo}{\dnCo}{\dnCz}{\dnCz}} & & \\
 \dnI{40} & \dnW: & {{\dnCo}{\dnCo}{\dnCz}{\dnCz}{\dnCz}{\dnCz}{\dnCz}{\dnCz}{\dnCz}{\dnCo}{\dnCz}{\dnCz}{\dnCo}{\dnCz}{\dnCz}{\dnCz}{\dnCz}{\dnCo}{\dnCo}{\dnCo}{\dnCz}{\dnCz}{\dnCz}{\dnCo}{\dnCo}{\dnCz}{\dnCz}{\dnCz}{\dnCz}{\dnCo}{\dnCz}{\dnCo}} & & \\
 \dnI{41} & \dnW: & {{\dnCz}{\dnCz}{\dnCz}{\dnCz}{\dnCz}{\dnCo}{\dnCo}{\dnCz}{\dnCo}{\dnCz}{\dnCz}{\dnCz}{\dnCz}{\dnCo}{\dnCo}{\dnCz}{\dnCo}{\dnCo}{\dnCo}{\dnCo}{\dnCz}{\dnCo}{\dnCz}{\dnCo}{\dnCz}{\dnCz}{\dnCo}{\dnCz}{\dnCz}{\dnCo}{\dnCo}{\dnCz}} & & \\
 \dnI{42} & \dnW: & {{\dnCz}{\dnCo}{\dnCz}{\dnCz}{\dnCo}{\dnCo}{\dnCo}{\dnCz}{\dnCo}{\dnCo}{\dnCz}{\dnCo}{\dnCo}{\dnCo}{\dnCz}{\dnCo}{\dnCo}{\dnCo}{\dnCo}{\dnCo}{\dnCo}{\dnCo}{\dnCo}{\dnCz}{\dnCo}{\dnCz}{\dnCz}{\dnCz}{\dnCz}{\dnCo}{\dnCo}{\dnCz}} & & \\
 \dnI{43} & \dnW: & {{\dnCz}{\dnCo}{\dnCz}{\dnCo}{\dnCz}{\dnCz}{\dnCz}{\dnCz}{\dnCz}{\dnCo}{\dnCo}{\dnCz}{\dnCz}{\dnCz}{\dnCo}{\dnCo}{\dnCo}{\dnCo}{\dnCz}{\dnCo}{\dnCz}{\dnCz}{\dnCz}{\dnCz}{\dnCz}{\dnCo}{\dnCo}{\dnCz}{\dnCo}{\dnCo}{\dnCz}{\dnCo}} & & \\
 \dnI{44} & \dnW: & {{\dnCo}{\dnCo}{\dnCo}{\dnCo}{\dnCo}{\dnCz}{\dnCz}{\dnCz}{\dnCz}{\dnCz}{\dnCz}{\dnCo}{\dnCz}{\dnCo}{\dnCo}{\dnCz}{\dnCo}{\dnCo}{\dnCo}{\dnCo}{\dnCz}{\dnCo}{\dnCo}{\dnCo}{\dnCz}{\dnCz}{\dnCz}{\dnCz}{\dnCo}{\dnCo}{\dnCz}{\dnCz}} & & \\
 \dnI{45} & \dnW: & {{\dnCo}{\dnCz}{\dnCz}{\dnCz}{\dnCo}{\dnCz}{\dnCo}{\dnCz}{\dnCo}{\dnCo}{\dnCz}{\dnCo}{\dnCo}{\dnCz}{\dnCo}{\dnCo}{\dnCz}{\dnCz}{\dnCo}{\dnCz}{\dnCo}{\dnCo}{\dnCz}{\dnCz}{\dnCz}{\dnCz}{\dnCz}{\dnCz}{\dnCz}{\dnCo}{\dnCz}{\dnCz}} & & \\
 \dnI{46} & \dnW: & {{\dnCo}{\dnCz}{\dnCz}{\dnCz}{\dnCz}{\dnCz}{\dnCo}{\dnCz}{\dnCo}{\dnCz}{\dnCz}{\dnCo}{\dnCo}{\dnCz}{\dnCz}{\dnCo}{\dnCz}{\dnCo}{\dnCz}{\dnCo}{\dnCo}{\dnCz}{\dnCz}{\dnCz}{\dnCo}{\dnCo}{\dnCz}{\dnCo}{\dnCo}{\dnCo}{\dnCz}{\dnCz}} & & \\
 \dnI{47} & \dnW: & {{\dnCo}{\dnCz}{\dnCz}{\dnCz}{\dnCz}{\dnCz}{\dnCz}{\dnCo}{\dnCo}{\dnCo}{\dnCo}{\dnCz}{\dnCz}{\dnCo}{\dnCz}{\dnCo}{\dnCo}{\dnCz}{\dnCo}{\dnCo}{\dnCz}{\dnCo}{\dnCz}{\dnCz}{\dnCo}{\dnCz}{\dnCo}{\dnCo}{\dnCo}{\dnCo}{\dnCz}{\dnCo}} & & \\
\hline
\end{tabular}
\hspace{10pt}$\Rightarrow$\hspace{10pt}
\begin{tabular}{|r|c|c|c|c|}
\hline
  $i$  &      &  $\nabla{S_{i,0}}$  &  $\nabla{S_{i,1}}$  &  $\nabla{S_{i,2}}$ \\
\hline
 \dnI{-4} & \dnW: & {{\dnCz}{\dnCo}{\dnCo}{\dnCz}{\dnCz}{\dnCo}{\dnCo}{\dnCo}{\dnCz}{\dnCo}{\dnCz}{\dnCz}{\dnCz}{\dnCo}{\dnCz}{\dnCo}{\dnCz}{\dnCz}{\dnCo}{\dnCz}{\dnCz}{\dnCz}{\dnCo}{\dnCo}{\dnCz}{\dnCz}{\dnCz}{\dnCz}{\dnCz}{\dnCz}{\dnCz}{\dnCo}} & & \\
 \dnI{-3} & \dnW: & {{\dnCz}{\dnCz}{\dnCz}{\dnCo}{\dnCz}{\dnCz}{\dnCz}{\dnCz}{\dnCz}{\dnCz}{\dnCo}{\dnCo}{\dnCz}{\dnCz}{\dnCo}{\dnCz}{\dnCz}{\dnCo}{\dnCz}{\dnCo}{\dnCz}{\dnCo}{\dnCz}{\dnCz}{\dnCz}{\dnCo}{\dnCo}{\dnCo}{\dnCz}{\dnCo}{\dnCo}{\dnCz}} & & \\
 \dnI{-2} & \dnW: & {{\dnCo}{\dnCz}{\dnCz}{\dnCo}{\dnCo}{\dnCz}{\dnCz}{\dnCz}{\dnCo}{\dnCz}{\dnCo}{\dnCo}{\dnCo}{\dnCz}{\dnCo}{\dnCz}{\dnCo}{\dnCo}{\dnCz}{\dnCo}{\dnCo}{\dnCo}{\dnCz}{\dnCz}{\dnCo}{\dnCo}{\dnCo}{\dnCo}{\dnCo}{\dnCo}{\dnCo}{\dnCz}} & & \\
 \dnI{-1} & \dnW: & {{\dnCo}{\dnCo}{\dnCo}{\dnCz}{\dnCo}{\dnCo}{\dnCo}{\dnCo}{\dnCo}{\dnCo}{\dnCz}{\dnCz}{\dnCo}{\dnCo}{\dnCz}{\dnCo}{\dnCo}{\dnCz}{\dnCo}{\dnCz}{\dnCo}{\dnCz}{\dnCo}{\dnCo}{\dnCo}{\dnCz}{\dnCz}{\dnCz}{\dnCo}{\dnCz}{\dnCz}{\dnCo}} & & \\
 \dnI{0} & \dnW: & {{\dnCz}{\dnCo}{\dnCo}{\dnCz}{\dnCo}{\dnCz}{\dnCo}{\dnCo}{\dnCo}{\dnCo}{\dnCz}{\dnCo}{\dnCz}{\dnCo}{\dnCz}{\dnCz}{\dnCo}{\dnCo}{\dnCo}{\dnCz}{\dnCz}{\dnCo}{\dnCz}{\dnCz}{\dnCz}{\dnCz}{\dnCz}{\dnCo}{\dnCz}{\dnCz}{\dnCo}{\dnCz}} & \dnW[W]{}: & {{\dnCz}{\dnCo}{\dnCz}{\dnCz}{\dnCo}{\dnCo}{\dnCz}{\dnCo}{\dnCz}{\dnCo}{\dnCo}{\dnCo}{\dnCo}{\dnCz}{\dnCo}{\dnCz}{\dnCo}{\dnCz}{\dnCz}{\dnCo}{\dnCo}{\dnCo}{\dnCz}{\dnCz}{\dnCo}{\dnCz}{\dnCz}{\dnCz}{\dnCz}{\dnCz}{\dnCo}{\dnCo}} \\
 \dnI{1} & \dnW: & {{\dnCz}{\dnCo}{\dnCo}{\dnCo}{\dnCz}{\dnCo}{\dnCo}{\dnCz}{\dnCz}{\dnCo}{\dnCz}{\dnCz}{\dnCo}{\dnCo}{\dnCo}{\dnCo}{\dnCo}{\dnCo}{\dnCo}{\dnCz}{\dnCo}{\dnCo}{\dnCo}{\dnCz}{\dnCz}{\dnCu}{\dnCo}{\dnCo}{\dnCz}{\dnCz}{\dnCz}{\dnCo}} & \dnW[W]{}: & {{\dnCu}{\dnCo}{\dnCz}{\dnCo}{\dnCz}{\dnCo}{\dnCo}{\dnCz}{\dnCo}{\dnCo}{\dnCz}{\dnCz}{\dnCo}{\dnCz}{\dnCo}{\dnCo}{\dnCo}{\dnCz}{\dnCz}{\dnCo}{\dnCz}{\dnCz}{\dnCo}{\dnCz}{\dnCz}{\dnCo}{\dnCo}{\dnCo}{\dnCo}{\dnCz}{\dnCo}{\dnCz}} \\
 \dnI{2} & \dnW: & {{\dnCo}{\dnCz}{\dnCo}{\dnCz}{\dnCo}{\dnCz}{\dnCo}{\dnCo}{\dnCz}{\dnCo}{\dnCz}{\dnCz}{\dnCz}{\dnCz}{\dnCz}{\dnCz}{\dnCz}{\dnCo}{\dnCo}{\dnCo}{\dnCz}{\dnCu}{\dnCz}{\dnCo}{\dnCn}{\dnCo}{\dnCo}{\dnCo}{\dnCz}{\dnCz}{\dnCo}{\dnCz}} & \dnW[W]{}: & {{\dnCn}{\dnCz}{\dnCo}{\dnCn}{\dnCo}{\dnCz}{\dnCz}{\dnCo}{\dnCo}{\dnCo}{\dnCz}{\dnCo}{\dnCz}{\dnCo}{\dnCz}{\dnCo}{\dnCo}{\dnCz}{\dnCo}{\dnCz}{\dnCz}{\dnCo}{\dnCz}{\dnCo}{\dnCz}{\dnCo}{\dnCo}{\dnCo}{\dnCo}{\dnCz}{\dnCz}{\dnCz}} \\
 \dnI{3} & \dnW: & {{\dnCo}{\dnCz}{\dnCo}{\dnCz}{\dnCo}{\dnCo}{\dnCu}{\dnCo}{\dnCz}{\dnCz}{\dnCo}{\dnCo}{\dnCo}{\dnCo}{\dnCz}{\dnCo}{\dnCz}{\dnCo}{\dnCz}{\dnCo}{\dnCz}{\dnCz}{\dnCo}{\dnCz}{\dnCz}{\dnCo}{\dnCz}{\dnCo}{\dnCz}{\dnCz}{\dnCz}{\dnCo}} & \dnW[W]{}: & {{\dnCz}{\dnCo}{\dnCz}{\dnCo}{\dnCz}{\dnCo}{\dnCo}{\dnCo}{\dnCo}{\dnCz}{\dnCo}{\dnCz}{\dnCz}{\dnCo}{\dnCo}{\dnCo}{\dnCo}{\dnCz}{\dnCo}{\dnCz}{\dnCz}{\dnCo}{\dnCz}{\dnCo}{\dnCo}{\dnCo}{\dnCo}{\dnCz}{\dnCo}{\dnCo}{\dnCo}{\dnCz}} \\
 \dnI{4} & \dnW: & {{\dnCz}{\dnCz}{\dnCo}{\dnCz}{\dnCo}{\dnCo}{\dnCz}{\dnCz}{\dnCz}{\dnCo}{\dnCo}{\dnCz}{\dnCz}{\dnCz}{\dnCo}{\dnCo}{\dnCz}{\dnCo}{\dnCz}{\dnCo}{\dnCz}{\dnCo}{\dnCz}{\dnCo}{\dnCo}{\dnCo}{\dnCo}{\dnCo}{\dnCz}{\dnCz}{\dnCo}{\dnCz}} & \dnW[W]{}: & {{\dnCo}{\dnCo}{\dnCz}{\dnCo}{\dnCo}{\dnCo}{\dnCo}{\dnCz}{\dnCz}{\dnCo}{\dnCo}{\dnCo}{\dnCz}{\dnCo}{\dnCz}{\dnCz}{\dnCo}{\dnCz}{\dnCz}{\dnCz}{\dnCo}{\dnCz}{\dnCo}{\dnCz}{\dnCz}{\dnCz}{\dnCo}{\dnCo}{\dnCo}{\dnCo}{\dnCz}{\dnCz}} \\
 \dnI{5} & \dnW: & {{\dnCz}{\dnCz}{\dnCz}{\dnCo}{\dnCo}{\dnCz}{\dnCo}{\dnCz}{\dnCz}{\dnCo}{\dnCo}{\dnCz}{\dnCz}{\dnCz}{\dnCo}{\dnCz}{\dnCo}{\dnCz}{\dnCu}{\dnCo}{\dnCo}{\dnCz}{\dnCo}{\dnCz}{\dnCz}{\dnCz}{\dnCz}{\dnCz}{\dnCz}{\dnCz}{\dnCz}{\dnCo}} & \dnW[W]{}: & {{\dnCo}{\dnCo}{\dnCz}{\dnCo}{\dnCo}{\dnCo}{\dnCz}{\dnCz}{\dnCo}{\dnCo}{\dnCz}{\dnCz}{\dnCz}{\dnCz}{\dnCo}{\dnCo}{\dnCz}{\dnCo}{\dnCo}{\dnCz}{\dnCz}{\dnCo}{\dnCo}{\dnCz}{\dnCo}{\dnCz}{\dnCo}{\dnCo}{\dnCz}{\dnCz}{\dnCo}{\dnCo}} \\
 \dnI{6} & \dnW: & {{\dnCz}{\dnCz}{\dnCz}{\dnCo}{\dnCo}{\dnCz}{\dnCo}{\dnCo}{\dnCz}{\dnCz}{\dnCu}{\dnCn}{\dnCu}{\dnCu}{\dnCo}{\dnCo}{\dnCz}{\dnCz}{\dnCz}{\dnCo}{\dnCz}{\dnCz}{\dnCz}{\dnCz}{\dnCz}{\dnCo}{\dnCo}{\dnCo}{\dnCo}{\dnCz}{\dnCo}{\dnCz}} & \dnW[W]{}: & {{\dnCo}{\dnCz}{\dnCo}{\dnCo}{\dnCz}{\dnCo}{\dnCo}{\dnCz}{\dnCo}{\dnCz}{\dnCz}{\dnCz}{\dnCz}{\dnCz}{\dnCo}{\dnCo}{\dnCo}{\dnCz}{\dnCo}{\dnCz}{\dnCz}{\dnCz}{\dnCz}{\dnCz}{\dnCz}{\dnCz}{\dnCo}{\dnCz}{\dnCz}{\dnCz}{\dnCz}{\dnCz}} \\
 \dnI{7} & \dnW: & {{\dnCz}{\dnCz}{\dnCo}{\dnCz}{\dnCo}{\dnCz}{\dnCo}{\dnCo}{\dnCo}{\dnCz}{\dnCz}{\dnCz}{\dnCz}{\dnCz}{\dnCz}{\dnCo}{\dnCz}{\dnCu}{\dnCn}{\dnCn}{\dnCz}{\dnCo}{\dnCo}{\dnCz}{\dnCz}{\dnCo}{\dnCz}{\dnCo}{\dnCz}{\dnCz}{\dnCz}{\dnCz}} & \dnW[W]{}: & {{\dnCz}{\dnCz}{\dnCo}{\dnCo}{\dnCo}{\dnCz}{\dnCo}{\dnCo}{\dnCz}{\dnCz}{\dnCo}{\dnCz}{\dnCo}{\dnCz}{\dnCo}{\dnCz}{\dnCz}{\dnCo}{\dnCz}{\dnCo}{\dnCo}{\dnCo}{\dnCz}{\dnCo}{\dnCo}{\dnCz}{\dnCz}{\dnCo}{\dnCo}{\dnCo}{\dnCo}{\dnCo}} \\
 \dnI{8} & \dnW: & {{\dnCz}{\dnCo}{\dnCo}{\dnCo}{\dnCz}{\dnCz}{\dnCo}{\dnCo}{\dnCz}{\dnCz}{\dnCo}{\dnCz}{\dnCz}{\dnCz}{\dnCo}{\dnCu}{\dnCo}{\dnCo}{\dnCo}{\dnCo}{\dnCo}{\dnCo}{\dnCo}{\dnCo}{\dnCo}{\dnCz}{\dnCo}{\dnCo}{\dnCz}{\dnCz}{\dnCz}{\dnCz}} & \dnW[W]{}: & {{\dnCo}{\dnCo}{\dnCz}{\dnCz}{\dnCz}{\dnCo}{\dnCo}{\dnCz}{\dnCo}{\dnCz}{\dnCz}{\dnCo}{\dnCo}{\dnCo}{\dnCz}{\dnCo}{\dnCz}{\dnCo}{\dnCo}{\dnCo}{\dnCz}{\dnCz}{\dnCz}{\dnCo}{\dnCo}{\dnCz}{\dnCo}{\dnCo}{\dnCz}{\dnCz}{\dnCo}{\dnCo}} \\
 \dnI{9} & \dnW: & {{\dnCo}{\dnCz}{\dnCo}{\dnCz}{\dnCo}{\dnCo}{\dnCn}{\dnCz}{\dnCo}{\dnCu}{\dnCn}{\dnCn}{\dnCu}{\dnCz}{\dnCz}{\dnCz}{\dnCo}{\dnCo}{\dnCo}{\dnCo}{\dnCo}{\dnCz}{\dnCz}{\dnCo}{\dnCo}{\dnCz}{\dnCz}{\dnCo}{\dnCo}{\dnCo}{\dnCo}{\dnCo}} & \dnW[W]{}: & {{\dnCo}{\dnCo}{\dnCo}{\dnCo}{\dnCo}{\dnCz}{\dnCz}{\dnCo}{\dnCo}{\dnCo}{\dnCo}{\dnCz}{\dnCo}{\dnCz}{\dnCz}{\dnCo}{\dnCo}{\dnCz}{\dnCz}{\dnCo}{\dnCz}{\dnCz}{\dnCz}{\dnCo}{\dnCo}{\dnCz}{\dnCz}{\dnCo}{\dnCo}{\dnCz}{\dnCz}{\dnCz}} \\
 \dnI{10} & \dnW: & {{\dnCo}{\dnCz}{\dnCn}{\dnCz}{\dnCz}{\dnCo}{\dnCz}{\dnCz}{\dnCo}{\dnCz}{\dnCz}{\dnCz}{\dnCz}{\dnCo}{\dnCz}{\dnCo}{\dnCz}{\dnCo}{\dnCz}{\dnCz}{\dnCz}{\dnCz}{\dnCz}{\dnCz}{\dnCo}{\dnCz}{\dnCo}{\dnCz}{\dnCo}{\dnCo}{\dnCo}{\dnCz}} & \dnW[W]{}: & {{\dnCo}{\dnCo}{\dnCz}{\dnCo}{\dnCz}{\dnCo}{\dnCo}{\dnCo}{\dnCo}{\dnCz}{\dnCz}{\dnCo}{\dnCo}{\dnCo}{\dnCo}{\dnCo}{\dnCo}{\dnCz}{\dnCz}{\dnCz}{\dnCz}{\dnCz}{\dnCz}{\dnCz}{\dnCz}{\dnCo}{\dnCz}{\dnCo}{\dnCo}{\dnCo}{\dnCo}{\dnCz}} \\
 \dnI{11} & \dnW: & {{\dnCu}{\dnCo}{\dnCz}{\dnCz}{\dnCz}{\dnCo}{\dnCo}{\dnCz}{\dnCo}{\dnCz}{\dnCo}{\dnCo}{\dnCz}{\dnCo}{\dnCo}{\dnCz}{\dnCz}{\dnCo}{\dnCz}{\dnCz}{\dnCo}{\dnCz}{\dnCo}{\dnCz}{\dnCo}{\dnCo}{\dnCo}{\dnCo}{\dnCo}{\dnCo}{\dnCo}{\dnCo}} & \dnW[W]{}: & {{\dnCo}{\dnCz}{\dnCo}{\dnCz}{\dnCz}{\dnCo}{\dnCo}{\dnCz}{\dnCz}{\dnCz}{\dnCo}{\dnCo}{\dnCo}{\dnCz}{\dnCo}{\dnCo}{\dnCo}{\dnCz}{\dnCo}{\dnCo}{\dnCz}{\dnCz}{\dnCo}{\dnCz}{\dnCo}{\dnCo}{\dnCo}{\dnCz}{\dnCo}{\dnCz}{\dnCz}{\dnCz}} \\
 \dnI{12} & \dnW: & {{\dnCz}{\dnCz}{\dnCo}{\dnCz}{\dnCo}{\dnCo}{\dnCu}{\dnCz}{\dnCz}{\dnCu}{\dnCo}{\dnCz}{\dnCo}{\dnCz}{\dnCo}{\dnCo}{\dnCo}{\dnCo}{\dnCo}{\dnCo}{\dnCo}{\dnCo}{\dnCz}{\dnCz}{\dnCz}{\dnCo}{\dnCo}{\dnCo}{\dnCo}{\dnCz}{\dnCo}{\dnCo}} & \dnW[W]{}: & {{\dnCz}{\dnCo}{\dnCz}{\dnCz}{\dnCz}{\dnCo}{\dnCz}{\dnCo}{\dnCo}{\dnCo}{\dnCz}{\dnCo}{\dnCo}{\dnCo}{\dnCz}{\dnCn}{\dnCo}{\dnCz}{\dnCz}{\dnCz}{\dnCo}{\dnCo}{\dnCo}{\dnCz}{\dnCz}{\dnCz}{\dnCo}{\dnCo}{\dnCz}{\dnCz}{\dnCz}{\dnCo}} \\
 \dnI{13} & \dnW: & {{\dnCo}{\dnCz}{\dnCu}{\dnCn}{\dnCo}{\dnCn}{\dnCz}{\dnCo}{\dnCz}{\dnCz}{\dnCo}{\dnCo}{\dnCz}{\dnCz}{\dnCz}{\dnCo}{\dnCz}{\dnCo}{\dnCz}{\dnCz}{\dnCz}{\dnCz}{\dnCz}{\dnCo}{\dnCo}{\dnCo}{\dnCo}{\dnCz}{\dnCz}{\dnCo}{\dnCz}{\dnCo}} & \dnW[W]{}: & {{\dnCo}{\dnCz}{\dnCz}{\dnCo}{\dnCz}{\dnCo}{\dnCo}{\dnCo}{\dnCo}{\dnCo}{\dnCo}{\dnCz}{\dnCz}{\dnCz}{\dnCo}{\dnCo}{\dnCz}{\dnCz}{\dnCz}{\dnCo}{\dnCo}{\dnCo}{\dnCo}{\dnCo}{\dnCo}{\dnCo}{\dnCo}{\dnCz}{\dnCz}{\dnCo}{\dnCz}{\dnCo}} \\
 \dnI{14} & \dnW: & {{\dnCz}{\dnCz}{\dnCz}{\dnCz}{\dnCo}{\dnCz}{\dnCo}{\dnCz}{\dnCz}{\dnCo}{\dnCz}{\dnCo}{\dnCz}{\dnCz}{\dnCz}{\dnCo}{\dnCo}{\dnCz}{\dnCz}{\dnCz}{\dnCo}{\dnCz}{\dnCz}{\dnCz}{\dnCo}{\dnCo}{\dnCz}{\dnCo}{\dnCz}{\dnCo}{\dnCo}{\dnCz}} & \dnW[W]{}: & {{\dnCz}{\dnCz}{\dnCo}{\dnCz}{\dnCz}{\dnCo}{\dnCo}{\dnCo}{\dnCo}{\dnCz}{\dnCz}{\dnCo}{\dnCz}{\dnCo}{\dnCz}{\dnCz}{\dnCo}{\dnCz}{\dnCo}{\dnCo}{\dnCo}{\dnCo}{\dnCo}{\dnCo}{\dnCz}{\dnCz}{\dnCz}{\dnCz}{\dnCo}{\dnCz}{\dnCz}{\dnCz}} \\
 \dnI{15} & \dnW: & {{\dnCz}{\dnCz}{\dnCz}{\dnCo}{\dnCo}{\dnCo}{\dnCo}{\dnCz}{\dnCo}{\dnCz}{\dnCo}{\dnCz}{\dnCo}{\dnCu}{\dnCz}{\dnCo}{\dnCz}{\dnCo}{\dnCo}{\dnCz}{\dnCz}{\dnCo}{\dnCo}{\dnCz}{\dnCo}{\dnCo}{\dnCz}{\dnCo}{\dnCz}{\dnCo}{\dnCz}{\dnCz}} & \dnW[W]{}: & {{\dnCo}{\dnCz}{\dnCo}{\dnCo}{\dnCo}{\dnCz}{\dnCz}{\dnCo}{\dnCo}{\dnCo}{\dnCo}{\dnCz}{\dnCo}{\dnCz}{\dnCz}{\dnCz}{\dnCo}{\dnCo}{\dnCz}{\dnCz}{\dnCz}{\dnCz}{\dnCo}{\dnCo}{\dnCo}{\dnCo}{\dnCo}{\dnCz}{\dnCo}{\dnCz}{\dnCz}{\dnCo}} \\
 \dnI{16} & \dnW: & {{\dnCn}{\dnCz}{\dnCz}{\dnCn}{\dnCz}{\dnCu}{\dnCn}{\dnCz}{\dnCo}{\dnCo}{\dnCz}{\dnCo}{\dnCz}{\dnCz}{\dnCo}{\dnCz}{\dnCo}{\dnCz}{\dnCz}{\dnCo}{\dnCo}{\dnCz}{\dnCo}{\dnCo}{\dnCz}{\dnCo}{\dnCz}{\dnCo}{\dnCo}{\dnCo}{\dnCo}{\dnCo}} & & \\
 \dnI{17} & \dnW: & {{\dnCz}{\dnCz}{\dnCz}{\dnCo}{\dnCo}{\dnCo}{\dnCo}{\dnCo}{\dnCz}{\dnCz}{\dnCo}{\dnCo}{\dnCo}{\dnCz}{\dnCo}{\dnCz}{\dnCz}{\dnCz}{\dnCz}{\dnCo}{\dnCz}{\dnCz}{\dnCo}{\dnCz}{\dnCz}{\dnCz}{\dnCz}{\dnCo}{\dnCo}{\dnCo}{\dnCo}{\dnCz}} & & \\
 \dnI{18} & \dnW: & {{\dnCz}{\dnCo}{\dnCz}{\dnCo}{\dnCz}{\dnCo}{\dnCo}{\dnCo}{\dnCz}{\dnCz}{\dnCz}{\dnCz}{\dnCo}{\dnCo}{\dnCz}{\dnCo}{\dnCz}{\dnCz}{\dnCz}{\dnCz}{\dnCz}{\dnCz}{\dnCz}{\dnCz}{\dnCo}{\dnCz}{\dnCz}{\dnCo}{\dnCz}{\dnCo}{\dnCz}{\dnCz}} & & \\
 \dnI{19} & \dnW: & {{\dnCu}{\dnCo}{\dnCn}{\dnCo}{\dnCz}{\dnCz}{\dnCz}{\dnCz}{\dnCz}{\dnCz}{\dnCz}{\dnCo}{\dnCz}{\dnCo}{\dnCo}{\dnCo}{\dnCo}{\dnCz}{\dnCz}{\dnCo}{\dnCo}{\dnCz}{\dnCo}{\dnCz}{\dnCo}{\dnCo}{\dnCz}{\dnCz}{\dnCz}{\dnCo}{\dnCz}{\dnCz}} & & \\
 \dnI{20} & \dnW: & {{\dnCn}{\dnCo}{\dnCu}{\dnCn}{\dnCo}{\dnCz}{\dnCz}{\dnCo}{\dnCo}{\dnCo}{\dnCo}{\dnCo}{\dnCo}{\dnCo}{\dnCz}{\dnCo}{\dnCo}{\dnCo}{\dnCz}{\dnCo}{\dnCz}{\dnCz}{\dnCz}{\dnCz}{\dnCz}{\dnCz}{\dnCo}{\dnCo}{\dnCz}{\dnCo}{\dnCz}{\dnCz}} & & \\
 \dnI{21} & \dnW: & {{\dnCo}{\dnCo}{\dnCo}{\dnCo}{\dnCz}{\dnCz}{\dnCo}{\dnCo}{\dnCo}{\dnCz}{\dnCo}{\dnCo}{\dnCz}{\dnCz}{\dnCz}{\dnCz}{\dnCz}{\dnCo}{\dnCz}{\dnCo}{\dnCo}{\dnCo}{\dnCo}{\dnCo}{\dnCo}{\dnCo}{\dnCz}{\dnCo}{\dnCz}{\dnCo}{\dnCz}{\dnCz}} & & \\
 \dnI{22} & \dnW: & {{\dnCz}{\dnCo}{\dnCz}{\dnCo}{\dnCo}{\dnCo}{\dnCz}{\dnCo}{\dnCo}{\dnCo}{\dnCz}{\dnCz}{\dnCo}{\dnCo}{\dnCz}{\dnCo}{\dnCz}{\dnCz}{\dnCo}{\dnCo}{\dnCz}{\dnCz}{\dnCo}{\dnCo}{\dnCz}{\dnCz}{\dnCo}{\dnCo}{\dnCo}{\dnCz}{\dnCo}{\dnCz}} & & \\
 \dnI{23} & \dnW: & {{\dnCz}{\dnCo}{\dnCz}{\dnCo}{\dnCz}{\dnCz}{\dnCz}{\dnCz}{\dnCo}{\dnCo}{\dnCo}{\dnCz}{\dnCo}{\dnCo}{\dnCo}{\dnCz}{\dnCo}{\dnCo}{\dnCz}{\dnCz}{\dnCz}{\dnCo}{\dnCo}{\dnCo}{\dnCo}{\dnCz}{\dnCz}{\dnCz}{\dnCo}{\dnCo}{\dnCo}{\dnCo}} & & \\
 \dnI{24} & \dnW: & {{\dnCz}{\dnCz}{\dnCz}{\dnCz}{\dnCz}{\dnCz}{\dnCo}{\dnCz}{\dnCz}{\dnCz}{\dnCz}{\dnCo}{\dnCz}{\dnCz}{\dnCo}{\dnCz}{\dnCz}{\dnCz}{\dnCo}{\dnCo}{\dnCz}{\dnCo}{\dnCo}{\dnCo}{\dnCz}{\dnCz}{\dnCz}{\dnCo}{\dnCo}{\dnCz}{\dnCo}{\dnCz}} & & \\
 \dnI{25} & \dnW: & {{\dnCo}{\dnCz}{\dnCo}{\dnCo}{\dnCz}{\dnCz}{\dnCz}{\dnCz}{\dnCo}{\dnCz}{\dnCz}{\dnCo}{\dnCz}{\dnCo}{\dnCo}{\dnCz}{\dnCz}{\dnCz}{\dnCz}{\dnCo}{\dnCz}{\dnCo}{\dnCz}{\dnCz}{\dnCo}{\dnCo}{\dnCo}{\dnCz}{\dnCo}{\dnCz}{\dnCo}{\dnCz}} & & \\
 \dnI{26} & \dnW: & {{\dnCz}{\dnCz}{\dnCz}{\dnCz}{\dnCo}{\dnCz}{\dnCo}{\dnCz}{\dnCo}{\dnCz}{\dnCz}{\dnCz}{\dnCo}{\dnCz}{\dnCz}{\dnCo}{\dnCz}{\dnCo}{\dnCo}{\dnCo}{\dnCz}{\dnCo}{\dnCo}{\dnCo}{\dnCz}{\dnCo}{\dnCz}{\dnCz}{\dnCz}{\dnCz}{\dnCz}{\dnCo}} & & \\
 \dnI{27} & \dnW: & {{\dnCz}{\dnCz}{\dnCz}{\dnCz}{\dnCz}{\dnCo}{\dnCo}{\dnCz}{\dnCo}{\dnCo}{\dnCo}{\dnCz}{\dnCo}{\dnCo}{\dnCo}{\dnCo}{\dnCz}{\dnCo}{\dnCz}{\dnCo}{\dnCo}{\dnCz}{\dnCo}{\dnCz}{\dnCo}{\dnCz}{\dnCo}{\dnCo}{\dnCz}{\dnCz}{\dnCo}{\dnCo}} & & \\
 \dnI{28} & \dnW: & {{\dnCo}{\dnCz}{\dnCo}{\dnCo}{\dnCz}{\dnCo}{\dnCo}{\dnCz}{\dnCz}{\dnCo}{\dnCz}{\dnCo}{\dnCo}{\dnCo}{\dnCz}{\dnCo}{\dnCz}{\dnCo}{\dnCo}{\dnCz}{\dnCo}{\dnCo}{\dnCz}{\dnCz}{\dnCz}{\dnCz}{\dnCo}{\dnCz}{\dnCz}{\dnCo}{\dnCz}{\dnCo}} & & \\
 \dnI{29} & \dnW: & {{\dnCo}{\dnCz}{\dnCo}{\dnCz}{\dnCz}{\dnCz}{\dnCo}{\dnCz}{\dnCz}{\dnCz}{\dnCz}{\dnCz}{\dnCo}{\dnCo}{\dnCz}{\dnCo}{\dnCz}{\dnCo}{\dnCz}{\dnCz}{\dnCo}{\dnCz}{\dnCz}{\dnCz}{\dnCz}{\dnCo}{\dnCo}{\dnCz}{\dnCo}{\dnCz}{\dnCz}{\dnCo}} & & \\
 \dnI{30} & \dnW: & {{\dnCz}{\dnCz}{\dnCo}{\dnCz}{\dnCo}{\dnCz}{\dnCz}{\dnCo}{\dnCo}{\dnCo}{\dnCz}{\dnCo}{\dnCz}{\dnCo}{\dnCo}{\dnCo}{\dnCo}{\dnCo}{\dnCz}{\dnCz}{\dnCz}{\dnCo}{\dnCo}{\dnCo}{\dnCz}{\dnCo}{\dnCo}{\dnCz}{\dnCz}{\dnCz}{\dnCo}{\dnCo}} & & \\
 \dnI{31} & \dnW: & {{\dnCo}{\dnCo}{\dnCo}{\dnCo}{\dnCo}{\dnCo}{\dnCz}{\dnCz}{\dnCo}{\dnCz}{\dnCz}{\dnCo}{\dnCz}{\dnCz}{\dnCo}{\dnCz}{\dnCo}{\dnCo}{\dnCz}{\dnCo}{\dnCz}{\dnCo}{\dnCo}{\dnCo}{\dnCo}{\dnCz}{\dnCo}{\dnCo}{\dnCz}{\dnCo}{\dnCo}{\dnCz}} & & \\
 \dnI{32} & \dnW: & {{\dnCz}{\dnCo}{\dnCz}{\dnCz}{\dnCo}{\dnCo}{\dnCo}{\dnCo}{\dnCo}{\dnCo}{\dnCz}{\dnCo}{\dnCz}{\dnCz}{\dnCo}{\dnCz}{\dnCz}{\dnCo}{\dnCo}{\dnCz}{\dnCo}{\dnCz}{\dnCz}{\dnCz}{\dnCz}{\dnCz}{\dnCo}{\dnCz}{\dnCo}{\dnCo}{\dnCo}{\dnCo}} & & \\
 \dnI{33} & \dnW: & {{\dnCz}{\dnCz}{\dnCo}{\dnCo}{\dnCo}{\dnCz}{\dnCz}{\dnCz}{\dnCz}{\dnCz}{\dnCo}{\dnCo}{\dnCo}{\dnCo}{\dnCz}{\dnCo}{\dnCz}{\dnCo}{\dnCo}{\dnCz}{\dnCo}{\dnCo}{\dnCo}{\dnCz}{\dnCo}{\dnCo}{\dnCo}{\dnCz}{\dnCz}{\dnCo}{\dnCz}{\dnCz}} & & \\
 \dnI{34} & \dnW: & {{\dnCz}{\dnCz}{\dnCo}{\dnCz}{\dnCz}{\dnCz}{\dnCz}{\dnCz}{\dnCz}{\dnCo}{\dnCo}{\dnCo}{\dnCz}{\dnCo}{\dnCz}{\dnCo}{\dnCo}{\dnCo}{\dnCo}{\dnCz}{\dnCo}{\dnCz}{\dnCz}{\dnCz}{\dnCz}{\dnCz}{\dnCz}{\dnCo}{\dnCz}{\dnCo}{\dnCz}{\dnCo}} & & \\
 \dnI{35} & \dnW: & {{\dnCn}{\dnCz}{\dnCo}{\dnCz}{\dnCz}{\dnCz}{\dnCz}{\dnCz}{\dnCz}{\dnCz}{\dnCo}{\dnCo}{\dnCz}{\dnCz}{\dnCo}{\dnCo}{\dnCz}{\dnCz}{\dnCz}{\dnCz}{\dnCz}{\dnCo}{\dnCz}{\dnCz}{\dnCz}{\dnCo}{\dnCo}{\dnCo}{\dnCz}{\dnCz}{\dnCo}{\dnCz}} & & \\
 \dnI{36} & \dnW: & {{\dnCn}{\dnCz}{\dnCz}{\dnCz}{\dnCz}{\dnCo}{\dnCo}{\dnCo}{\dnCo}{\dnCo}{\dnCo}{\dnCz}{\dnCo}{\dnCz}{\dnCo}{\dnCo}{\dnCo}{\dnCo}{\dnCz}{\dnCo}{\dnCo}{\dnCo}{\dnCo}{\dnCz}{\dnCz}{\dnCo}{\dnCz}{\dnCo}{\dnCo}{\dnCz}{\dnCz}{\dnCo}} & & \\
 \dnI{37} & \dnW: & {{\dnCo}{\dnCo}{\dnCz}{\dnCz}{\dnCo}{\dnCz}{\dnCz}{\dnCz}{\dnCz}{\dnCz}{\dnCz}{\dnCo}{\dnCo}{\dnCz}{\dnCo}{\dnCz}{\dnCz}{\dnCo}{\dnCz}{\dnCz}{\dnCz}{\dnCz}{\dnCo}{\dnCo}{\dnCz}{\dnCz}{\dnCz}{\dnCz}{\dnCo}{\dnCo}{\dnCz}{\dnCz}} & & \\
 \dnI{38} & \dnW: & {{\dnCo}{\dnCz}{\dnCo}{\dnCo}{\dnCz}{\dnCz}{\dnCz}{\dnCz}{\dnCz}{\dnCo}{\dnCo}{\dnCz}{\dnCz}{\dnCo}{\dnCo}{\dnCo}{\dnCo}{\dnCo}{\dnCo}{\dnCz}{\dnCo}{\dnCz}{\dnCz}{\dnCo}{\dnCo}{\dnCz}{\dnCo}{\dnCz}{\dnCo}{\dnCo}{\dnCz}{\dnCz}} & & \\
 \dnI{39} & \dnW: & {{\dnCz}{\dnCz}{\dnCz}{\dnCo}{\dnCz}{\dnCz}{\dnCo}{\dnCz}{\dnCz}{\dnCz}{\dnCz}{\dnCz}{\dnCo}{\dnCz}{\dnCo}{\dnCz}{\dnCz}{\dnCz}{\dnCz}{\dnCo}{\dnCo}{\dnCz}{\dnCo}{\dnCo}{\dnCz}{\dnCz}{\dnCz}{\dnCo}{\dnCo}{\dnCo}{\dnCz}{\dnCz}} & & \\
 \dnI{40} & \dnW: & {{\dnCo}{\dnCo}{\dnCz}{\dnCz}{\dnCz}{\dnCz}{\dnCz}{\dnCz}{\dnCz}{\dnCo}{\dnCz}{\dnCz}{\dnCo}{\dnCz}{\dnCz}{\dnCz}{\dnCz}{\dnCo}{\dnCo}{\dnCo}{\dnCz}{\dnCz}{\dnCz}{\dnCo}{\dnCo}{\dnCz}{\dnCz}{\dnCz}{\dnCz}{\dnCo}{\dnCz}{\dnCo}} & & \\
 \dnI{41} & \dnW: & {{\dnCz}{\dnCz}{\dnCz}{\dnCz}{\dnCz}{\dnCo}{\dnCo}{\dnCz}{\dnCo}{\dnCz}{\dnCz}{\dnCz}{\dnCz}{\dnCo}{\dnCo}{\dnCz}{\dnCo}{\dnCo}{\dnCo}{\dnCo}{\dnCz}{\dnCo}{\dnCz}{\dnCo}{\dnCz}{\dnCz}{\dnCo}{\dnCz}{\dnCz}{\dnCo}{\dnCo}{\dnCz}} & & \\
 \dnI{42} & \dnW: & {{\dnCz}{\dnCo}{\dnCz}{\dnCz}{\dnCo}{\dnCo}{\dnCo}{\dnCz}{\dnCo}{\dnCo}{\dnCz}{\dnCo}{\dnCo}{\dnCo}{\dnCz}{\dnCo}{\dnCo}{\dnCo}{\dnCo}{\dnCo}{\dnCo}{\dnCo}{\dnCo}{\dnCz}{\dnCo}{\dnCz}{\dnCz}{\dnCz}{\dnCz}{\dnCo}{\dnCo}{\dnCz}} & & \\
 \dnI{43} & \dnW: & {{\dnCz}{\dnCo}{\dnCz}{\dnCo}{\dnCz}{\dnCz}{\dnCz}{\dnCz}{\dnCz}{\dnCo}{\dnCo}{\dnCz}{\dnCz}{\dnCz}{\dnCo}{\dnCo}{\dnCo}{\dnCo}{\dnCz}{\dnCo}{\dnCz}{\dnCz}{\dnCz}{\dnCz}{\dnCz}{\dnCo}{\dnCo}{\dnCz}{\dnCo}{\dnCo}{\dnCz}{\dnCo}} & & \\
 \dnI{44} & \dnW: & {{\dnCo}{\dnCo}{\dnCo}{\dnCo}{\dnCo}{\dnCz}{\dnCz}{\dnCz}{\dnCz}{\dnCz}{\dnCz}{\dnCo}{\dnCz}{\dnCo}{\dnCo}{\dnCz}{\dnCo}{\dnCo}{\dnCo}{\dnCo}{\dnCz}{\dnCo}{\dnCo}{\dnCo}{\dnCz}{\dnCz}{\dnCz}{\dnCz}{\dnCo}{\dnCo}{\dnCz}{\dnCz}} & & \\
 \dnI{45} & \dnW: & {{\dnCo}{\dnCz}{\dnCz}{\dnCz}{\dnCo}{\dnCz}{\dnCo}{\dnCz}{\dnCo}{\dnCo}{\dnCz}{\dnCo}{\dnCo}{\dnCz}{\dnCo}{\dnCo}{\dnCz}{\dnCz}{\dnCo}{\dnCz}{\dnCo}{\dnCo}{\dnCz}{\dnCz}{\dnCz}{\dnCz}{\dnCz}{\dnCz}{\dnCz}{\dnCo}{\dnCz}{\dnCz}} & & \\
 \dnI{46} & \dnW: & {{\dnCo}{\dnCz}{\dnCz}{\dnCz}{\dnCz}{\dnCz}{\dnCo}{\dnCz}{\dnCo}{\dnCz}{\dnCz}{\dnCo}{\dnCo}{\dnCz}{\dnCz}{\dnCo}{\dnCz}{\dnCo}{\dnCz}{\dnCo}{\dnCo}{\dnCz}{\dnCz}{\dnCz}{\dnCo}{\dnCo}{\dnCz}{\dnCo}{\dnCo}{\dnCo}{\dnCz}{\dnCz}} & & \\
 \dnI{47} & \dnW: & {{\dnCo}{\dnCz}{\dnCz}{\dnCz}{\dnCz}{\dnCz}{\dnCz}{\dnCo}{\dnCo}{\dnCo}{\dnCo}{\dnCz}{\dnCz}{\dnCo}{\dnCz}{\dnCo}{\dnCo}{\dnCz}{\dnCo}{\dnCo}{\dnCz}{\dnCo}{\dnCz}{\dnCz}{\dnCo}{\dnCz}{\dnCo}{\dnCo}{\dnCo}{\dnCo}{\dnCz}{\dnCo}} & & \\
\hline
\end{tabular}
}
\caption[%
  One of the original MD4 collisions by Wang et al,
  with a few bits underspecified and followingly propagated%
]{%
  One of the original MD4 collisions by Wang et al,
  with a few bits underspecified (left) and propagated values (right).
  The question marks indicate that any bit configuration for the two bits
  are possible. Dashes indicate that the bits have the same configuration
  in both instances, but the value itself is unknown.
  However, it turns out the description with missing values in iteration~6
  (message) and iterations 8--19 is complete enough such that missing values
  can be deduced by other values the description of the algorithm.
  Therefore the collision They propagates through the intermediate
  values until the differences cancel out after round~36.
}
\label{tab:wang-collision-propagated}
\end{center}
\end{table}
}

\renewcommand*\chappic{img/satisfiability.pdf}
\renewcommand*\chapquote{What idiot called them logic errors rather than bool shit?}
\renewcommand*\chapquotesrc{Unknown}
\chapter{Satisfiability}
\label{ch:sat}
%
Boolean algebra allows us to describe functions over two-valued variables.
Satisfiability is the question for an assignment such that a function
evaluates to true. Satisfiability problems are solved by SAT~solvers.
We discuss the basic theory behind satisfiability. We will learn that any
computation can be represented as satisfiability problem. In Chapter~\ref{ch:enc}
we will represent a differential cryptanalysis problem such that it is
solvable iff the corresponding SAT~problem is satisfiable.

\section{Basic notation and definitions}
\label{sec:sat-intro}
%
\index{Boolean function}
\begin{defi}[Boolean function]
  A \emph{Boolean function} is a mapping $h: X \to Y$ with $X = \left\{0,1\right\}^n$
  for $n \in \mathbb N_{\geq 1}$ and $Y = \left\{0,1\right\}$.
\end{defi}

\index{Assignment}
\begin{defi}[Assignment]
  A \emph{$k$-assignment} is an element of $\left\{0,1\right\}^k$.

  \noindent
  Let $f$ be some $k$-ary Boolean function.
  An \emph{assignment for function $f$} is any $k$-assignment.
\end{defi}

\index{Truth table}
\begin{defi}[Truth table]
  Let $f$ be some $k$-ary Boolean function.
  The \emph{truth table of Boolean function $f$} assigns
  truth value $0$ or $1$ to any assignment of $f$.
\end{defi}

Boolean functions are characterized by their corresponding truth table.

\begin{table}[pt]
  \centering
  \subfloat[\boolf{AND}]{%
    \begin{tabular}{cc|c}
      $x_1$ & $x_2$ & $f(x_1, x_2)$ \\
     \hline
      $1$ & $1$ & $1$ \\
      $1$ & $0$ & $0$ \\
      $0$ & $1$ & $0$ \\
      $0$ & $0$ & $0$
    \end{tabular}
  }
  ~
  \subfloat[\boolf{OR}]{%
    \begin{tabular}{cc|c}
      $x_1$ & $x_2$ & $f(x_1, x_2)$ \\
     \hline
      $1$ & $1$ & $1$ \\
      $1$ & $0$ & $1$ \\
      $0$ & $1$ & $1$ \\
      $0$ & $0$ & $0$
    \end{tabular}
  }
  ~
  \raisebox{13.6pt}{%
    \subfloat[\boolf{NOT}]{%
      \begin{tabular}{c|c}
        $v$ & $f(v)$ \\
       \hline
        $1$ & $0$ \\
        $0$ & $1$
      \end{tabular}
    }%
  }%
  \caption{Truth tables for \boolf{AND}, \boolf{OR} and \boolf{NOT}}
  \label{tab:andornot-truthtables}
\end{table}

Table~\ref{tab:andornot-truthtables} shows example truth tables for
the Boolean \boolf{AND}, \boolf{OR} and \boolf{NOT} functions.
A different definition of the three functions is given the following way:

\index{AND (Boolean function)}
\index{OR (Boolean function)}
\index{NOT (Boolean function)}
\begin{defi}
  Let \boolf{AND}, \boolf{OR} and \boolf{NOT} be three Boolean functions.
  \begin{itemize}[noitemsep,topsep=0pt]
    \item
      \boolf{AND} maps $X = \left\{0,1\right\}^2$
      to $1$ if all values of $X$ are $1$.
    \item
      \boolf{OR} maps $X = \left\{0,1\right\}^2$
      to $1$ if any value of $X$ is $1$.
    \item
      \boolf{NOT} maps $X = \left\{0,1\right\}^1$
      to $1$ if the single value of $X$ is $0$.
  \end{itemize}
  All functions return $0$ in the other case.

  Those functions are denoted $a_0 \land a_1$, $a_0 \lor a_1$
  and $\neg a_0$ respectively, for input parameters $a_0$ and $a_1$.
\end{defi}

It is interesting to observe, that any Boolean function can be represented
using only these three operators. This can be proven by complete induction
over the number of arguments $k$ of the function.

Let $k = 1$. Then we consider any possible $2$-assignment for one input
variable $x_1$ and one value of $f(x_1)$. Then four truth tables are possible
listed in Table~\ref{tab:unary-f}. The description shows the corresponding
definition of $f$ using \boolf{AND}, \boolf{OR} and \boolf{NOT} only.

Now let $g$ be some $k$-ary function. Let $(a_0, a_1, \ldots, a_k)$ be the
$k$ input arguments to $g$ and $x_1 \coloneqq g(a_0, a_1, \ldots, a_k)$.
Then we can again look at Table~\ref{tab:unary-f} to discover that 4 cases
are possible: 2 cases where the return value of our new $(k+1)$-ary function
depends on value $x_1$ and 2 cases where the return value is constant.

This completes our proof.

\begin{table}[ht]
  \centering
  \subfloat[$f: x \mapsto 1$]{%
    \begin{tabular}{cc}
      $x_1$ & $f(x_1)$ \\
     \hline
      $1$ & $1$ \\
      $0$ & $1$
    \end{tabular}
  }
  ~
  \subfloat[$f: x \mapsto x$]{%
    \begin{tabular}{cc}
      $x_1$ & $f(x_1)$ \\
     \hline
      $1$ & $1$ \\
      $0$ & $0$
    \end{tabular}
  }
  ~
  \subfloat[$f: x \mapsto \neg x$]{%
    \begin{tabular}{cc}
      $x_1$ & $f(x_1)$ \\
     \hline
      $1$ & $0$ \\
      $0$ & $1$
    \end{tabular}
  }
  ~
  \subfloat[$f: x \mapsto 0$]{%
    \begin{tabular}{cc}
      $x_1$ & $f(x_1)$ \\
     \hline
      $1$ & $0$ \\
      $0$ & $0$
    \end{tabular}
  }%
  \caption{Unary $f$ and its four possible cases}
  \label{tab:unary-f}
\end{table}

Boolean functions have an important property which is described
in the following definition:

\index{Satisfiability}
\index{Assignment}
\index{Model}
\begin{defi}
  A Boolean function $f$ is \emph{satisfiable} iff there exists at least one
  input $x \in X$ such that $f(x) = 1$.
  Every input $x \in X$ satisfying this property is called \emph{model}.
\end{defi}

The corresponding tool to determine satisfiability is defined as follows:

\index{SAT solver}
\begin{defi}
  A \emph{SAT solver} is a tool to determine satisfiability (SAT or UNSAT)
  of a Boolean function. If satisfiability is given, it returns some model.
\end{defi}

\subsection{Computational considerations}
\label{sec:sat-complexity}
%
The generic complexity of SAT determination is given by $2^n$ for $n$ Boolean variables.

Let $n$ be the number of variables of a Boolean function.
No known algorithm exists to determine satisfiability in polynomial runtime.
This means no algorithm solves the SAT problem with runtime behavior
which depends polynomially on the growth of $n$.

This is known as the famous {\cP $\overset{?}{\neq}$ \cNP} problem.

\index{Unit propagation}
However, SAT solver can take advantage of the problem's description.
For example consider function $f$ in Display~\ref{eq:3f}.
\begin{align} f(x_0, x_1, x_2) &= x_0 \land (\neg x_1 \lor x_2) \label{eq:3f} \end{align}
Instead of trying all possible 8~cases for 3~Boolean variables,
we can immediately see that $x_0$ is required to be $1$.
So we don't need to test $x_0 = 0$ and can skip 4~cases.
This particular strategy is called \emph{unit propagation}.

\subsection{SAT competitions}
\label{sec:sat-competitions}
%
SAT research is heavily concerned with finding good heuristics to find some model
for a given SAT problem as fast as possible. Biyearly
\href{http://satcompetition.org/}{SAT competitions} take place to challenge
SAT solvers in a set of benchmarks. The committee evaluates the most successful
SAT solvers solving the most problems within a given time frame.

SAT~2016 is currently ongoing, but in 2014 lingeling by Armin Biere has won first prize in
the Application benchmarks track and second prize in the Hard Combinatorial benchmarks
track for SAT and UNSAT instances respectively. Its parallelized sibling plingeling
and Cube \& Conquer sibling treengeling have won prizes in parallel settings.

In chapter~\ref{ch:results} we will look at runtime results shown by (but not limited to)
those SAT solvers.

%\section{Satisfiability of hash algorithm states}
%\label{sec:intro-algo-sat}
%%
%We discussed Boolean functions and satisfiability. At the same time we looked
%at basic properties of hash algorithms. But the question remains how we can
%link those areas together? This section is dedicated to this question.

%\index{Algorithm}
%\index{I/O Algorithm}
%\begin{defi}
%  An \emph{algorithm} is a step-wise set of instructions to solve a problem.
%  An \emph{I/O algorithm} transforms given input values to output values.
%\end{defi}

%Hash algorithms are one example of I/O algorithms.

%I/O algorithms can be implemented as a sequence of instructions for computers.
%At the same time I/O algorithms can be represented as combination of Boolean
%functions. This claim is backed in more detail in Section~\ref{sec:sat-dimacs}
%with Theorem~\ref{thm:all-cnf}. It follows immediately that we can represent
%I/O algorithms such as hash algorithms entirely as Boolean function.

%\begin{theorem}
%  Every algorithm can be represented as Boolean function.
%\end{theorem}

%\index{Least significant bit}
%\index{Most significant bit}
%We consider 2-bit addition as small example. Let $a_{i}$ be the first argument
%where $i$ denotes the binary position. If $i=0$, the \emph{least significant bit}
%(LSB) is considered. If $i=1$, the \emph{most significant bit} (MSB) is considered.

%Let $b_{i}$ be the second argument and $s_{i}$ be the output value.
%Furthermore $c_{i}$ is the carry bit, where $c_1$ is left out, because
%it is not used in 2-bit addition. This model of 2-bit addition as Boolean
%function can be seen in Figure~\ref{fig:intro-2bit-addition}.

%\begin{figure}
%  \begin{center}
%    \begin{tabular}{lrcc}
%      1st arg: &   & $a_{1}$ & $a_{0}$ \\
%      2nd arg: & + & $b_{1}$ & $b_{0}$ \\
%    \hline
%      carry:   &   &         & $c_0$ \\
%      sum:     & = & $s_{1}$ & $s_0$
%    \end{tabular}
%    \hspace{10pt}~$\leadsto$\hspace{10pt}%
%    \begin{minipage}{50pt}%
%      \begin{flalign*}
%        s_0 &= \boolf{XOR}(a_0, b_0) &\\
%        c_0 &= a_0 \land b_0 &\\
%        s_1 &= \boolf{XOR}(a_0, b_0, c_0) &
%      \end{flalign*}
%    \end{minipage}
%  \end{center}
%  \caption{Modelling 2bit addition (left) as Boolean function (right)}
%  \label{fig:intro-2bit-addition}
%\end{figure}


\section{The DIMACS de-facto standard}
\label{sec:sat-dimacs}
%
\index{Conjunction}
\index{Disjunction}
\index{Literal}
\index{Positive literal}
\index{Negative literal}
\index{Conjunctive Normal Form}
\begin{defi}
  A \emph{conjunction} is a sequence of Boolean functions combined using
  a logical OR. A \emph{disjunction} is a sequence of Boolean functions
  combined using a logical AND. A \emph{literal} is a Boolean variable
  (\emph{positive}) or its negation (\emph{negative}).

  A SAT problem is given in \emph{Conjunctive Normal Form} (CNF) if
  the problem is defined as conjunction of disjunctions of literals.
\end{defi}

A simple example for a SAT problem in CNF is the exclusive OR (XOR).
It takes two Boolean values $a$ and $b$ as arguments and returns true
if and only if the two arguments differ.
{
\setlength{\abovedisplayskip}{5pt}
\setlength{\belowdisplayskip}{5pt}
\setlength{\abovedisplayshortskip}{0pt}
\setlength{\belowdisplayshortskip}{0pt}
\begin{align} (a \lor b) \land (\neg a \lor \neg b) \label{eq:xor}\end{align}
}
Display~\ref{eq:xor} shows one conjunction (denoted $\land$) of two disjunctions
(denoted $\lor$) of literals (denoted $a$ and $b$ where prefix $\neg$ represents
negation). This structure constitutes a CNF.

\index{Disjunctive Normal Form}
Analogously we define a \emph{Disjunctive Normal Form} (DNF) as disjunction
of conjunctions of literals. The negation of a CNF is in DNF, because literals
are negated and conjunctions become disjunctions, vice versa.

\begin{theorem}
  \label{thm:all-cnf}
  Every Boolean function can be represented as CNF.
\end{theorem}

Theorem~\ref{thm:all-cnf} is easy to prove.
% Consider null-ary functions as
% induction base. A (null-ary) Boolean function is a set of associations.
% An association maps an assignment for input arguments to one output value.
% True can be represented as $a \lor \neg a$ and false can be
% represented as $a \land \neg a$ for some free Boolean variable $a$ and using
% only negation, conjunctions and disjunctions. Recognize that those representations
% are given in DNF as well as CNF. Hence all null-ary Boolean functions can be
% represented as CNF and DNF.
%
% Let $f$ be a Boolean function represented in DNF with $k$ input arguments.
% Extend every conjunction with a Boolean variable $v_{k+1}$.
Consider the truth table of an arbitrary Boolean function $f$ with $k$ input arguments
and $j$ rows of output value false. We represent $f$ as CNF.

Consider Boolean variables $b_{i,l}$ with $0 \leq i \leq j$ and $0 \leq l \leq k$.
For every row $i$ of the truth table with assignment $(r_i)$, add one disjunction to the CNF.
This disjunction contains $b_{i,l}$ if $r_{i,l}$ is false.
The disjunction contains $b_{i,l}$ if $r_{i,l}$ is true.

As far as $f$ is an arbitrary $k$-ary Boolean function, we have proven that
any function can be represented as CNF.

SAT problems are usually represented in the DIMACS de-facto standard.
Consider a SAT problem in CNF with \emph{nbclauses} clauses and
enumerate all variables from 1 to \emph{nbvars}. A DIMACS file is an ASCII text
file. Lines starting with \enquote{\texttt{c}} are skipped (comment lines).
The first remaining line has to begin with \enquote{\texttt{p cnf}} followed by
\emph{nbclauses} and \emph{nbvars} separated by spaces (header line).
All following non-comment lines are space-separated indices of Boolean variables
optionally prefixed by a minus symbol. Then one line represents one clause and
must be terminated with a zero symbol after a space. All lines are conjuncted
to form a CNF.

Variations of the DIMACS de-facto standard also allow multiline clauses (the
zero symbol constitutes the end of a clause) or arbitrary whitespace instead of
spaces. The syntactical details are individually published on a per competition
basis.

\renewcommand{\lstlistingname}{Listing}  % TODO otherwise it's Japanese, WTF
\begin{lstlisting}[caption={Display~\ref{eq:xor} represented in DIMACS format}]
p cnf 2 2
a b
-a -b
\end{lstlisting}

\section{Terminology}
\label{sec:sat-terminology}
%
Given a conjunctive structure of disjunctions, we can define define terms
related to this structure:

\index{Clause}
\index{$k$-clause}
\index{Unit clause}
\index{Horn clause}
\index{Definite clause}
\begin{defi}
  A \emph{clause} is a disjunction of literals.
  A $k$-\emph{clause} is a clause consisting of exactly $k$ literals.
  A \emph{unit clause} is a $1$-clause.
  A \emph{Horn clause} is a clause with at most one positive literal.
  A \emph{definite clause} is a clause with exactly one positive literal.
  A \emph{goal clause} is a clause with no positive literal.
\end{defi}

\index{Negated literal}
\index{Existential literal}
\index{Used variable}
\begin{defi}
  Given a literal, its \emph{negated literal} is the literal with its sign negated.
  A literal is \emph{positive}, if its sign is positive. A literal is \emph{negative} if its sign is negative.
  An \emph{existential literal} is a literal which occurs exactly once and
  its negation does not occur. A \emph{used variable} is a variable which
  occurs at least once in the CNF.

  The \emph{literal frequency} is the number of occurences of a literal in the CNF divided by the number of clauses declared. Equivalently \emph{variable frequency} defines the number of variable occurences divided by the number of clauses declared.
\end{defi}

\index{Clause length}
\index{Tautological clause}
\begin{defi}
  The \emph{clause length of a clause} is the number of literals contained.
  A clause is called \emph{tautological} if a literal and its negated literal occurs in it.
\end{defi}

\section{Basic SAT solving techniques}
\label{sec:sat-solving}
%
\index{Equisatisfiability}
\begin{defi}
  Given two CNFs $A$ and $B$, they are called \emph{equisatisfiable} if and only
  if $A$ is satisfiable if and only if $B$.
\end{defi}

\subsection{Boolean constraint propagation (BCP)}
\label{sec:sat-bcp}
%
\index{Boolean constraint propagation}
One of the most basic techniques to SAT solving is \emph{Boolean Constraint Propagation},
also called \emph{unit propagation}.
It is so fundamental that SATzilla, introduced in Section~\ref{sec:features-related},
applies it immediately before looking at SAT features.

Let $l$ be the literal of a unit clause in a CNF. Remove any clause containing
$l$ and replace any occurences of $-l$ from the CNF. It is easy to see, that
the resulting CNF is equisatisfiable, because due to the unit clause $l$ must
be true. So any clause containing $l$ is satisfied and $-l$ yields false,
where $A \lor \bot$ is equivalent to $A$ for any Boolean function $A$.

\subsection{Watched literals}
\label{sec:sat-wl}
%
\index{Watched Literals}
Watched Literals are another fundamental concept in SAT solving. It is very
expensive to check satisfiability of all clauses for every value of a literal.
Watched Literals is a neat technique to reduce the number of checks.

Consider the solver assigns a value for literal $l$. Instead of looking at
all clauses and testing whether the clause is falsified by $l$, only clauses
containing $l$ are checked if $l$ is one of the watched literals of the clause.
This empirical approach was established with the Chaff and zChaff SAT solvers
and has proven useful in various variants.

\subsection{Remark}
\label{sec:sat-remark}
%
The previous two techniques shall illustrate basic approaches, but actual SAT
solving research requires decades of development to tune individual SAT solvers.
Memory models and concurrency strategies lead to fundamentally different runtime
behaviours of SAT solvers.

As such an initial idea to initiate an individual SAT solver specifically designed for
solving problems in differential cryptanalysis was dropped, because development time
is expected too long for a master thesis to be fruitful. As such we focused on popular
and established SAT solvers of the SAT community.

\section{SAT solvers in use}
\label{sec:sat-solvers}
%
\index{lingeling}
In this thesis we consider the several SAT solvers.
They have been selected either by their popularity
or their good results at previous SAT competitions:
\begin{itemize}
  \item MiniSat 2.2.0
  \item treengeling, lingeling and plingeling, in versions:
    \begin{itemize}
      \item lingeling ats1
      \item lingeling ats1o1
      \item lingeling ats1o2
      \item lingeling ats1o4
      \item lingeling baz
    \end{itemize}
  \item CryptoMiniSat 4.5
  \item CryptoMiniSat 5
  \item glucose syrup
\end{itemize}

Specifically this means the hash collision attacks we looked have run with
these SAT solvers. The results are discussed in Section~\ref{ch:results}
and provided in Appendix~\ref{app:runtimes}.

lingeling ats1o1, ats1o2 and ats1o4 are non-public releases of lingeling.
They have been developed in private communication with Armin Biere.
Our main goal was to achieve a separation between two sets of variables.
First all variables of the first need to be assigned in the best possible
way. Afterwards the second set of variables is considered.
Specifically variables modelling the differences between the two hash
algorithm instances should constitute the first set.

ats1o1 implements that difference variables are guessed with false first and
usual heuristics apply for all other variables.
Our intermediate results with incomplete CNF files showed a high number of restarts.
Therefore ats1o2 disables backjumping and therefore skips decisions for important variables.
Finally ats1o4 is not expected to distinguish from ats1o2. It only provides further debugging information.

The SAT solvers have generally been run without any special options, except for
\begin{itemize}
  \item MiniSat was run with \texttt{\textendash{}\textendash{}pre=once} as it is generally recommended to run with the builtin preprocessor.
  \item Lingeling has been generally run with \texttt{\textendash{}\textendash{}phase=-1} to prefer false as initial assignment to literals.
    However, lingeling ats1o1 implements this with more forceful strategy.
\end{itemize}

Preprocessing is a difficult topic on its own. Sometimes preprocessing can provide a speedup,
before actually solving the problem, but mostly SAT solvers implement preprocessing strategies
themselves and run them repeatedly when solving the problem.

TODO: testcases with lingeling have been run 5 times and mean was taken

TODO: glucose-syrup is glucose in many parallel threads

\renewcommand*\chappic{img/satfeatures.pdf}
\renewcommand*\chapquote{To be usable effectively [\dots] these features must be related to instance hardness and relatively cheap to compute}
\renewcommand*\chapquotesrc{SATzilla}
\chapter{SAT features}
\label{ch:features}
%
At the very beginning I was very intrigued by the question
\enquote{What is an \enquote{average} SAT problem?}. Answers to this
question can help to optimize SAT solver memory layouts.
Specifically for this thesis I wanted to find out whether
our problems distinguish from \enquote{average} problems in any
way such that we can use this distinction for runtime optimization.

I came up with 8 questions related to basic properties of SAT problems
we will discuss in depth in this section. We will characterize an
average SAT problem at Section~\ref{sec:features-average}:
\begin{enumerate}
\item Given an arbitrary literal. What is the percentage it is positive?
\item What is the variables / clauses ratio?
\item How many literals occur only once either positive or negative?
\item What is the average and longest clause length among CNF benchmarks?
\item How many Horn clauses exist in a CNF?
\item Are there any tautological clauses?
\item Are there any CNF files with more than one connected variable component?
\item How many variables of a CNF are covered by unit clauses?
\end{enumerate}

We will now define the terms used in those questions.

\section{Definition}
\label{sec:features-definition}
%
\index{SAT feature}
\index{Feature value}
\begin{defi}[SAT feature]
  A \emph{SAT feature} is a statistical value (named \emph{feature value})
  retrievable from some given SAT problem.
\end{defi}

The most basic example of a SAT feature is the number of variables and clauses
of a given SAT problem. This SAT feature is stored in the CNF header of a SAT
problem encoded in the DIMACS format.

The general goal is to write a tool which evaluates several SAT features at the same
time and retrieve them for comparison with other problems. Therefore it should be
computationally easy to evaluate SAT features of a given SAT problem. A suggested
computational limit is given with polynomial complexity in terms of number of
variables and number of clauses for memory as well as runtime for evaluation algorithms.
Sticking to this convention implies that evaluation of satisfiability must not be
necessary to evaluate a SAT feature under the assumption that \cPneqNP. Hence the
number of valid models cannot be a SAT feature as far as satisfiability needs to
be determined. But no actual hard computational limit is defined.

\section{Related work}
\label{sec:features-related}
%
The most similar resource I found looking at SAT features was the
SATzilla project~\cite{satzilla2004,satzilla2008} in 2012. The authors systematically defined
138 SAT~features categorized in 12 groups. Some features are only evaluated conditionally.
The features themselves are not defined formally, but an implementation is provided bundled
with example data. The following list provides an excerpt of the features:

\begin{description}
\item[\satfeature{nvarsOrig}] number of variables defined in the CNF header
\item[\satfeature{nvars}] number of active variables
\item[\satfeature{reducedVars}] \satfeature{nvarsOrig} - \satfeature{nvars}, divided by \satfeature{nvars}
\item[\satfeature{vars-clauses-ratio}] nvars divided by number of active clauses
\item[\satfeature{POSNEG-RATIO-CLAUSE-mean}] mean of $2 \cdot \left\| 0.5 - {\text{pos}}/{\text{length}}\right\|$ where $\text{pos}$ is the number of positive literals and $\text{length}$ clause length of a specific clause
\item[\satfeature{POSNEG-RATIO-CLAUSE-entropy}] like \satfeature{POSNEG-RATIO-CLAUSE-mean} but entropy
\item[\satfeature{TRINARY+}] number of clauses with clause length 1, 2 or 3 divided by number of active clauses
\item[\satfeature{HORNY-VAR-min}] minimum number of times a variable occurs in a Horn clause
\item[\satfeature{cluster-coeff-mean}]
   let neighbors of a clause be all clauses containing any literal negated
   and let clauses $c_1$ and $c_2$ be conflicting if $c_1$ contains literal $l$ and $c_2$ contains $-l$,
   then return the mean of $2$ times the number of conflicting neighbors of a clause $c$
   divided by the number of unordered pairs of neighbors,
   returned iff computable within 20 seconds for all clauses
\end{description}

Please recognize that active clauses are the unsatisfied clauses after BCP has been applied.
Equivalently active variables are remaining variables after application of BCP.

Many SAT solvers collect feature values to improve algorithm selection,
restart strategies and estimate problem sizes. Recent trends to apply Machine
Learning to SAT solving imply feature evaluation. SAT features and the resulting
satisfiability runtime are used as training data for Machine Learning. One example
using SAT features for algorithm selection is ASlib~\cite{aslib}.

The SAT solvers of Section~\ref{sec:sat-solvers} we use also compute features they use
when computing a solution. For example CryptoMiniSat 4.5.3 prints the following lines

\begin{verbatim}
c [features] numVars 56118, numClauses 358991, var_cl_ratio 0.156,
binary 0.019, trinary 0.520, horn 0.387, horn_mean 0.000, horn_std
0.000, horn_min 0.000, horn_max 0.000, horn_spread 0.000,
vcg_var_mean 0.000, vcg_var_std 0.902, vcg_var_min 0.000,
vcg_var_max 0.000, vcg_var_spread -0.000, vcg_cls_mean 0.000,
...
\end{verbatim}

Even though we will partially use equivalent features (like Horn clauses),
many are actually related to the current state of evaluation like decisions
per conflicts. We consider this as a property of the evaluation and not the
SAT problem itself.

\section{Statistical features}
\label{sec:features-stats}
%
For our SAT~features we need to define some basic statistical terminology.
Let $x_1, x_2, \ldots, x_n$ be a finite sequence of numbers ($n \in \mathbb N$).
\begin{description}
  \item[Arithmetic mean] (or short: mean)
    is defined as
    \[ \overline{x} = \frac1n \sum_{i=1}^n x_i \]
  \item[Standard deviation] (or short: sd)
    with mean $\overline{x}$ is defined as
    \[ \sigma(x) = \sqrt{\frac1n \sum_{i=1}^n (x_i - \overline{x})} \]
  \item[Median]
    with $x_1 \leq x_2 \leq \ldots \leq x_n$
    (i.e. sorted ascendingly) is defined as
    \[
       m = \begin{cases}
         x_{\lceil \text{mid}\rceil} & \text{if } n \text{ odd} \\
         \frac{x_{\text{mid}} + x_{\text{mid} + 1}}{2} & \text{if } n \text{ even}
       \end{cases}
       \quad\text{ with } \text{mid} = \frac{n}{2}
    \]
    and often considered more \enquote{robust} than the arithmetic mean.
  \item[Entropy]
    is defined according to Claude Shannon's information theory:
    \[ H(x) = -\sum_{i=1}^n x_i \cdot \log_2(x_i) \]
    where $0 \cdot \log_2(0) \coloneqq 0$.
\end{description}

Furthermore \emph{count} refers to the number of elements $n$,
\emph{largest} refers to the maximum element $\max_{1 \leq i \leq n}(x_i)$
and \emph{smallest} refers to the minimum element $\min_{1 \leq i \leq n}(x_i)$.

\section{Suggested SAT features}
\label{sec:features-suggested}
%
\index{cnf-analysis}
We wrote a tool called cnf-analysis. The evaluated features are partially inspired
by SATzilla and lingeling. The latter prints basic statistics for every CNF it
evaluates.

A summary of our suggested SAT features is given:

\begin{description}
\item[\satfeature{clause\_variables\_sd\_mean}] \hfill{} \\
  mean of sd of variables in a clause
\item[\satfeature{clauses\_length\_(largest, smallest, mean, median, sd)}] \hfill{} \\
  statistics related to the clause length
\item[\satfeature{connected\_(literal, variable)\_components\_count}] \hfill{} \\
  two literals (variables) are connected if they occur in some clause together,
  count the number of connected components
\item[\satfeature{definite\_clauses\_count}] \hfill{} \\
  number of definite clauses in the CNF
\item[\satfeature{existential\_literals\_count}] \hfill{} \\
  number of existential literals in the CNF
\item[\satfeature{existential\_positive\_literals\_count}] \hfill{} \\
  number of positive, existential literals in the CNF
\item[\satfeature{(false, true)\_trivial}] \hfill{} \\
  is the CNF satisfied if all variables are claimed to be false (true)?
\item[\satfeature{goal\_clauses\_count}] \hfill{} \\
  number of goal clauses in the CNF
\item[\satfeature{literals\_count}] \hfill{} \\
  number of literals in the CNF (i.e. sum of clause lengths)
\item[\satfeature{literals\_frequency\_$k$\_to\_$k+5$}] \hfill{} \\
  let $n_l$ be the literal frequency of literal $l$,
  count the number of $n_l$ satisfying $\frac{k}{100} \leq n_l < \frac{k+5}{100}$
  where $k$ is a variable in $\{0,5,10,\ldots,90,95\}$ and $k=95$ counts
  $\frac{k}{100} \leq n_l \leq \frac{k+5}{100}$.
\item[\satfeature{literals\_frequency\_(largest, smallest, mean, median, sd)\_entropy}] \hfill{} \\
  statistics related to literal frequencies
\item[\satfeature{literals\_occurence\_one\_count}] \hfill{} \\
  number of literals with occurence $1$
\item[\satfeature{nbclauses}, \satfeature{nbvars}] \hfill{}
  number of clauses (variables) as defined in the CNF header
\item[\satfeature{negative\_literals\_in\_clause\_(smallest, largest, mean)}] \hfill{} \\
  statistics related to number of negative literals in clauses
\item[\satfeature{(positive, negative)\_unit\_clause\_count}] \hfill{} \\
  number of unit clauses with a positive (negative) literal
\item[\satfeature{positive\_literals\_count}] \hfill{} \\
  number of positive literals in CNF
\item[\satfeature{positive\_literals\_in\_clause\_(largest, smallest, mean, median, sd)}] \hfill{} \\
  statistics related to number of positive literals in clauses
\item[\satfeature{positive\_negative\_literals\_in\_clause\_ratio\_(mean, entropy)}] \hfill{} \\
  let $r_c$ be the number of positive literals divided by clause length of clause $c$,
  mean and related of all $r_c$
\item[\satfeature{positive\_negative\_literals\_in\_clause\_ratio\_mean}] \hfill{} \\
  mean of all $r_c$
\item[\satfeature{tautological\_literals\_count}] \hfill{} \\
  number of clauses which contain a tautological literal
\item[\satfeature{two\_literals\_clause\_count}] \hfill{} \\
  number of clauses with two literals
\item[\satfeature{variables\_frequency\_$k$\_to\_$k+5$}] \hfill{} \\
  same as \satfeature{literals\_frequency\_$k$\_to\_$k+5$} but for variables
\item[\satfeature{variables\_frequency\_(largest, smallest, mean, median, sd, entropy)}] \hfill{} \\
  same as \satfeature{literals\_frequency} but for variables
\item[\satfeature{variables\_used\_count}] \hfill{} \\
  number of variables with occurence greater 0
\end{description}

\section{Evaluation efficiency}
\label{sec:features-efficiency}
%
The resource requirements of those features have been classified:
\begin{description}
  \item[Type 1] read the files as bytestring, a DIMACS CNF parser is not necessary, constant memory is used
  \item[Type 2] features understand what a clause is, but still need constant memory
  \item[Type 3] subquadratic runtime and linear memory
  \item[Type 4] unrestricted
\end{description}
%
Memory and runtime is always considered in comparison with the filesize.

This classification should support future considerations regarding feature evaluation tools.
The suggested SAT features above have been explicitly selected to avoid Type 4 implementations to limit the time to compute features.
Furthermore tools evaluating only a subset of features (like Type~2 features) with better performance characteristics.

The Python implementation triggered MemoryErrors on a computer with 4~GB RAM for a 770~MB CNF file.
Followingly a much more efficient Go implementation was implemented which requires much less memory and is much faster.
\texttt{bench\_573.smt2.cnf} took 1 second in Go instead of 2 minutes in Python.
However, the data evaluated is less accurate compared to Python, because Python unlike Go provides a nice implementation of statistical tools in the standard library. Go algorithms were written on our own.

In the following section we define the dataset we consider.

\section{CNF dataset}
\label{sec:features-dataset}
%
To evaluate CNF features of a representative set of CNF files, it was necessary to identify equivalent CNF files in the best possible way.
Therefore we defined a hashing algorithm standardizing the CNF input provided to a SHA1 instance. Every CNF file is identifiable by its
\enquote{cnfhash 2.0.0} hash value.

In the next step a complete set of CNF files of previous SAT competitions was collected.
The following CNF file collections have been considered:

\begin{center}
  \begin{minipage}{0.48\linewidth}
    \begin{itemize}
      \itemsep2pt
      \item SAT Race 2008
      \item SAT09 Competition
      \item SAT-Race 2010
      \item SAT11 Competition
      \item SAT Challenge 2012
    \end{itemize}
  \end{minipage}
  \begin{minipage}{0.48\linewidth}
    \begin{itemize}
      \itemsep2pt
      \item SAT Competition 2013
      \item SAT Competition 2014
      \item SAT-Race 2015
      \item SAT Competition 2016
      \item SATlib
    \end{itemize}
  \end{minipage}
\end{center}

The benchmarks are mostly contributed by the participants of the associated conferences.
Others are reused from previous years. Individual projects allow to generate CNF files
for specific problems in a selectable problem size; such as CNFgen~\cite{cnfgen} by
Massimo Lauria.

Some files turned out to be problematic. In SATlib, 3 gzipped files couldn't be decompressed and several files
contain empty clauses. Empty clauses are assumed to immediately falsify the CNF and are therefore pointless.
I removed trailing zeros in CNFs. Variants of the DIMACS standard also expect lines with a percent symbol to
terminate the CNF. Besides those minor issues documented as part of the cnf-analysis project,
175 gigabytes of CNF files have been evaluated with a total of 68,069 CNF files (62,251 unique CNF files).

\section{The average SAT problem}
\label{sec:features-average}
%
\begin{prop}
  The set of public benchmarks in SAT competitions between 2008 and 2015
  represent average SAT problems
\end{prop}

It is important to point out that public benchmark files are specifically chosen
to be evaluated before a conference is held. Hence they are expected to terminate
within a given time frame and are therefore not oversized. On the one hand this
ensures that the problems are feasible, but on the other hand they might be a
biased selection. At this point no better data set is available and therefore
we proceeded with this dataset.

According to my results, an average SAT problem consists of:
\begin{itemize}
  \item 83,542 clauses in average ranging from 21 to 53,616,734 (median 430, $\sigma$ 848388)
  \item The longest clause we found had 61,473~literals, but the longest clause of an average CNF covers 20 literals.
  \item The number of total literals in a CNF ranges from 60 up to 150,609,758.
  \item The clause-variables ratio lies between 1.22 and 27,720 with mean 9.54 and $\sigma$ 139.
  \item The average length of a clause is expected to be 3. The largest clause length mean we found was 19.58 compared to 2.83 as the smallest clause length mean.
  \item Surprisingly, in average a CNF file has 205 connected literal components and 53 connected variable components. However both corresponding medians are $1$ meaning that at least have of the problems still have only one component.
  \item In average 32,787 clauses are definite clauses and 35,094 clauses are goal clauses.
  \item In average a literal occurs in 1.4~\% of the clauses of the CNF.
  \item 47~\% of literals in a clause are positive.
  \item The arithmetic mean tells 124 unit clauses per CNF file can be expected, but the median tells it is mostly 0.
  \item The largest variable found was 13,842,706 and 13,829,558~variables were used at most.
\end{itemize}

\section{Benford's law in CNF files}
\label{sec:benford}
%
Given this huge set of CNF files and therefore integers, we evaluated whether Benford's law holds.

A paper by Theodore P. Hill~\cite{hill1998first} characterizes Benford's Law
as the following conjecture:
\begin{theorem}
  Consider arbitrary data from tables, listings or other sources in publications, newspaper
  and so on and so forth, the digit $1$ occurs in about 30~\% of the time.
  The first digit is $1$ about $30$~percent of the time
  and $9$ only about $4.6$~percent of the time.
\end{theorem}
We evaluated that the leading digits 1 to 9 occur with probabilities
28.76~\%, 19.09~\%, 13.56~\%, 10.97~\%, 7.48~\%, 5.76~\%, 5.12~\%, 4.79~\%
and 4.46~\% respectively. We concluded that:
\begin{quote}
  Considering a deviation of 1.34~\% for digit $1$ (expected $30.1$, actual $28.76$) and
  0.12~\% for digit $9$ (expected $4.58$, actual $4.46$),
  we claimed the CNF~files considered satisfy Benford's law.
\end{quote}

\renewcommand*\chappic{img/encoding.pdf}
\renewcommand*\chapquote{There is concensus that encoding techniques usually have a dramatic impact on the efficiency of the SAT solver}
\renewcommand*\chapquotesrc{Magnus Bj\"ork}
\chapter{Problem encoding}
\label{ch:enc}

We already discussed how SAT solvers work and which input they take.
We also sketched how hash algorithm properties got broken using
differential cryptanalysis. In this section we combine those
subjects and describe how we designed an attack setting.

Every section represents a major approach whereas subsections
represent derivatives of this approach with minor tweaks.

\section{STP approach}
\label{sec:enc-stp}
%
Our first approach started with STP~\cite{stp} initially written by Vijay Ganesh and David L. Dill.
It is currently maintained by Mate Soos.

First we wrote an implementation using the CVC language to model the MD4 hash algorithm.
Reimplementing hash algorithms in CVC language (i.e. generating the corresponding code)
seemed cumbersome and we switched to the Python binding. With little modifications to
a working pure-Python MD4 implementation, the prototype was working.

However, STP was not fruitful for us, because we needed good control over the SAT encoding
which we expected to have a major influence on the performance. We used minisat as SAT solver
in the backend, but STP allows to exchange it for CryptoMiniSat which is a more modern and
versatile SAT solver. In the end this straight-forward approach worked nicely for MD4,
but we needed a different toolchain for SHA-256.

\section{algotocnf}
\label{sec:enc-algotocnf}
%
We implemented our own library \emph{algotocnf} to achieve greater flexibility in SAT design.

\subsection{Two instances and its difference}
\label{sec:enc-original}
%
\index{algotocnf}
algotocnf generated a CNF for a given hash algorithm implementation which is fed with
a symbolic variable. Only integer operations have been implemented.

Our tool \emph{algotocnf} implements the following strategy:
\begin{enumerate}
  \item Take a differential characteristic and the hash algorithm as input.
  \item For two instances of the hash algorithm, initialize symbolic variables
    (implemented as bitvectors) representing the bytes used to compute
    the hash value.

    Every bit is therefore represented as a boolean variable.
    If you apply addition, operator overloading in python will ensure
    that clauses are generated to describe the addition consisting of
    \boolf{XOR}s and \boolf{MAJ}s. Every operation is modelled as assignment.
    Hence an operation operating on a few Boolean variables is equivalent
    to a single variable which represents the result.
    Equivalently other operations related to integers are implemented as well.
  \item Constants used in the implementation are automatically converted
    to bitvectors with unit clauses.
  \item After running the hash algorithm with bitvectors per instance,
    all constraints related to the hash algorithm are added.
  \item Followingly the differential characteristic is read. Values such as $A_i$
    represent intermediate states of bitvectors. Therefore the corresponding bitvectors
    are looked up and equivalences with temporary bitvectors are added.
    Those temporary bitvectors are initialized with all constraints resulting from
    the bit conditions of this bitvector.
    In conclusion all constraints resulting from the differential characteristic are added.
  \item Finally the SAT solver is called. The CNF was mostly solved on a cluster
    specified in Appendix~\ref{app:setup}.
  \item Afterwards the program is run again to
    create the exactly same problem instance and the solver's solution replaces
    symbolic values with actual Boolean values. The resulting differential
    characteristic is parsed backed and printed out as differential characteristic
    where expectedly many bit conditions have been strengthened.
\end{enumerate}

We think \emph{algotocnf} mainly differs from other SAT abstraction tools because
of its implementation of differential logic. When adding clauses resulting from
the differential characteristic as constraints, the question arises how those
bit conditions are encoded. Essentially, we have only boolean values available,
but bit conditions tell constraints such as \enquote{a difference is given,
but the actual value is unknown}.

It seemed trivial to add a \emph{difference variable} for every pair of boolean
values representing a bit in the two instances. Furthermore the difference
variable $\Delta x$ is connected by a XOR with the variables of the pair $(x, x')$.
\[ \Delta x = x \oplus x' \]
Therefore it should be trivial for a preprocessor to simplify the formula
appropriately or actually we don't expect runtime differences for the larger
amount of variables.

And finally we expect the CNF to inherit a property of hash functions.
Inputs are provided into the hash algorithm and strongly intermingled
with other values. This should result in a high diffusion and almost every
variable is expected to share a clause with another variable.

The design decisions of this encoding are fundamental to the
resulting runtime discussed in chapter~\ref{ch:results}.

\subsection{Adding the differential description}
\label{sec:enc-diff-desc}
%
Using the approach in the previous section, we were able to find actual MD4 collisions
using a SAT solver. A SHA256 implementation followed which obviously lead to worse
runtime results, because the internal state of SHA-256 is much larger (by a factor
of at least 2). Can we further improve the runtime of the SAT solver?

Because we work with bitvectors and apply high-level operations like \boolf{MAJ} or addition,
we can additionally implement how differences in those operations propagate.
Magnus Daum's thesis on \enquote{Cryptanalysis of Hash Functions of the
MD4-Family}~\cite[Table 4.4]{daum2005cryptanalysis} discusses how differences propagate in \boolf{MAJ} and
\boolf{ITE} functions. Trivially, \boolf{XOR}s propagate differences the way they are~\footnote{
A difference in the arguments of two \boolf{XOR} instances remains the same difference
after applying \boolf{XOR} to each instance}.

This approach explicitly models differentiable behavior, which should be deducible
by the SAT solver itself. However, this lead to a major speedup which can be observed
in the runtime results of chapter~\ref{ch:results}.

\newpage
\subsection{Evaluating difference variables first}
\label{sec:enc-diff-desc-ocnf}
%
\begin{prop}
  Deriving difference values first, followed by actual bit values for the two instances,
  leads to a speedup.
\end{prop}

This proposed principle is fundamental to differential cryptanalysis. A previous tool
at our department implements propagation of hash algorithm values without a SAT solver
and this strategy is essential to good performance. This strategy was introduced
in the very early days of differential cryptanalysis and was also used by Wang et
al.~\cite{wang2004} to find their hash collisions.

So basically in terms of SAT solvers we want to guess values for differential variables
first and once all have been assigned, we try to find actual values for the two
hash algorithm instances.

It is important to point out that DIMACS does not specify a way to annotate Boolean
variables. As such that SAT solver cannot distinguish between difference variables
and variables of the instances. Therefore implementing this approach requires a custom
SAT solver which is given with lingeling~ats1o1.

Another proposition is important for this approach:

\begin{prop}
  Guessing difference values false first, followed by true,
  should solve hash collision problems faster.
\end{prop}

This proposition is justified by the desire to find as little differences as possible
in a hash collision to increase the probability of values cancelling each other out
in the later rounds.

\subsection{A light-weight approach}
\label{sec:enc-ocnf}
%
In this tweak we made a step back and considered the ideas of the previous
section but neglected the differential description. This approach was interesting
to see the effect of differential descriptions.

\subsection{Influencing the evaluation order}
\label{sec:enc-diff-desc-eo}
%
We took to idea to influence the evaluation order to the next level,
by applying the following design:

Let $\Delta x$ be the difference variable of pair $(x, x')$. We introduce a new boolean
variable $x^*$. We add clause
\[ x^* = (\Delta x \land x) \]
and explicitly tell the SAT solver to guess on $x^*$ before guessing on $\Delta x$, $x$ or $x'$.

The SAT solver will assign $x'=0$ first, because of the evaluation order. So either $\Delta x$
or $x$ must be false. $\Delta x$ is assigned false, because as difference variable it has a higher
priority over $x$. Equivalently for $x'=1$ we have $\Delta x$ needs to be true. So we actually
achieve an early guess on the difference variable.

This another approach evaluates in the results chapter.

\renewcommand*\chappic{img/megosat.pdf}
\renewcommand*\chapquote{}
\renewcommand*\chapquotesrc{}
\chapter{Results}
\label{ch:results}
%

\section{Benchmark results}

\section{Related work}

\section{Conclusion}
\label{sec:conclusion}




%\section{The distinction between concurrency and parallelism}

%\section{The characteristics of Go routines}

%CSP, Hoare paper

\renewcommand*\chappic{}
\renewcommand*\chapquote{}
\chapter{Summary and Future Work}
\label{ch:summary}
%
\section{Related work}
\label{sec:results-related}
%
Already in my bachelor thesis~\cite{bach} we tried to integrate a SAT solver into
a dedicated, automated tool finding collisions in hash functions using differential
cryptanalysis. This approach was not very successful as restarts between hash
algorithm rounds implied that intermediate results by the SAT solver got lost.
This was one motivation for this master thesis: We represent all details in CNF.

Research was already done by Ilya Mironov and Lintao Zhang~\cite{mironov2006applications}
to apply SAT solvers to differential cryptanalysis specifically to find hash collisions
in MD4. Their approach corresponds to our basic approach applied to Testcase~A
(compare with Section~\ref{sec:tcA}), where all differences are assigned.
Therefore our results in Table~\ref{tab:tcA-results} within one minute are
comparable with their evaluations within ten minutes.
We can clearly see how SAT technology and CPU performance
has progressed over 10 years since publication of this paper.

Finding hash collisions in Testcases~B, C, 21, 23 and 24 seems
to be a novel result of this master thesis and has not been found
in related work.

\section{Conclusion}
\label{sec:conclusion}
%
We successfully found full-round hash collisions for MD4
using SAT solvers mentioned in Section~\ref{sec:sat-solvers}.
We modified the lingeling SAT solver to improve our
runtime results further and found 24-round hash collisions
for SHA-256. Our attack starting points for MD4 ---
Testcases~\ref{fig:tcA}, \ref{fig:tcB} and \ref{fig:tcC} ---
are based on the work by Yusuke Naito, Yu Sasaki, Noboru Kunihiro and
Kazuo Ohta~\cite{sasaki2007new}. Our starting points for SHA-256
--- Testcases~\ref{fig:tc18}, \ref{fig:tc21}, \ref{fig:tc23}
and \ref{fig:tc24} --- are based on the work by Ivica Nikoli{\'c}
and Alex Biryukov~\cite{nikolic2008collisions}.

\section{Contributions}
\label{sec:contributions}
%
To encourage future work, the source code and data resulting from this
thesis is available online.
It allows the reader to run the experiments again and verify our claims.
We did our best to describe our hardware setup as accurately as possible.
At the following website, any results part of this project are collected:

\begin{quote}
  \url{http://lukas-prokop.at/proj/megosat/}
\end{quote}

% \noindent
% Several subprojects are part of this master thesis:
% \begin{description}
% \item[algotocnf]\hfill{} \\
%   A python library implementing the encoding described in Chapter~\ref{ch:enc}. \\[4pt]
%   \textbf{Python3 library and program:} \url{https://github.com/prokls/algotocnf}
% \item[cnf-hash]\hfill{} \\
%   A standardized way to produce a hash value for CNF files \\[4pt]
%   \textbf{Go implementation:} \url{https://github.com/prokls/cnf-hash-go} \\
%   \textbf{Python3 implementation:} \url{https://github.com/prokls/cnf-hash-py} \\
%   \textbf{Testsuite:} \url{https://github.com/prokls/cnf-hash-tests2}
% \item[cnf-analysis]\hfill{} \\
%   Evaluate SAT features for a given CNF file. \\[4pt]
%   \textbf{Go implementation:} \url{https://github.com/prokls/cnf-analysis-go} \\
%   \textbf{Python3 implementation:} \url{https://github.com/prokls/cnf-analysis-py} \\
%   \textbf{Testsuite:} \url{https://github.com/prokls/cnf-analysis-tests}
% \end{description}

\section{Future work}
\label{sec:future}
%
Future work might want to consider our design decisions made in
Chapter~\ref{ch:enc}.

In general, it would be interesting to generate our testcase for a broader
set of differential characteristics. Probably we can come up with empirical
results showing a relation between the size of unspecified areas in the
differential characteristic and the evaluated runtimes.

Furthermore, many SAT-related effort could be put to thoroughly discuss
why differential description provides such a significant performance improvement.
This necessarily means the SAT solver is generally not capable of deriving
the useful clauses, differential description provides. The resulting
numbers of restarts could also be subject of further research.

Lingeling ats1o1 was experimentally modified to define a separate set of
variables to be evaluated first. This approach seemed promising and
should be subject to future SAT research.

As far as cnf-analysis is concerned, the project aims to extend to
a larger set of SAT features and feedback by SAT solver developers
is appreciated. Our main contribution is a search
interface to search for SAT features given the cnfhash of a CNF file.
We hope to get in touch with new SAT feature ideas and SAT benchmark files.

\begin{appendices}
\chapter{Hardware setup}
\label{app:setup}
%
We introduce two hardwarde setups used in this master thesis.
%In the following we introduce two hardware setups used for evaluating
%data of this master thesis. The first setup is referred to as
%\enquote{Thinkpad x220} throughout this thesis whereas the second
%setup is referred to as \enquote{cluster}.
Thinkpad x220 was used
to simplify CNF files and evaluate SAT features. The cluster was used
for running SAT solvers and retrieving runtime results.

%\subsection{Thinkpad x220 Tablet}
%
\newcommand\caches[1][L3]{L1--#1 cache~sizes}
\begin{table}[!ht]
  \begin{center}
    \begin{tabular}{rc}
      \hline
      \emph{Model type}            & Thinkpad Lenovo x220 tablet, 4299-2P6 \\
      \emph{Processor}             & Intel i5-2520M, 2.50~GHz, dual-core, Hyperthreaded \\
      \emph{RAM}                   & 16~GB (extension to common retail setup) \\
      \emph{\caches{}}             & 32~KB / 256~KB / 3,072~KB \\
      \hline
    \end{tabular}
    \caption{Thinkpad x220 Tablet specification~\cite{thinkpadx220}}
    \label{setup:thinkpadx220}
  \end{center}
\end{table}

\begin{table}[!ht]
  \begin{center}
    \begin{tabular}{rl}
      \hline
        \multicolumn{2}{c}{\emph{Cluster node nehalem192g0 specification~\cite{intelX5690}}} \\
      \hline
        \emph{Processor}            & Intel Xeon X5690, 3.47~GHz, 6 cores, Hyperthreaded \\
        \emph{RAM}                  & 192~GB \\
        \emph{\caches{}}            & 32~KB / 256~KB / 12,288~KB \\
      \hline
        \multicolumn{2}{c}{\emph{Cluster node nehalem72g0 specification}} \\
      \hline
        \emph{Processor}            & Intel Xeon X5550, 2.67~GHz, 4 cores, Hyperthreaded \\
        \emph{RAM}                  & 72~GB \\
        \emph{\caches{}}            & 32~KB / 256~KB / 8,192~KB \\
      \hline
        \multicolumn{2}{c}{\emph{Cluster node xeon64g* specification}} \\
      \hline
        \emph{Processor}            & Intel Xeon E5430, 2.66~GHz, 4 cores, Hyperthreaded \\
        \emph{RAM}                  & 72~GB \\
        \emph{\caches{}}            & 32~KB / 6144~KB \\
      \hline
    \end{tabular}
    \caption{One node nehalem192g0, one node nehalem72g0 and four nodes xeon64g* were used for evaluating runtimes}
    \label{setup:cluster}
  \end{center}
\end{table}

\chapter{Testcases}
\label{app:tc}
%
\newcommand\cmp[1]{\emph{Please compare with Figure~\ref{fig:#1}.} \\ }
%
\section{MD4 testcase A}
\label{sec:tcA}
\cmp{tcA}
We can clearly see that all difference variables are defined.
Either they are true (bit condition {\dnCx}) or false (bit condition {\dnCh}),
but no variable has an undeciable state like {\dnCq}. At the top we can see
bit conditions {\dnCz} and {\dnCo} encoding the MD4 initial vector defined
by the hash algorithm. The differences {\dnCx} introduce the hash collision
and with round~47 being set of {\dnCh} only, the output is forced to be
equal between both hash algorithm instances. This testcase
is a trivial version of the differential characteristic described
in~\cite{sasaki2007new}.

\section{MD4 testcase B}
\label{sec:tcB}
\cmp{tcB}
In this testcase we have less knowledge about the state than in testcase A
because many values are encoded with {\dnCq} meaning that neither their
difference nor their actual values are known.
However, of course some {\dnCx} exists to introduce a hash collision
and the last round only consists of dashes to assert no difference
in the output.
So unlike testcase A, the SAT solver needs to figure out the difference
variables in rounds 0--11 increasing its overall runtime in all SAT solver
implementations.

\section{MD4 testcase C}
\label{sec:tcC}
\cmp{tcC}
This testcases introduces a hash collision which is expected to cancel out
at round 32. At the same time no information is provided about the intermediate
state of the hash algorithm in rounds 0--20.

\section{SHA-256 testcase 18 rounds}
\label{sec:tc18}
\cmp{tc18}
All SHA-256 testcases got derived from a paper by Ivica Nikolić and Alex Biryukov~\cite{nikolic2008collisions}.
Recall that message input of one block only fills W0 to W15. W16 or later
are generated based on the previous message words. This differential characteristic
depicts a hash collision introduced in round 3. When solving this differential characteristic
differences will especially occur in the most-significant bit of words with unspecified
difference. Differences cancel out after round~7.

\section{SHA-256 testcase 21 rounds}
\label{sec:tc21}
\cmp{tc21}
Unlike Testcase~\ref{sec:tc18} word A5 is underspecified with question marks.
This certainly makes the problem harder for a SAT solver.
At the same time differences of the message words W9 and W10 are specified.
This makes the problem easier. At the same time the whole collision covers
21~rounds, unlike Testcase~\ref{sec:tc18} with 18~rounds.

\section{SHA-256 testcase 23 rounds}
\label{sec:tc23}
\cmp{tc23}
Message word W9 specifies no difference, but the collision is extended
to 23~rounds.

\section{SHA-256 testcase 24 rounds}
\label{sec:tc24}
\cmp{tc24}
No measures were taken to simplify the problem for the SAT solver,
but the hash collision needs to be found for 24~rounds, which
makes the problem hard for the SAT solver because of its increased
state size.




\texttt{
\begin{figure}[b]
\begin{center}
\scalebox{.55}{
\begin{tabular}{|r|c|c|c|c|}
\hline
  $i$  &      &  A  &   &  W \\
\hline
\dnI{-4} & \dnW: & {{\dnCz}{\dnCo}{\dnCo}{\dnCz}{\dnCz}{\dnCo}{\dnCo}{\dnCo}{\dnCz}{\dnCo}{\dnCz}{\dnCz}{\dnCz}{\dnCo}{\dnCz}{\dnCo}{\dnCz}{\dnCz}{\dnCo}{\dnCz}{\dnCz}{\dnCz}{\dnCo}{\dnCo}{\dnCz}{\dnCz}{\dnCz}{\dnCz}{\dnCz}{\dnCz}{\dnCz}{\dnCo}} & & \\
\dnI{-3} & \dnW: & {{\dnCz}{\dnCz}{\dnCz}{\dnCo}{\dnCz}{\dnCz}{\dnCz}{\dnCz}{\dnCz}{\dnCz}{\dnCo}{\dnCo}{\dnCz}{\dnCz}{\dnCo}{\dnCz}{\dnCz}{\dnCo}{\dnCz}{\dnCo}{\dnCz}{\dnCo}{\dnCz}{\dnCz}{\dnCz}{\dnCo}{\dnCo}{\dnCo}{\dnCz}{\dnCo}{\dnCo}{\dnCz}} & & \\
\dnI{-2} & \dnW: & {{\dnCo}{\dnCz}{\dnCz}{\dnCo}{\dnCo}{\dnCz}{\dnCz}{\dnCz}{\dnCo}{\dnCz}{\dnCo}{\dnCo}{\dnCo}{\dnCz}{\dnCo}{\dnCz}{\dnCo}{\dnCo}{\dnCz}{\dnCo}{\dnCo}{\dnCo}{\dnCz}{\dnCz}{\dnCo}{\dnCo}{\dnCo}{\dnCo}{\dnCo}{\dnCo}{\dnCo}{\dnCz}} & & \\
\dnI{-1} & \dnW: & {{\dnCo}{\dnCo}{\dnCo}{\dnCz}{\dnCo}{\dnCo}{\dnCo}{\dnCo}{\dnCo}{\dnCo}{\dnCz}{\dnCz}{\dnCo}{\dnCo}{\dnCz}{\dnCo}{\dnCo}{\dnCz}{\dnCo}{\dnCz}{\dnCo}{\dnCz}{\dnCo}{\dnCo}{\dnCo}{\dnCz}{\dnCz}{\dnCz}{\dnCo}{\dnCz}{\dnCz}{\dnCo}} & & \\
\dnI{0} & \dnW: & {{\dnCx}{\dnCh}{\dnCh}{\dnCh}{\dnCh}{\dnCh}{\dnCh}{\dnCh}{\dnCh}{\dnCh}{\dnCh}{\dnCh}{\dnCh}{\dnCh}{\dnCh}{\dnCh}{\dnCh}{\dnCh}{\dnCh}{\dnCh}{\dnCh}{\dnCh}{\dnCh}{\dnCh}{\dnCh}{\dnCh}{\dnCh}{\dnCh}{\dnCh}{\dnCh}{\dnCh}{\dnCh}} & \dnW[W]{}: & {{\dnCh}{\dnCh}{\dnCh}{\dnCx}{\dnCh}{\dnCh}{\dnCh}{\dnCh}{\dnCh}{\dnCh}{\dnCh}{\dnCh}{\dnCh}{\dnCh}{\dnCh}{\dnCh}{\dnCh}{\dnCh}{\dnCh}{\dnCh}{\dnCh}{\dnCh}{\dnCh}{\dnCh}{\dnCh}{\dnCh}{\dnCh}{\dnCh}{\dnCh}{\dnCh}{\dnCh}{\dnCh}} \\
\dnI{1} & \dnW: & {{\dnCh}{\dnCh}{\dnCh}{\dnCh}{\dnCh}{\dnCh}{\dnCh}{\dnCh}{\dnCh}{\dnCh}{\dnCh}{\dnCh}{\dnCh}{\dnCh}{\dnCh}{\dnCh}{\dnCh}{\dnCh}{\dnCh}{\dnCh}{\dnCh}{\dnCh}{\dnCh}{\dnCh}{\dnCh}{\dnCh}{\dnCh}{\dnCh}{\dnCh}{\dnCh}{\dnCh}{\dnCh}} & \dnW[W]{}: & {{\dnCh}{\dnCh}{\dnCh}{\dnCh}{\dnCh}{\dnCh}{\dnCh}{\dnCh}{\dnCh}{\dnCh}{\dnCh}{\dnCh}{\dnCh}{\dnCh}{\dnCh}{\dnCh}{\dnCh}{\dnCh}{\dnCh}{\dnCh}{\dnCh}{\dnCh}{\dnCh}{\dnCh}{\dnCh}{\dnCh}{\dnCh}{\dnCh}{\dnCh}{\dnCh}{\dnCh}{\dnCh}} \\
\dnI{2} & \dnW: & {{\dnCh}{\dnCh}{\dnCh}{\dnCh}{\dnCh}{\dnCh}{\dnCh}{\dnCh}{\dnCh}{\dnCh}{\dnCh}{\dnCh}{\dnCh}{\dnCh}{\dnCh}{\dnCh}{\dnCh}{\dnCh}{\dnCh}{\dnCh}{\dnCh}{\dnCx}{\dnCh}{\dnCh}{\dnCh}{\dnCh}{\dnCh}{\dnCh}{\dnCh}{\dnCh}{\dnCh}{\dnCh}} & \dnW[W]{}: & {{\dnCx}{\dnCh}{\dnCh}{\dnCh}{\dnCh}{\dnCh}{\dnCh}{\dnCh}{\dnCh}{\dnCh}{\dnCh}{\dnCh}{\dnCh}{\dnCh}{\dnCh}{\dnCh}{\dnCh}{\dnCh}{\dnCh}{\dnCh}{\dnCh}{\dnCh}{\dnCh}{\dnCh}{\dnCh}{\dnCh}{\dnCh}{\dnCh}{\dnCh}{\dnCh}{\dnCh}{\dnCh}} \\
\dnI{3} & \dnW: & {{\dnCx}{\dnCx}{\dnCx}{\dnCh}{\dnCh}{\dnCh}{\dnCh}{\dnCh}{\dnCh}{\dnCh}{\dnCh}{\dnCh}{\dnCh}{\dnCh}{\dnCh}{\dnCh}{\dnCh}{\dnCh}{\dnCh}{\dnCh}{\dnCh}{\dnCh}{\dnCh}{\dnCh}{\dnCh}{\dnCh}{\dnCh}{\dnCh}{\dnCh}{\dnCh}{\dnCh}{\dnCh}} & \dnW[W]{}: & {{\dnCh}{\dnCh}{\dnCh}{\dnCh}{\dnCh}{\dnCh}{\dnCh}{\dnCh}{\dnCh}{\dnCh}{\dnCh}{\dnCh}{\dnCh}{\dnCh}{\dnCh}{\dnCh}{\dnCh}{\dnCh}{\dnCh}{\dnCh}{\dnCh}{\dnCh}{\dnCh}{\dnCh}{\dnCh}{\dnCh}{\dnCh}{\dnCh}{\dnCh}{\dnCh}{\dnCh}{\dnCh}} \\
\dnI{4} & \dnW: & {{\dnCh}{\dnCh}{\dnCh}{\dnCh}{\dnCh}{\dnCh}{\dnCh}{\dnCh}{\dnCh}{\dnCh}{\dnCh}{\dnCh}{\dnCh}{\dnCh}{\dnCh}{\dnCh}{\dnCh}{\dnCh}{\dnCh}{\dnCh}{\dnCh}{\dnCh}{\dnCh}{\dnCh}{\dnCh}{\dnCh}{\dnCh}{\dnCh}{\dnCh}{\dnCh}{\dnCx}{\dnCx}} & \dnW[W]{}: & {{\dnCx}{\dnCh}{\dnCh}{\dnCh}{\dnCh}{\dnCh}{\dnCh}{\dnCh}{\dnCh}{\dnCh}{\dnCh}{\dnCh}{\dnCh}{\dnCh}{\dnCh}{\dnCh}{\dnCh}{\dnCh}{\dnCh}{\dnCh}{\dnCh}{\dnCh}{\dnCh}{\dnCh}{\dnCh}{\dnCh}{\dnCh}{\dnCh}{\dnCh}{\dnCh}{\dnCh}{\dnCh}} \\
\dnI{5} & \dnW: & {{\dnCh}{\dnCh}{\dnCh}{\dnCh}{\dnCh}{\dnCh}{\dnCh}{\dnCx}{\dnCx}{\dnCx}{\dnCx}{\dnCx}{\dnCx}{\dnCx}{\dnCx}{\dnCx}{\dnCx}{\dnCx}{\dnCx}{\dnCx}{\dnCx}{\dnCx}{\dnCx}{\dnCh}{\dnCx}{\dnCh}{\dnCh}{\dnCh}{\dnCh}{\dnCh}{\dnCh}{\dnCh}} & \dnW[W]{}: & {{\dnCh}{\dnCh}{\dnCh}{\dnCh}{\dnCh}{\dnCh}{\dnCh}{\dnCh}{\dnCh}{\dnCh}{\dnCh}{\dnCh}{\dnCh}{\dnCh}{\dnCh}{\dnCh}{\dnCh}{\dnCh}{\dnCh}{\dnCh}{\dnCh}{\dnCh}{\dnCh}{\dnCh}{\dnCh}{\dnCh}{\dnCh}{\dnCh}{\dnCh}{\dnCh}{\dnCh}{\dnCh}} \\
\dnI{6} & \dnW: & {{\dnCx}{\dnCh}{\dnCh}{\dnCh}{\dnCh}{\dnCh}{\dnCh}{\dnCh}{\dnCh}{\dnCx}{\dnCh}{\dnCh}{\dnCh}{\dnCh}{\dnCh}{\dnCh}{\dnCh}{\dnCh}{\dnCh}{\dnCh}{\dnCh}{\dnCx}{\dnCh}{\dnCx}{\dnCh}{\dnCx}{\dnCx}{\dnCx}{\dnCx}{\dnCh}{\dnCh}{\dnCx}} & \dnW[W]{}: & {{\dnCh}{\dnCh}{\dnCh}{\dnCh}{\dnCh}{\dnCh}{\dnCh}{\dnCh}{\dnCh}{\dnCh}{\dnCh}{\dnCh}{\dnCh}{\dnCh}{\dnCh}{\dnCh}{\dnCh}{\dnCh}{\dnCh}{\dnCh}{\dnCh}{\dnCh}{\dnCh}{\dnCh}{\dnCh}{\dnCh}{\dnCh}{\dnCh}{\dnCh}{\dnCh}{\dnCh}{\dnCh}} \\
\dnI{7} & \dnW: & {{\dnCh}{\dnCh}{\dnCh}{\dnCh}{\dnCh}{\dnCh}{\dnCh}{\dnCh}{\dnCh}{\dnCx}{\dnCh}{\dnCh}{\dnCh}{\dnCh}{\dnCh}{\dnCx}{\dnCh}{\dnCh}{\dnCh}{\dnCh}{\dnCh}{\dnCh}{\dnCh}{\dnCh}{\dnCh}{\dnCh}{\dnCh}{\dnCh}{\dnCh}{\dnCh}{\dnCh}{\dnCh}} & \dnW[W]{}: & {{\dnCh}{\dnCh}{\dnCh}{\dnCh}{\dnCh}{\dnCh}{\dnCh}{\dnCh}{\dnCh}{\dnCh}{\dnCh}{\dnCh}{\dnCh}{\dnCh}{\dnCh}{\dnCh}{\dnCh}{\dnCh}{\dnCh}{\dnCh}{\dnCh}{\dnCh}{\dnCh}{\dnCh}{\dnCh}{\dnCh}{\dnCh}{\dnCh}{\dnCh}{\dnCh}{\dnCh}{\dnCh}} \\
\dnI{8} & \dnW: & {{\dnCh}{\dnCh}{\dnCh}{\dnCh}{\dnCh}{\dnCh}{\dnCh}{\dnCh}{\dnCh}{\dnCx}{\dnCh}{\dnCh}{\dnCh}{\dnCh}{\dnCh}{\dnCh}{\dnCh}{\dnCx}{\dnCh}{\dnCx}{\dnCh}{\dnCx}{\dnCh}{\dnCh}{\dnCh}{\dnCh}{\dnCh}{\dnCh}{\dnCh}{\dnCh}{\dnCh}{\dnCh}} & \dnW[W]{}: & {{\dnCx}{\dnCh}{\dnCh}{\dnCh}{\dnCh}{\dnCh}{\dnCh}{\dnCh}{\dnCh}{\dnCh}{\dnCh}{\dnCh}{\dnCh}{\dnCh}{\dnCh}{\dnCh}{\dnCh}{\dnCh}{\dnCh}{\dnCh}{\dnCh}{\dnCh}{\dnCh}{\dnCh}{\dnCh}{\dnCh}{\dnCh}{\dnCh}{\dnCh}{\dnCh}{\dnCh}{\dnCh}} \\
\dnI{9} & \dnW: & {{\dnCh}{\dnCh}{\dnCh}{\dnCh}{\dnCh}{\dnCh}{\dnCh}{\dnCh}{\dnCh}{\dnCh}{\dnCh}{\dnCh}{\dnCh}{\dnCh}{\dnCh}{\dnCh}{\dnCh}{\dnCx}{\dnCh}{\dnCh}{\dnCh}{\dnCx}{\dnCh}{\dnCh}{\dnCh}{\dnCh}{\dnCh}{\dnCh}{\dnCh}{\dnCh}{\dnCh}{\dnCh}} & \dnW[W]{}: & {{\dnCh}{\dnCh}{\dnCh}{\dnCh}{\dnCh}{\dnCh}{\dnCh}{\dnCh}{\dnCh}{\dnCh}{\dnCh}{\dnCh}{\dnCh}{\dnCh}{\dnCh}{\dnCh}{\dnCh}{\dnCh}{\dnCh}{\dnCh}{\dnCh}{\dnCh}{\dnCh}{\dnCh}{\dnCh}{\dnCh}{\dnCh}{\dnCh}{\dnCh}{\dnCh}{\dnCh}{\dnCh}} \\
\dnI{10} & \dnW: & {{\dnCh}{\dnCh}{\dnCh}{\dnCh}{\dnCh}{\dnCh}{\dnCx}{\dnCh}{\dnCh}{\dnCh}{\dnCh}{\dnCh}{\dnCx}{\dnCh}{\dnCh}{\dnCx}{\dnCx}{\dnCx}{\dnCh}{\dnCx}{\dnCx}{\dnCx}{\dnCh}{\dnCh}{\dnCh}{\dnCh}{\dnCh}{\dnCh}{\dnCh}{\dnCh}{\dnCh}{\dnCh}} & \dnW[W]{}: & {{\dnCh}{\dnCh}{\dnCh}{\dnCh}{\dnCh}{\dnCh}{\dnCh}{\dnCh}{\dnCh}{\dnCh}{\dnCh}{\dnCh}{\dnCh}{\dnCh}{\dnCh}{\dnCh}{\dnCh}{\dnCh}{\dnCh}{\dnCh}{\dnCh}{\dnCh}{\dnCh}{\dnCh}{\dnCh}{\dnCh}{\dnCh}{\dnCh}{\dnCh}{\dnCh}{\dnCh}{\dnCh}} \\
\dnI{11} & \dnW: & {{\dnCx}{\dnCh}{\dnCh}{\dnCh}{\dnCh}{\dnCh}{\dnCh}{\dnCh}{\dnCh}{\dnCh}{\dnCh}{\dnCh}{\dnCh}{\dnCh}{\dnCh}{\dnCh}{\dnCh}{\dnCx}{\dnCx}{\dnCx}{\dnCh}{\dnCx}{\dnCh}{\dnCh}{\dnCh}{\dnCh}{\dnCh}{\dnCh}{\dnCh}{\dnCh}{\dnCh}{\dnCh}} & \dnW[W]{}: & {{\dnCh}{\dnCh}{\dnCh}{\dnCh}{\dnCh}{\dnCh}{\dnCh}{\dnCh}{\dnCh}{\dnCh}{\dnCh}{\dnCh}{\dnCh}{\dnCh}{\dnCh}{\dnCh}{\dnCh}{\dnCh}{\dnCh}{\dnCh}{\dnCh}{\dnCh}{\dnCh}{\dnCh}{\dnCh}{\dnCh}{\dnCh}{\dnCh}{\dnCh}{\dnCh}{\dnCh}{\dnCh}} \\
\dnI{12} & \dnW: & {{\dnCh}{\dnCh}{\dnCh}{\dnCx}{\dnCh}{\dnCh}{\dnCx}{\dnCh}{\dnCh}{\dnCh}{\dnCh}{\dnCh}{\dnCh}{\dnCh}{\dnCh}{\dnCh}{\dnCh}{\dnCh}{\dnCh}{\dnCh}{\dnCh}{\dnCh}{\dnCh}{\dnCh}{\dnCh}{\dnCh}{\dnCh}{\dnCh}{\dnCh}{\dnCh}{\dnCh}{\dnCh}} & \dnW[W]{}: & {{\dnCx}{\dnCh}{\dnCh}{\dnCh}{\dnCh}{\dnCh}{\dnCh}{\dnCh}{\dnCh}{\dnCh}{\dnCh}{\dnCh}{\dnCh}{\dnCh}{\dnCh}{\dnCh}{\dnCh}{\dnCh}{\dnCh}{\dnCh}{\dnCh}{\dnCh}{\dnCh}{\dnCh}{\dnCh}{\dnCh}{\dnCh}{\dnCh}{\dnCh}{\dnCh}{\dnCh}{\dnCh}} \\
\dnI{13} & \dnW: & {{\dnCh}{\dnCh}{\dnCh}{\dnCh}{\dnCh}{\dnCh}{\dnCh}{\dnCh}{\dnCh}{\dnCh}{\dnCh}{\dnCh}{\dnCh}{\dnCh}{\dnCh}{\dnCh}{\dnCh}{\dnCh}{\dnCh}{\dnCh}{\dnCh}{\dnCh}{\dnCh}{\dnCh}{\dnCh}{\dnCh}{\dnCh}{\dnCh}{\dnCh}{\dnCh}{\dnCh}{\dnCh}} & \dnW[W]{}: & {{\dnCh}{\dnCh}{\dnCh}{\dnCh}{\dnCh}{\dnCh}{\dnCh}{\dnCh}{\dnCh}{\dnCh}{\dnCh}{\dnCh}{\dnCh}{\dnCh}{\dnCh}{\dnCh}{\dnCh}{\dnCh}{\dnCh}{\dnCh}{\dnCh}{\dnCh}{\dnCh}{\dnCh}{\dnCh}{\dnCh}{\dnCh}{\dnCh}{\dnCh}{\dnCh}{\dnCh}{\dnCh}} \\
\dnI{14} & \dnW: & {{\dnCh}{\dnCx}{\dnCh}{\dnCh}{\dnCh}{\dnCh}{\dnCh}{\dnCh}{\dnCh}{\dnCh}{\dnCh}{\dnCh}{\dnCh}{\dnCh}{\dnCh}{\dnCh}{\dnCh}{\dnCh}{\dnCh}{\dnCh}{\dnCh}{\dnCh}{\dnCh}{\dnCh}{\dnCh}{\dnCh}{\dnCh}{\dnCh}{\dnCh}{\dnCh}{\dnCh}{\dnCh}} & \dnW[W]{}: & {{\dnCh}{\dnCh}{\dnCh}{\dnCh}{\dnCh}{\dnCh}{\dnCh}{\dnCh}{\dnCh}{\dnCh}{\dnCh}{\dnCh}{\dnCh}{\dnCh}{\dnCh}{\dnCh}{\dnCh}{\dnCh}{\dnCh}{\dnCh}{\dnCh}{\dnCh}{\dnCh}{\dnCh}{\dnCh}{\dnCh}{\dnCh}{\dnCh}{\dnCh}{\dnCh}{\dnCh}{\dnCh}} \\
\dnI{15} & \dnW: & {{\dnCx}{\dnCh}{\dnCx}{\dnCh}{\dnCh}{\dnCh}{\dnCh}{\dnCh}{\dnCh}{\dnCh}{\dnCh}{\dnCh}{\dnCh}{\dnCx}{\dnCh}{\dnCh}{\dnCh}{\dnCh}{\dnCh}{\dnCh}{\dnCh}{\dnCh}{\dnCh}{\dnCh}{\dnCh}{\dnCh}{\dnCh}{\dnCh}{\dnCh}{\dnCh}{\dnCh}{\dnCh}} & \dnW[W]{}: & {{\dnCh}{\dnCh}{\dnCh}{\dnCh}{\dnCh}{\dnCh}{\dnCh}{\dnCh}{\dnCh}{\dnCh}{\dnCh}{\dnCh}{\dnCh}{\dnCh}{\dnCh}{\dnCh}{\dnCh}{\dnCh}{\dnCh}{\dnCh}{\dnCh}{\dnCh}{\dnCh}{\dnCh}{\dnCh}{\dnCh}{\dnCh}{\dnCh}{\dnCh}{\dnCh}{\dnCh}{\dnCh}} \\
\dnI{16} & \dnW: & {{\dnCh}{\dnCx}{\dnCx}{\dnCx}{\dnCh}{\dnCh}{\dnCh}{\dnCh}{\dnCh}{\dnCh}{\dnCh}{\dnCh}{\dnCh}{\dnCh}{\dnCh}{\dnCh}{\dnCh}{\dnCh}{\dnCh}{\dnCh}{\dnCh}{\dnCh}{\dnCh}{\dnCh}{\dnCh}{\dnCh}{\dnCh}{\dnCh}{\dnCh}{\dnCh}{\dnCh}{\dnCh}} & & \\
\dnI{17} & \dnW: & {{\dnCh}{\dnCh}{\dnCh}{\dnCh}{\dnCh}{\dnCh}{\dnCh}{\dnCh}{\dnCh}{\dnCh}{\dnCh}{\dnCh}{\dnCh}{\dnCh}{\dnCh}{\dnCh}{\dnCh}{\dnCh}{\dnCh}{\dnCh}{\dnCh}{\dnCh}{\dnCh}{\dnCh}{\dnCh}{\dnCh}{\dnCh}{\dnCh}{\dnCh}{\dnCh}{\dnCh}{\dnCh}} & & \\
\dnI{18} & \dnW: & {{\dnCh}{\dnCh}{\dnCh}{\dnCh}{\dnCh}{\dnCh}{\dnCh}{\dnCh}{\dnCh}{\dnCh}{\dnCh}{\dnCh}{\dnCh}{\dnCh}{\dnCh}{\dnCh}{\dnCh}{\dnCh}{\dnCh}{\dnCh}{\dnCh}{\dnCh}{\dnCh}{\dnCh}{\dnCh}{\dnCh}{\dnCh}{\dnCh}{\dnCh}{\dnCh}{\dnCh}{\dnCh}} & & \\
\dnI{19} & \dnW: & {{\dnCx}{\dnCh}{\dnCh}{\dnCh}{\dnCh}{\dnCh}{\dnCh}{\dnCh}{\dnCh}{\dnCh}{\dnCh}{\dnCh}{\dnCh}{\dnCh}{\dnCh}{\dnCh}{\dnCh}{\dnCh}{\dnCh}{\dnCh}{\dnCh}{\dnCh}{\dnCh}{\dnCh}{\dnCh}{\dnCh}{\dnCh}{\dnCh}{\dnCh}{\dnCh}{\dnCh}{\dnCh}} & & \\
\dnI{20} & \dnW: & {{\dnCx}{\dnCh}{\dnCh}{\dnCh}{\dnCh}{\dnCh}{\dnCh}{\dnCh}{\dnCh}{\dnCh}{\dnCh}{\dnCh}{\dnCh}{\dnCh}{\dnCh}{\dnCh}{\dnCh}{\dnCh}{\dnCh}{\dnCh}{\dnCh}{\dnCh}{\dnCh}{\dnCh}{\dnCh}{\dnCh}{\dnCh}{\dnCh}{\dnCh}{\dnCh}{\dnCh}{\dnCh}} & & \\
\dnI{21} & \dnW: & {{\dnCh}{\dnCh}{\dnCh}{\dnCh}{\dnCh}{\dnCh}{\dnCh}{\dnCh}{\dnCh}{\dnCh}{\dnCh}{\dnCh}{\dnCh}{\dnCh}{\dnCh}{\dnCh}{\dnCh}{\dnCh}{\dnCh}{\dnCh}{\dnCh}{\dnCh}{\dnCh}{\dnCh}{\dnCh}{\dnCh}{\dnCh}{\dnCh}{\dnCh}{\dnCh}{\dnCh}{\dnCh}} & & \\
\dnI{22} & \dnW: & {{\dnCh}{\dnCh}{\dnCh}{\dnCh}{\dnCh}{\dnCh}{\dnCh}{\dnCh}{\dnCh}{\dnCh}{\dnCh}{\dnCh}{\dnCh}{\dnCh}{\dnCh}{\dnCh}{\dnCh}{\dnCh}{\dnCh}{\dnCh}{\dnCh}{\dnCh}{\dnCh}{\dnCh}{\dnCh}{\dnCh}{\dnCh}{\dnCh}{\dnCh}{\dnCh}{\dnCh}{\dnCh}} & & \\
\dnI{23} & \dnW: & {{\dnCh}{\dnCh}{\dnCh}{\dnCh}{\dnCh}{\dnCh}{\dnCh}{\dnCh}{\dnCh}{\dnCh}{\dnCh}{\dnCh}{\dnCh}{\dnCh}{\dnCh}{\dnCh}{\dnCh}{\dnCh}{\dnCh}{\dnCh}{\dnCh}{\dnCh}{\dnCh}{\dnCh}{\dnCh}{\dnCh}{\dnCh}{\dnCh}{\dnCh}{\dnCh}{\dnCh}{\dnCh}} & & \\
\dnI{24} & \dnW: & {{\dnCh}{\dnCh}{\dnCh}{\dnCh}{\dnCh}{\dnCh}{\dnCh}{\dnCh}{\dnCh}{\dnCh}{\dnCh}{\dnCh}{\dnCh}{\dnCh}{\dnCh}{\dnCh}{\dnCh}{\dnCh}{\dnCh}{\dnCh}{\dnCh}{\dnCh}{\dnCh}{\dnCh}{\dnCh}{\dnCh}{\dnCh}{\dnCh}{\dnCh}{\dnCh}{\dnCh}{\dnCh}} & & \\
\dnI{25} & \dnW: & {{\dnCh}{\dnCh}{\dnCh}{\dnCh}{\dnCh}{\dnCh}{\dnCh}{\dnCh}{\dnCh}{\dnCh}{\dnCh}{\dnCh}{\dnCh}{\dnCh}{\dnCh}{\dnCh}{\dnCh}{\dnCh}{\dnCh}{\dnCh}{\dnCh}{\dnCh}{\dnCh}{\dnCh}{\dnCh}{\dnCh}{\dnCh}{\dnCh}{\dnCh}{\dnCh}{\dnCh}{\dnCh}} & & \\
\dnI{26} & \dnW: & {{\dnCh}{\dnCh}{\dnCh}{\dnCh}{\dnCh}{\dnCh}{\dnCh}{\dnCh}{\dnCh}{\dnCh}{\dnCh}{\dnCh}{\dnCh}{\dnCh}{\dnCh}{\dnCh}{\dnCh}{\dnCh}{\dnCh}{\dnCh}{\dnCh}{\dnCh}{\dnCh}{\dnCh}{\dnCh}{\dnCh}{\dnCh}{\dnCh}{\dnCh}{\dnCh}{\dnCh}{\dnCh}} & & \\
\dnI{27} & \dnW: & {{\dnCh}{\dnCh}{\dnCh}{\dnCh}{\dnCh}{\dnCh}{\dnCh}{\dnCh}{\dnCh}{\dnCh}{\dnCh}{\dnCh}{\dnCh}{\dnCh}{\dnCh}{\dnCh}{\dnCh}{\dnCh}{\dnCh}{\dnCh}{\dnCh}{\dnCh}{\dnCh}{\dnCh}{\dnCh}{\dnCh}{\dnCh}{\dnCh}{\dnCh}{\dnCh}{\dnCh}{\dnCh}} & & \\
\dnI{28} & \dnW: & {{\dnCh}{\dnCh}{\dnCh}{\dnCh}{\dnCh}{\dnCh}{\dnCh}{\dnCh}{\dnCh}{\dnCh}{\dnCh}{\dnCh}{\dnCh}{\dnCh}{\dnCh}{\dnCh}{\dnCh}{\dnCh}{\dnCh}{\dnCh}{\dnCh}{\dnCh}{\dnCh}{\dnCh}{\dnCh}{\dnCh}{\dnCh}{\dnCh}{\dnCh}{\dnCh}{\dnCh}{\dnCh}} & & \\
\dnI{29} & \dnW: & {{\dnCh}{\dnCh}{\dnCh}{\dnCh}{\dnCh}{\dnCh}{\dnCh}{\dnCh}{\dnCh}{\dnCh}{\dnCh}{\dnCh}{\dnCh}{\dnCh}{\dnCh}{\dnCh}{\dnCh}{\dnCh}{\dnCh}{\dnCh}{\dnCh}{\dnCh}{\dnCh}{\dnCh}{\dnCh}{\dnCh}{\dnCh}{\dnCh}{\dnCh}{\dnCh}{\dnCh}{\dnCh}} & & \\
\dnI{30} & \dnW: & {{\dnCh}{\dnCh}{\dnCh}{\dnCh}{\dnCh}{\dnCh}{\dnCh}{\dnCh}{\dnCh}{\dnCh}{\dnCh}{\dnCh}{\dnCh}{\dnCh}{\dnCh}{\dnCh}{\dnCh}{\dnCh}{\dnCh}{\dnCh}{\dnCh}{\dnCh}{\dnCh}{\dnCh}{\dnCh}{\dnCh}{\dnCh}{\dnCh}{\dnCh}{\dnCh}{\dnCh}{\dnCh}} & & \\
\dnI{31} & \dnW: & {{\dnCh}{\dnCh}{\dnCh}{\dnCh}{\dnCh}{\dnCh}{\dnCh}{\dnCh}{\dnCh}{\dnCh}{\dnCh}{\dnCh}{\dnCh}{\dnCh}{\dnCh}{\dnCh}{\dnCh}{\dnCh}{\dnCh}{\dnCh}{\dnCh}{\dnCh}{\dnCh}{\dnCh}{\dnCh}{\dnCh}{\dnCh}{\dnCh}{\dnCh}{\dnCh}{\dnCh}{\dnCh}} & & \\
\dnI{32} & \dnW: & {{\dnCx}{\dnCh}{\dnCh}{\dnCh}{\dnCh}{\dnCh}{\dnCh}{\dnCh}{\dnCh}{\dnCh}{\dnCh}{\dnCh}{\dnCh}{\dnCh}{\dnCh}{\dnCh}{\dnCh}{\dnCh}{\dnCh}{\dnCh}{\dnCh}{\dnCh}{\dnCh}{\dnCh}{\dnCh}{\dnCh}{\dnCh}{\dnCh}{\dnCh}{\dnCh}{\dnCh}{\dnCh}} & & \\
\dnI{33} & \dnW: & {{\dnCh}{\dnCh}{\dnCh}{\dnCh}{\dnCh}{\dnCh}{\dnCh}{\dnCh}{\dnCh}{\dnCh}{\dnCh}{\dnCh}{\dnCh}{\dnCh}{\dnCh}{\dnCh}{\dnCh}{\dnCh}{\dnCh}{\dnCh}{\dnCh}{\dnCh}{\dnCh}{\dnCh}{\dnCh}{\dnCh}{\dnCh}{\dnCh}{\dnCh}{\dnCh}{\dnCh}{\dnCh}} & & \\
\dnI{34} & \dnW: & {{\dnCh}{\dnCh}{\dnCh}{\dnCh}{\dnCh}{\dnCh}{\dnCh}{\dnCh}{\dnCh}{\dnCh}{\dnCh}{\dnCh}{\dnCh}{\dnCh}{\dnCh}{\dnCh}{\dnCh}{\dnCh}{\dnCh}{\dnCh}{\dnCh}{\dnCh}{\dnCh}{\dnCh}{\dnCh}{\dnCh}{\dnCh}{\dnCh}{\dnCh}{\dnCh}{\dnCh}{\dnCh}} & & \\
\dnI{35} & \dnW: & {{\dnCh}{\dnCh}{\dnCh}{\dnCh}{\dnCh}{\dnCh}{\dnCh}{\dnCh}{\dnCh}{\dnCh}{\dnCh}{\dnCh}{\dnCh}{\dnCh}{\dnCh}{\dnCh}{\dnCh}{\dnCh}{\dnCh}{\dnCh}{\dnCh}{\dnCh}{\dnCh}{\dnCh}{\dnCh}{\dnCh}{\dnCh}{\dnCh}{\dnCh}{\dnCh}{\dnCh}{\dnCh}} & & \\
\dnI{36} & \dnW: & {{\dnCh}{\dnCh}{\dnCh}{\dnCh}{\dnCh}{\dnCh}{\dnCh}{\dnCh}{\dnCh}{\dnCh}{\dnCh}{\dnCh}{\dnCh}{\dnCh}{\dnCh}{\dnCh}{\dnCh}{\dnCh}{\dnCh}{\dnCh}{\dnCh}{\dnCh}{\dnCh}{\dnCh}{\dnCh}{\dnCh}{\dnCh}{\dnCh}{\dnCh}{\dnCh}{\dnCh}{\dnCh}} & & \\
\dnI{37} & \dnW: & {{\dnCh}{\dnCh}{\dnCh}{\dnCh}{\dnCh}{\dnCh}{\dnCh}{\dnCh}{\dnCh}{\dnCh}{\dnCh}{\dnCh}{\dnCh}{\dnCh}{\dnCh}{\dnCh}{\dnCh}{\dnCh}{\dnCh}{\dnCh}{\dnCh}{\dnCh}{\dnCh}{\dnCh}{\dnCh}{\dnCh}{\dnCh}{\dnCh}{\dnCh}{\dnCh}{\dnCh}{\dnCh}} & & \\
\dnI{38} & \dnW: & {{\dnCh}{\dnCh}{\dnCh}{\dnCh}{\dnCh}{\dnCh}{\dnCh}{\dnCh}{\dnCh}{\dnCh}{\dnCh}{\dnCh}{\dnCh}{\dnCh}{\dnCh}{\dnCh}{\dnCh}{\dnCh}{\dnCh}{\dnCh}{\dnCh}{\dnCh}{\dnCh}{\dnCh}{\dnCh}{\dnCh}{\dnCh}{\dnCh}{\dnCh}{\dnCh}{\dnCh}{\dnCh}} & & \\
\dnI{39} & \dnW: & {{\dnCh}{\dnCh}{\dnCh}{\dnCh}{\dnCh}{\dnCh}{\dnCh}{\dnCh}{\dnCh}{\dnCh}{\dnCh}{\dnCh}{\dnCh}{\dnCh}{\dnCh}{\dnCh}{\dnCh}{\dnCh}{\dnCh}{\dnCh}{\dnCh}{\dnCh}{\dnCh}{\dnCh}{\dnCh}{\dnCh}{\dnCh}{\dnCh}{\dnCh}{\dnCh}{\dnCh}{\dnCh}} & & \\
\dnI{40} & \dnW: & {{\dnCh}{\dnCh}{\dnCh}{\dnCh}{\dnCh}{\dnCh}{\dnCh}{\dnCh}{\dnCh}{\dnCh}{\dnCh}{\dnCh}{\dnCh}{\dnCh}{\dnCh}{\dnCh}{\dnCh}{\dnCh}{\dnCh}{\dnCh}{\dnCh}{\dnCh}{\dnCh}{\dnCh}{\dnCh}{\dnCh}{\dnCh}{\dnCh}{\dnCh}{\dnCh}{\dnCh}{\dnCh}} & & \\
\dnI{41} & \dnW: & {{\dnCh}{\dnCh}{\dnCh}{\dnCh}{\dnCh}{\dnCh}{\dnCh}{\dnCh}{\dnCh}{\dnCh}{\dnCh}{\dnCh}{\dnCh}{\dnCh}{\dnCh}{\dnCh}{\dnCh}{\dnCh}{\dnCh}{\dnCh}{\dnCh}{\dnCh}{\dnCh}{\dnCh}{\dnCh}{\dnCh}{\dnCh}{\dnCh}{\dnCh}{\dnCh}{\dnCh}{\dnCh}} & & \\
\dnI{42} & \dnW: & {{\dnCh}{\dnCh}{\dnCh}{\dnCh}{\dnCh}{\dnCh}{\dnCh}{\dnCh}{\dnCh}{\dnCh}{\dnCh}{\dnCh}{\dnCh}{\dnCh}{\dnCh}{\dnCh}{\dnCh}{\dnCh}{\dnCh}{\dnCh}{\dnCh}{\dnCh}{\dnCh}{\dnCh}{\dnCh}{\dnCh}{\dnCh}{\dnCh}{\dnCh}{\dnCh}{\dnCh}{\dnCh}} & & \\
\dnI{43} & \dnW: & {{\dnCh}{\dnCh}{\dnCh}{\dnCh}{\dnCh}{\dnCh}{\dnCh}{\dnCh}{\dnCh}{\dnCh}{\dnCh}{\dnCh}{\dnCh}{\dnCh}{\dnCh}{\dnCh}{\dnCh}{\dnCh}{\dnCh}{\dnCh}{\dnCh}{\dnCh}{\dnCh}{\dnCh}{\dnCh}{\dnCh}{\dnCh}{\dnCh}{\dnCh}{\dnCh}{\dnCh}{\dnCh}} & & \\
\dnI{44} & \dnW: & {{\dnCh}{\dnCh}{\dnCh}{\dnCh}{\dnCh}{\dnCh}{\dnCh}{\dnCh}{\dnCh}{\dnCh}{\dnCh}{\dnCh}{\dnCh}{\dnCh}{\dnCh}{\dnCh}{\dnCh}{\dnCh}{\dnCh}{\dnCh}{\dnCh}{\dnCh}{\dnCh}{\dnCh}{\dnCh}{\dnCh}{\dnCh}{\dnCh}{\dnCh}{\dnCh}{\dnCh}{\dnCh}} & & \\
\dnI{45} & \dnW: & {{\dnCh}{\dnCh}{\dnCh}{\dnCh}{\dnCh}{\dnCh}{\dnCh}{\dnCh}{\dnCh}{\dnCh}{\dnCh}{\dnCh}{\dnCh}{\dnCh}{\dnCh}{\dnCh}{\dnCh}{\dnCh}{\dnCh}{\dnCh}{\dnCh}{\dnCh}{\dnCh}{\dnCh}{\dnCh}{\dnCh}{\dnCh}{\dnCh}{\dnCh}{\dnCh}{\dnCh}{\dnCh}} & & \\
\dnI{46} & \dnW: & {{\dnCh}{\dnCh}{\dnCh}{\dnCh}{\dnCh}{\dnCh}{\dnCh}{\dnCh}{\dnCh}{\dnCh}{\dnCh}{\dnCh}{\dnCh}{\dnCh}{\dnCh}{\dnCh}{\dnCh}{\dnCh}{\dnCh}{\dnCh}{\dnCh}{\dnCh}{\dnCh}{\dnCh}{\dnCh}{\dnCh}{\dnCh}{\dnCh}{\dnCh}{\dnCh}{\dnCh}{\dnCh}} & & \\
\dnI{47} & \dnW: & {{\dnCh}{\dnCh}{\dnCh}{\dnCh}{\dnCh}{\dnCh}{\dnCh}{\dnCh}{\dnCh}{\dnCh}{\dnCh}{\dnCh}{\dnCh}{\dnCh}{\dnCh}{\dnCh}{\dnCh}{\dnCh}{\dnCh}{\dnCh}{\dnCh}{\dnCh}{\dnCh}{\dnCh}{\dnCh}{\dnCh}{\dnCh}{\dnCh}{\dnCh}{\dnCh}{\dnCh}{\dnCh}} & & \\
\hline
\end{tabular}
}
% internal testcase name: 48steps_m0m2m4m8m12_3.xml
\caption[MD4 testcase A]{Differential characteristic of MD4 testcase A}
\label{fig:tcA}
\end{center}
\end{figure}
}

\texttt{
\begin{figure}[!ht]
\begin{center}
\scalebox{.55}{
\begin{tabular}{|r|c|c|c|c|}
\hline
  $i$  &      &  A   &     &  W \\
\hline
\dnI{-4} & \dnW: & {{\dnCz}{\dnCo}{\dnCo}{\dnCz}{\dnCz}{\dnCo}{\dnCo}{\dnCo}{\dnCz}{\dnCo}{\dnCz}{\dnCz}{\dnCz}{\dnCo}{\dnCz}{\dnCo}{\dnCz}{\dnCz}{\dnCo}{\dnCz}{\dnCz}{\dnCz}{\dnCo}{\dnCo}{\dnCz}{\dnCz}{\dnCz}{\dnCz}{\dnCz}{\dnCz}{\dnCz}{\dnCo}} & & \\
\dnI{-3} & \dnW: & {{\dnCz}{\dnCz}{\dnCz}{\dnCo}{\dnCz}{\dnCz}{\dnCz}{\dnCz}{\dnCz}{\dnCz}{\dnCo}{\dnCo}{\dnCz}{\dnCz}{\dnCo}{\dnCz}{\dnCz}{\dnCo}{\dnCz}{\dnCo}{\dnCz}{\dnCo}{\dnCz}{\dnCz}{\dnCz}{\dnCo}{\dnCo}{\dnCo}{\dnCz}{\dnCo}{\dnCo}{\dnCz}} & & \\
\dnI{-2} & \dnW: & {{\dnCo}{\dnCz}{\dnCz}{\dnCo}{\dnCo}{\dnCz}{\dnCz}{\dnCz}{\dnCo}{\dnCz}{\dnCo}{\dnCo}{\dnCo}{\dnCz}{\dnCo}{\dnCz}{\dnCo}{\dnCo}{\dnCz}{\dnCo}{\dnCo}{\dnCo}{\dnCz}{\dnCz}{\dnCo}{\dnCo}{\dnCo}{\dnCo}{\dnCo}{\dnCo}{\dnCo}{\dnCz}} & & \\
\dnI{-1} & \dnW: & {{\dnCo}{\dnCo}{\dnCo}{\dnCz}{\dnCo}{\dnCo}{\dnCo}{\dnCo}{\dnCo}{\dnCo}{\dnCz}{\dnCz}{\dnCo}{\dnCo}{\dnCz}{\dnCo}{\dnCo}{\dnCz}{\dnCo}{\dnCz}{\dnCo}{\dnCz}{\dnCo}{\dnCo}{\dnCo}{\dnCz}{\dnCz}{\dnCz}{\dnCo}{\dnCz}{\dnCz}{\dnCo}} & & \\
\dnI{0} & \dnW: & {{\dnCq}{\dnCq}{\dnCq}{\dnCq}{\dnCq}{\dnCq}{\dnCq}{\dnCq}{\dnCq}{\dnCq}{\dnCq}{\dnCq}{\dnCq}{\dnCq}{\dnCq}{\dnCq}{\dnCq}{\dnCq}{\dnCq}{\dnCq}{\dnCq}{\dnCq}{\dnCq}{\dnCq}{\dnCq}{\dnCq}{\dnCq}{\dnCq}{\dnCq}{\dnCq}{\dnCq}{\dnCq}} & \dnW[W]{}: & {{\dnCh}{\dnCh}{\dnCh}{\dnCx}{\dnCh}{\dnCh}{\dnCh}{\dnCh}{\dnCh}{\dnCh}{\dnCh}{\dnCh}{\dnCh}{\dnCh}{\dnCh}{\dnCh}{\dnCh}{\dnCh}{\dnCh}{\dnCh}{\dnCh}{\dnCh}{\dnCh}{\dnCh}{\dnCh}{\dnCh}{\dnCh}{\dnCh}{\dnCh}{\dnCh}{\dnCh}{\dnCh}} \\
\dnI{1} & \dnW: & {{\dnCq}{\dnCq}{\dnCq}{\dnCq}{\dnCq}{\dnCq}{\dnCq}{\dnCq}{\dnCq}{\dnCq}{\dnCq}{\dnCq}{\dnCq}{\dnCq}{\dnCq}{\dnCq}{\dnCq}{\dnCq}{\dnCq}{\dnCq}{\dnCq}{\dnCq}{\dnCq}{\dnCq}{\dnCq}{\dnCq}{\dnCq}{\dnCq}{\dnCq}{\dnCq}{\dnCq}{\dnCq}} & \dnW[W]{}: & {{\dnCh}{\dnCh}{\dnCh}{\dnCh}{\dnCh}{\dnCh}{\dnCh}{\dnCh}{\dnCh}{\dnCh}{\dnCh}{\dnCh}{\dnCh}{\dnCh}{\dnCh}{\dnCh}{\dnCh}{\dnCh}{\dnCh}{\dnCh}{\dnCh}{\dnCh}{\dnCh}{\dnCh}{\dnCh}{\dnCh}{\dnCh}{\dnCh}{\dnCh}{\dnCh}{\dnCh}{\dnCh}} \\
\dnI{2} & \dnW: & {{\dnCq}{\dnCq}{\dnCq}{\dnCq}{\dnCq}{\dnCq}{\dnCq}{\dnCq}{\dnCq}{\dnCq}{\dnCq}{\dnCq}{\dnCq}{\dnCq}{\dnCq}{\dnCq}{\dnCq}{\dnCq}{\dnCq}{\dnCq}{\dnCq}{\dnCq}{\dnCq}{\dnCq}{\dnCq}{\dnCq}{\dnCq}{\dnCq}{\dnCq}{\dnCq}{\dnCq}{\dnCq}} & \dnW[W]{}: & {{\dnCx}{\dnCh}{\dnCh}{\dnCh}{\dnCh}{\dnCh}{\dnCh}{\dnCh}{\dnCh}{\dnCh}{\dnCh}{\dnCh}{\dnCh}{\dnCh}{\dnCh}{\dnCh}{\dnCh}{\dnCh}{\dnCh}{\dnCh}{\dnCh}{\dnCh}{\dnCh}{\dnCh}{\dnCh}{\dnCh}{\dnCh}{\dnCh}{\dnCh}{\dnCh}{\dnCh}{\dnCh}} \\
\dnI{3} & \dnW: & {{\dnCq}{\dnCq}{\dnCq}{\dnCq}{\dnCq}{\dnCq}{\dnCq}{\dnCq}{\dnCq}{\dnCq}{\dnCq}{\dnCq}{\dnCq}{\dnCq}{\dnCq}{\dnCq}{\dnCq}{\dnCq}{\dnCq}{\dnCq}{\dnCq}{\dnCq}{\dnCq}{\dnCq}{\dnCq}{\dnCq}{\dnCq}{\dnCq}{\dnCq}{\dnCq}{\dnCq}{\dnCq}} & \dnW[W]{}: & {{\dnCh}{\dnCh}{\dnCh}{\dnCh}{\dnCh}{\dnCh}{\dnCh}{\dnCh}{\dnCh}{\dnCh}{\dnCh}{\dnCh}{\dnCh}{\dnCh}{\dnCh}{\dnCh}{\dnCh}{\dnCh}{\dnCh}{\dnCh}{\dnCh}{\dnCh}{\dnCh}{\dnCh}{\dnCh}{\dnCh}{\dnCh}{\dnCh}{\dnCh}{\dnCh}{\dnCh}{\dnCh}} \\
\dnI{4} & \dnW: & {{\dnCq}{\dnCq}{\dnCq}{\dnCq}{\dnCq}{\dnCq}{\dnCq}{\dnCq}{\dnCq}{\dnCq}{\dnCq}{\dnCq}{\dnCq}{\dnCq}{\dnCq}{\dnCq}{\dnCq}{\dnCq}{\dnCq}{\dnCq}{\dnCq}{\dnCq}{\dnCq}{\dnCq}{\dnCq}{\dnCq}{\dnCq}{\dnCq}{\dnCq}{\dnCq}{\dnCq}{\dnCq}} & \dnW[W]{}: & {{\dnCx}{\dnCh}{\dnCh}{\dnCh}{\dnCh}{\dnCh}{\dnCh}{\dnCh}{\dnCh}{\dnCh}{\dnCh}{\dnCh}{\dnCh}{\dnCh}{\dnCh}{\dnCh}{\dnCh}{\dnCh}{\dnCh}{\dnCh}{\dnCh}{\dnCh}{\dnCh}{\dnCh}{\dnCh}{\dnCh}{\dnCh}{\dnCh}{\dnCh}{\dnCh}{\dnCh}{\dnCh}} \\
\dnI{5} & \dnW: & {{\dnCq}{\dnCq}{\dnCq}{\dnCq}{\dnCq}{\dnCq}{\dnCq}{\dnCq}{\dnCq}{\dnCq}{\dnCq}{\dnCq}{\dnCq}{\dnCq}{\dnCq}{\dnCq}{\dnCq}{\dnCq}{\dnCq}{\dnCq}{\dnCq}{\dnCq}{\dnCq}{\dnCq}{\dnCq}{\dnCq}{\dnCq}{\dnCq}{\dnCq}{\dnCq}{\dnCq}{\dnCq}} & \dnW[W]{}: & {{\dnCh}{\dnCh}{\dnCh}{\dnCh}{\dnCh}{\dnCh}{\dnCh}{\dnCh}{\dnCh}{\dnCh}{\dnCh}{\dnCh}{\dnCh}{\dnCh}{\dnCh}{\dnCh}{\dnCh}{\dnCh}{\dnCh}{\dnCh}{\dnCh}{\dnCh}{\dnCh}{\dnCh}{\dnCh}{\dnCh}{\dnCh}{\dnCh}{\dnCh}{\dnCh}{\dnCh}{\dnCh}} \\
\dnI{6} & \dnW: & {{\dnCq}{\dnCq}{\dnCq}{\dnCq}{\dnCq}{\dnCq}{\dnCq}{\dnCq}{\dnCq}{\dnCq}{\dnCq}{\dnCq}{\dnCq}{\dnCq}{\dnCq}{\dnCq}{\dnCq}{\dnCq}{\dnCq}{\dnCq}{\dnCq}{\dnCq}{\dnCq}{\dnCq}{\dnCq}{\dnCq}{\dnCq}{\dnCq}{\dnCq}{\dnCq}{\dnCq}{\dnCq}} & \dnW[W]{}: & {{\dnCh}{\dnCh}{\dnCh}{\dnCh}{\dnCh}{\dnCh}{\dnCh}{\dnCh}{\dnCh}{\dnCh}{\dnCh}{\dnCh}{\dnCh}{\dnCh}{\dnCh}{\dnCh}{\dnCh}{\dnCh}{\dnCh}{\dnCh}{\dnCh}{\dnCh}{\dnCh}{\dnCh}{\dnCh}{\dnCh}{\dnCh}{\dnCh}{\dnCh}{\dnCh}{\dnCh}{\dnCh}} \\
\dnI{7} & \dnW: & {{\dnCq}{\dnCq}{\dnCq}{\dnCq}{\dnCq}{\dnCq}{\dnCq}{\dnCq}{\dnCq}{\dnCq}{\dnCq}{\dnCq}{\dnCq}{\dnCq}{\dnCq}{\dnCq}{\dnCq}{\dnCq}{\dnCq}{\dnCq}{\dnCq}{\dnCq}{\dnCq}{\dnCq}{\dnCq}{\dnCq}{\dnCq}{\dnCq}{\dnCq}{\dnCq}{\dnCq}{\dnCq}} & \dnW[W]{}: & {{\dnCh}{\dnCh}{\dnCh}{\dnCh}{\dnCh}{\dnCh}{\dnCh}{\dnCh}{\dnCh}{\dnCh}{\dnCh}{\dnCh}{\dnCh}{\dnCh}{\dnCh}{\dnCh}{\dnCh}{\dnCh}{\dnCh}{\dnCh}{\dnCh}{\dnCh}{\dnCh}{\dnCh}{\dnCh}{\dnCh}{\dnCh}{\dnCh}{\dnCh}{\dnCh}{\dnCh}{\dnCh}} \\
\dnI{8} & \dnW: & {{\dnCq}{\dnCq}{\dnCq}{\dnCq}{\dnCq}{\dnCq}{\dnCq}{\dnCq}{\dnCq}{\dnCq}{\dnCq}{\dnCq}{\dnCq}{\dnCq}{\dnCq}{\dnCq}{\dnCq}{\dnCq}{\dnCq}{\dnCq}{\dnCq}{\dnCq}{\dnCq}{\dnCq}{\dnCq}{\dnCq}{\dnCq}{\dnCq}{\dnCq}{\dnCq}{\dnCq}{\dnCq}} & \dnW[W]{}: & {{\dnCx}{\dnCh}{\dnCh}{\dnCh}{\dnCh}{\dnCh}{\dnCh}{\dnCh}{\dnCh}{\dnCh}{\dnCh}{\dnCh}{\dnCh}{\dnCh}{\dnCh}{\dnCh}{\dnCh}{\dnCh}{\dnCh}{\dnCh}{\dnCh}{\dnCh}{\dnCh}{\dnCh}{\dnCh}{\dnCh}{\dnCh}{\dnCh}{\dnCh}{\dnCh}{\dnCh}{\dnCh}} \\
\dnI{9} & \dnW: & {{\dnCq}{\dnCq}{\dnCq}{\dnCq}{\dnCq}{\dnCq}{\dnCq}{\dnCq}{\dnCq}{\dnCq}{\dnCq}{\dnCq}{\dnCq}{\dnCq}{\dnCq}{\dnCq}{\dnCq}{\dnCq}{\dnCq}{\dnCq}{\dnCq}{\dnCq}{\dnCq}{\dnCq}{\dnCq}{\dnCq}{\dnCq}{\dnCq}{\dnCq}{\dnCq}{\dnCq}{\dnCq}} & \dnW[W]{}: & {{\dnCh}{\dnCh}{\dnCh}{\dnCh}{\dnCh}{\dnCh}{\dnCh}{\dnCh}{\dnCh}{\dnCh}{\dnCh}{\dnCh}{\dnCh}{\dnCh}{\dnCh}{\dnCh}{\dnCh}{\dnCh}{\dnCh}{\dnCh}{\dnCh}{\dnCh}{\dnCh}{\dnCh}{\dnCh}{\dnCh}{\dnCh}{\dnCh}{\dnCh}{\dnCh}{\dnCh}{\dnCh}} \\
\dnI{10} & \dnW: & {{\dnCq}{\dnCq}{\dnCq}{\dnCq}{\dnCq}{\dnCq}{\dnCq}{\dnCq}{\dnCq}{\dnCq}{\dnCq}{\dnCq}{\dnCq}{\dnCq}{\dnCq}{\dnCq}{\dnCq}{\dnCq}{\dnCq}{\dnCq}{\dnCq}{\dnCq}{\dnCq}{\dnCq}{\dnCq}{\dnCq}{\dnCq}{\dnCq}{\dnCq}{\dnCq}{\dnCq}{\dnCq}} & \dnW[W]{}: & {{\dnCh}{\dnCh}{\dnCh}{\dnCh}{\dnCh}{\dnCh}{\dnCh}{\dnCh}{\dnCh}{\dnCh}{\dnCh}{\dnCh}{\dnCh}{\dnCh}{\dnCh}{\dnCh}{\dnCh}{\dnCh}{\dnCh}{\dnCh}{\dnCh}{\dnCh}{\dnCh}{\dnCh}{\dnCh}{\dnCh}{\dnCh}{\dnCh}{\dnCh}{\dnCh}{\dnCh}{\dnCh}} \\
\dnI{11} & \dnW: & {{\dnCq}{\dnCq}{\dnCq}{\dnCq}{\dnCq}{\dnCq}{\dnCq}{\dnCq}{\dnCq}{\dnCq}{\dnCq}{\dnCq}{\dnCq}{\dnCq}{\dnCq}{\dnCq}{\dnCq}{\dnCq}{\dnCq}{\dnCq}{\dnCq}{\dnCq}{\dnCq}{\dnCq}{\dnCq}{\dnCq}{\dnCq}{\dnCq}{\dnCq}{\dnCq}{\dnCq}{\dnCq}} & \dnW[W]{}: & {{\dnCh}{\dnCh}{\dnCh}{\dnCh}{\dnCh}{\dnCh}{\dnCh}{\dnCh}{\dnCh}{\dnCh}{\dnCh}{\dnCh}{\dnCh}{\dnCh}{\dnCh}{\dnCh}{\dnCh}{\dnCh}{\dnCh}{\dnCh}{\dnCh}{\dnCh}{\dnCh}{\dnCh}{\dnCh}{\dnCh}{\dnCh}{\dnCh}{\dnCh}{\dnCh}{\dnCh}{\dnCh}} \\
\dnI{12} & \dnW: & {{\dnCq}{\dnCq}{\dnCq}{\dnCq}{\dnCq}{\dnCq}{\dnCq}{\dnCq}{\dnCq}{\dnCq}{\dnCq}{\dnCq}{\dnCq}{\dnCq}{\dnCh}{\dnCh}{\dnCh}{\dnCh}{\dnCh}{\dnCh}{\dnCh}{\dnCh}{\dnCh}{\dnCh}{\dnCh}{\dnCh}{\dnCh}{\dnCh}{\dnCh}{\dnCh}{\dnCh}{\dnCh}} & \dnW[W]{}: & {{\dnCx}{\dnCh}{\dnCh}{\dnCh}{\dnCh}{\dnCh}{\dnCh}{\dnCh}{\dnCh}{\dnCh}{\dnCh}{\dnCh}{\dnCh}{\dnCh}{\dnCh}{\dnCh}{\dnCh}{\dnCh}{\dnCh}{\dnCh}{\dnCh}{\dnCh}{\dnCh}{\dnCh}{\dnCh}{\dnCh}{\dnCh}{\dnCh}{\dnCh}{\dnCh}{\dnCh}{\dnCh}} \\
\dnI{13} & \dnW: & {{\dnCq}{\dnCq}{\dnCq}{\dnCq}{\dnCq}{\dnCq}{\dnCq}{\dnCq}{\dnCq}{\dnCq}{\dnCq}{\dnCq}{\dnCq}{\dnCq}{\dnCh}{\dnCh}{\dnCh}{\dnCh}{\dnCh}{\dnCh}{\dnCh}{\dnCh}{\dnCh}{\dnCh}{\dnCh}{\dnCh}{\dnCh}{\dnCh}{\dnCh}{\dnCh}{\dnCh}{\dnCh}} & \dnW[W]{}: & {{\dnCh}{\dnCh}{\dnCh}{\dnCh}{\dnCh}{\dnCh}{\dnCh}{\dnCh}{\dnCh}{\dnCh}{\dnCh}{\dnCh}{\dnCh}{\dnCh}{\dnCh}{\dnCh}{\dnCh}{\dnCh}{\dnCh}{\dnCh}{\dnCh}{\dnCh}{\dnCh}{\dnCh}{\dnCh}{\dnCh}{\dnCh}{\dnCh}{\dnCh}{\dnCh}{\dnCh}{\dnCh}} \\
\dnI{14} & \dnW: & {{\dnCq}{\dnCq}{\dnCq}{\dnCq}{\dnCq}{\dnCq}{\dnCq}{\dnCq}{\dnCq}{\dnCq}{\dnCq}{\dnCq}{\dnCq}{\dnCq}{\dnCh}{\dnCh}{\dnCh}{\dnCh}{\dnCh}{\dnCh}{\dnCh}{\dnCh}{\dnCh}{\dnCh}{\dnCh}{\dnCh}{\dnCh}{\dnCh}{\dnCh}{\dnCh}{\dnCh}{\dnCh}} & \dnW[W]{}: & {{\dnCh}{\dnCh}{\dnCh}{\dnCh}{\dnCh}{\dnCh}{\dnCh}{\dnCh}{\dnCh}{\dnCh}{\dnCh}{\dnCh}{\dnCh}{\dnCh}{\dnCh}{\dnCh}{\dnCh}{\dnCh}{\dnCh}{\dnCh}{\dnCh}{\dnCh}{\dnCh}{\dnCh}{\dnCh}{\dnCh}{\dnCh}{\dnCh}{\dnCh}{\dnCh}{\dnCh}{\dnCh}} \\
\dnI{15} & \dnW: & {{\dnCq}{\dnCq}{\dnCq}{\dnCq}{\dnCq}{\dnCq}{\dnCq}{\dnCq}{\dnCq}{\dnCq}{\dnCq}{\dnCq}{\dnCq}{\dnCq}{\dnCh}{\dnCh}{\dnCh}{\dnCh}{\dnCh}{\dnCh}{\dnCh}{\dnCh}{\dnCh}{\dnCh}{\dnCh}{\dnCh}{\dnCh}{\dnCh}{\dnCh}{\dnCh}{\dnCh}{\dnCh}} & \dnW[W]{}: & {{\dnCh}{\dnCh}{\dnCh}{\dnCh}{\dnCh}{\dnCh}{\dnCh}{\dnCh}{\dnCh}{\dnCh}{\dnCh}{\dnCh}{\dnCh}{\dnCh}{\dnCh}{\dnCh}{\dnCh}{\dnCh}{\dnCh}{\dnCh}{\dnCh}{\dnCh}{\dnCh}{\dnCh}{\dnCh}{\dnCh}{\dnCh}{\dnCh}{\dnCh}{\dnCh}{\dnCh}{\dnCh}} \\
\dnI{16} & \dnW: & {{\dnCq}{\dnCq}{\dnCq}{\dnCx}{\dnCh}{\dnCh}{\dnCh}{\dnCh}{\dnCh}{\dnCh}{\dnCh}{\dnCh}{\dnCh}{\dnCh}{\dnCh}{\dnCh}{\dnCh}{\dnCh}{\dnCh}{\dnCh}{\dnCh}{\dnCh}{\dnCh}{\dnCh}{\dnCh}{\dnCh}{\dnCh}{\dnCh}{\dnCh}{\dnCh}{\dnCh}{\dnCh}} & & \\
\dnI{17} & \dnW: & {{\dnCq}{\dnCh}{\dnCh}{\dnCh}{\dnCh}{\dnCh}{\dnCh}{\dnCh}{\dnCh}{\dnCh}{\dnCh}{\dnCh}{\dnCh}{\dnCh}{\dnCh}{\dnCh}{\dnCh}{\dnCh}{\dnCh}{\dnCh}{\dnCh}{\dnCh}{\dnCh}{\dnCh}{\dnCh}{\dnCh}{\dnCh}{\dnCh}{\dnCh}{\dnCh}{\dnCh}{\dnCh}} & & \\
\dnI{18} & \dnW: & {{\dnCq}{\dnCh}{\dnCh}{\dnCh}{\dnCh}{\dnCh}{\dnCh}{\dnCh}{\dnCh}{\dnCh}{\dnCh}{\dnCh}{\dnCh}{\dnCh}{\dnCh}{\dnCh}{\dnCh}{\dnCh}{\dnCh}{\dnCh}{\dnCh}{\dnCh}{\dnCh}{\dnCh}{\dnCh}{\dnCh}{\dnCh}{\dnCh}{\dnCh}{\dnCh}{\dnCh}{\dnCh}} & & \\
\dnI{19} & \dnW: & {{\dnCq}{\dnCh}{\dnCh}{\dnCh}{\dnCh}{\dnCh}{\dnCh}{\dnCh}{\dnCh}{\dnCh}{\dnCh}{\dnCh}{\dnCh}{\dnCh}{\dnCh}{\dnCh}{\dnCh}{\dnCh}{\dnCh}{\dnCh}{\dnCh}{\dnCh}{\dnCh}{\dnCh}{\dnCh}{\dnCh}{\dnCh}{\dnCh}{\dnCh}{\dnCh}{\dnCh}{\dnCh}} & & \\
\dnI{20} & \dnW: & {{\dnCx}{\dnCh}{\dnCh}{\dnCh}{\dnCh}{\dnCh}{\dnCh}{\dnCh}{\dnCh}{\dnCh}{\dnCh}{\dnCh}{\dnCh}{\dnCh}{\dnCh}{\dnCh}{\dnCh}{\dnCh}{\dnCh}{\dnCh}{\dnCh}{\dnCh}{\dnCh}{\dnCh}{\dnCh}{\dnCh}{\dnCh}{\dnCh}{\dnCh}{\dnCh}{\dnCh}{\dnCh}} & & \\
\dnI{21} & \dnW: & {{\dnCh}{\dnCh}{\dnCh}{\dnCh}{\dnCh}{\dnCh}{\dnCh}{\dnCh}{\dnCh}{\dnCh}{\dnCh}{\dnCh}{\dnCh}{\dnCh}{\dnCh}{\dnCh}{\dnCh}{\dnCh}{\dnCh}{\dnCh}{\dnCh}{\dnCh}{\dnCh}{\dnCh}{\dnCh}{\dnCh}{\dnCh}{\dnCh}{\dnCh}{\dnCh}{\dnCh}{\dnCh}} & & \\
\dnI{22} & \dnW: & {{\dnCh}{\dnCh}{\dnCh}{\dnCh}{\dnCh}{\dnCh}{\dnCh}{\dnCh}{\dnCh}{\dnCh}{\dnCh}{\dnCh}{\dnCh}{\dnCh}{\dnCh}{\dnCh}{\dnCh}{\dnCh}{\dnCh}{\dnCh}{\dnCh}{\dnCh}{\dnCh}{\dnCh}{\dnCh}{\dnCh}{\dnCh}{\dnCh}{\dnCh}{\dnCh}{\dnCh}{\dnCh}} & & \\
\dnI{23} & \dnW: & {{\dnCh}{\dnCh}{\dnCh}{\dnCh}{\dnCh}{\dnCh}{\dnCh}{\dnCh}{\dnCh}{\dnCh}{\dnCh}{\dnCh}{\dnCh}{\dnCh}{\dnCh}{\dnCh}{\dnCh}{\dnCh}{\dnCh}{\dnCh}{\dnCh}{\dnCh}{\dnCh}{\dnCh}{\dnCh}{\dnCh}{\dnCh}{\dnCh}{\dnCh}{\dnCh}{\dnCh}{\dnCh}} & & \\
\dnI{24} & \dnW: & {{\dnCh}{\dnCh}{\dnCh}{\dnCh}{\dnCh}{\dnCh}{\dnCh}{\dnCh}{\dnCh}{\dnCh}{\dnCh}{\dnCh}{\dnCh}{\dnCh}{\dnCh}{\dnCh}{\dnCh}{\dnCh}{\dnCh}{\dnCh}{\dnCh}{\dnCh}{\dnCh}{\dnCh}{\dnCh}{\dnCh}{\dnCh}{\dnCh}{\dnCh}{\dnCh}{\dnCh}{\dnCh}} & & \\
\dnI{25} & \dnW: & {{\dnCh}{\dnCh}{\dnCh}{\dnCh}{\dnCh}{\dnCh}{\dnCh}{\dnCh}{\dnCh}{\dnCh}{\dnCh}{\dnCh}{\dnCh}{\dnCh}{\dnCh}{\dnCh}{\dnCh}{\dnCh}{\dnCh}{\dnCh}{\dnCh}{\dnCh}{\dnCh}{\dnCh}{\dnCh}{\dnCh}{\dnCh}{\dnCh}{\dnCh}{\dnCh}{\dnCh}{\dnCh}} & & \\
\dnI{26} & \dnW: & {{\dnCh}{\dnCh}{\dnCh}{\dnCh}{\dnCh}{\dnCh}{\dnCh}{\dnCh}{\dnCh}{\dnCh}{\dnCh}{\dnCh}{\dnCh}{\dnCh}{\dnCh}{\dnCh}{\dnCh}{\dnCh}{\dnCh}{\dnCh}{\dnCh}{\dnCh}{\dnCh}{\dnCh}{\dnCh}{\dnCh}{\dnCh}{\dnCh}{\dnCh}{\dnCh}{\dnCh}{\dnCh}} & & \\
\dnI{27} & \dnW: & {{\dnCh}{\dnCh}{\dnCh}{\dnCh}{\dnCh}{\dnCh}{\dnCh}{\dnCh}{\dnCh}{\dnCh}{\dnCh}{\dnCh}{\dnCh}{\dnCh}{\dnCh}{\dnCh}{\dnCh}{\dnCh}{\dnCh}{\dnCh}{\dnCh}{\dnCh}{\dnCh}{\dnCh}{\dnCh}{\dnCh}{\dnCh}{\dnCh}{\dnCh}{\dnCh}{\dnCh}{\dnCh}} & & \\
\dnI{28} & \dnW: & {{\dnCh}{\dnCh}{\dnCh}{\dnCh}{\dnCh}{\dnCh}{\dnCh}{\dnCh}{\dnCh}{\dnCh}{\dnCh}{\dnCh}{\dnCh}{\dnCh}{\dnCh}{\dnCh}{\dnCh}{\dnCh}{\dnCh}{\dnCh}{\dnCh}{\dnCh}{\dnCh}{\dnCh}{\dnCh}{\dnCh}{\dnCh}{\dnCh}{\dnCh}{\dnCh}{\dnCh}{\dnCh}} & & \\
\dnI{29} & \dnW: & {{\dnCh}{\dnCh}{\dnCh}{\dnCh}{\dnCh}{\dnCh}{\dnCh}{\dnCh}{\dnCh}{\dnCh}{\dnCh}{\dnCh}{\dnCh}{\dnCh}{\dnCh}{\dnCh}{\dnCh}{\dnCh}{\dnCh}{\dnCh}{\dnCh}{\dnCh}{\dnCh}{\dnCh}{\dnCh}{\dnCh}{\dnCh}{\dnCh}{\dnCh}{\dnCh}{\dnCh}{\dnCh}} & & \\
\dnI{30} & \dnW: & {{\dnCh}{\dnCh}{\dnCh}{\dnCh}{\dnCh}{\dnCh}{\dnCh}{\dnCh}{\dnCh}{\dnCh}{\dnCh}{\dnCh}{\dnCh}{\dnCh}{\dnCh}{\dnCh}{\dnCh}{\dnCh}{\dnCh}{\dnCh}{\dnCh}{\dnCh}{\dnCh}{\dnCh}{\dnCh}{\dnCh}{\dnCh}{\dnCh}{\dnCh}{\dnCh}{\dnCh}{\dnCh}} & & \\
\dnI{31} & \dnW: & {{\dnCh}{\dnCh}{\dnCh}{\dnCh}{\dnCh}{\dnCh}{\dnCh}{\dnCh}{\dnCh}{\dnCh}{\dnCh}{\dnCh}{\dnCh}{\dnCh}{\dnCh}{\dnCh}{\dnCh}{\dnCh}{\dnCh}{\dnCh}{\dnCh}{\dnCh}{\dnCh}{\dnCh}{\dnCh}{\dnCh}{\dnCh}{\dnCh}{\dnCh}{\dnCh}{\dnCh}{\dnCh}} & & \\
\dnI{32} & \dnW: & {{\dnCx}{\dnCh}{\dnCh}{\dnCh}{\dnCh}{\dnCh}{\dnCh}{\dnCh}{\dnCh}{\dnCh}{\dnCh}{\dnCh}{\dnCh}{\dnCh}{\dnCh}{\dnCh}{\dnCh}{\dnCh}{\dnCh}{\dnCh}{\dnCh}{\dnCh}{\dnCh}{\dnCh}{\dnCh}{\dnCh}{\dnCh}{\dnCh}{\dnCh}{\dnCh}{\dnCh}{\dnCh}} & & \\
\dnI{33} & \dnW: & {{\dnCh}{\dnCh}{\dnCh}{\dnCh}{\dnCh}{\dnCh}{\dnCh}{\dnCh}{\dnCh}{\dnCh}{\dnCh}{\dnCh}{\dnCh}{\dnCh}{\dnCh}{\dnCh}{\dnCh}{\dnCh}{\dnCh}{\dnCh}{\dnCh}{\dnCh}{\dnCh}{\dnCh}{\dnCh}{\dnCh}{\dnCh}{\dnCh}{\dnCh}{\dnCh}{\dnCh}{\dnCh}} & & \\
\dnI{34} & \dnW: & {{\dnCh}{\dnCh}{\dnCh}{\dnCh}{\dnCh}{\dnCh}{\dnCh}{\dnCh}{\dnCh}{\dnCh}{\dnCh}{\dnCh}{\dnCh}{\dnCh}{\dnCh}{\dnCh}{\dnCh}{\dnCh}{\dnCh}{\dnCh}{\dnCh}{\dnCh}{\dnCh}{\dnCh}{\dnCh}{\dnCh}{\dnCh}{\dnCh}{\dnCh}{\dnCh}{\dnCh}{\dnCh}} & & \\
\dnI{35} & \dnW: & {{\dnCh}{\dnCh}{\dnCh}{\dnCh}{\dnCh}{\dnCh}{\dnCh}{\dnCh}{\dnCh}{\dnCh}{\dnCh}{\dnCh}{\dnCh}{\dnCh}{\dnCh}{\dnCh}{\dnCh}{\dnCh}{\dnCh}{\dnCh}{\dnCh}{\dnCh}{\dnCh}{\dnCh}{\dnCh}{\dnCh}{\dnCh}{\dnCh}{\dnCh}{\dnCh}{\dnCh}{\dnCh}} & & \\
\dnI{36} & \dnW: & {{\dnCh}{\dnCh}{\dnCh}{\dnCh}{\dnCh}{\dnCh}{\dnCh}{\dnCh}{\dnCh}{\dnCh}{\dnCh}{\dnCh}{\dnCh}{\dnCh}{\dnCh}{\dnCh}{\dnCh}{\dnCh}{\dnCh}{\dnCh}{\dnCh}{\dnCh}{\dnCh}{\dnCh}{\dnCh}{\dnCh}{\dnCh}{\dnCh}{\dnCh}{\dnCh}{\dnCh}{\dnCh}} & & \\
\dnI{37} & \dnW: & {{\dnCh}{\dnCh}{\dnCh}{\dnCh}{\dnCh}{\dnCh}{\dnCh}{\dnCh}{\dnCh}{\dnCh}{\dnCh}{\dnCh}{\dnCh}{\dnCh}{\dnCh}{\dnCh}{\dnCh}{\dnCh}{\dnCh}{\dnCh}{\dnCh}{\dnCh}{\dnCh}{\dnCh}{\dnCh}{\dnCh}{\dnCh}{\dnCh}{\dnCh}{\dnCh}{\dnCh}{\dnCh}} & & \\
\dnI{38} & \dnW: & {{\dnCh}{\dnCh}{\dnCh}{\dnCh}{\dnCh}{\dnCh}{\dnCh}{\dnCh}{\dnCh}{\dnCh}{\dnCh}{\dnCh}{\dnCh}{\dnCh}{\dnCh}{\dnCh}{\dnCh}{\dnCh}{\dnCh}{\dnCh}{\dnCh}{\dnCh}{\dnCh}{\dnCh}{\dnCh}{\dnCh}{\dnCh}{\dnCh}{\dnCh}{\dnCh}{\dnCh}{\dnCh}} & & \\
\dnI{39} & \dnW: & {{\dnCh}{\dnCh}{\dnCh}{\dnCh}{\dnCh}{\dnCh}{\dnCh}{\dnCh}{\dnCh}{\dnCh}{\dnCh}{\dnCh}{\dnCh}{\dnCh}{\dnCh}{\dnCh}{\dnCh}{\dnCh}{\dnCh}{\dnCh}{\dnCh}{\dnCh}{\dnCh}{\dnCh}{\dnCh}{\dnCh}{\dnCh}{\dnCh}{\dnCh}{\dnCh}{\dnCh}{\dnCh}} & & \\
\dnI{40} & \dnW: & {{\dnCh}{\dnCh}{\dnCh}{\dnCh}{\dnCh}{\dnCh}{\dnCh}{\dnCh}{\dnCh}{\dnCh}{\dnCh}{\dnCh}{\dnCh}{\dnCh}{\dnCh}{\dnCh}{\dnCh}{\dnCh}{\dnCh}{\dnCh}{\dnCh}{\dnCh}{\dnCh}{\dnCh}{\dnCh}{\dnCh}{\dnCh}{\dnCh}{\dnCh}{\dnCh}{\dnCh}{\dnCh}} & & \\
\dnI{41} & \dnW: & {{\dnCh}{\dnCh}{\dnCh}{\dnCh}{\dnCh}{\dnCh}{\dnCh}{\dnCh}{\dnCh}{\dnCh}{\dnCh}{\dnCh}{\dnCh}{\dnCh}{\dnCh}{\dnCh}{\dnCh}{\dnCh}{\dnCh}{\dnCh}{\dnCh}{\dnCh}{\dnCh}{\dnCh}{\dnCh}{\dnCh}{\dnCh}{\dnCh}{\dnCh}{\dnCh}{\dnCh}{\dnCh}} & & \\
\dnI{42} & \dnW: & {{\dnCh}{\dnCh}{\dnCh}{\dnCh}{\dnCh}{\dnCh}{\dnCh}{\dnCh}{\dnCh}{\dnCh}{\dnCh}{\dnCh}{\dnCh}{\dnCh}{\dnCh}{\dnCh}{\dnCh}{\dnCh}{\dnCh}{\dnCh}{\dnCh}{\dnCh}{\dnCh}{\dnCh}{\dnCh}{\dnCh}{\dnCh}{\dnCh}{\dnCh}{\dnCh}{\dnCh}{\dnCh}} & & \\
\dnI{43} & \dnW: & {{\dnCh}{\dnCh}{\dnCh}{\dnCh}{\dnCh}{\dnCh}{\dnCh}{\dnCh}{\dnCh}{\dnCh}{\dnCh}{\dnCh}{\dnCh}{\dnCh}{\dnCh}{\dnCh}{\dnCh}{\dnCh}{\dnCh}{\dnCh}{\dnCh}{\dnCh}{\dnCh}{\dnCh}{\dnCh}{\dnCh}{\dnCh}{\dnCh}{\dnCh}{\dnCh}{\dnCh}{\dnCh}} & & \\
\dnI{44} & \dnW: & {{\dnCh}{\dnCh}{\dnCh}{\dnCh}{\dnCh}{\dnCh}{\dnCh}{\dnCh}{\dnCh}{\dnCh}{\dnCh}{\dnCh}{\dnCh}{\dnCh}{\dnCh}{\dnCh}{\dnCh}{\dnCh}{\dnCh}{\dnCh}{\dnCh}{\dnCh}{\dnCh}{\dnCh}{\dnCh}{\dnCh}{\dnCh}{\dnCh}{\dnCh}{\dnCh}{\dnCh}{\dnCh}} & & \\
\dnI{45} & \dnW: & {{\dnCh}{\dnCh}{\dnCh}{\dnCh}{\dnCh}{\dnCh}{\dnCh}{\dnCh}{\dnCh}{\dnCh}{\dnCh}{\dnCh}{\dnCh}{\dnCh}{\dnCh}{\dnCh}{\dnCh}{\dnCh}{\dnCh}{\dnCh}{\dnCh}{\dnCh}{\dnCh}{\dnCh}{\dnCh}{\dnCh}{\dnCh}{\dnCh}{\dnCh}{\dnCh}{\dnCh}{\dnCh}} & & \\
\dnI{46} & \dnW: & {{\dnCh}{\dnCh}{\dnCh}{\dnCh}{\dnCh}{\dnCh}{\dnCh}{\dnCh}{\dnCh}{\dnCh}{\dnCh}{\dnCh}{\dnCh}{\dnCh}{\dnCh}{\dnCh}{\dnCh}{\dnCh}{\dnCh}{\dnCh}{\dnCh}{\dnCh}{\dnCh}{\dnCh}{\dnCh}{\dnCh}{\dnCh}{\dnCh}{\dnCh}{\dnCh}{\dnCh}{\dnCh}} & & \\
\dnI{47} & \dnW: & {{\dnCh}{\dnCh}{\dnCh}{\dnCh}{\dnCh}{\dnCh}{\dnCh}{\dnCh}{\dnCh}{\dnCh}{\dnCh}{\dnCh}{\dnCh}{\dnCh}{\dnCh}{\dnCh}{\dnCh}{\dnCh}{\dnCh}{\dnCh}{\dnCh}{\dnCh}{\dnCh}{\dnCh}{\dnCh}{\dnCh}{\dnCh}{\dnCh}{\dnCh}{\dnCh}{\dnCh}{\dnCh}} & & \\
\hline
\end{tabular}
}
% internal testcase name: 48steps_m0m2m4m8m12_2.xml
\caption[MD4 B testcase]{Differential characteristic of MD4 testcase B}
\label{fig:tcB}
\end{center}
\end{figure}
}

\texttt{
\begin{figure}[!ht]
\begin{center}
\scalebox{.55}{
\begin{tabular}{|r|c|c|c|c|}
\hline
  $i$  &      &  A  &    &  W \\
\hline
\dnI{-4} & \dnW: & {{\dnCz}{\dnCo}{\dnCo}{\dnCz}{\dnCz}{\dnCo}{\dnCo}{\dnCo}{\dnCz}{\dnCo}{\dnCz}{\dnCz}{\dnCz}{\dnCo}{\dnCz}{\dnCo}{\dnCz}{\dnCz}{\dnCo}{\dnCz}{\dnCz}{\dnCz}{\dnCo}{\dnCo}{\dnCz}{\dnCz}{\dnCz}{\dnCz}{\dnCz}{\dnCz}{\dnCz}{\dnCo}} & & \\
\dnI{-3} & \dnW: & {{\dnCz}{\dnCz}{\dnCz}{\dnCo}{\dnCz}{\dnCz}{\dnCz}{\dnCz}{\dnCz}{\dnCz}{\dnCo}{\dnCo}{\dnCz}{\dnCz}{\dnCo}{\dnCz}{\dnCz}{\dnCo}{\dnCz}{\dnCo}{\dnCz}{\dnCo}{\dnCz}{\dnCz}{\dnCz}{\dnCo}{\dnCo}{\dnCo}{\dnCz}{\dnCo}{\dnCo}{\dnCz}} & & \\
\dnI{-2} & \dnW: & {{\dnCo}{\dnCz}{\dnCz}{\dnCo}{\dnCo}{\dnCz}{\dnCz}{\dnCz}{\dnCo}{\dnCz}{\dnCo}{\dnCo}{\dnCo}{\dnCz}{\dnCo}{\dnCz}{\dnCo}{\dnCo}{\dnCz}{\dnCo}{\dnCo}{\dnCo}{\dnCz}{\dnCz}{\dnCo}{\dnCo}{\dnCo}{\dnCo}{\dnCo}{\dnCo}{\dnCo}{\dnCz}} & & \\
\dnI{-1} & \dnW: & {{\dnCo}{\dnCo}{\dnCo}{\dnCz}{\dnCo}{\dnCo}{\dnCo}{\dnCo}{\dnCo}{\dnCo}{\dnCz}{\dnCz}{\dnCo}{\dnCo}{\dnCz}{\dnCo}{\dnCo}{\dnCz}{\dnCo}{\dnCz}{\dnCo}{\dnCz}{\dnCo}{\dnCo}{\dnCo}{\dnCz}{\dnCz}{\dnCz}{\dnCo}{\dnCz}{\dnCz}{\dnCo}} & & \\
\dnI{0} & \dnW: & {{\dnCq}{\dnCq}{\dnCq}{\dnCq}{\dnCq}{\dnCq}{\dnCq}{\dnCq}{\dnCq}{\dnCq}{\dnCq}{\dnCq}{\dnCq}{\dnCq}{\dnCq}{\dnCq}{\dnCq}{\dnCq}{\dnCq}{\dnCq}{\dnCq}{\dnCq}{\dnCq}{\dnCq}{\dnCq}{\dnCq}{\dnCq}{\dnCq}{\dnCq}{\dnCq}{\dnCq}{\dnCq}} & \dnW[W]{}: & {{\dnCh}{\dnCh}{\dnCh}{\dnCx}{\dnCh}{\dnCh}{\dnCh}{\dnCh}{\dnCh}{\dnCh}{\dnCh}{\dnCh}{\dnCh}{\dnCh}{\dnCh}{\dnCh}{\dnCh}{\dnCh}{\dnCh}{\dnCh}{\dnCh}{\dnCh}{\dnCh}{\dnCh}{\dnCh}{\dnCh}{\dnCh}{\dnCh}{\dnCh}{\dnCh}{\dnCh}{\dnCh}} \\
\dnI{1} & \dnW: & {{\dnCq}{\dnCq}{\dnCq}{\dnCq}{\dnCq}{\dnCq}{\dnCq}{\dnCq}{\dnCq}{\dnCq}{\dnCq}{\dnCq}{\dnCq}{\dnCq}{\dnCq}{\dnCq}{\dnCq}{\dnCq}{\dnCq}{\dnCq}{\dnCq}{\dnCq}{\dnCq}{\dnCq}{\dnCq}{\dnCq}{\dnCq}{\dnCq}{\dnCq}{\dnCq}{\dnCq}{\dnCq}} & \dnW[W]{}: & {{\dnCh}{\dnCh}{\dnCh}{\dnCh}{\dnCh}{\dnCh}{\dnCh}{\dnCh}{\dnCh}{\dnCh}{\dnCh}{\dnCh}{\dnCh}{\dnCh}{\dnCh}{\dnCh}{\dnCh}{\dnCh}{\dnCh}{\dnCh}{\dnCh}{\dnCh}{\dnCh}{\dnCh}{\dnCh}{\dnCh}{\dnCh}{\dnCh}{\dnCh}{\dnCh}{\dnCh}{\dnCh}} \\
\dnI{2} & \dnW: & {{\dnCq}{\dnCq}{\dnCq}{\dnCq}{\dnCq}{\dnCq}{\dnCq}{\dnCq}{\dnCq}{\dnCq}{\dnCq}{\dnCq}{\dnCq}{\dnCq}{\dnCq}{\dnCq}{\dnCq}{\dnCq}{\dnCq}{\dnCq}{\dnCq}{\dnCq}{\dnCq}{\dnCq}{\dnCq}{\dnCq}{\dnCq}{\dnCq}{\dnCq}{\dnCq}{\dnCq}{\dnCq}} & \dnW[W]{}: & {{\dnCx}{\dnCh}{\dnCh}{\dnCh}{\dnCh}{\dnCh}{\dnCh}{\dnCh}{\dnCh}{\dnCh}{\dnCh}{\dnCh}{\dnCh}{\dnCh}{\dnCh}{\dnCh}{\dnCh}{\dnCh}{\dnCh}{\dnCh}{\dnCh}{\dnCh}{\dnCh}{\dnCh}{\dnCh}{\dnCh}{\dnCh}{\dnCh}{\dnCh}{\dnCh}{\dnCh}{\dnCh}} \\
\dnI{3} & \dnW: & {{\dnCq}{\dnCq}{\dnCq}{\dnCq}{\dnCq}{\dnCq}{\dnCq}{\dnCq}{\dnCq}{\dnCq}{\dnCq}{\dnCq}{\dnCq}{\dnCq}{\dnCq}{\dnCq}{\dnCq}{\dnCq}{\dnCq}{\dnCq}{\dnCq}{\dnCq}{\dnCq}{\dnCq}{\dnCq}{\dnCq}{\dnCq}{\dnCq}{\dnCq}{\dnCq}{\dnCq}{\dnCq}} & \dnW[W]{}: & {{\dnCh}{\dnCh}{\dnCh}{\dnCh}{\dnCh}{\dnCh}{\dnCh}{\dnCh}{\dnCh}{\dnCh}{\dnCh}{\dnCh}{\dnCh}{\dnCh}{\dnCh}{\dnCh}{\dnCh}{\dnCh}{\dnCh}{\dnCh}{\dnCh}{\dnCh}{\dnCh}{\dnCh}{\dnCh}{\dnCh}{\dnCh}{\dnCh}{\dnCh}{\dnCh}{\dnCh}{\dnCh}} \\
\dnI{4} & \dnW: & {{\dnCq}{\dnCq}{\dnCq}{\dnCq}{\dnCq}{\dnCq}{\dnCq}{\dnCq}{\dnCq}{\dnCq}{\dnCq}{\dnCq}{\dnCq}{\dnCq}{\dnCq}{\dnCq}{\dnCq}{\dnCq}{\dnCq}{\dnCq}{\dnCq}{\dnCq}{\dnCq}{\dnCq}{\dnCq}{\dnCq}{\dnCq}{\dnCq}{\dnCq}{\dnCq}{\dnCq}{\dnCq}} & \dnW[W]{}: & {{\dnCx}{\dnCh}{\dnCh}{\dnCh}{\dnCh}{\dnCh}{\dnCh}{\dnCh}{\dnCh}{\dnCh}{\dnCh}{\dnCh}{\dnCh}{\dnCh}{\dnCh}{\dnCh}{\dnCh}{\dnCh}{\dnCh}{\dnCh}{\dnCh}{\dnCh}{\dnCh}{\dnCh}{\dnCh}{\dnCh}{\dnCh}{\dnCh}{\dnCh}{\dnCh}{\dnCh}{\dnCh}} \\
\dnI{5} & \dnW: & {{\dnCq}{\dnCq}{\dnCq}{\dnCq}{\dnCq}{\dnCq}{\dnCq}{\dnCq}{\dnCq}{\dnCq}{\dnCq}{\dnCq}{\dnCq}{\dnCq}{\dnCq}{\dnCq}{\dnCq}{\dnCq}{\dnCq}{\dnCq}{\dnCq}{\dnCq}{\dnCq}{\dnCq}{\dnCq}{\dnCq}{\dnCq}{\dnCq}{\dnCq}{\dnCq}{\dnCq}{\dnCq}} & \dnW[W]{}: & {{\dnCh}{\dnCh}{\dnCh}{\dnCh}{\dnCh}{\dnCh}{\dnCh}{\dnCh}{\dnCh}{\dnCh}{\dnCh}{\dnCh}{\dnCh}{\dnCh}{\dnCh}{\dnCh}{\dnCh}{\dnCh}{\dnCh}{\dnCh}{\dnCh}{\dnCh}{\dnCh}{\dnCh}{\dnCh}{\dnCh}{\dnCh}{\dnCh}{\dnCh}{\dnCh}{\dnCh}{\dnCh}} \\
\dnI{6} & \dnW: & {{\dnCq}{\dnCq}{\dnCq}{\dnCq}{\dnCq}{\dnCq}{\dnCq}{\dnCq}{\dnCq}{\dnCq}{\dnCq}{\dnCq}{\dnCq}{\dnCq}{\dnCq}{\dnCq}{\dnCq}{\dnCq}{\dnCq}{\dnCq}{\dnCq}{\dnCq}{\dnCq}{\dnCq}{\dnCq}{\dnCq}{\dnCq}{\dnCq}{\dnCq}{\dnCq}{\dnCq}{\dnCq}} & \dnW[W]{}: & {{\dnCh}{\dnCh}{\dnCh}{\dnCh}{\dnCh}{\dnCh}{\dnCh}{\dnCh}{\dnCh}{\dnCh}{\dnCh}{\dnCh}{\dnCh}{\dnCh}{\dnCh}{\dnCh}{\dnCh}{\dnCh}{\dnCh}{\dnCh}{\dnCh}{\dnCh}{\dnCh}{\dnCh}{\dnCh}{\dnCh}{\dnCh}{\dnCh}{\dnCh}{\dnCh}{\dnCh}{\dnCh}} \\
\dnI{7} & \dnW: & {{\dnCq}{\dnCq}{\dnCq}{\dnCq}{\dnCq}{\dnCq}{\dnCq}{\dnCq}{\dnCq}{\dnCq}{\dnCq}{\dnCq}{\dnCq}{\dnCq}{\dnCq}{\dnCq}{\dnCq}{\dnCq}{\dnCq}{\dnCq}{\dnCq}{\dnCq}{\dnCq}{\dnCq}{\dnCq}{\dnCq}{\dnCq}{\dnCq}{\dnCq}{\dnCq}{\dnCq}{\dnCq}} & \dnW[W]{}: & {{\dnCh}{\dnCh}{\dnCh}{\dnCh}{\dnCh}{\dnCh}{\dnCh}{\dnCh}{\dnCh}{\dnCh}{\dnCh}{\dnCh}{\dnCh}{\dnCh}{\dnCh}{\dnCh}{\dnCh}{\dnCh}{\dnCh}{\dnCh}{\dnCh}{\dnCh}{\dnCh}{\dnCh}{\dnCh}{\dnCh}{\dnCh}{\dnCh}{\dnCh}{\dnCh}{\dnCh}{\dnCh}} \\
\dnI{8} & \dnW: & {{\dnCq}{\dnCq}{\dnCq}{\dnCq}{\dnCq}{\dnCq}{\dnCq}{\dnCq}{\dnCq}{\dnCq}{\dnCq}{\dnCq}{\dnCq}{\dnCq}{\dnCq}{\dnCq}{\dnCq}{\dnCq}{\dnCq}{\dnCq}{\dnCq}{\dnCq}{\dnCq}{\dnCq}{\dnCq}{\dnCq}{\dnCq}{\dnCq}{\dnCq}{\dnCq}{\dnCq}{\dnCq}} & \dnW[W]{}: & {{\dnCx}{\dnCh}{\dnCh}{\dnCh}{\dnCh}{\dnCh}{\dnCh}{\dnCh}{\dnCh}{\dnCh}{\dnCh}{\dnCh}{\dnCh}{\dnCh}{\dnCh}{\dnCh}{\dnCh}{\dnCh}{\dnCh}{\dnCh}{\dnCh}{\dnCh}{\dnCh}{\dnCh}{\dnCh}{\dnCh}{\dnCh}{\dnCh}{\dnCh}{\dnCh}{\dnCh}{\dnCh}} \\
\dnI{9} & \dnW: & {{\dnCq}{\dnCq}{\dnCq}{\dnCq}{\dnCq}{\dnCq}{\dnCq}{\dnCq}{\dnCq}{\dnCq}{\dnCq}{\dnCq}{\dnCq}{\dnCq}{\dnCq}{\dnCq}{\dnCq}{\dnCq}{\dnCq}{\dnCq}{\dnCq}{\dnCq}{\dnCq}{\dnCq}{\dnCq}{\dnCq}{\dnCq}{\dnCq}{\dnCq}{\dnCq}{\dnCq}{\dnCq}} & \dnW[W]{}: & {{\dnCh}{\dnCh}{\dnCh}{\dnCh}{\dnCh}{\dnCh}{\dnCh}{\dnCh}{\dnCh}{\dnCh}{\dnCh}{\dnCh}{\dnCh}{\dnCh}{\dnCh}{\dnCh}{\dnCh}{\dnCh}{\dnCh}{\dnCh}{\dnCh}{\dnCh}{\dnCh}{\dnCh}{\dnCh}{\dnCh}{\dnCh}{\dnCh}{\dnCh}{\dnCh}{\dnCh}{\dnCh}} \\
\dnI{10} & \dnW: & {{\dnCq}{\dnCq}{\dnCq}{\dnCq}{\dnCq}{\dnCq}{\dnCq}{\dnCq}{\dnCq}{\dnCq}{\dnCq}{\dnCq}{\dnCq}{\dnCq}{\dnCq}{\dnCq}{\dnCq}{\dnCq}{\dnCq}{\dnCq}{\dnCq}{\dnCq}{\dnCq}{\dnCq}{\dnCq}{\dnCq}{\dnCq}{\dnCq}{\dnCq}{\dnCq}{\dnCq}{\dnCq}} & \dnW[W]{}: & {{\dnCh}{\dnCh}{\dnCh}{\dnCh}{\dnCh}{\dnCh}{\dnCh}{\dnCh}{\dnCh}{\dnCh}{\dnCh}{\dnCh}{\dnCh}{\dnCh}{\dnCh}{\dnCh}{\dnCh}{\dnCh}{\dnCh}{\dnCh}{\dnCh}{\dnCh}{\dnCh}{\dnCh}{\dnCh}{\dnCh}{\dnCh}{\dnCh}{\dnCh}{\dnCh}{\dnCh}{\dnCh}} \\
\dnI{11} & \dnW: & {{\dnCq}{\dnCq}{\dnCq}{\dnCq}{\dnCq}{\dnCq}{\dnCq}{\dnCq}{\dnCq}{\dnCq}{\dnCq}{\dnCq}{\dnCq}{\dnCq}{\dnCq}{\dnCq}{\dnCq}{\dnCq}{\dnCq}{\dnCq}{\dnCq}{\dnCq}{\dnCq}{\dnCq}{\dnCq}{\dnCq}{\dnCq}{\dnCq}{\dnCq}{\dnCq}{\dnCq}{\dnCq}} & \dnW[W]{}: & {{\dnCh}{\dnCh}{\dnCh}{\dnCh}{\dnCh}{\dnCh}{\dnCh}{\dnCh}{\dnCh}{\dnCh}{\dnCh}{\dnCh}{\dnCh}{\dnCh}{\dnCh}{\dnCh}{\dnCh}{\dnCh}{\dnCh}{\dnCh}{\dnCh}{\dnCh}{\dnCh}{\dnCh}{\dnCh}{\dnCh}{\dnCh}{\dnCh}{\dnCh}{\dnCh}{\dnCh}{\dnCh}} \\
\dnI{12} & \dnW: & {{\dnCq}{\dnCq}{\dnCq}{\dnCq}{\dnCq}{\dnCq}{\dnCq}{\dnCq}{\dnCq}{\dnCq}{\dnCq}{\dnCq}{\dnCq}{\dnCq}{\dnCq}{\dnCq}{\dnCq}{\dnCq}{\dnCq}{\dnCq}{\dnCq}{\dnCq}{\dnCq}{\dnCq}{\dnCq}{\dnCq}{\dnCq}{\dnCq}{\dnCq}{\dnCq}{\dnCq}{\dnCq}} & \dnW[W]{}: & {{\dnCx}{\dnCh}{\dnCh}{\dnCh}{\dnCh}{\dnCh}{\dnCh}{\dnCh}{\dnCh}{\dnCh}{\dnCh}{\dnCh}{\dnCh}{\dnCh}{\dnCh}{\dnCh}{\dnCh}{\dnCh}{\dnCh}{\dnCh}{\dnCh}{\dnCh}{\dnCh}{\dnCh}{\dnCh}{\dnCh}{\dnCh}{\dnCh}{\dnCh}{\dnCh}{\dnCh}{\dnCh}} \\
\dnI{13} & \dnW: & {{\dnCq}{\dnCq}{\dnCq}{\dnCq}{\dnCq}{\dnCq}{\dnCq}{\dnCq}{\dnCq}{\dnCq}{\dnCq}{\dnCq}{\dnCq}{\dnCq}{\dnCq}{\dnCq}{\dnCq}{\dnCq}{\dnCq}{\dnCq}{\dnCq}{\dnCq}{\dnCq}{\dnCq}{\dnCq}{\dnCq}{\dnCq}{\dnCq}{\dnCq}{\dnCq}{\dnCq}{\dnCq}} & \dnW[W]{}: & {{\dnCh}{\dnCh}{\dnCh}{\dnCh}{\dnCh}{\dnCh}{\dnCh}{\dnCh}{\dnCh}{\dnCh}{\dnCh}{\dnCh}{\dnCh}{\dnCh}{\dnCh}{\dnCh}{\dnCh}{\dnCh}{\dnCh}{\dnCh}{\dnCh}{\dnCh}{\dnCh}{\dnCh}{\dnCh}{\dnCh}{\dnCh}{\dnCh}{\dnCh}{\dnCh}{\dnCh}{\dnCh}} \\
\dnI{14} & \dnW: & {{\dnCq}{\dnCq}{\dnCq}{\dnCq}{\dnCq}{\dnCq}{\dnCq}{\dnCq}{\dnCq}{\dnCq}{\dnCq}{\dnCq}{\dnCq}{\dnCq}{\dnCq}{\dnCq}{\dnCq}{\dnCq}{\dnCq}{\dnCq}{\dnCq}{\dnCq}{\dnCq}{\dnCq}{\dnCq}{\dnCq}{\dnCq}{\dnCq}{\dnCq}{\dnCq}{\dnCq}{\dnCq}} & \dnW[W]{}: & {{\dnCh}{\dnCh}{\dnCh}{\dnCh}{\dnCh}{\dnCh}{\dnCh}{\dnCh}{\dnCh}{\dnCh}{\dnCh}{\dnCh}{\dnCh}{\dnCh}{\dnCh}{\dnCh}{\dnCh}{\dnCh}{\dnCh}{\dnCh}{\dnCh}{\dnCh}{\dnCh}{\dnCh}{\dnCh}{\dnCh}{\dnCh}{\dnCh}{\dnCh}{\dnCh}{\dnCh}{\dnCh}} \\
\dnI{15} & \dnW: & {{\dnCq}{\dnCq}{\dnCq}{\dnCq}{\dnCq}{\dnCq}{\dnCq}{\dnCq}{\dnCq}{\dnCq}{\dnCq}{\dnCq}{\dnCq}{\dnCq}{\dnCq}{\dnCq}{\dnCq}{\dnCq}{\dnCq}{\dnCq}{\dnCq}{\dnCq}{\dnCq}{\dnCq}{\dnCq}{\dnCq}{\dnCq}{\dnCq}{\dnCq}{\dnCq}{\dnCq}{\dnCq}} & \dnW[W]{}: & {{\dnCh}{\dnCh}{\dnCh}{\dnCh}{\dnCh}{\dnCh}{\dnCh}{\dnCh}{\dnCh}{\dnCh}{\dnCh}{\dnCh}{\dnCh}{\dnCh}{\dnCh}{\dnCh}{\dnCh}{\dnCh}{\dnCh}{\dnCh}{\dnCh}{\dnCh}{\dnCh}{\dnCh}{\dnCh}{\dnCh}{\dnCh}{\dnCh}{\dnCh}{\dnCh}{\dnCh}{\dnCh}} \\
\dnI{16} & \dnW: & {{\dnCq}{\dnCq}{\dnCq}{\dnCq}{\dnCq}{\dnCq}{\dnCq}{\dnCq}{\dnCq}{\dnCq}{\dnCq}{\dnCq}{\dnCq}{\dnCq}{\dnCq}{\dnCq}{\dnCq}{\dnCq}{\dnCq}{\dnCq}{\dnCq}{\dnCq}{\dnCq}{\dnCq}{\dnCq}{\dnCq}{\dnCq}{\dnCq}{\dnCq}{\dnCq}{\dnCq}{\dnCq}} & & \\
\dnI{17} & \dnW: & {{\dnCq}{\dnCq}{\dnCq}{\dnCq}{\dnCq}{\dnCq}{\dnCq}{\dnCq}{\dnCq}{\dnCq}{\dnCq}{\dnCq}{\dnCq}{\dnCq}{\dnCq}{\dnCq}{\dnCq}{\dnCq}{\dnCq}{\dnCq}{\dnCq}{\dnCq}{\dnCq}{\dnCq}{\dnCq}{\dnCq}{\dnCq}{\dnCq}{\dnCq}{\dnCq}{\dnCq}{\dnCq}} & & \\
\dnI{18} & \dnW: & {{\dnCq}{\dnCq}{\dnCq}{\dnCq}{\dnCq}{\dnCq}{\dnCq}{\dnCq}{\dnCq}{\dnCq}{\dnCq}{\dnCq}{\dnCq}{\dnCq}{\dnCq}{\dnCq}{\dnCq}{\dnCq}{\dnCq}{\dnCq}{\dnCq}{\dnCq}{\dnCq}{\dnCq}{\dnCq}{\dnCq}{\dnCq}{\dnCq}{\dnCq}{\dnCq}{\dnCq}{\dnCq}} & & \\
\dnI{19} & \dnW: & {{\dnCq}{\dnCq}{\dnCq}{\dnCq}{\dnCq}{\dnCq}{\dnCq}{\dnCq}{\dnCq}{\dnCq}{\dnCq}{\dnCq}{\dnCq}{\dnCq}{\dnCq}{\dnCq}{\dnCq}{\dnCq}{\dnCq}{\dnCq}{\dnCq}{\dnCq}{\dnCq}{\dnCq}{\dnCq}{\dnCq}{\dnCq}{\dnCq}{\dnCq}{\dnCq}{\dnCq}{\dnCq}} & & \\
\dnI{20} & \dnW: & {{\dnCq}{\dnCq}{\dnCq}{\dnCq}{\dnCq}{\dnCq}{\dnCq}{\dnCq}{\dnCq}{\dnCq}{\dnCq}{\dnCq}{\dnCq}{\dnCq}{\dnCq}{\dnCq}{\dnCq}{\dnCq}{\dnCq}{\dnCq}{\dnCq}{\dnCq}{\dnCq}{\dnCq}{\dnCq}{\dnCq}{\dnCq}{\dnCq}{\dnCq}{\dnCq}{\dnCq}{\dnCq}} & & \\
\dnI{21} & \dnW: & {{\dnCh}{\dnCh}{\dnCh}{\dnCh}{\dnCh}{\dnCh}{\dnCh}{\dnCh}{\dnCh}{\dnCh}{\dnCh}{\dnCh}{\dnCh}{\dnCh}{\dnCh}{\dnCh}{\dnCh}{\dnCh}{\dnCh}{\dnCh}{\dnCh}{\dnCh}{\dnCh}{\dnCh}{\dnCh}{\dnCh}{\dnCh}{\dnCh}{\dnCh}{\dnCh}{\dnCh}{\dnCh}} & & \\
\dnI{22} & \dnW: & {{\dnCh}{\dnCh}{\dnCh}{\dnCh}{\dnCh}{\dnCh}{\dnCh}{\dnCh}{\dnCh}{\dnCh}{\dnCh}{\dnCh}{\dnCh}{\dnCh}{\dnCh}{\dnCh}{\dnCh}{\dnCh}{\dnCh}{\dnCh}{\dnCh}{\dnCh}{\dnCh}{\dnCh}{\dnCh}{\dnCh}{\dnCh}{\dnCh}{\dnCh}{\dnCh}{\dnCh}{\dnCh}} & & \\
\dnI{23} & \dnW: & {{\dnCh}{\dnCh}{\dnCh}{\dnCh}{\dnCh}{\dnCh}{\dnCh}{\dnCh}{\dnCh}{\dnCh}{\dnCh}{\dnCh}{\dnCh}{\dnCh}{\dnCh}{\dnCh}{\dnCh}{\dnCh}{\dnCh}{\dnCh}{\dnCh}{\dnCh}{\dnCh}{\dnCh}{\dnCh}{\dnCh}{\dnCh}{\dnCh}{\dnCh}{\dnCh}{\dnCh}{\dnCh}} & & \\
\dnI{24} & \dnW: & {{\dnCh}{\dnCh}{\dnCh}{\dnCh}{\dnCh}{\dnCh}{\dnCh}{\dnCh}{\dnCh}{\dnCh}{\dnCh}{\dnCh}{\dnCh}{\dnCh}{\dnCh}{\dnCh}{\dnCh}{\dnCh}{\dnCh}{\dnCh}{\dnCh}{\dnCh}{\dnCh}{\dnCh}{\dnCh}{\dnCh}{\dnCh}{\dnCh}{\dnCh}{\dnCh}{\dnCh}{\dnCh}} & & \\
\dnI{25} & \dnW: & {{\dnCh}{\dnCh}{\dnCh}{\dnCh}{\dnCh}{\dnCh}{\dnCh}{\dnCh}{\dnCh}{\dnCh}{\dnCh}{\dnCh}{\dnCh}{\dnCh}{\dnCh}{\dnCh}{\dnCh}{\dnCh}{\dnCh}{\dnCh}{\dnCh}{\dnCh}{\dnCh}{\dnCh}{\dnCh}{\dnCh}{\dnCh}{\dnCh}{\dnCh}{\dnCh}{\dnCh}{\dnCh}} & & \\
\dnI{26} & \dnW: & {{\dnCh}{\dnCh}{\dnCh}{\dnCh}{\dnCh}{\dnCh}{\dnCh}{\dnCh}{\dnCh}{\dnCh}{\dnCh}{\dnCh}{\dnCh}{\dnCh}{\dnCh}{\dnCh}{\dnCh}{\dnCh}{\dnCh}{\dnCh}{\dnCh}{\dnCh}{\dnCh}{\dnCh}{\dnCh}{\dnCh}{\dnCh}{\dnCh}{\dnCh}{\dnCh}{\dnCh}{\dnCh}} & & \\
\dnI{27} & \dnW: & {{\dnCh}{\dnCh}{\dnCh}{\dnCh}{\dnCh}{\dnCh}{\dnCh}{\dnCh}{\dnCh}{\dnCh}{\dnCh}{\dnCh}{\dnCh}{\dnCh}{\dnCh}{\dnCh}{\dnCh}{\dnCh}{\dnCh}{\dnCh}{\dnCh}{\dnCh}{\dnCh}{\dnCh}{\dnCh}{\dnCh}{\dnCh}{\dnCh}{\dnCh}{\dnCh}{\dnCh}{\dnCh}} & & \\
\dnI{28} & \dnW: & {{\dnCh}{\dnCh}{\dnCh}{\dnCh}{\dnCh}{\dnCh}{\dnCh}{\dnCh}{\dnCh}{\dnCh}{\dnCh}{\dnCh}{\dnCh}{\dnCh}{\dnCh}{\dnCh}{\dnCh}{\dnCh}{\dnCh}{\dnCh}{\dnCh}{\dnCh}{\dnCh}{\dnCh}{\dnCh}{\dnCh}{\dnCh}{\dnCh}{\dnCh}{\dnCh}{\dnCh}{\dnCh}} & & \\
\dnI{29} & \dnW: & {{\dnCh}{\dnCh}{\dnCh}{\dnCh}{\dnCh}{\dnCh}{\dnCh}{\dnCh}{\dnCh}{\dnCh}{\dnCh}{\dnCh}{\dnCh}{\dnCh}{\dnCh}{\dnCh}{\dnCh}{\dnCh}{\dnCh}{\dnCh}{\dnCh}{\dnCh}{\dnCh}{\dnCh}{\dnCh}{\dnCh}{\dnCh}{\dnCh}{\dnCh}{\dnCh}{\dnCh}{\dnCh}} & & \\
\dnI{30} & \dnW: & {{\dnCh}{\dnCh}{\dnCh}{\dnCh}{\dnCh}{\dnCh}{\dnCh}{\dnCh}{\dnCh}{\dnCh}{\dnCh}{\dnCh}{\dnCh}{\dnCh}{\dnCh}{\dnCh}{\dnCh}{\dnCh}{\dnCh}{\dnCh}{\dnCh}{\dnCh}{\dnCh}{\dnCh}{\dnCh}{\dnCh}{\dnCh}{\dnCh}{\dnCh}{\dnCh}{\dnCh}{\dnCh}} & & \\
\dnI{31} & \dnW: & {{\dnCh}{\dnCh}{\dnCh}{\dnCh}{\dnCh}{\dnCh}{\dnCh}{\dnCh}{\dnCh}{\dnCh}{\dnCh}{\dnCh}{\dnCh}{\dnCh}{\dnCh}{\dnCh}{\dnCh}{\dnCh}{\dnCh}{\dnCh}{\dnCh}{\dnCh}{\dnCh}{\dnCh}{\dnCh}{\dnCh}{\dnCh}{\dnCh}{\dnCh}{\dnCh}{\dnCh}{\dnCh}} & & \\
\dnI{32} & \dnW: & {{\dnCx}{\dnCh}{\dnCh}{\dnCh}{\dnCh}{\dnCh}{\dnCh}{\dnCh}{\dnCh}{\dnCh}{\dnCh}{\dnCh}{\dnCh}{\dnCh}{\dnCh}{\dnCh}{\dnCh}{\dnCh}{\dnCh}{\dnCh}{\dnCh}{\dnCh}{\dnCh}{\dnCh}{\dnCh}{\dnCh}{\dnCh}{\dnCh}{\dnCh}{\dnCh}{\dnCh}{\dnCh}} & & \\
\dnI{33} & \dnW: & {{\dnCh}{\dnCh}{\dnCh}{\dnCh}{\dnCh}{\dnCh}{\dnCh}{\dnCh}{\dnCh}{\dnCh}{\dnCh}{\dnCh}{\dnCh}{\dnCh}{\dnCh}{\dnCh}{\dnCh}{\dnCh}{\dnCh}{\dnCh}{\dnCh}{\dnCh}{\dnCh}{\dnCh}{\dnCh}{\dnCh}{\dnCh}{\dnCh}{\dnCh}{\dnCh}{\dnCh}{\dnCh}} & & \\
\dnI{34} & \dnW: & {{\dnCh}{\dnCh}{\dnCh}{\dnCh}{\dnCh}{\dnCh}{\dnCh}{\dnCh}{\dnCh}{\dnCh}{\dnCh}{\dnCh}{\dnCh}{\dnCh}{\dnCh}{\dnCh}{\dnCh}{\dnCh}{\dnCh}{\dnCh}{\dnCh}{\dnCh}{\dnCh}{\dnCh}{\dnCh}{\dnCh}{\dnCh}{\dnCh}{\dnCh}{\dnCh}{\dnCh}{\dnCh}} & & \\
\dnI{35} & \dnW: & {{\dnCh}{\dnCh}{\dnCh}{\dnCh}{\dnCh}{\dnCh}{\dnCh}{\dnCh}{\dnCh}{\dnCh}{\dnCh}{\dnCh}{\dnCh}{\dnCh}{\dnCh}{\dnCh}{\dnCh}{\dnCh}{\dnCh}{\dnCh}{\dnCh}{\dnCh}{\dnCh}{\dnCh}{\dnCh}{\dnCh}{\dnCh}{\dnCh}{\dnCh}{\dnCh}{\dnCh}{\dnCh}} & & \\
\dnI{36} & \dnW: & {{\dnCh}{\dnCh}{\dnCh}{\dnCh}{\dnCh}{\dnCh}{\dnCh}{\dnCh}{\dnCh}{\dnCh}{\dnCh}{\dnCh}{\dnCh}{\dnCh}{\dnCh}{\dnCh}{\dnCh}{\dnCh}{\dnCh}{\dnCh}{\dnCh}{\dnCh}{\dnCh}{\dnCh}{\dnCh}{\dnCh}{\dnCh}{\dnCh}{\dnCh}{\dnCh}{\dnCh}{\dnCh}} & & \\
\dnI{37} & \dnW: & {{\dnCh}{\dnCh}{\dnCh}{\dnCh}{\dnCh}{\dnCh}{\dnCh}{\dnCh}{\dnCh}{\dnCh}{\dnCh}{\dnCh}{\dnCh}{\dnCh}{\dnCh}{\dnCh}{\dnCh}{\dnCh}{\dnCh}{\dnCh}{\dnCh}{\dnCh}{\dnCh}{\dnCh}{\dnCh}{\dnCh}{\dnCh}{\dnCh}{\dnCh}{\dnCh}{\dnCh}{\dnCh}} & & \\
\dnI{38} & \dnW: & {{\dnCh}{\dnCh}{\dnCh}{\dnCh}{\dnCh}{\dnCh}{\dnCh}{\dnCh}{\dnCh}{\dnCh}{\dnCh}{\dnCh}{\dnCh}{\dnCh}{\dnCh}{\dnCh}{\dnCh}{\dnCh}{\dnCh}{\dnCh}{\dnCh}{\dnCh}{\dnCh}{\dnCh}{\dnCh}{\dnCh}{\dnCh}{\dnCh}{\dnCh}{\dnCh}{\dnCh}{\dnCh}} & & \\
\dnI{39} & \dnW: & {{\dnCh}{\dnCh}{\dnCh}{\dnCh}{\dnCh}{\dnCh}{\dnCh}{\dnCh}{\dnCh}{\dnCh}{\dnCh}{\dnCh}{\dnCh}{\dnCh}{\dnCh}{\dnCh}{\dnCh}{\dnCh}{\dnCh}{\dnCh}{\dnCh}{\dnCh}{\dnCh}{\dnCh}{\dnCh}{\dnCh}{\dnCh}{\dnCh}{\dnCh}{\dnCh}{\dnCh}{\dnCh}} & & \\
\dnI{40} & \dnW: & {{\dnCh}{\dnCh}{\dnCh}{\dnCh}{\dnCh}{\dnCh}{\dnCh}{\dnCh}{\dnCh}{\dnCh}{\dnCh}{\dnCh}{\dnCh}{\dnCh}{\dnCh}{\dnCh}{\dnCh}{\dnCh}{\dnCh}{\dnCh}{\dnCh}{\dnCh}{\dnCh}{\dnCh}{\dnCh}{\dnCh}{\dnCh}{\dnCh}{\dnCh}{\dnCh}{\dnCh}{\dnCh}} & & \\
\dnI{41} & \dnW: & {{\dnCh}{\dnCh}{\dnCh}{\dnCh}{\dnCh}{\dnCh}{\dnCh}{\dnCh}{\dnCh}{\dnCh}{\dnCh}{\dnCh}{\dnCh}{\dnCh}{\dnCh}{\dnCh}{\dnCh}{\dnCh}{\dnCh}{\dnCh}{\dnCh}{\dnCh}{\dnCh}{\dnCh}{\dnCh}{\dnCh}{\dnCh}{\dnCh}{\dnCh}{\dnCh}{\dnCh}{\dnCh}} & & \\
\dnI{42} & \dnW: & {{\dnCh}{\dnCh}{\dnCh}{\dnCh}{\dnCh}{\dnCh}{\dnCh}{\dnCh}{\dnCh}{\dnCh}{\dnCh}{\dnCh}{\dnCh}{\dnCh}{\dnCh}{\dnCh}{\dnCh}{\dnCh}{\dnCh}{\dnCh}{\dnCh}{\dnCh}{\dnCh}{\dnCh}{\dnCh}{\dnCh}{\dnCh}{\dnCh}{\dnCh}{\dnCh}{\dnCh}{\dnCh}} & & \\
\dnI{43} & \dnW: & {{\dnCh}{\dnCh}{\dnCh}{\dnCh}{\dnCh}{\dnCh}{\dnCh}{\dnCh}{\dnCh}{\dnCh}{\dnCh}{\dnCh}{\dnCh}{\dnCh}{\dnCh}{\dnCh}{\dnCh}{\dnCh}{\dnCh}{\dnCh}{\dnCh}{\dnCh}{\dnCh}{\dnCh}{\dnCh}{\dnCh}{\dnCh}{\dnCh}{\dnCh}{\dnCh}{\dnCh}{\dnCh}} & & \\
\dnI{44} & \dnW: & {{\dnCh}{\dnCh}{\dnCh}{\dnCh}{\dnCh}{\dnCh}{\dnCh}{\dnCh}{\dnCh}{\dnCh}{\dnCh}{\dnCh}{\dnCh}{\dnCh}{\dnCh}{\dnCh}{\dnCh}{\dnCh}{\dnCh}{\dnCh}{\dnCh}{\dnCh}{\dnCh}{\dnCh}{\dnCh}{\dnCh}{\dnCh}{\dnCh}{\dnCh}{\dnCh}{\dnCh}{\dnCh}} & & \\
\dnI{45} & \dnW: & {{\dnCh}{\dnCh}{\dnCh}{\dnCh}{\dnCh}{\dnCh}{\dnCh}{\dnCh}{\dnCh}{\dnCh}{\dnCh}{\dnCh}{\dnCh}{\dnCh}{\dnCh}{\dnCh}{\dnCh}{\dnCh}{\dnCh}{\dnCh}{\dnCh}{\dnCh}{\dnCh}{\dnCh}{\dnCh}{\dnCh}{\dnCh}{\dnCh}{\dnCh}{\dnCh}{\dnCh}{\dnCh}} & & \\
\dnI{46} & \dnW: & {{\dnCh}{\dnCh}{\dnCh}{\dnCh}{\dnCh}{\dnCh}{\dnCh}{\dnCh}{\dnCh}{\dnCh}{\dnCh}{\dnCh}{\dnCh}{\dnCh}{\dnCh}{\dnCh}{\dnCh}{\dnCh}{\dnCh}{\dnCh}{\dnCh}{\dnCh}{\dnCh}{\dnCh}{\dnCh}{\dnCh}{\dnCh}{\dnCh}{\dnCh}{\dnCh}{\dnCh}{\dnCh}} & & \\
\dnI{47} & \dnW: & {{\dnCh}{\dnCh}{\dnCh}{\dnCh}{\dnCh}{\dnCh}{\dnCh}{\dnCh}{\dnCh}{\dnCh}{\dnCh}{\dnCh}{\dnCh}{\dnCh}{\dnCh}{\dnCh}{\dnCh}{\dnCh}{\dnCh}{\dnCh}{\dnCh}{\dnCh}{\dnCh}{\dnCh}{\dnCh}{\dnCh}{\dnCh}{\dnCh}{\dnCh}{\dnCh}{\dnCh}{\dnCh}} & & \\
\hline
\end{tabular}
}
% internal testcase name: 48steps_m0m2m4m8m12_1.xml
\caption[MD4 testcase C]{Differential characteristic of MD4 testcase C}
\label{fig:tcC}
\end{center}
\end{figure}
}

\texttt{
\begin{figure}[!ht]
\begin{center}
\scalebox{.55}{
\begin{tabular}{|r|c|c|c|c|c|c|}
\hline
  $i$  &      &  A  &  &  E &  & W \\
\hline
 \dnI{-4} & \dnW: & {{\dnCh}{\dnCh}{\dnCh}{\dnCh}{\dnCh}{\dnCh}{\dnCh}{\dnCh}{\dnCh}{\dnCh}{\dnCh}{\dnCh}{\dnCh}{\dnCh}{\dnCh}{\dnCh}{\dnCh}{\dnCh}{\dnCh}{\dnCh}{\dnCh}{\dnCh}{\dnCh}{\dnCh}{\dnCh}{\dnCh}{\dnCh}{\dnCh}{\dnCh}{\dnCh}{\dnCh}{\dnCh}} & \dnW[E]{}: & {{\dnCh}{\dnCh}{\dnCh}{\dnCh}{\dnCh}{\dnCh}{\dnCh}{\dnCh}{\dnCh}{\dnCh}{\dnCh}{\dnCh}{\dnCh}{\dnCh}{\dnCh}{\dnCh}{\dnCh}{\dnCh}{\dnCh}{\dnCh}{\dnCh}{\dnCh}{\dnCh}{\dnCh}{\dnCh}{\dnCh}{\dnCh}{\dnCh}{\dnCh}{\dnCh}{\dnCh}{\dnCh}} & &     \\
 \dnI{-3} & \dnW: & {{\dnCh}{\dnCh}{\dnCh}{\dnCh}{\dnCh}{\dnCh}{\dnCh}{\dnCh}{\dnCh}{\dnCh}{\dnCh}{\dnCh}{\dnCh}{\dnCh}{\dnCh}{\dnCh}{\dnCh}{\dnCh}{\dnCh}{\dnCh}{\dnCh}{\dnCh}{\dnCh}{\dnCh}{\dnCh}{\dnCh}{\dnCh}{\dnCh}{\dnCh}{\dnCh}{\dnCh}{\dnCh}} & \dnW[E]{}: & {{\dnCh}{\dnCh}{\dnCh}{\dnCh}{\dnCh}{\dnCh}{\dnCh}{\dnCh}{\dnCh}{\dnCh}{\dnCh}{\dnCh}{\dnCh}{\dnCh}{\dnCh}{\dnCh}{\dnCh}{\dnCh}{\dnCh}{\dnCh}{\dnCh}{\dnCh}{\dnCh}{\dnCh}{\dnCh}{\dnCh}{\dnCh}{\dnCh}{\dnCh}{\dnCh}{\dnCh}{\dnCh}} & &     \\
 \dnI{-2} & \dnW: & {{\dnCh}{\dnCh}{\dnCh}{\dnCh}{\dnCh}{\dnCh}{\dnCh}{\dnCh}{\dnCh}{\dnCh}{\dnCh}{\dnCh}{\dnCh}{\dnCh}{\dnCh}{\dnCh}{\dnCh}{\dnCh}{\dnCh}{\dnCh}{\dnCh}{\dnCh}{\dnCh}{\dnCh}{\dnCh}{\dnCh}{\dnCh}{\dnCh}{\dnCh}{\dnCh}{\dnCh}{\dnCh}} & \dnW[E]{}: & {{\dnCh}{\dnCh}{\dnCh}{\dnCh}{\dnCh}{\dnCh}{\dnCh}{\dnCh}{\dnCh}{\dnCh}{\dnCh}{\dnCh}{\dnCh}{\dnCh}{\dnCh}{\dnCh}{\dnCh}{\dnCh}{\dnCh}{\dnCh}{\dnCh}{\dnCh}{\dnCh}{\dnCh}{\dnCh}{\dnCh}{\dnCh}{\dnCh}{\dnCh}{\dnCh}{\dnCh}{\dnCh}} & &     \\
 \dnI{-1} & \dnW: & {{\dnCh}{\dnCh}{\dnCh}{\dnCh}{\dnCh}{\dnCh}{\dnCh}{\dnCh}{\dnCh}{\dnCh}{\dnCh}{\dnCh}{\dnCh}{\dnCh}{\dnCh}{\dnCh}{\dnCh}{\dnCh}{\dnCh}{\dnCh}{\dnCh}{\dnCh}{\dnCh}{\dnCh}{\dnCh}{\dnCh}{\dnCh}{\dnCh}{\dnCh}{\dnCh}{\dnCh}{\dnCh}} & \dnW[E]{}: & {{\dnCh}{\dnCh}{\dnCh}{\dnCh}{\dnCh}{\dnCh}{\dnCh}{\dnCh}{\dnCh}{\dnCh}{\dnCh}{\dnCh}{\dnCh}{\dnCh}{\dnCh}{\dnCh}{\dnCh}{\dnCh}{\dnCh}{\dnCh}{\dnCh}{\dnCh}{\dnCh}{\dnCh}{\dnCh}{\dnCh}{\dnCh}{\dnCh}{\dnCh}{\dnCh}{\dnCh}{\dnCh}} & &     \\
 \dnI{0} & \dnW: & {{\dnCh}{\dnCh}{\dnCh}{\dnCh}{\dnCh}{\dnCh}{\dnCh}{\dnCh}{\dnCh}{\dnCh}{\dnCh}{\dnCh}{\dnCh}{\dnCh}{\dnCh}{\dnCh}{\dnCh}{\dnCh}{\dnCh}{\dnCh}{\dnCh}{\dnCh}{\dnCh}{\dnCh}{\dnCh}{\dnCh}{\dnCh}{\dnCh}{\dnCh}{\dnCh}{\dnCh}{\dnCh}} & \dnW[E]{}: & {{\dnCh}{\dnCh}{\dnCh}{\dnCh}{\dnCh}{\dnCh}{\dnCh}{\dnCh}{\dnCh}{\dnCh}{\dnCh}{\dnCh}{\dnCh}{\dnCh}{\dnCh}{\dnCh}{\dnCh}{\dnCh}{\dnCh}{\dnCh}{\dnCh}{\dnCh}{\dnCh}{\dnCh}{\dnCh}{\dnCh}{\dnCh}{\dnCh}{\dnCh}{\dnCh}{\dnCh}{\dnCh}} & \dnW[W]{}: & {{\dnCh}{\dnCh}{\dnCh}{\dnCh}{\dnCh}{\dnCh}{\dnCh}{\dnCh}{\dnCh}{\dnCh}{\dnCh}{\dnCh}{\dnCh}{\dnCh}{\dnCh}{\dnCh}{\dnCh}{\dnCh}{\dnCh}{\dnCh}{\dnCh}{\dnCh}{\dnCh}{\dnCh}{\dnCh}{\dnCh}{\dnCh}{\dnCh}{\dnCh}{\dnCh}{\dnCh}{\dnCh}} \\
 \dnI{1} & \dnW: & {{\dnCh}{\dnCh}{\dnCh}{\dnCh}{\dnCh}{\dnCh}{\dnCh}{\dnCh}{\dnCh}{\dnCh}{\dnCh}{\dnCh}{\dnCh}{\dnCh}{\dnCh}{\dnCh}{\dnCh}{\dnCh}{\dnCh}{\dnCh}{\dnCh}{\dnCh}{\dnCh}{\dnCh}{\dnCh}{\dnCh}{\dnCh}{\dnCh}{\dnCh}{\dnCh}{\dnCh}{\dnCh}} & \dnW[E]{}: & {{\dnCh}{\dnCh}{\dnCh}{\dnCh}{\dnCh}{\dnCh}{\dnCh}{\dnCh}{\dnCh}{\dnCh}{\dnCh}{\dnCh}{\dnCh}{\dnCh}{\dnCh}{\dnCh}{\dnCh}{\dnCh}{\dnCh}{\dnCh}{\dnCh}{\dnCh}{\dnCh}{\dnCh}{\dnCh}{\dnCh}{\dnCh}{\dnCh}{\dnCh}{\dnCh}{\dnCh}{\dnCh}} & \dnW[W]{}: & {{\dnCh}{\dnCh}{\dnCh}{\dnCh}{\dnCh}{\dnCh}{\dnCh}{\dnCh}{\dnCh}{\dnCh}{\dnCh}{\dnCh}{\dnCh}{\dnCh}{\dnCh}{\dnCh}{\dnCh}{\dnCh}{\dnCh}{\dnCh}{\dnCh}{\dnCh}{\dnCh}{\dnCh}{\dnCh}{\dnCh}{\dnCh}{\dnCh}{\dnCh}{\dnCh}{\dnCh}{\dnCh}} \\
 \dnI{2} & \dnW: & {{\dnCh}{\dnCh}{\dnCh}{\dnCh}{\dnCh}{\dnCh}{\dnCh}{\dnCh}{\dnCh}{\dnCh}{\dnCh}{\dnCh}{\dnCh}{\dnCh}{\dnCh}{\dnCh}{\dnCh}{\dnCh}{\dnCh}{\dnCh}{\dnCh}{\dnCh}{\dnCh}{\dnCh}{\dnCh}{\dnCh}{\dnCh}{\dnCh}{\dnCh}{\dnCh}{\dnCh}{\dnCh}} & \dnW[E]{}: & {{\dnCh}{\dnCh}{\dnCh}{\dnCh}{\dnCh}{\dnCh}{\dnCh}{\dnCh}{\dnCh}{\dnCh}{\dnCh}{\dnCh}{\dnCh}{\dnCh}{\dnCh}{\dnCh}{\dnCh}{\dnCh}{\dnCh}{\dnCh}{\dnCh}{\dnCh}{\dnCh}{\dnCh}{\dnCh}{\dnCh}{\dnCh}{\dnCh}{\dnCh}{\dnCh}{\dnCh}{\dnCh}} & \dnW[W]{}: & {{\dnCh}{\dnCh}{\dnCh}{\dnCh}{\dnCh}{\dnCh}{\dnCh}{\dnCh}{\dnCh}{\dnCh}{\dnCh}{\dnCh}{\dnCh}{\dnCh}{\dnCh}{\dnCh}{\dnCh}{\dnCh}{\dnCh}{\dnCh}{\dnCh}{\dnCh}{\dnCh}{\dnCh}{\dnCh}{\dnCh}{\dnCh}{\dnCh}{\dnCh}{\dnCh}{\dnCh}{\dnCh}} \\
 \dnI{3} & \dnW: & {{\dnCx}{\dnCh}{\dnCh}{\dnCh}{\dnCh}{\dnCh}{\dnCh}{\dnCh}{\dnCh}{\dnCh}{\dnCh}{\dnCh}{\dnCh}{\dnCh}{\dnCh}{\dnCh}{\dnCh}{\dnCh}{\dnCh}{\dnCh}{\dnCh}{\dnCh}{\dnCh}{\dnCh}{\dnCh}{\dnCh}{\dnCh}{\dnCh}{\dnCh}{\dnCh}{\dnCh}{\dnCh}} & \dnW[E]{}: & {{\dnCq}{\dnCq}{\dnCq}{\dnCq}{\dnCq}{\dnCq}{\dnCq}{\dnCq}{\dnCq}{\dnCq}{\dnCq}{\dnCq}{\dnCq}{\dnCq}{\dnCq}{\dnCq}{\dnCq}{\dnCq}{\dnCq}{\dnCq}{\dnCq}{\dnCq}{\dnCq}{\dnCq}{\dnCq}{\dnCq}{\dnCq}{\dnCq}{\dnCq}{\dnCq}{\dnCq}{\dnCq}} & \dnW[W]{}: & {{\dnCq}{\dnCq}{\dnCq}{\dnCq}{\dnCq}{\dnCq}{\dnCq}{\dnCq}{\dnCq}{\dnCq}{\dnCq}{\dnCq}{\dnCq}{\dnCq}{\dnCq}{\dnCq}{\dnCq}{\dnCq}{\dnCq}{\dnCq}{\dnCq}{\dnCq}{\dnCq}{\dnCq}{\dnCq}{\dnCq}{\dnCq}{\dnCq}{\dnCq}{\dnCq}{\dnCq}{\dnCq}} \\
 \dnI{4} & \dnW: & {{\dnCh}{\dnCh}{\dnCh}{\dnCh}{\dnCh}{\dnCh}{\dnCh}{\dnCh}{\dnCh}{\dnCh}{\dnCh}{\dnCh}{\dnCh}{\dnCh}{\dnCh}{\dnCh}{\dnCh}{\dnCh}{\dnCh}{\dnCh}{\dnCh}{\dnCh}{\dnCh}{\dnCh}{\dnCh}{\dnCh}{\dnCh}{\dnCh}{\dnCh}{\dnCh}{\dnCh}{\dnCh}} & \dnW[E]{}: & {{\dnCq}{\dnCq}{\dnCq}{\dnCq}{\dnCq}{\dnCq}{\dnCq}{\dnCq}{\dnCq}{\dnCq}{\dnCq}{\dnCq}{\dnCq}{\dnCq}{\dnCq}{\dnCq}{\dnCq}{\dnCq}{\dnCq}{\dnCq}{\dnCq}{\dnCq}{\dnCq}{\dnCq}{\dnCq}{\dnCq}{\dnCq}{\dnCq}{\dnCq}{\dnCq}{\dnCq}{\dnCq}} & \dnW[W]{}: & {{\dnCq}{\dnCq}{\dnCq}{\dnCq}{\dnCq}{\dnCq}{\dnCq}{\dnCq}{\dnCq}{\dnCq}{\dnCq}{\dnCq}{\dnCq}{\dnCq}{\dnCq}{\dnCq}{\dnCq}{\dnCq}{\dnCq}{\dnCq}{\dnCq}{\dnCq}{\dnCq}{\dnCq}{\dnCq}{\dnCq}{\dnCq}{\dnCq}{\dnCq}{\dnCq}{\dnCq}{\dnCq}} \\
 \dnI{5} & \dnW: & {{\dnCh}{\dnCh}{\dnCh}{\dnCh}{\dnCh}{\dnCh}{\dnCh}{\dnCh}{\dnCh}{\dnCh}{\dnCh}{\dnCh}{\dnCh}{\dnCh}{\dnCh}{\dnCh}{\dnCh}{\dnCh}{\dnCh}{\dnCh}{\dnCh}{\dnCh}{\dnCh}{\dnCh}{\dnCh}{\dnCh}{\dnCh}{\dnCh}{\dnCh}{\dnCh}{\dnCh}{\dnCh}} & \dnW[E]{}: & {{\dnCq}{\dnCq}{\dnCq}{\dnCq}{\dnCq}{\dnCq}{\dnCq}{\dnCq}{\dnCq}{\dnCq}{\dnCq}{\dnCq}{\dnCq}{\dnCq}{\dnCq}{\dnCq}{\dnCq}{\dnCq}{\dnCq}{\dnCq}{\dnCq}{\dnCq}{\dnCq}{\dnCq}{\dnCq}{\dnCq}{\dnCq}{\dnCq}{\dnCq}{\dnCq}{\dnCq}{\dnCq}} & \dnW[W]{}: & {{\dnCq}{\dnCq}{\dnCq}{\dnCq}{\dnCq}{\dnCq}{\dnCq}{\dnCq}{\dnCq}{\dnCq}{\dnCq}{\dnCq}{\dnCq}{\dnCq}{\dnCq}{\dnCq}{\dnCq}{\dnCq}{\dnCq}{\dnCq}{\dnCq}{\dnCq}{\dnCq}{\dnCq}{\dnCq}{\dnCq}{\dnCq}{\dnCq}{\dnCq}{\dnCq}{\dnCq}{\dnCq}} \\
 \dnI{6} & \dnW: & {{\dnCh}{\dnCh}{\dnCh}{\dnCh}{\dnCh}{\dnCh}{\dnCh}{\dnCh}{\dnCh}{\dnCh}{\dnCh}{\dnCh}{\dnCh}{\dnCh}{\dnCh}{\dnCh}{\dnCh}{\dnCh}{\dnCh}{\dnCh}{\dnCh}{\dnCh}{\dnCh}{\dnCh}{\dnCh}{\dnCh}{\dnCh}{\dnCh}{\dnCh}{\dnCh}{\dnCh}{\dnCh}} & \dnW[E]{}: & {{\dnCq}{\dnCq}{\dnCq}{\dnCq}{\dnCq}{\dnCq}{\dnCq}{\dnCq}{\dnCq}{\dnCq}{\dnCq}{\dnCq}{\dnCq}{\dnCq}{\dnCq}{\dnCq}{\dnCq}{\dnCq}{\dnCq}{\dnCq}{\dnCq}{\dnCq}{\dnCq}{\dnCq}{\dnCq}{\dnCq}{\dnCq}{\dnCq}{\dnCq}{\dnCq}{\dnCq}{\dnCq}} & \dnW[W]{}: & {{\dnCq}{\dnCq}{\dnCq}{\dnCq}{\dnCq}{\dnCq}{\dnCq}{\dnCq}{\dnCq}{\dnCq}{\dnCq}{\dnCq}{\dnCq}{\dnCq}{\dnCq}{\dnCq}{\dnCq}{\dnCq}{\dnCq}{\dnCq}{\dnCq}{\dnCq}{\dnCq}{\dnCq}{\dnCq}{\dnCq}{\dnCq}{\dnCq}{\dnCq}{\dnCq}{\dnCq}{\dnCq}} \\
 \dnI{7} & \dnW: & {{\dnCh}{\dnCh}{\dnCh}{\dnCh}{\dnCh}{\dnCh}{\dnCh}{\dnCh}{\dnCh}{\dnCh}{\dnCh}{\dnCh}{\dnCh}{\dnCh}{\dnCh}{\dnCh}{\dnCh}{\dnCh}{\dnCh}{\dnCh}{\dnCh}{\dnCh}{\dnCh}{\dnCh}{\dnCh}{\dnCh}{\dnCh}{\dnCh}{\dnCh}{\dnCh}{\dnCh}{\dnCh}} & \dnW[E]{}: & {{\dnCq}{\dnCq}{\dnCq}{\dnCq}{\dnCq}{\dnCq}{\dnCq}{\dnCq}{\dnCq}{\dnCq}{\dnCq}{\dnCq}{\dnCq}{\dnCq}{\dnCq}{\dnCq}{\dnCq}{\dnCq}{\dnCq}{\dnCq}{\dnCq}{\dnCq}{\dnCq}{\dnCq}{\dnCq}{\dnCq}{\dnCq}{\dnCq}{\dnCq}{\dnCq}{\dnCq}{\dnCq}} & \dnW[W]{}: & {{\dnCq}{\dnCq}{\dnCq}{\dnCq}{\dnCq}{\dnCq}{\dnCq}{\dnCq}{\dnCq}{\dnCq}{\dnCq}{\dnCq}{\dnCq}{\dnCq}{\dnCq}{\dnCq}{\dnCq}{\dnCq}{\dnCq}{\dnCq}{\dnCq}{\dnCq}{\dnCq}{\dnCq}{\dnCq}{\dnCq}{\dnCq}{\dnCq}{\dnCq}{\dnCq}{\dnCq}{\dnCq}} \\
 \dnI{8} & \dnW: & {{\dnCh}{\dnCh}{\dnCh}{\dnCh}{\dnCh}{\dnCh}{\dnCh}{\dnCh}{\dnCh}{\dnCh}{\dnCh}{\dnCh}{\dnCh}{\dnCh}{\dnCh}{\dnCh}{\dnCh}{\dnCh}{\dnCh}{\dnCh}{\dnCh}{\dnCh}{\dnCh}{\dnCh}{\dnCh}{\dnCh}{\dnCh}{\dnCh}{\dnCh}{\dnCh}{\dnCh}{\dnCh}} & \dnW[E]{}: & {{\dnCh}{\dnCh}{\dnCh}{\dnCh}{\dnCh}{\dnCh}{\dnCh}{\dnCh}{\dnCh}{\dnCh}{\dnCh}{\dnCh}{\dnCh}{\dnCh}{\dnCh}{\dnCh}{\dnCh}{\dnCh}{\dnCh}{\dnCh}{\dnCh}{\dnCh}{\dnCh}{\dnCh}{\dnCh}{\dnCh}{\dnCh}{\dnCh}{\dnCh}{\dnCh}{\dnCh}{\dnCh}} & \dnW[W]{}: & {{\dnCq}{\dnCq}{\dnCq}{\dnCq}{\dnCq}{\dnCq}{\dnCq}{\dnCq}{\dnCq}{\dnCq}{\dnCq}{\dnCq}{\dnCq}{\dnCq}{\dnCq}{\dnCq}{\dnCq}{\dnCq}{\dnCq}{\dnCq}{\dnCq}{\dnCq}{\dnCq}{\dnCq}{\dnCq}{\dnCq}{\dnCq}{\dnCq}{\dnCq}{\dnCq}{\dnCq}{\dnCq}} \\
 \dnI{9} & \dnW: & {{\dnCh}{\dnCh}{\dnCh}{\dnCh}{\dnCh}{\dnCh}{\dnCh}{\dnCh}{\dnCh}{\dnCh}{\dnCh}{\dnCh}{\dnCh}{\dnCh}{\dnCh}{\dnCh}{\dnCh}{\dnCh}{\dnCh}{\dnCh}{\dnCh}{\dnCh}{\dnCh}{\dnCh}{\dnCh}{\dnCh}{\dnCh}{\dnCh}{\dnCh}{\dnCh}{\dnCh}{\dnCh}} & \dnW[E]{}: & {{\dnCh}{\dnCh}{\dnCh}{\dnCh}{\dnCh}{\dnCh}{\dnCh}{\dnCh}{\dnCh}{\dnCh}{\dnCh}{\dnCh}{\dnCh}{\dnCh}{\dnCh}{\dnCh}{\dnCh}{\dnCh}{\dnCh}{\dnCh}{\dnCh}{\dnCh}{\dnCh}{\dnCh}{\dnCh}{\dnCh}{\dnCh}{\dnCh}{\dnCh}{\dnCh}{\dnCh}{\dnCh}} & \dnW[W]{}: & {{\dnCh}{\dnCh}{\dnCh}{\dnCh}{\dnCh}{\dnCh}{\dnCh}{\dnCh}{\dnCh}{\dnCh}{\dnCh}{\dnCh}{\dnCh}{\dnCh}{\dnCh}{\dnCh}{\dnCh}{\dnCh}{\dnCh}{\dnCh}{\dnCh}{\dnCh}{\dnCh}{\dnCh}{\dnCh}{\dnCh}{\dnCh}{\dnCh}{\dnCh}{\dnCh}{\dnCh}{\dnCh}} \\
 \dnI{10} & \dnW: & {{\dnCh}{\dnCh}{\dnCh}{\dnCh}{\dnCh}{\dnCh}{\dnCh}{\dnCh}{\dnCh}{\dnCh}{\dnCh}{\dnCh}{\dnCh}{\dnCh}{\dnCh}{\dnCh}{\dnCh}{\dnCh}{\dnCh}{\dnCh}{\dnCh}{\dnCh}{\dnCh}{\dnCh}{\dnCh}{\dnCh}{\dnCh}{\dnCh}{\dnCh}{\dnCh}{\dnCh}{\dnCh}} & \dnW[E]{}: & {{\dnCh}{\dnCh}{\dnCh}{\dnCh}{\dnCh}{\dnCh}{\dnCh}{\dnCh}{\dnCh}{\dnCh}{\dnCh}{\dnCh}{\dnCh}{\dnCh}{\dnCh}{\dnCh}{\dnCh}{\dnCh}{\dnCh}{\dnCh}{\dnCh}{\dnCh}{\dnCh}{\dnCh}{\dnCh}{\dnCh}{\dnCh}{\dnCh}{\dnCh}{\dnCh}{\dnCh}{\dnCh}} & \dnW[W]{}: & {{\dnCh}{\dnCh}{\dnCh}{\dnCh}{\dnCh}{\dnCh}{\dnCh}{\dnCh}{\dnCh}{\dnCh}{\dnCh}{\dnCh}{\dnCh}{\dnCh}{\dnCh}{\dnCh}{\dnCh}{\dnCh}{\dnCh}{\dnCh}{\dnCh}{\dnCh}{\dnCh}{\dnCh}{\dnCh}{\dnCh}{\dnCh}{\dnCh}{\dnCh}{\dnCh}{\dnCh}{\dnCh}} \\
 \dnI{11} & \dnW: & {{\dnCh}{\dnCh}{\dnCh}{\dnCh}{\dnCh}{\dnCh}{\dnCh}{\dnCh}{\dnCh}{\dnCh}{\dnCh}{\dnCh}{\dnCh}{\dnCh}{\dnCh}{\dnCh}{\dnCh}{\dnCh}{\dnCh}{\dnCh}{\dnCh}{\dnCh}{\dnCh}{\dnCh}{\dnCh}{\dnCh}{\dnCh}{\dnCh}{\dnCh}{\dnCh}{\dnCh}{\dnCh}} & \dnW[E]{}: & {{\dnCh}{\dnCh}{\dnCh}{\dnCh}{\dnCh}{\dnCh}{\dnCh}{\dnCh}{\dnCh}{\dnCh}{\dnCh}{\dnCh}{\dnCh}{\dnCh}{\dnCh}{\dnCh}{\dnCh}{\dnCh}{\dnCh}{\dnCh}{\dnCh}{\dnCh}{\dnCh}{\dnCh}{\dnCh}{\dnCh}{\dnCh}{\dnCh}{\dnCh}{\dnCh}{\dnCh}{\dnCh}} & \dnW[W]{}: & {{\dnCq}{\dnCq}{\dnCq}{\dnCq}{\dnCq}{\dnCq}{\dnCq}{\dnCq}{\dnCq}{\dnCq}{\dnCq}{\dnCq}{\dnCq}{\dnCq}{\dnCq}{\dnCq}{\dnCq}{\dnCq}{\dnCq}{\dnCq}{\dnCq}{\dnCq}{\dnCq}{\dnCq}{\dnCq}{\dnCq}{\dnCq}{\dnCq}{\dnCq}{\dnCq}{\dnCq}{\dnCq}} \\
 \dnI{12} & \dnW: & {{\dnCh}{\dnCh}{\dnCh}{\dnCh}{\dnCh}{\dnCh}{\dnCh}{\dnCh}{\dnCh}{\dnCh}{\dnCh}{\dnCh}{\dnCh}{\dnCh}{\dnCh}{\dnCh}{\dnCh}{\dnCh}{\dnCh}{\dnCh}{\dnCh}{\dnCh}{\dnCh}{\dnCh}{\dnCh}{\dnCh}{\dnCh}{\dnCh}{\dnCh}{\dnCh}{\dnCh}{\dnCh}} & \dnW[E]{}: & {{\dnCh}{\dnCh}{\dnCh}{\dnCh}{\dnCh}{\dnCh}{\dnCh}{\dnCh}{\dnCh}{\dnCh}{\dnCh}{\dnCh}{\dnCh}{\dnCh}{\dnCh}{\dnCh}{\dnCh}{\dnCh}{\dnCh}{\dnCh}{\dnCh}{\dnCh}{\dnCh}{\dnCh}{\dnCh}{\dnCh}{\dnCh}{\dnCh}{\dnCh}{\dnCh}{\dnCh}{\dnCh}} & \dnW[W]{}: & {{\dnCh}{\dnCh}{\dnCh}{\dnCh}{\dnCh}{\dnCh}{\dnCh}{\dnCh}{\dnCh}{\dnCh}{\dnCh}{\dnCh}{\dnCh}{\dnCh}{\dnCh}{\dnCh}{\dnCh}{\dnCh}{\dnCh}{\dnCh}{\dnCh}{\dnCh}{\dnCh}{\dnCh}{\dnCh}{\dnCh}{\dnCh}{\dnCh}{\dnCh}{\dnCh}{\dnCh}{\dnCh}} \\
 \dnI{13} & \dnW: & {{\dnCh}{\dnCh}{\dnCh}{\dnCh}{\dnCh}{\dnCh}{\dnCh}{\dnCh}{\dnCh}{\dnCh}{\dnCh}{\dnCh}{\dnCh}{\dnCh}{\dnCh}{\dnCh}{\dnCh}{\dnCh}{\dnCh}{\dnCh}{\dnCh}{\dnCh}{\dnCh}{\dnCh}{\dnCh}{\dnCh}{\dnCh}{\dnCh}{\dnCh}{\dnCh}{\dnCh}{\dnCh}} & \dnW[E]{}: & {{\dnCh}{\dnCh}{\dnCh}{\dnCh}{\dnCh}{\dnCh}{\dnCh}{\dnCh}{\dnCh}{\dnCh}{\dnCh}{\dnCh}{\dnCh}{\dnCh}{\dnCh}{\dnCh}{\dnCh}{\dnCh}{\dnCh}{\dnCh}{\dnCh}{\dnCh}{\dnCh}{\dnCh}{\dnCh}{\dnCh}{\dnCh}{\dnCh}{\dnCh}{\dnCh}{\dnCh}{\dnCh}} & \dnW[W]{}: & {{\dnCh}{\dnCh}{\dnCh}{\dnCh}{\dnCh}{\dnCh}{\dnCh}{\dnCh}{\dnCh}{\dnCh}{\dnCh}{\dnCh}{\dnCh}{\dnCh}{\dnCh}{\dnCh}{\dnCh}{\dnCh}{\dnCh}{\dnCh}{\dnCh}{\dnCh}{\dnCh}{\dnCh}{\dnCh}{\dnCh}{\dnCh}{\dnCh}{\dnCh}{\dnCh}{\dnCh}{\dnCh}} \\
 \dnI{14} & \dnW: & {{\dnCh}{\dnCh}{\dnCh}{\dnCh}{\dnCh}{\dnCh}{\dnCh}{\dnCh}{\dnCh}{\dnCh}{\dnCh}{\dnCh}{\dnCh}{\dnCh}{\dnCh}{\dnCh}{\dnCh}{\dnCh}{\dnCh}{\dnCh}{\dnCh}{\dnCh}{\dnCh}{\dnCh}{\dnCh}{\dnCh}{\dnCh}{\dnCh}{\dnCh}{\dnCh}{\dnCh}{\dnCh}} & \dnW[E]{}: & {{\dnCh}{\dnCh}{\dnCh}{\dnCh}{\dnCh}{\dnCh}{\dnCh}{\dnCh}{\dnCh}{\dnCh}{\dnCh}{\dnCh}{\dnCh}{\dnCh}{\dnCh}{\dnCh}{\dnCh}{\dnCh}{\dnCh}{\dnCh}{\dnCh}{\dnCh}{\dnCh}{\dnCh}{\dnCh}{\dnCh}{\dnCh}{\dnCh}{\dnCh}{\dnCh}{\dnCh}{\dnCh}} & \dnW[W]{}: & {{\dnCh}{\dnCh}{\dnCh}{\dnCh}{\dnCh}{\dnCh}{\dnCh}{\dnCh}{\dnCh}{\dnCh}{\dnCh}{\dnCh}{\dnCh}{\dnCh}{\dnCh}{\dnCh}{\dnCh}{\dnCh}{\dnCh}{\dnCh}{\dnCh}{\dnCh}{\dnCh}{\dnCh}{\dnCh}{\dnCh}{\dnCh}{\dnCh}{\dnCh}{\dnCh}{\dnCh}{\dnCh}} \\
 \dnI{15} & \dnW: & {{\dnCh}{\dnCh}{\dnCh}{\dnCh}{\dnCh}{\dnCh}{\dnCh}{\dnCh}{\dnCh}{\dnCh}{\dnCh}{\dnCh}{\dnCh}{\dnCh}{\dnCh}{\dnCh}{\dnCh}{\dnCh}{\dnCh}{\dnCh}{\dnCh}{\dnCh}{\dnCh}{\dnCh}{\dnCh}{\dnCh}{\dnCh}{\dnCh}{\dnCh}{\dnCh}{\dnCh}{\dnCh}} & \dnW[E]{}: & {{\dnCh}{\dnCh}{\dnCh}{\dnCh}{\dnCh}{\dnCh}{\dnCh}{\dnCh}{\dnCh}{\dnCh}{\dnCh}{\dnCh}{\dnCh}{\dnCh}{\dnCh}{\dnCh}{\dnCh}{\dnCh}{\dnCh}{\dnCh}{\dnCh}{\dnCh}{\dnCh}{\dnCh}{\dnCh}{\dnCh}{\dnCh}{\dnCh}{\dnCh}{\dnCh}{\dnCh}{\dnCh}} & \dnW[W]{}: & {{\dnCh}{\dnCh}{\dnCh}{\dnCh}{\dnCh}{\dnCh}{\dnCh}{\dnCh}{\dnCh}{\dnCh}{\dnCh}{\dnCh}{\dnCh}{\dnCh}{\dnCh}{\dnCh}{\dnCh}{\dnCh}{\dnCh}{\dnCh}{\dnCh}{\dnCh}{\dnCh}{\dnCh}{\dnCh}{\dnCh}{\dnCh}{\dnCh}{\dnCh}{\dnCh}{\dnCh}{\dnCh}} \\
 \dnI{16} & \dnW: & {{\dnCh}{\dnCh}{\dnCh}{\dnCh}{\dnCh}{\dnCh}{\dnCh}{\dnCh}{\dnCh}{\dnCh}{\dnCh}{\dnCh}{\dnCh}{\dnCh}{\dnCh}{\dnCh}{\dnCh}{\dnCh}{\dnCh}{\dnCh}{\dnCh}{\dnCh}{\dnCh}{\dnCh}{\dnCh}{\dnCh}{\dnCh}{\dnCh}{\dnCh}{\dnCh}{\dnCh}{\dnCh}} & \dnW[E]{}: & {{\dnCh}{\dnCh}{\dnCh}{\dnCh}{\dnCh}{\dnCh}{\dnCh}{\dnCh}{\dnCh}{\dnCh}{\dnCh}{\dnCh}{\dnCh}{\dnCh}{\dnCh}{\dnCh}{\dnCh}{\dnCh}{\dnCh}{\dnCh}{\dnCh}{\dnCh}{\dnCh}{\dnCh}{\dnCh}{\dnCh}{\dnCh}{\dnCh}{\dnCh}{\dnCh}{\dnCh}{\dnCh}} & \dnW[W]{}: & {{\dnCh}{\dnCh}{\dnCh}{\dnCh}{\dnCh}{\dnCh}{\dnCh}{\dnCh}{\dnCh}{\dnCh}{\dnCh}{\dnCh}{\dnCh}{\dnCh}{\dnCh}{\dnCh}{\dnCh}{\dnCh}{\dnCh}{\dnCh}{\dnCh}{\dnCh}{\dnCh}{\dnCh}{\dnCh}{\dnCh}{\dnCh}{\dnCh}{\dnCh}{\dnCh}{\dnCh}{\dnCh}} \\
 \dnI{17} & \dnW: & {{\dnCh}{\dnCh}{\dnCh}{\dnCh}{\dnCh}{\dnCh}{\dnCh}{\dnCh}{\dnCh}{\dnCh}{\dnCh}{\dnCh}{\dnCh}{\dnCh}{\dnCh}{\dnCh}{\dnCh}{\dnCh}{\dnCh}{\dnCh}{\dnCh}{\dnCh}{\dnCh}{\dnCh}{\dnCh}{\dnCh}{\dnCh}{\dnCh}{\dnCh}{\dnCh}{\dnCh}{\dnCh}} & \dnW[E]{}: & {{\dnCh}{\dnCh}{\dnCh}{\dnCh}{\dnCh}{\dnCh}{\dnCh}{\dnCh}{\dnCh}{\dnCh}{\dnCh}{\dnCh}{\dnCh}{\dnCh}{\dnCh}{\dnCh}{\dnCh}{\dnCh}{\dnCh}{\dnCh}{\dnCh}{\dnCh}{\dnCh}{\dnCh}{\dnCh}{\dnCh}{\dnCh}{\dnCh}{\dnCh}{\dnCh}{\dnCh}{\dnCh}} & \dnW[W]{}: & {{\dnCh}{\dnCh}{\dnCh}{\dnCh}{\dnCh}{\dnCh}{\dnCh}{\dnCh}{\dnCh}{\dnCh}{\dnCh}{\dnCh}{\dnCh}{\dnCh}{\dnCh}{\dnCh}{\dnCh}{\dnCh}{\dnCh}{\dnCh}{\dnCh}{\dnCh}{\dnCh}{\dnCh}{\dnCh}{\dnCh}{\dnCh}{\dnCh}{\dnCh}{\dnCh}{\dnCh}{\dnCh}} \\
\hline
\end{tabular}
}
\caption{SHA256 hash collision over 18 rounds}
\label{fig:tc18}
\end{center}
\end{figure}
}

\texttt{
\begin{figure}[!ht]
\begin{center}
\scalebox{.55}{
\begin{tabular}{|r|c|c|c|c|c|c|}
\hline
  $i$  &      &  A  &    &  E &  & W \\
\hline
 \dnI{-4} & \dnW: & {{\dnCh}{\dnCh}{\dnCh}{\dnCh}{\dnCh}{\dnCh}{\dnCh}{\dnCh}{\dnCh}{\dnCh}{\dnCh}{\dnCh}{\dnCh}{\dnCh}{\dnCh}{\dnCh}{\dnCh}{\dnCh}{\dnCh}{\dnCh}{\dnCh}{\dnCh}{\dnCh}{\dnCh}{\dnCh}{\dnCh}{\dnCh}{\dnCh}{\dnCh}{\dnCh}{\dnCh}{\dnCh}} & \dnW[E]{}: & {{\dnCh}{\dnCh}{\dnCh}{\dnCh}{\dnCh}{\dnCh}{\dnCh}{\dnCh}{\dnCh}{\dnCh}{\dnCh}{\dnCh}{\dnCh}{\dnCh}{\dnCh}{\dnCh}{\dnCh}{\dnCh}{\dnCh}{\dnCh}{\dnCh}{\dnCh}{\dnCh}{\dnCh}{\dnCh}{\dnCh}{\dnCh}{\dnCh}{\dnCh}{\dnCh}{\dnCh}{\dnCh}} & &     \\
 \dnI{-3} & \dnW: & {{\dnCh}{\dnCh}{\dnCh}{\dnCh}{\dnCh}{\dnCh}{\dnCh}{\dnCh}{\dnCh}{\dnCh}{\dnCh}{\dnCh}{\dnCh}{\dnCh}{\dnCh}{\dnCh}{\dnCh}{\dnCh}{\dnCh}{\dnCh}{\dnCh}{\dnCh}{\dnCh}{\dnCh}{\dnCh}{\dnCh}{\dnCh}{\dnCh}{\dnCh}{\dnCh}{\dnCh}{\dnCh}} & \dnW[E]{}: & {{\dnCh}{\dnCh}{\dnCh}{\dnCh}{\dnCh}{\dnCh}{\dnCh}{\dnCh}{\dnCh}{\dnCh}{\dnCh}{\dnCh}{\dnCh}{\dnCh}{\dnCh}{\dnCh}{\dnCh}{\dnCh}{\dnCh}{\dnCh}{\dnCh}{\dnCh}{\dnCh}{\dnCh}{\dnCh}{\dnCh}{\dnCh}{\dnCh}{\dnCh}{\dnCh}{\dnCh}{\dnCh}} & &     \\
 \dnI{-2} & \dnW: & {{\dnCh}{\dnCh}{\dnCh}{\dnCh}{\dnCh}{\dnCh}{\dnCh}{\dnCh}{\dnCh}{\dnCh}{\dnCh}{\dnCh}{\dnCh}{\dnCh}{\dnCh}{\dnCh}{\dnCh}{\dnCh}{\dnCh}{\dnCh}{\dnCh}{\dnCh}{\dnCh}{\dnCh}{\dnCh}{\dnCh}{\dnCh}{\dnCh}{\dnCh}{\dnCh}{\dnCh}{\dnCh}} & \dnW[E]{}: & {{\dnCh}{\dnCh}{\dnCh}{\dnCh}{\dnCh}{\dnCh}{\dnCh}{\dnCh}{\dnCh}{\dnCh}{\dnCh}{\dnCh}{\dnCh}{\dnCh}{\dnCh}{\dnCh}{\dnCh}{\dnCh}{\dnCh}{\dnCh}{\dnCh}{\dnCh}{\dnCh}{\dnCh}{\dnCh}{\dnCh}{\dnCh}{\dnCh}{\dnCh}{\dnCh}{\dnCh}{\dnCh}} & &     \\
 \dnI{-1} & \dnW: & {{\dnCh}{\dnCh}{\dnCh}{\dnCh}{\dnCh}{\dnCh}{\dnCh}{\dnCh}{\dnCh}{\dnCh}{\dnCh}{\dnCh}{\dnCh}{\dnCh}{\dnCh}{\dnCh}{\dnCh}{\dnCh}{\dnCh}{\dnCh}{\dnCh}{\dnCh}{\dnCh}{\dnCh}{\dnCh}{\dnCh}{\dnCh}{\dnCh}{\dnCh}{\dnCh}{\dnCh}{\dnCh}} & \dnW[E]{}: & {{\dnCh}{\dnCh}{\dnCh}{\dnCh}{\dnCh}{\dnCh}{\dnCh}{\dnCh}{\dnCh}{\dnCh}{\dnCh}{\dnCh}{\dnCh}{\dnCh}{\dnCh}{\dnCh}{\dnCh}{\dnCh}{\dnCh}{\dnCh}{\dnCh}{\dnCh}{\dnCh}{\dnCh}{\dnCh}{\dnCh}{\dnCh}{\dnCh}{\dnCh}{\dnCh}{\dnCh}{\dnCh}} & &     \\
 \dnI{0} & \dnW: & {{\dnCh}{\dnCh}{\dnCh}{\dnCh}{\dnCh}{\dnCh}{\dnCh}{\dnCh}{\dnCh}{\dnCh}{\dnCh}{\dnCh}{\dnCh}{\dnCh}{\dnCh}{\dnCh}{\dnCh}{\dnCh}{\dnCh}{\dnCh}{\dnCh}{\dnCh}{\dnCh}{\dnCh}{\dnCh}{\dnCh}{\dnCh}{\dnCh}{\dnCh}{\dnCh}{\dnCh}{\dnCh}} & \dnW[E]{}: & {{\dnCh}{\dnCh}{\dnCh}{\dnCh}{\dnCh}{\dnCh}{\dnCh}{\dnCh}{\dnCh}{\dnCh}{\dnCh}{\dnCh}{\dnCh}{\dnCh}{\dnCh}{\dnCh}{\dnCh}{\dnCh}{\dnCh}{\dnCh}{\dnCh}{\dnCh}{\dnCh}{\dnCh}{\dnCh}{\dnCh}{\dnCh}{\dnCh}{\dnCh}{\dnCh}{\dnCh}{\dnCh}} & \dnW[W]{}: & {{\dnCh}{\dnCh}{\dnCh}{\dnCh}{\dnCh}{\dnCh}{\dnCh}{\dnCh}{\dnCh}{\dnCh}{\dnCh}{\dnCh}{\dnCh}{\dnCh}{\dnCh}{\dnCh}{\dnCh}{\dnCh}{\dnCh}{\dnCh}{\dnCh}{\dnCh}{\dnCh}{\dnCh}{\dnCh}{\dnCh}{\dnCh}{\dnCh}{\dnCh}{\dnCh}{\dnCh}{\dnCh}} \\
 \dnI{1} & \dnW: & {{\dnCh}{\dnCh}{\dnCh}{\dnCh}{\dnCh}{\dnCh}{\dnCh}{\dnCh}{\dnCh}{\dnCh}{\dnCh}{\dnCh}{\dnCh}{\dnCh}{\dnCh}{\dnCh}{\dnCh}{\dnCh}{\dnCh}{\dnCh}{\dnCh}{\dnCh}{\dnCh}{\dnCh}{\dnCh}{\dnCh}{\dnCh}{\dnCh}{\dnCh}{\dnCh}{\dnCh}{\dnCh}} & \dnW[E]{}: & {{\dnCh}{\dnCh}{\dnCh}{\dnCh}{\dnCh}{\dnCh}{\dnCh}{\dnCh}{\dnCh}{\dnCh}{\dnCh}{\dnCh}{\dnCh}{\dnCh}{\dnCh}{\dnCh}{\dnCh}{\dnCh}{\dnCh}{\dnCh}{\dnCh}{\dnCh}{\dnCh}{\dnCh}{\dnCh}{\dnCh}{\dnCh}{\dnCh}{\dnCh}{\dnCh}{\dnCh}{\dnCh}} & \dnW[W]{}: & {{\dnCh}{\dnCh}{\dnCh}{\dnCh}{\dnCh}{\dnCh}{\dnCh}{\dnCh}{\dnCh}{\dnCh}{\dnCh}{\dnCh}{\dnCh}{\dnCh}{\dnCh}{\dnCh}{\dnCh}{\dnCh}{\dnCh}{\dnCh}{\dnCh}{\dnCh}{\dnCh}{\dnCh}{\dnCh}{\dnCh}{\dnCh}{\dnCh}{\dnCh}{\dnCh}{\dnCh}{\dnCh}} \\
 \dnI{2} & \dnW: & {{\dnCh}{\dnCh}{\dnCh}{\dnCh}{\dnCh}{\dnCh}{\dnCh}{\dnCh}{\dnCh}{\dnCh}{\dnCh}{\dnCh}{\dnCh}{\dnCh}{\dnCh}{\dnCh}{\dnCh}{\dnCh}{\dnCh}{\dnCh}{\dnCh}{\dnCh}{\dnCh}{\dnCh}{\dnCh}{\dnCh}{\dnCh}{\dnCh}{\dnCh}{\dnCh}{\dnCh}{\dnCh}} & \dnW[E]{}: & {{\dnCh}{\dnCh}{\dnCh}{\dnCh}{\dnCh}{\dnCh}{\dnCh}{\dnCh}{\dnCh}{\dnCh}{\dnCh}{\dnCh}{\dnCh}{\dnCh}{\dnCh}{\dnCh}{\dnCh}{\dnCh}{\dnCh}{\dnCh}{\dnCh}{\dnCh}{\dnCh}{\dnCh}{\dnCh}{\dnCh}{\dnCh}{\dnCh}{\dnCh}{\dnCh}{\dnCh}{\dnCh}} & \dnW[W]{}: & {{\dnCh}{\dnCh}{\dnCh}{\dnCh}{\dnCh}{\dnCh}{\dnCh}{\dnCh}{\dnCh}{\dnCh}{\dnCh}{\dnCh}{\dnCh}{\dnCh}{\dnCh}{\dnCh}{\dnCh}{\dnCh}{\dnCh}{\dnCh}{\dnCh}{\dnCh}{\dnCh}{\dnCh}{\dnCh}{\dnCh}{\dnCh}{\dnCh}{\dnCh}{\dnCh}{\dnCh}{\dnCh}} \\
 \dnI{3} & \dnW: & {{\dnCh}{\dnCh}{\dnCh}{\dnCh}{\dnCh}{\dnCh}{\dnCh}{\dnCh}{\dnCh}{\dnCh}{\dnCh}{\dnCh}{\dnCh}{\dnCh}{\dnCh}{\dnCh}{\dnCh}{\dnCh}{\dnCh}{\dnCh}{\dnCh}{\dnCh}{\dnCh}{\dnCh}{\dnCh}{\dnCh}{\dnCh}{\dnCh}{\dnCh}{\dnCh}{\dnCh}{\dnCh}} & \dnW[E]{}: & {{\dnCh}{\dnCh}{\dnCh}{\dnCh}{\dnCh}{\dnCh}{\dnCh}{\dnCh}{\dnCh}{\dnCh}{\dnCh}{\dnCh}{\dnCh}{\dnCh}{\dnCh}{\dnCh}{\dnCh}{\dnCh}{\dnCh}{\dnCh}{\dnCh}{\dnCh}{\dnCh}{\dnCh}{\dnCh}{\dnCh}{\dnCh}{\dnCh}{\dnCh}{\dnCh}{\dnCh}{\dnCh}} & \dnW[W]{}: & {{\dnCh}{\dnCh}{\dnCh}{\dnCh}{\dnCh}{\dnCh}{\dnCh}{\dnCh}{\dnCh}{\dnCh}{\dnCh}{\dnCh}{\dnCh}{\dnCh}{\dnCh}{\dnCh}{\dnCh}{\dnCh}{\dnCh}{\dnCh}{\dnCh}{\dnCh}{\dnCh}{\dnCh}{\dnCh}{\dnCh}{\dnCh}{\dnCh}{\dnCh}{\dnCh}{\dnCh}{\dnCh}} \\
 \dnI{4} & \dnW: & {{\dnCh}{\dnCh}{\dnCh}{\dnCh}{\dnCh}{\dnCh}{\dnCh}{\dnCh}{\dnCh}{\dnCh}{\dnCh}{\dnCh}{\dnCh}{\dnCh}{\dnCh}{\dnCh}{\dnCh}{\dnCh}{\dnCh}{\dnCh}{\dnCh}{\dnCh}{\dnCh}{\dnCh}{\dnCh}{\dnCh}{\dnCh}{\dnCh}{\dnCh}{\dnCh}{\dnCh}{\dnCh}} & \dnW[E]{}: & {{\dnCh}{\dnCh}{\dnCh}{\dnCh}{\dnCh}{\dnCh}{\dnCh}{\dnCh}{\dnCh}{\dnCh}{\dnCh}{\dnCh}{\dnCh}{\dnCh}{\dnCh}{\dnCh}{\dnCh}{\dnCh}{\dnCh}{\dnCh}{\dnCh}{\dnCh}{\dnCh}{\dnCh}{\dnCh}{\dnCh}{\dnCh}{\dnCh}{\dnCh}{\dnCh}{\dnCh}{\dnCh}} & \dnW[W]{}: & {{\dnCh}{\dnCh}{\dnCh}{\dnCh}{\dnCh}{\dnCh}{\dnCh}{\dnCh}{\dnCh}{\dnCh}{\dnCh}{\dnCh}{\dnCh}{\dnCh}{\dnCh}{\dnCh}{\dnCh}{\dnCh}{\dnCh}{\dnCh}{\dnCh}{\dnCh}{\dnCh}{\dnCh}{\dnCh}{\dnCh}{\dnCh}{\dnCh}{\dnCh}{\dnCh}{\dnCh}{\dnCh}} \\
 \dnI{5} & \dnW: & {{\dnCx}{\dnCq}{\dnCq}{\dnCq}{\dnCq}{\dnCq}{\dnCq}{\dnCq}{\dnCq}{\dnCq}{\dnCq}{\dnCq}{\dnCq}{\dnCq}{\dnCq}{\dnCq}{\dnCq}{\dnCq}{\dnCq}{\dnCq}{\dnCq}{\dnCq}{\dnCq}{\dnCq}{\dnCq}{\dnCq}{\dnCq}{\dnCq}{\dnCq}{\dnCq}{\dnCq}{\dnCq}} & \dnW[E]{}: & {{\dnCq}{\dnCq}{\dnCq}{\dnCq}{\dnCq}{\dnCq}{\dnCq}{\dnCq}{\dnCq}{\dnCq}{\dnCq}{\dnCq}{\dnCq}{\dnCq}{\dnCq}{\dnCq}{\dnCq}{\dnCq}{\dnCq}{\dnCq}{\dnCq}{\dnCq}{\dnCq}{\dnCq}{\dnCq}{\dnCq}{\dnCq}{\dnCq}{\dnCq}{\dnCq}{\dnCq}{\dnCq}} & \dnW[W]{}: & {{\dnCq}{\dnCq}{\dnCq}{\dnCq}{\dnCq}{\dnCq}{\dnCq}{\dnCq}{\dnCq}{\dnCq}{\dnCq}{\dnCq}{\dnCq}{\dnCq}{\dnCq}{\dnCq}{\dnCq}{\dnCq}{\dnCq}{\dnCq}{\dnCq}{\dnCq}{\dnCq}{\dnCq}{\dnCq}{\dnCq}{\dnCq}{\dnCq}{\dnCq}{\dnCq}{\dnCq}{\dnCq}} \\
 \dnI{6} & \dnW: & {{\dnCh}{\dnCh}{\dnCh}{\dnCh}{\dnCh}{\dnCh}{\dnCh}{\dnCh}{\dnCh}{\dnCh}{\dnCh}{\dnCh}{\dnCh}{\dnCh}{\dnCh}{\dnCh}{\dnCh}{\dnCh}{\dnCh}{\dnCh}{\dnCh}{\dnCh}{\dnCh}{\dnCh}{\dnCh}{\dnCh}{\dnCh}{\dnCh}{\dnCh}{\dnCh}{\dnCh}{\dnCh}} & \dnW[E]{}: & {{\dnCq}{\dnCq}{\dnCq}{\dnCq}{\dnCq}{\dnCq}{\dnCq}{\dnCq}{\dnCq}{\dnCq}{\dnCq}{\dnCq}{\dnCq}{\dnCq}{\dnCq}{\dnCq}{\dnCq}{\dnCq}{\dnCq}{\dnCq}{\dnCq}{\dnCq}{\dnCq}{\dnCq}{\dnCq}{\dnCq}{\dnCq}{\dnCq}{\dnCq}{\dnCq}{\dnCq}{\dnCq}} & \dnW[W]{}: & {{\dnCq}{\dnCq}{\dnCq}{\dnCq}{\dnCq}{\dnCq}{\dnCq}{\dnCq}{\dnCq}{\dnCq}{\dnCq}{\dnCq}{\dnCq}{\dnCq}{\dnCq}{\dnCq}{\dnCq}{\dnCq}{\dnCq}{\dnCq}{\dnCq}{\dnCq}{\dnCq}{\dnCq}{\dnCq}{\dnCq}{\dnCq}{\dnCq}{\dnCq}{\dnCq}{\dnCq}{\dnCq}} \\
 \dnI{7} & \dnW: & {{\dnCh}{\dnCh}{\dnCh}{\dnCh}{\dnCh}{\dnCh}{\dnCh}{\dnCh}{\dnCh}{\dnCh}{\dnCh}{\dnCh}{\dnCh}{\dnCh}{\dnCh}{\dnCh}{\dnCh}{\dnCh}{\dnCh}{\dnCh}{\dnCh}{\dnCh}{\dnCh}{\dnCh}{\dnCh}{\dnCh}{\dnCh}{\dnCh}{\dnCh}{\dnCh}{\dnCh}{\dnCh}} & \dnW[E]{}: & {{\dnCq}{\dnCq}{\dnCq}{\dnCq}{\dnCq}{\dnCq}{\dnCq}{\dnCq}{\dnCq}{\dnCq}{\dnCq}{\dnCq}{\dnCq}{\dnCq}{\dnCq}{\dnCq}{\dnCq}{\dnCq}{\dnCq}{\dnCq}{\dnCq}{\dnCq}{\dnCq}{\dnCq}{\dnCq}{\dnCq}{\dnCq}{\dnCq}{\dnCq}{\dnCq}{\dnCq}{\dnCq}} & \dnW[W]{}: & {{\dnCq}{\dnCq}{\dnCq}{\dnCq}{\dnCq}{\dnCq}{\dnCq}{\dnCq}{\dnCq}{\dnCq}{\dnCq}{\dnCq}{\dnCq}{\dnCq}{\dnCq}{\dnCq}{\dnCq}{\dnCq}{\dnCq}{\dnCq}{\dnCq}{\dnCq}{\dnCq}{\dnCq}{\dnCq}{\dnCq}{\dnCq}{\dnCq}{\dnCq}{\dnCq}{\dnCq}{\dnCq}} \\
 \dnI{8} & \dnW: & {{\dnCh}{\dnCh}{\dnCh}{\dnCh}{\dnCh}{\dnCh}{\dnCh}{\dnCh}{\dnCh}{\dnCh}{\dnCh}{\dnCh}{\dnCh}{\dnCh}{\dnCh}{\dnCh}{\dnCh}{\dnCh}{\dnCh}{\dnCh}{\dnCh}{\dnCh}{\dnCh}{\dnCh}{\dnCh}{\dnCh}{\dnCh}{\dnCh}{\dnCh}{\dnCh}{\dnCh}{\dnCh}} & \dnW[E]{}: & {{\dnCq}{\dnCq}{\dnCq}{\dnCq}{\dnCq}{\dnCq}{\dnCq}{\dnCq}{\dnCq}{\dnCq}{\dnCq}{\dnCq}{\dnCq}{\dnCq}{\dnCq}{\dnCq}{\dnCq}{\dnCq}{\dnCq}{\dnCq}{\dnCq}{\dnCq}{\dnCq}{\dnCq}{\dnCq}{\dnCq}{\dnCq}{\dnCq}{\dnCq}{\dnCq}{\dnCq}{\dnCq}} & \dnW[W]{}: & {{\dnCq}{\dnCq}{\dnCq}{\dnCq}{\dnCq}{\dnCq}{\dnCq}{\dnCq}{\dnCq}{\dnCq}{\dnCq}{\dnCq}{\dnCq}{\dnCq}{\dnCq}{\dnCq}{\dnCq}{\dnCq}{\dnCq}{\dnCq}{\dnCq}{\dnCq}{\dnCq}{\dnCq}{\dnCq}{\dnCq}{\dnCq}{\dnCq}{\dnCq}{\dnCq}{\dnCq}{\dnCq}} \\
 \dnI{9} & \dnW: & {{\dnCh}{\dnCh}{\dnCh}{\dnCh}{\dnCh}{\dnCh}{\dnCh}{\dnCh}{\dnCh}{\dnCh}{\dnCh}{\dnCh}{\dnCh}{\dnCh}{\dnCh}{\dnCh}{\dnCh}{\dnCh}{\dnCh}{\dnCh}{\dnCh}{\dnCh}{\dnCh}{\dnCh}{\dnCh}{\dnCh}{\dnCh}{\dnCh}{\dnCh}{\dnCh}{\dnCh}{\dnCh}} & \dnW[E]{}: & {{\dnCq}{\dnCq}{\dnCq}{\dnCq}{\dnCq}{\dnCq}{\dnCq}{\dnCq}{\dnCq}{\dnCq}{\dnCq}{\dnCq}{\dnCq}{\dnCq}{\dnCq}{\dnCq}{\dnCq}{\dnCq}{\dnCq}{\dnCq}{\dnCq}{\dnCq}{\dnCq}{\dnCq}{\dnCq}{\dnCq}{\dnCq}{\dnCq}{\dnCq}{\dnCq}{\dnCq}{\dnCq}} & \dnW[W]{}: & {{\dnCh}{\dnCh}{\dnCh}{\dnCh}{\dnCh}{\dnCh}{\dnCh}{\dnCh}{\dnCh}{\dnCh}{\dnCh}{\dnCh}{\dnCh}{\dnCh}{\dnCh}{\dnCh}{\dnCh}{\dnCh}{\dnCh}{\dnCh}{\dnCh}{\dnCh}{\dnCh}{\dnCh}{\dnCh}{\dnCh}{\dnCh}{\dnCh}{\dnCh}{\dnCh}{\dnCh}{\dnCh}} \\
 \dnI{10} & \dnW: & {{\dnCh}{\dnCh}{\dnCh}{\dnCh}{\dnCh}{\dnCh}{\dnCh}{\dnCh}{\dnCh}{\dnCh}{\dnCh}{\dnCh}{\dnCh}{\dnCh}{\dnCh}{\dnCh}{\dnCh}{\dnCh}{\dnCh}{\dnCh}{\dnCh}{\dnCh}{\dnCh}{\dnCh}{\dnCh}{\dnCh}{\dnCh}{\dnCh}{\dnCh}{\dnCh}{\dnCh}{\dnCh}} & \dnW[E]{}: & {{\dnCh}{\dnCh}{\dnCh}{\dnCh}{\dnCh}{\dnCh}{\dnCh}{\dnCh}{\dnCh}{\dnCh}{\dnCh}{\dnCh}{\dnCh}{\dnCh}{\dnCh}{\dnCh}{\dnCh}{\dnCh}{\dnCh}{\dnCh}{\dnCh}{\dnCh}{\dnCh}{\dnCh}{\dnCh}{\dnCh}{\dnCh}{\dnCh}{\dnCh}{\dnCh}{\dnCh}{\dnCh}} & \dnW[W]{}: & {{\dnCh}{\dnCh}{\dnCh}{\dnCh}{\dnCh}{\dnCh}{\dnCh}{\dnCh}{\dnCh}{\dnCh}{\dnCh}{\dnCh}{\dnCh}{\dnCh}{\dnCh}{\dnCh}{\dnCh}{\dnCh}{\dnCh}{\dnCh}{\dnCh}{\dnCh}{\dnCh}{\dnCh}{\dnCh}{\dnCh}{\dnCh}{\dnCh}{\dnCh}{\dnCh}{\dnCh}{\dnCh}} \\
 \dnI{11} & \dnW: & {{\dnCh}{\dnCh}{\dnCh}{\dnCh}{\dnCh}{\dnCh}{\dnCh}{\dnCh}{\dnCh}{\dnCh}{\dnCh}{\dnCh}{\dnCh}{\dnCh}{\dnCh}{\dnCh}{\dnCh}{\dnCh}{\dnCh}{\dnCh}{\dnCh}{\dnCh}{\dnCh}{\dnCh}{\dnCh}{\dnCh}{\dnCh}{\dnCh}{\dnCh}{\dnCh}{\dnCh}{\dnCh}} & \dnW[E]{}: & {{\dnCh}{\dnCh}{\dnCh}{\dnCh}{\dnCh}{\dnCh}{\dnCh}{\dnCh}{\dnCh}{\dnCh}{\dnCh}{\dnCh}{\dnCh}{\dnCh}{\dnCh}{\dnCh}{\dnCh}{\dnCh}{\dnCh}{\dnCh}{\dnCh}{\dnCh}{\dnCh}{\dnCh}{\dnCh}{\dnCh}{\dnCh}{\dnCh}{\dnCh}{\dnCh}{\dnCh}{\dnCh}} & \dnW[W]{}: & {{\dnCh}{\dnCh}{\dnCh}{\dnCh}{\dnCh}{\dnCh}{\dnCh}{\dnCh}{\dnCh}{\dnCh}{\dnCh}{\dnCh}{\dnCh}{\dnCh}{\dnCh}{\dnCh}{\dnCh}{\dnCh}{\dnCh}{\dnCh}{\dnCh}{\dnCh}{\dnCh}{\dnCh}{\dnCh}{\dnCh}{\dnCh}{\dnCh}{\dnCh}{\dnCh}{\dnCh}{\dnCh}} \\
 \dnI{12} & \dnW: & {{\dnCh}{\dnCh}{\dnCh}{\dnCh}{\dnCh}{\dnCh}{\dnCh}{\dnCh}{\dnCh}{\dnCh}{\dnCh}{\dnCh}{\dnCh}{\dnCh}{\dnCh}{\dnCh}{\dnCh}{\dnCh}{\dnCh}{\dnCh}{\dnCh}{\dnCh}{\dnCh}{\dnCh}{\dnCh}{\dnCh}{\dnCh}{\dnCh}{\dnCh}{\dnCh}{\dnCh}{\dnCh}} & \dnW[E]{}: & {{\dnCh}{\dnCh}{\dnCh}{\dnCh}{\dnCh}{\dnCh}{\dnCh}{\dnCh}{\dnCh}{\dnCh}{\dnCh}{\dnCh}{\dnCh}{\dnCh}{\dnCh}{\dnCh}{\dnCh}{\dnCh}{\dnCh}{\dnCh}{\dnCh}{\dnCh}{\dnCh}{\dnCh}{\dnCh}{\dnCh}{\dnCh}{\dnCh}{\dnCh}{\dnCh}{\dnCh}{\dnCh}} & \dnW[W]{}: & {{\dnCh}{\dnCh}{\dnCh}{\dnCh}{\dnCh}{\dnCh}{\dnCh}{\dnCh}{\dnCh}{\dnCh}{\dnCh}{\dnCh}{\dnCh}{\dnCh}{\dnCh}{\dnCh}{\dnCh}{\dnCh}{\dnCh}{\dnCh}{\dnCh}{\dnCh}{\dnCh}{\dnCh}{\dnCh}{\dnCh}{\dnCh}{\dnCh}{\dnCh}{\dnCh}{\dnCh}{\dnCh}} \\
 \dnI{13} & \dnW: & {{\dnCh}{\dnCh}{\dnCh}{\dnCh}{\dnCh}{\dnCh}{\dnCh}{\dnCh}{\dnCh}{\dnCh}{\dnCh}{\dnCh}{\dnCh}{\dnCh}{\dnCh}{\dnCh}{\dnCh}{\dnCh}{\dnCh}{\dnCh}{\dnCh}{\dnCh}{\dnCh}{\dnCh}{\dnCh}{\dnCh}{\dnCh}{\dnCh}{\dnCh}{\dnCh}{\dnCh}{\dnCh}} & \dnW[E]{}: & {{\dnCh}{\dnCh}{\dnCh}{\dnCh}{\dnCh}{\dnCh}{\dnCh}{\dnCh}{\dnCh}{\dnCh}{\dnCh}{\dnCh}{\dnCh}{\dnCh}{\dnCh}{\dnCh}{\dnCh}{\dnCh}{\dnCh}{\dnCh}{\dnCh}{\dnCh}{\dnCh}{\dnCh}{\dnCh}{\dnCh}{\dnCh}{\dnCh}{\dnCh}{\dnCh}{\dnCh}{\dnCh}} & \dnW[W]{}: & {{\dnCq}{\dnCq}{\dnCq}{\dnCq}{\dnCq}{\dnCq}{\dnCq}{\dnCq}{\dnCq}{\dnCq}{\dnCq}{\dnCq}{\dnCq}{\dnCq}{\dnCq}{\dnCq}{\dnCq}{\dnCq}{\dnCq}{\dnCq}{\dnCq}{\dnCq}{\dnCq}{\dnCq}{\dnCq}{\dnCq}{\dnCq}{\dnCq}{\dnCq}{\dnCq}{\dnCq}{\dnCq}} \\
 \dnI{14} & \dnW: & {{\dnCh}{\dnCh}{\dnCh}{\dnCh}{\dnCh}{\dnCh}{\dnCh}{\dnCh}{\dnCh}{\dnCh}{\dnCh}{\dnCh}{\dnCh}{\dnCh}{\dnCh}{\dnCh}{\dnCh}{\dnCh}{\dnCh}{\dnCh}{\dnCh}{\dnCh}{\dnCh}{\dnCh}{\dnCh}{\dnCh}{\dnCh}{\dnCh}{\dnCh}{\dnCh}{\dnCh}{\dnCh}} & \dnW[E]{}: & {{\dnCh}{\dnCh}{\dnCh}{\dnCh}{\dnCh}{\dnCh}{\dnCh}{\dnCh}{\dnCh}{\dnCh}{\dnCh}{\dnCh}{\dnCh}{\dnCh}{\dnCh}{\dnCh}{\dnCh}{\dnCh}{\dnCh}{\dnCh}{\dnCh}{\dnCh}{\dnCh}{\dnCh}{\dnCh}{\dnCh}{\dnCh}{\dnCh}{\dnCh}{\dnCh}{\dnCh}{\dnCh}} & \dnW[W]{}: & {{\dnCh}{\dnCh}{\dnCh}{\dnCh}{\dnCh}{\dnCh}{\dnCh}{\dnCh}{\dnCh}{\dnCh}{\dnCh}{\dnCh}{\dnCh}{\dnCh}{\dnCh}{\dnCh}{\dnCh}{\dnCh}{\dnCh}{\dnCh}{\dnCh}{\dnCh}{\dnCh}{\dnCh}{\dnCh}{\dnCh}{\dnCh}{\dnCh}{\dnCh}{\dnCh}{\dnCh}{\dnCh}} \\
 \dnI{15} & \dnW: & {{\dnCh}{\dnCh}{\dnCh}{\dnCh}{\dnCh}{\dnCh}{\dnCh}{\dnCh}{\dnCh}{\dnCh}{\dnCh}{\dnCh}{\dnCh}{\dnCh}{\dnCh}{\dnCh}{\dnCh}{\dnCh}{\dnCh}{\dnCh}{\dnCh}{\dnCh}{\dnCh}{\dnCh}{\dnCh}{\dnCh}{\dnCh}{\dnCh}{\dnCh}{\dnCh}{\dnCh}{\dnCh}} & \dnW[E]{}: & {{\dnCh}{\dnCh}{\dnCh}{\dnCh}{\dnCh}{\dnCh}{\dnCh}{\dnCh}{\dnCh}{\dnCh}{\dnCh}{\dnCh}{\dnCh}{\dnCh}{\dnCh}{\dnCh}{\dnCh}{\dnCh}{\dnCh}{\dnCh}{\dnCh}{\dnCh}{\dnCh}{\dnCh}{\dnCh}{\dnCh}{\dnCh}{\dnCh}{\dnCh}{\dnCh}{\dnCh}{\dnCh}} & \dnW[W]{}: & {{\dnCh}{\dnCh}{\dnCh}{\dnCh}{\dnCh}{\dnCh}{\dnCh}{\dnCh}{\dnCh}{\dnCh}{\dnCh}{\dnCh}{\dnCh}{\dnCh}{\dnCh}{\dnCh}{\dnCh}{\dnCh}{\dnCh}{\dnCh}{\dnCh}{\dnCh}{\dnCh}{\dnCh}{\dnCh}{\dnCh}{\dnCh}{\dnCh}{\dnCh}{\dnCh}{\dnCh}{\dnCh}} \\
 \dnI{16} & \dnW: & {{\dnCh}{\dnCh}{\dnCh}{\dnCh}{\dnCh}{\dnCh}{\dnCh}{\dnCh}{\dnCh}{\dnCh}{\dnCh}{\dnCh}{\dnCh}{\dnCh}{\dnCh}{\dnCh}{\dnCh}{\dnCh}{\dnCh}{\dnCh}{\dnCh}{\dnCh}{\dnCh}{\dnCh}{\dnCh}{\dnCh}{\dnCh}{\dnCh}{\dnCh}{\dnCh}{\dnCh}{\dnCh}} & \dnW[E]{}: & {{\dnCh}{\dnCh}{\dnCh}{\dnCh}{\dnCh}{\dnCh}{\dnCh}{\dnCh}{\dnCh}{\dnCh}{\dnCh}{\dnCh}{\dnCh}{\dnCh}{\dnCh}{\dnCh}{\dnCh}{\dnCh}{\dnCh}{\dnCh}{\dnCh}{\dnCh}{\dnCh}{\dnCh}{\dnCh}{\dnCh}{\dnCh}{\dnCh}{\dnCh}{\dnCh}{\dnCh}{\dnCh}} & \dnW[W]{}: & {{\dnCh}{\dnCh}{\dnCh}{\dnCh}{\dnCh}{\dnCh}{\dnCh}{\dnCh}{\dnCh}{\dnCh}{\dnCh}{\dnCh}{\dnCh}{\dnCh}{\dnCh}{\dnCh}{\dnCh}{\dnCh}{\dnCh}{\dnCh}{\dnCh}{\dnCh}{\dnCh}{\dnCh}{\dnCh}{\dnCh}{\dnCh}{\dnCh}{\dnCh}{\dnCh}{\dnCh}{\dnCh}} \\
 \dnI{17} & \dnW: & {{\dnCh}{\dnCh}{\dnCh}{\dnCh}{\dnCh}{\dnCh}{\dnCh}{\dnCh}{\dnCh}{\dnCh}{\dnCh}{\dnCh}{\dnCh}{\dnCh}{\dnCh}{\dnCh}{\dnCh}{\dnCh}{\dnCh}{\dnCh}{\dnCh}{\dnCh}{\dnCh}{\dnCh}{\dnCh}{\dnCh}{\dnCh}{\dnCh}{\dnCh}{\dnCh}{\dnCh}{\dnCh}} & \dnW[E]{}: & {{\dnCh}{\dnCh}{\dnCh}{\dnCh}{\dnCh}{\dnCh}{\dnCh}{\dnCh}{\dnCh}{\dnCh}{\dnCh}{\dnCh}{\dnCh}{\dnCh}{\dnCh}{\dnCh}{\dnCh}{\dnCh}{\dnCh}{\dnCh}{\dnCh}{\dnCh}{\dnCh}{\dnCh}{\dnCh}{\dnCh}{\dnCh}{\dnCh}{\dnCh}{\dnCh}{\dnCh}{\dnCh}} & \dnW[W]{}: & {{\dnCh}{\dnCh}{\dnCh}{\dnCh}{\dnCh}{\dnCh}{\dnCh}{\dnCh}{\dnCh}{\dnCh}{\dnCh}{\dnCh}{\dnCh}{\dnCh}{\dnCh}{\dnCh}{\dnCh}{\dnCh}{\dnCh}{\dnCh}{\dnCh}{\dnCh}{\dnCh}{\dnCh}{\dnCh}{\dnCh}{\dnCh}{\dnCh}{\dnCh}{\dnCh}{\dnCh}{\dnCh}} \\
 \dnI{18} & \dnW: & {{\dnCh}{\dnCh}{\dnCh}{\dnCh}{\dnCh}{\dnCh}{\dnCh}{\dnCh}{\dnCh}{\dnCh}{\dnCh}{\dnCh}{\dnCh}{\dnCh}{\dnCh}{\dnCh}{\dnCh}{\dnCh}{\dnCh}{\dnCh}{\dnCh}{\dnCh}{\dnCh}{\dnCh}{\dnCh}{\dnCh}{\dnCh}{\dnCh}{\dnCh}{\dnCh}{\dnCh}{\dnCh}} & \dnW[E]{}: & {{\dnCh}{\dnCh}{\dnCh}{\dnCh}{\dnCh}{\dnCh}{\dnCh}{\dnCh}{\dnCh}{\dnCh}{\dnCh}{\dnCh}{\dnCh}{\dnCh}{\dnCh}{\dnCh}{\dnCh}{\dnCh}{\dnCh}{\dnCh}{\dnCh}{\dnCh}{\dnCh}{\dnCh}{\dnCh}{\dnCh}{\dnCh}{\dnCh}{\dnCh}{\dnCh}{\dnCh}{\dnCh}} & \dnW[W]{}: & {{\dnCh}{\dnCh}{\dnCh}{\dnCh}{\dnCh}{\dnCh}{\dnCh}{\dnCh}{\dnCh}{\dnCh}{\dnCh}{\dnCh}{\dnCh}{\dnCh}{\dnCh}{\dnCh}{\dnCh}{\dnCh}{\dnCh}{\dnCh}{\dnCh}{\dnCh}{\dnCh}{\dnCh}{\dnCh}{\dnCh}{\dnCh}{\dnCh}{\dnCh}{\dnCh}{\dnCh}{\dnCh}} \\
 \dnI{19} & \dnW: & {{\dnCh}{\dnCh}{\dnCh}{\dnCh}{\dnCh}{\dnCh}{\dnCh}{\dnCh}{\dnCh}{\dnCh}{\dnCh}{\dnCh}{\dnCh}{\dnCh}{\dnCh}{\dnCh}{\dnCh}{\dnCh}{\dnCh}{\dnCh}{\dnCh}{\dnCh}{\dnCh}{\dnCh}{\dnCh}{\dnCh}{\dnCh}{\dnCh}{\dnCh}{\dnCh}{\dnCh}{\dnCh}} & \dnW[E]{}: & {{\dnCh}{\dnCh}{\dnCh}{\dnCh}{\dnCh}{\dnCh}{\dnCh}{\dnCh}{\dnCh}{\dnCh}{\dnCh}{\dnCh}{\dnCh}{\dnCh}{\dnCh}{\dnCh}{\dnCh}{\dnCh}{\dnCh}{\dnCh}{\dnCh}{\dnCh}{\dnCh}{\dnCh}{\dnCh}{\dnCh}{\dnCh}{\dnCh}{\dnCh}{\dnCh}{\dnCh}{\dnCh}} & \dnW[W]{}: & {{\dnCh}{\dnCh}{\dnCh}{\dnCh}{\dnCh}{\dnCh}{\dnCh}{\dnCh}{\dnCh}{\dnCh}{\dnCh}{\dnCh}{\dnCh}{\dnCh}{\dnCh}{\dnCh}{\dnCh}{\dnCh}{\dnCh}{\dnCh}{\dnCh}{\dnCh}{\dnCh}{\dnCh}{\dnCh}{\dnCh}{\dnCh}{\dnCh}{\dnCh}{\dnCh}{\dnCh}{\dnCh}} \\
 \dnI{20} & \dnW: & {{\dnCh}{\dnCh}{\dnCh}{\dnCh}{\dnCh}{\dnCh}{\dnCh}{\dnCh}{\dnCh}{\dnCh}{\dnCh}{\dnCh}{\dnCh}{\dnCh}{\dnCh}{\dnCh}{\dnCh}{\dnCh}{\dnCh}{\dnCh}{\dnCh}{\dnCh}{\dnCh}{\dnCh}{\dnCh}{\dnCh}{\dnCh}{\dnCh}{\dnCh}{\dnCh}{\dnCh}{\dnCh}} & \dnW[E]{}: & {{\dnCh}{\dnCh}{\dnCh}{\dnCh}{\dnCh}{\dnCh}{\dnCh}{\dnCh}{\dnCh}{\dnCh}{\dnCh}{\dnCh}{\dnCh}{\dnCh}{\dnCh}{\dnCh}{\dnCh}{\dnCh}{\dnCh}{\dnCh}{\dnCh}{\dnCh}{\dnCh}{\dnCh}{\dnCh}{\dnCh}{\dnCh}{\dnCh}{\dnCh}{\dnCh}{\dnCh}{\dnCh}} & \dnW[W]{}: & {{\dnCh}{\dnCh}{\dnCh}{\dnCh}{\dnCh}{\dnCh}{\dnCh}{\dnCh}{\dnCh}{\dnCh}{\dnCh}{\dnCh}{\dnCh}{\dnCh}{\dnCh}{\dnCh}{\dnCh}{\dnCh}{\dnCh}{\dnCh}{\dnCh}{\dnCh}{\dnCh}{\dnCh}{\dnCh}{\dnCh}{\dnCh}{\dnCh}{\dnCh}{\dnCh}{\dnCh}{\dnCh}} \\
\hline
\end{tabular}
}
\caption{SHA-256 hash collision over 21 rounds}
\label{fig:tc21}
\end{center}
\end{figure}
}

\texttt{
\begin{figure}[!ht]
\begin{center}
\scalebox{.55}{
\begin{tabular}{|r|c|c|c|c|c|c|}
\hline
  $i$  &      &  A  &  &  E &  & W \\
\hline
 \dnI{-4} & \dnW: & {{\dnCh}{\dnCh}{\dnCh}{\dnCh}{\dnCh}{\dnCh}{\dnCh}{\dnCh}{\dnCh}{\dnCh}{\dnCh}{\dnCh}{\dnCh}{\dnCh}{\dnCh}{\dnCh}{\dnCh}{\dnCh}{\dnCh}{\dnCh}{\dnCh}{\dnCh}{\dnCh}{\dnCh}{\dnCh}{\dnCh}{\dnCh}{\dnCh}{\dnCh}{\dnCh}{\dnCh}{\dnCh}} & \dnW[E]{}: & {{\dnCh}{\dnCh}{\dnCh}{\dnCh}{\dnCh}{\dnCh}{\dnCh}{\dnCh}{\dnCh}{\dnCh}{\dnCh}{\dnCh}{\dnCh}{\dnCh}{\dnCh}{\dnCh}{\dnCh}{\dnCh}{\dnCh}{\dnCh}{\dnCh}{\dnCh}{\dnCh}{\dnCh}{\dnCh}{\dnCh}{\dnCh}{\dnCh}{\dnCh}{\dnCh}{\dnCh}{\dnCh}} & &     \\
 \dnI{-3} & \dnW: & {{\dnCh}{\dnCh}{\dnCh}{\dnCh}{\dnCh}{\dnCh}{\dnCh}{\dnCh}{\dnCh}{\dnCh}{\dnCh}{\dnCh}{\dnCh}{\dnCh}{\dnCh}{\dnCh}{\dnCh}{\dnCh}{\dnCh}{\dnCh}{\dnCh}{\dnCh}{\dnCh}{\dnCh}{\dnCh}{\dnCh}{\dnCh}{\dnCh}{\dnCh}{\dnCh}{\dnCh}{\dnCh}} & \dnW[E]{}: & {{\dnCh}{\dnCh}{\dnCh}{\dnCh}{\dnCh}{\dnCh}{\dnCh}{\dnCh}{\dnCh}{\dnCh}{\dnCh}{\dnCh}{\dnCh}{\dnCh}{\dnCh}{\dnCh}{\dnCh}{\dnCh}{\dnCh}{\dnCh}{\dnCh}{\dnCh}{\dnCh}{\dnCh}{\dnCh}{\dnCh}{\dnCh}{\dnCh}{\dnCh}{\dnCh}{\dnCh}{\dnCh}} & &     \\
 \dnI{-2} & \dnW: & {{\dnCh}{\dnCh}{\dnCh}{\dnCh}{\dnCh}{\dnCh}{\dnCh}{\dnCh}{\dnCh}{\dnCh}{\dnCh}{\dnCh}{\dnCh}{\dnCh}{\dnCh}{\dnCh}{\dnCh}{\dnCh}{\dnCh}{\dnCh}{\dnCh}{\dnCh}{\dnCh}{\dnCh}{\dnCh}{\dnCh}{\dnCh}{\dnCh}{\dnCh}{\dnCh}{\dnCh}{\dnCh}} & \dnW[E]{}: & {{\dnCh}{\dnCh}{\dnCh}{\dnCh}{\dnCh}{\dnCh}{\dnCh}{\dnCh}{\dnCh}{\dnCh}{\dnCh}{\dnCh}{\dnCh}{\dnCh}{\dnCh}{\dnCh}{\dnCh}{\dnCh}{\dnCh}{\dnCh}{\dnCh}{\dnCh}{\dnCh}{\dnCh}{\dnCh}{\dnCh}{\dnCh}{\dnCh}{\dnCh}{\dnCh}{\dnCh}{\dnCh}} & &     \\
 \dnI{-1} & \dnW: & {{\dnCh}{\dnCh}{\dnCh}{\dnCh}{\dnCh}{\dnCh}{\dnCh}{\dnCh}{\dnCh}{\dnCh}{\dnCh}{\dnCh}{\dnCh}{\dnCh}{\dnCh}{\dnCh}{\dnCh}{\dnCh}{\dnCh}{\dnCh}{\dnCh}{\dnCh}{\dnCh}{\dnCh}{\dnCh}{\dnCh}{\dnCh}{\dnCh}{\dnCh}{\dnCh}{\dnCh}{\dnCh}} & \dnW[E]{}: & {{\dnCh}{\dnCh}{\dnCh}{\dnCh}{\dnCh}{\dnCh}{\dnCh}{\dnCh}{\dnCh}{\dnCh}{\dnCh}{\dnCh}{\dnCh}{\dnCh}{\dnCh}{\dnCh}{\dnCh}{\dnCh}{\dnCh}{\dnCh}{\dnCh}{\dnCh}{\dnCh}{\dnCh}{\dnCh}{\dnCh}{\dnCh}{\dnCh}{\dnCh}{\dnCh}{\dnCh}{\dnCh}} & &     \\
 \dnI{0} & \dnW: & {{\dnCh}{\dnCh}{\dnCh}{\dnCh}{\dnCh}{\dnCh}{\dnCh}{\dnCh}{\dnCh}{\dnCh}{\dnCh}{\dnCh}{\dnCh}{\dnCh}{\dnCh}{\dnCh}{\dnCh}{\dnCh}{\dnCh}{\dnCh}{\dnCh}{\dnCh}{\dnCh}{\dnCh}{\dnCh}{\dnCh}{\dnCh}{\dnCh}{\dnCh}{\dnCh}{\dnCh}{\dnCh}} & \dnW[E]{}: & {{\dnCh}{\dnCh}{\dnCh}{\dnCh}{\dnCh}{\dnCh}{\dnCh}{\dnCh}{\dnCh}{\dnCh}{\dnCh}{\dnCh}{\dnCh}{\dnCh}{\dnCh}{\dnCh}{\dnCh}{\dnCh}{\dnCh}{\dnCh}{\dnCh}{\dnCh}{\dnCh}{\dnCh}{\dnCh}{\dnCh}{\dnCh}{\dnCh}{\dnCh}{\dnCh}{\dnCh}{\dnCh}} & \dnW[W]{}: & {{\dnCh}{\dnCh}{\dnCh}{\dnCh}{\dnCh}{\dnCh}{\dnCh}{\dnCh}{\dnCh}{\dnCh}{\dnCh}{\dnCh}{\dnCh}{\dnCh}{\dnCh}{\dnCh}{\dnCh}{\dnCh}{\dnCh}{\dnCh}{\dnCh}{\dnCh}{\dnCh}{\dnCh}{\dnCh}{\dnCh}{\dnCh}{\dnCh}{\dnCh}{\dnCh}{\dnCh}{\dnCh}} \\
 \dnI{1} & \dnW: & {{\dnCh}{\dnCh}{\dnCh}{\dnCh}{\dnCh}{\dnCh}{\dnCh}{\dnCh}{\dnCh}{\dnCh}{\dnCh}{\dnCh}{\dnCh}{\dnCh}{\dnCh}{\dnCh}{\dnCh}{\dnCh}{\dnCh}{\dnCh}{\dnCh}{\dnCh}{\dnCh}{\dnCh}{\dnCh}{\dnCh}{\dnCh}{\dnCh}{\dnCh}{\dnCh}{\dnCh}{\dnCh}} & \dnW[E]{}: & {{\dnCh}{\dnCh}{\dnCh}{\dnCh}{\dnCh}{\dnCh}{\dnCh}{\dnCh}{\dnCh}{\dnCh}{\dnCh}{\dnCh}{\dnCh}{\dnCh}{\dnCh}{\dnCh}{\dnCh}{\dnCh}{\dnCh}{\dnCh}{\dnCh}{\dnCh}{\dnCh}{\dnCh}{\dnCh}{\dnCh}{\dnCh}{\dnCh}{\dnCh}{\dnCh}{\dnCh}{\dnCh}} & \dnW[W]{}: & {{\dnCh}{\dnCh}{\dnCh}{\dnCh}{\dnCh}{\dnCh}{\dnCh}{\dnCh}{\dnCh}{\dnCh}{\dnCh}{\dnCh}{\dnCh}{\dnCh}{\dnCh}{\dnCh}{\dnCh}{\dnCh}{\dnCh}{\dnCh}{\dnCh}{\dnCh}{\dnCh}{\dnCh}{\dnCh}{\dnCh}{\dnCh}{\dnCh}{\dnCh}{\dnCh}{\dnCh}{\dnCh}} \\
 \dnI{2} & \dnW: & {{\dnCh}{\dnCh}{\dnCh}{\dnCh}{\dnCh}{\dnCh}{\dnCh}{\dnCh}{\dnCh}{\dnCh}{\dnCh}{\dnCh}{\dnCh}{\dnCh}{\dnCh}{\dnCh}{\dnCh}{\dnCh}{\dnCh}{\dnCh}{\dnCh}{\dnCh}{\dnCh}{\dnCh}{\dnCh}{\dnCh}{\dnCh}{\dnCh}{\dnCh}{\dnCh}{\dnCh}{\dnCh}} & \dnW[E]{}: & {{\dnCh}{\dnCh}{\dnCh}{\dnCh}{\dnCh}{\dnCh}{\dnCh}{\dnCh}{\dnCh}{\dnCh}{\dnCh}{\dnCh}{\dnCh}{\dnCh}{\dnCh}{\dnCh}{\dnCh}{\dnCh}{\dnCh}{\dnCh}{\dnCh}{\dnCh}{\dnCh}{\dnCh}{\dnCh}{\dnCh}{\dnCh}{\dnCh}{\dnCh}{\dnCh}{\dnCh}{\dnCh}} & \dnW[W]{}: & {{\dnCh}{\dnCh}{\dnCh}{\dnCh}{\dnCh}{\dnCh}{\dnCh}{\dnCh}{\dnCh}{\dnCh}{\dnCh}{\dnCh}{\dnCh}{\dnCh}{\dnCh}{\dnCh}{\dnCh}{\dnCh}{\dnCh}{\dnCh}{\dnCh}{\dnCh}{\dnCh}{\dnCh}{\dnCh}{\dnCh}{\dnCh}{\dnCh}{\dnCh}{\dnCh}{\dnCh}{\dnCh}} \\
 \dnI{3} & \dnW: & {{\dnCh}{\dnCh}{\dnCh}{\dnCh}{\dnCh}{\dnCh}{\dnCh}{\dnCh}{\dnCh}{\dnCh}{\dnCh}{\dnCh}{\dnCh}{\dnCh}{\dnCh}{\dnCh}{\dnCh}{\dnCh}{\dnCh}{\dnCh}{\dnCh}{\dnCh}{\dnCh}{\dnCh}{\dnCh}{\dnCh}{\dnCh}{\dnCh}{\dnCh}{\dnCh}{\dnCh}{\dnCh}} & \dnW[E]{}: & {{\dnCh}{\dnCh}{\dnCh}{\dnCh}{\dnCh}{\dnCh}{\dnCh}{\dnCh}{\dnCh}{\dnCh}{\dnCh}{\dnCh}{\dnCh}{\dnCh}{\dnCh}{\dnCh}{\dnCh}{\dnCh}{\dnCh}{\dnCh}{\dnCh}{\dnCh}{\dnCh}{\dnCh}{\dnCh}{\dnCh}{\dnCh}{\dnCh}{\dnCh}{\dnCh}{\dnCh}{\dnCh}} & \dnW[W]{}: & {{\dnCh}{\dnCh}{\dnCh}{\dnCh}{\dnCh}{\dnCh}{\dnCh}{\dnCh}{\dnCh}{\dnCh}{\dnCh}{\dnCh}{\dnCh}{\dnCh}{\dnCh}{\dnCh}{\dnCh}{\dnCh}{\dnCh}{\dnCh}{\dnCh}{\dnCh}{\dnCh}{\dnCh}{\dnCh}{\dnCh}{\dnCh}{\dnCh}{\dnCh}{\dnCh}{\dnCh}{\dnCh}} \\
 \dnI{4} & \dnW: & {{\dnCh}{\dnCh}{\dnCh}{\dnCh}{\dnCh}{\dnCh}{\dnCh}{\dnCh}{\dnCh}{\dnCh}{\dnCh}{\dnCh}{\dnCh}{\dnCh}{\dnCh}{\dnCh}{\dnCh}{\dnCh}{\dnCh}{\dnCh}{\dnCh}{\dnCh}{\dnCh}{\dnCh}{\dnCh}{\dnCh}{\dnCh}{\dnCh}{\dnCh}{\dnCh}{\dnCh}{\dnCh}} & \dnW[E]{}: & {{\dnCh}{\dnCh}{\dnCh}{\dnCh}{\dnCh}{\dnCh}{\dnCh}{\dnCh}{\dnCh}{\dnCh}{\dnCh}{\dnCh}{\dnCh}{\dnCh}{\dnCh}{\dnCh}{\dnCh}{\dnCh}{\dnCh}{\dnCh}{\dnCh}{\dnCh}{\dnCh}{\dnCh}{\dnCh}{\dnCh}{\dnCh}{\dnCh}{\dnCh}{\dnCh}{\dnCh}{\dnCh}} & \dnW[W]{}: & {{\dnCh}{\dnCh}{\dnCh}{\dnCh}{\dnCh}{\dnCh}{\dnCh}{\dnCh}{\dnCh}{\dnCh}{\dnCh}{\dnCh}{\dnCh}{\dnCh}{\dnCh}{\dnCh}{\dnCh}{\dnCh}{\dnCh}{\dnCh}{\dnCh}{\dnCh}{\dnCh}{\dnCh}{\dnCh}{\dnCh}{\dnCh}{\dnCh}{\dnCh}{\dnCh}{\dnCh}{\dnCh}} \\
 \dnI{5} & \dnW: & {{\dnCh}{\dnCh}{\dnCh}{\dnCh}{\dnCh}{\dnCh}{\dnCh}{\dnCh}{\dnCh}{\dnCh}{\dnCh}{\dnCh}{\dnCh}{\dnCh}{\dnCh}{\dnCh}{\dnCh}{\dnCh}{\dnCh}{\dnCh}{\dnCh}{\dnCh}{\dnCh}{\dnCh}{\dnCh}{\dnCh}{\dnCh}{\dnCh}{\dnCh}{\dnCh}{\dnCh}{\dnCh}} & \dnW[E]{}: & {{\dnCh}{\dnCh}{\dnCh}{\dnCh}{\dnCh}{\dnCh}{\dnCh}{\dnCh}{\dnCh}{\dnCh}{\dnCh}{\dnCh}{\dnCh}{\dnCh}{\dnCh}{\dnCh}{\dnCh}{\dnCh}{\dnCh}{\dnCh}{\dnCh}{\dnCh}{\dnCh}{\dnCh}{\dnCh}{\dnCh}{\dnCh}{\dnCh}{\dnCh}{\dnCh}{\dnCh}{\dnCh}} & \dnW[W]{}: & {{\dnCh}{\dnCh}{\dnCh}{\dnCh}{\dnCh}{\dnCh}{\dnCh}{\dnCh}{\dnCh}{\dnCh}{\dnCh}{\dnCh}{\dnCh}{\dnCh}{\dnCh}{\dnCh}{\dnCh}{\dnCh}{\dnCh}{\dnCh}{\dnCh}{\dnCh}{\dnCh}{\dnCh}{\dnCh}{\dnCh}{\dnCh}{\dnCh}{\dnCh}{\dnCh}{\dnCh}{\dnCh}} \\
 \dnI{6} & \dnW: & {{\dnCh}{\dnCh}{\dnCh}{\dnCh}{\dnCh}{\dnCh}{\dnCh}{\dnCh}{\dnCh}{\dnCh}{\dnCh}{\dnCh}{\dnCh}{\dnCh}{\dnCh}{\dnCh}{\dnCh}{\dnCh}{\dnCh}{\dnCh}{\dnCh}{\dnCh}{\dnCh}{\dnCh}{\dnCh}{\dnCh}{\dnCh}{\dnCh}{\dnCh}{\dnCh}{\dnCh}{\dnCh}} & \dnW[E]{}: & {{\dnCh}{\dnCh}{\dnCh}{\dnCh}{\dnCh}{\dnCh}{\dnCh}{\dnCh}{\dnCh}{\dnCh}{\dnCh}{\dnCh}{\dnCh}{\dnCh}{\dnCh}{\dnCh}{\dnCh}{\dnCh}{\dnCh}{\dnCh}{\dnCh}{\dnCh}{\dnCh}{\dnCh}{\dnCh}{\dnCh}{\dnCh}{\dnCh}{\dnCh}{\dnCh}{\dnCh}{\dnCh}} & \dnW[W]{}: & {{\dnCh}{\dnCh}{\dnCh}{\dnCh}{\dnCh}{\dnCh}{\dnCh}{\dnCh}{\dnCh}{\dnCh}{\dnCh}{\dnCh}{\dnCh}{\dnCh}{\dnCh}{\dnCh}{\dnCh}{\dnCh}{\dnCh}{\dnCh}{\dnCh}{\dnCh}{\dnCh}{\dnCh}{\dnCh}{\dnCh}{\dnCh}{\dnCh}{\dnCh}{\dnCh}{\dnCh}{\dnCh}} \\
 \dnI{7} & \dnW: & {{\dnCx}{\dnCq}{\dnCq}{\dnCq}{\dnCq}{\dnCq}{\dnCq}{\dnCq}{\dnCq}{\dnCq}{\dnCq}{\dnCq}{\dnCq}{\dnCq}{\dnCq}{\dnCq}{\dnCq}{\dnCq}{\dnCq}{\dnCq}{\dnCq}{\dnCq}{\dnCq}{\dnCq}{\dnCq}{\dnCq}{\dnCq}{\dnCq}{\dnCq}{\dnCq}{\dnCq}{\dnCq}} & \dnW[E]{}: & {{\dnCq}{\dnCq}{\dnCq}{\dnCq}{\dnCq}{\dnCq}{\dnCq}{\dnCq}{\dnCq}{\dnCq}{\dnCq}{\dnCq}{\dnCq}{\dnCq}{\dnCq}{\dnCq}{\dnCq}{\dnCq}{\dnCq}{\dnCq}{\dnCq}{\dnCq}{\dnCq}{\dnCq}{\dnCq}{\dnCq}{\dnCq}{\dnCq}{\dnCq}{\dnCq}{\dnCq}{\dnCq}} & \dnW[W]{}: & {{\dnCq}{\dnCq}{\dnCq}{\dnCq}{\dnCq}{\dnCq}{\dnCq}{\dnCq}{\dnCq}{\dnCq}{\dnCq}{\dnCq}{\dnCq}{\dnCq}{\dnCq}{\dnCq}{\dnCq}{\dnCq}{\dnCq}{\dnCq}{\dnCq}{\dnCq}{\dnCq}{\dnCq}{\dnCq}{\dnCq}{\dnCq}{\dnCq}{\dnCq}{\dnCq}{\dnCq}{\dnCq}} \\
 \dnI{8} & \dnW: & {{\dnCh}{\dnCh}{\dnCh}{\dnCh}{\dnCh}{\dnCh}{\dnCh}{\dnCh}{\dnCh}{\dnCh}{\dnCh}{\dnCh}{\dnCh}{\dnCh}{\dnCh}{\dnCh}{\dnCh}{\dnCh}{\dnCh}{\dnCh}{\dnCh}{\dnCh}{\dnCh}{\dnCh}{\dnCh}{\dnCh}{\dnCh}{\dnCh}{\dnCh}{\dnCh}{\dnCh}{\dnCh}} & \dnW[E]{}: & {{\dnCq}{\dnCq}{\dnCq}{\dnCq}{\dnCq}{\dnCq}{\dnCq}{\dnCq}{\dnCq}{\dnCq}{\dnCq}{\dnCq}{\dnCq}{\dnCq}{\dnCq}{\dnCq}{\dnCq}{\dnCq}{\dnCq}{\dnCq}{\dnCq}{\dnCq}{\dnCq}{\dnCq}{\dnCq}{\dnCq}{\dnCq}{\dnCq}{\dnCq}{\dnCq}{\dnCq}{\dnCq}} & \dnW[W]{}: & {{\dnCq}{\dnCq}{\dnCq}{\dnCq}{\dnCq}{\dnCq}{\dnCq}{\dnCq}{\dnCq}{\dnCq}{\dnCq}{\dnCq}{\dnCq}{\dnCq}{\dnCq}{\dnCq}{\dnCq}{\dnCq}{\dnCq}{\dnCq}{\dnCq}{\dnCq}{\dnCq}{\dnCq}{\dnCq}{\dnCq}{\dnCq}{\dnCq}{\dnCq}{\dnCq}{\dnCq}{\dnCq}} \\
 \dnI{9} & \dnW: & {{\dnCh}{\dnCh}{\dnCh}{\dnCh}{\dnCh}{\dnCh}{\dnCh}{\dnCh}{\dnCh}{\dnCh}{\dnCh}{\dnCh}{\dnCh}{\dnCh}{\dnCh}{\dnCh}{\dnCh}{\dnCh}{\dnCh}{\dnCh}{\dnCh}{\dnCh}{\dnCh}{\dnCh}{\dnCh}{\dnCh}{\dnCh}{\dnCh}{\dnCh}{\dnCh}{\dnCh}{\dnCh}} & \dnW[E]{}: & {{\dnCq}{\dnCq}{\dnCq}{\dnCq}{\dnCq}{\dnCq}{\dnCq}{\dnCq}{\dnCq}{\dnCq}{\dnCq}{\dnCq}{\dnCq}{\dnCq}{\dnCq}{\dnCq}{\dnCq}{\dnCq}{\dnCq}{\dnCq}{\dnCq}{\dnCq}{\dnCq}{\dnCq}{\dnCq}{\dnCq}{\dnCq}{\dnCq}{\dnCq}{\dnCq}{\dnCq}{\dnCq}} & \dnW[W]{}: & {{\dnCh}{\dnCh}{\dnCh}{\dnCh}{\dnCh}{\dnCh}{\dnCh}{\dnCh}{\dnCh}{\dnCh}{\dnCh}{\dnCh}{\dnCh}{\dnCh}{\dnCh}{\dnCh}{\dnCh}{\dnCh}{\dnCh}{\dnCh}{\dnCh}{\dnCh}{\dnCh}{\dnCh}{\dnCh}{\dnCh}{\dnCh}{\dnCh}{\dnCh}{\dnCh}{\dnCh}{\dnCh}} \\
 \dnI{10} & \dnW: & {{\dnCh}{\dnCh}{\dnCh}{\dnCh}{\dnCh}{\dnCh}{\dnCh}{\dnCh}{\dnCh}{\dnCh}{\dnCh}{\dnCh}{\dnCh}{\dnCh}{\dnCh}{\dnCh}{\dnCh}{\dnCh}{\dnCh}{\dnCh}{\dnCh}{\dnCh}{\dnCh}{\dnCh}{\dnCh}{\dnCh}{\dnCh}{\dnCh}{\dnCh}{\dnCh}{\dnCh}{\dnCh}} & \dnW[E]{}: & {{\dnCq}{\dnCq}{\dnCq}{\dnCq}{\dnCq}{\dnCq}{\dnCq}{\dnCq}{\dnCq}{\dnCq}{\dnCq}{\dnCq}{\dnCq}{\dnCq}{\dnCq}{\dnCq}{\dnCq}{\dnCq}{\dnCq}{\dnCq}{\dnCq}{\dnCq}{\dnCq}{\dnCq}{\dnCq}{\dnCq}{\dnCq}{\dnCq}{\dnCq}{\dnCq}{\dnCq}{\dnCq}} & \dnW[W]{}: & {{\dnCq}{\dnCq}{\dnCq}{\dnCq}{\dnCq}{\dnCq}{\dnCq}{\dnCq}{\dnCq}{\dnCq}{\dnCq}{\dnCq}{\dnCq}{\dnCq}{\dnCq}{\dnCq}{\dnCq}{\dnCq}{\dnCq}{\dnCq}{\dnCq}{\dnCq}{\dnCq}{\dnCq}{\dnCq}{\dnCq}{\dnCq}{\dnCq}{\dnCq}{\dnCq}{\dnCq}{\dnCq}} \\
 \dnI{11} & \dnW: & {{\dnCh}{\dnCh}{\dnCh}{\dnCh}{\dnCh}{\dnCh}{\dnCh}{\dnCh}{\dnCh}{\dnCh}{\dnCh}{\dnCh}{\dnCh}{\dnCh}{\dnCh}{\dnCh}{\dnCh}{\dnCh}{\dnCh}{\dnCh}{\dnCh}{\dnCh}{\dnCh}{\dnCh}{\dnCh}{\dnCh}{\dnCh}{\dnCh}{\dnCh}{\dnCh}{\dnCh}{\dnCh}} & \dnW[E]{}: & {{\dnCq}{\dnCq}{\dnCq}{\dnCq}{\dnCq}{\dnCq}{\dnCq}{\dnCq}{\dnCq}{\dnCq}{\dnCq}{\dnCq}{\dnCq}{\dnCq}{\dnCq}{\dnCq}{\dnCq}{\dnCq}{\dnCq}{\dnCq}{\dnCq}{\dnCq}{\dnCq}{\dnCq}{\dnCq}{\dnCq}{\dnCq}{\dnCq}{\dnCq}{\dnCq}{\dnCq}{\dnCq}} & \dnW[W]{}: & {{\dnCh}{\dnCh}{\dnCh}{\dnCh}{\dnCh}{\dnCh}{\dnCh}{\dnCh}{\dnCh}{\dnCh}{\dnCh}{\dnCh}{\dnCh}{\dnCh}{\dnCh}{\dnCh}{\dnCh}{\dnCh}{\dnCh}{\dnCh}{\dnCh}{\dnCh}{\dnCh}{\dnCh}{\dnCh}{\dnCh}{\dnCh}{\dnCh}{\dnCh}{\dnCh}{\dnCh}{\dnCh}} \\
 \dnI{12} & \dnW: & {{\dnCh}{\dnCh}{\dnCh}{\dnCh}{\dnCh}{\dnCh}{\dnCh}{\dnCh}{\dnCh}{\dnCh}{\dnCh}{\dnCh}{\dnCh}{\dnCh}{\dnCh}{\dnCh}{\dnCh}{\dnCh}{\dnCh}{\dnCh}{\dnCh}{\dnCh}{\dnCh}{\dnCh}{\dnCh}{\dnCh}{\dnCh}{\dnCh}{\dnCh}{\dnCh}{\dnCh}{\dnCh}} & \dnW[E]{}: & {{\dnCh}{\dnCh}{\dnCh}{\dnCh}{\dnCh}{\dnCh}{\dnCh}{\dnCh}{\dnCh}{\dnCh}{\dnCh}{\dnCh}{\dnCh}{\dnCh}{\dnCh}{\dnCh}{\dnCh}{\dnCh}{\dnCh}{\dnCh}{\dnCh}{\dnCh}{\dnCh}{\dnCh}{\dnCh}{\dnCh}{\dnCh}{\dnCh}{\dnCh}{\dnCh}{\dnCh}{\dnCh}} & \dnW[W]{}: & {{\dnCh}{\dnCh}{\dnCh}{\dnCh}{\dnCh}{\dnCh}{\dnCh}{\dnCh}{\dnCh}{\dnCh}{\dnCh}{\dnCh}{\dnCh}{\dnCh}{\dnCh}{\dnCh}{\dnCh}{\dnCh}{\dnCh}{\dnCh}{\dnCh}{\dnCh}{\dnCh}{\dnCh}{\dnCh}{\dnCh}{\dnCh}{\dnCh}{\dnCh}{\dnCh}{\dnCh}{\dnCh}} \\
 \dnI{13} & \dnW: & {{\dnCh}{\dnCh}{\dnCh}{\dnCh}{\dnCh}{\dnCh}{\dnCh}{\dnCh}{\dnCh}{\dnCh}{\dnCh}{\dnCh}{\dnCh}{\dnCh}{\dnCh}{\dnCh}{\dnCh}{\dnCh}{\dnCh}{\dnCh}{\dnCh}{\dnCh}{\dnCh}{\dnCh}{\dnCh}{\dnCh}{\dnCh}{\dnCh}{\dnCh}{\dnCh}{\dnCh}{\dnCh}} & \dnW[E]{}: & {{\dnCh}{\dnCh}{\dnCh}{\dnCh}{\dnCh}{\dnCh}{\dnCh}{\dnCh}{\dnCh}{\dnCh}{\dnCh}{\dnCh}{\dnCh}{\dnCh}{\dnCh}{\dnCh}{\dnCh}{\dnCh}{\dnCh}{\dnCh}{\dnCh}{\dnCh}{\dnCh}{\dnCh}{\dnCh}{\dnCh}{\dnCh}{\dnCh}{\dnCh}{\dnCh}{\dnCh}{\dnCh}} & \dnW[W]{}: & {{\dnCh}{\dnCh}{\dnCh}{\dnCh}{\dnCh}{\dnCh}{\dnCh}{\dnCh}{\dnCh}{\dnCh}{\dnCh}{\dnCh}{\dnCh}{\dnCh}{\dnCh}{\dnCh}{\dnCh}{\dnCh}{\dnCh}{\dnCh}{\dnCh}{\dnCh}{\dnCh}{\dnCh}{\dnCh}{\dnCh}{\dnCh}{\dnCh}{\dnCh}{\dnCh}{\dnCh}{\dnCh}} \\
 \dnI{14} & \dnW: & {{\dnCh}{\dnCh}{\dnCh}{\dnCh}{\dnCh}{\dnCh}{\dnCh}{\dnCh}{\dnCh}{\dnCh}{\dnCh}{\dnCh}{\dnCh}{\dnCh}{\dnCh}{\dnCh}{\dnCh}{\dnCh}{\dnCh}{\dnCh}{\dnCh}{\dnCh}{\dnCh}{\dnCh}{\dnCh}{\dnCh}{\dnCh}{\dnCh}{\dnCh}{\dnCh}{\dnCh}{\dnCh}} & \dnW[E]{}: & {{\dnCh}{\dnCh}{\dnCh}{\dnCh}{\dnCh}{\dnCh}{\dnCh}{\dnCh}{\dnCh}{\dnCh}{\dnCh}{\dnCh}{\dnCh}{\dnCh}{\dnCh}{\dnCh}{\dnCh}{\dnCh}{\dnCh}{\dnCh}{\dnCh}{\dnCh}{\dnCh}{\dnCh}{\dnCh}{\dnCh}{\dnCh}{\dnCh}{\dnCh}{\dnCh}{\dnCh}{\dnCh}} & \dnW[W]{}: & {{\dnCh}{\dnCh}{\dnCh}{\dnCh}{\dnCh}{\dnCh}{\dnCh}{\dnCh}{\dnCh}{\dnCh}{\dnCh}{\dnCh}{\dnCh}{\dnCh}{\dnCh}{\dnCh}{\dnCh}{\dnCh}{\dnCh}{\dnCh}{\dnCh}{\dnCh}{\dnCh}{\dnCh}{\dnCh}{\dnCh}{\dnCh}{\dnCh}{\dnCh}{\dnCh}{\dnCh}{\dnCh}} \\
 \dnI{15} & \dnW: & {{\dnCh}{\dnCh}{\dnCh}{\dnCh}{\dnCh}{\dnCh}{\dnCh}{\dnCh}{\dnCh}{\dnCh}{\dnCh}{\dnCh}{\dnCh}{\dnCh}{\dnCh}{\dnCh}{\dnCh}{\dnCh}{\dnCh}{\dnCh}{\dnCh}{\dnCh}{\dnCh}{\dnCh}{\dnCh}{\dnCh}{\dnCh}{\dnCh}{\dnCh}{\dnCh}{\dnCh}{\dnCh}} & \dnW[E]{}: & {{\dnCh}{\dnCh}{\dnCh}{\dnCh}{\dnCh}{\dnCh}{\dnCh}{\dnCh}{\dnCh}{\dnCh}{\dnCh}{\dnCh}{\dnCh}{\dnCh}{\dnCh}{\dnCh}{\dnCh}{\dnCh}{\dnCh}{\dnCh}{\dnCh}{\dnCh}{\dnCh}{\dnCh}{\dnCh}{\dnCh}{\dnCh}{\dnCh}{\dnCh}{\dnCh}{\dnCh}{\dnCh}} & \dnW[W]{}: & {{\dnCq}{\dnCq}{\dnCq}{\dnCq}{\dnCq}{\dnCq}{\dnCq}{\dnCq}{\dnCq}{\dnCq}{\dnCq}{\dnCq}{\dnCq}{\dnCq}{\dnCq}{\dnCq}{\dnCq}{\dnCq}{\dnCq}{\dnCq}{\dnCq}{\dnCq}{\dnCq}{\dnCq}{\dnCq}{\dnCq}{\dnCq}{\dnCq}{\dnCq}{\dnCq}{\dnCq}{\dnCq}} \\
 \dnI{16} & \dnW: & {{\dnCh}{\dnCh}{\dnCh}{\dnCh}{\dnCh}{\dnCh}{\dnCh}{\dnCh}{\dnCh}{\dnCh}{\dnCh}{\dnCh}{\dnCh}{\dnCh}{\dnCh}{\dnCh}{\dnCh}{\dnCh}{\dnCh}{\dnCh}{\dnCh}{\dnCh}{\dnCh}{\dnCh}{\dnCh}{\dnCh}{\dnCh}{\dnCh}{\dnCh}{\dnCh}{\dnCh}{\dnCh}} & \dnW[E]{}: & {{\dnCh}{\dnCh}{\dnCh}{\dnCh}{\dnCh}{\dnCh}{\dnCh}{\dnCh}{\dnCh}{\dnCh}{\dnCh}{\dnCh}{\dnCh}{\dnCh}{\dnCh}{\dnCh}{\dnCh}{\dnCh}{\dnCh}{\dnCh}{\dnCh}{\dnCh}{\dnCh}{\dnCh}{\dnCh}{\dnCh}{\dnCh}{\dnCh}{\dnCh}{\dnCh}{\dnCh}{\dnCh}} & \dnW[W]{}: & {{\dnCh}{\dnCh}{\dnCh}{\dnCh}{\dnCh}{\dnCh}{\dnCh}{\dnCh}{\dnCh}{\dnCh}{\dnCh}{\dnCh}{\dnCh}{\dnCh}{\dnCh}{\dnCh}{\dnCh}{\dnCh}{\dnCh}{\dnCh}{\dnCh}{\dnCh}{\dnCh}{\dnCh}{\dnCh}{\dnCh}{\dnCh}{\dnCh}{\dnCh}{\dnCh}{\dnCh}{\dnCh}} \\
 \dnI{17} & \dnW: & {{\dnCh}{\dnCh}{\dnCh}{\dnCh}{\dnCh}{\dnCh}{\dnCh}{\dnCh}{\dnCh}{\dnCh}{\dnCh}{\dnCh}{\dnCh}{\dnCh}{\dnCh}{\dnCh}{\dnCh}{\dnCh}{\dnCh}{\dnCh}{\dnCh}{\dnCh}{\dnCh}{\dnCh}{\dnCh}{\dnCh}{\dnCh}{\dnCh}{\dnCh}{\dnCh}{\dnCh}{\dnCh}} & \dnW[E]{}: & {{\dnCh}{\dnCh}{\dnCh}{\dnCh}{\dnCh}{\dnCh}{\dnCh}{\dnCh}{\dnCh}{\dnCh}{\dnCh}{\dnCh}{\dnCh}{\dnCh}{\dnCh}{\dnCh}{\dnCh}{\dnCh}{\dnCh}{\dnCh}{\dnCh}{\dnCh}{\dnCh}{\dnCh}{\dnCh}{\dnCh}{\dnCh}{\dnCh}{\dnCh}{\dnCh}{\dnCh}{\dnCh}} & \dnW[W]{}: & {{\dnCh}{\dnCh}{\dnCh}{\dnCh}{\dnCh}{\dnCh}{\dnCh}{\dnCh}{\dnCh}{\dnCh}{\dnCh}{\dnCh}{\dnCh}{\dnCh}{\dnCh}{\dnCh}{\dnCh}{\dnCh}{\dnCh}{\dnCh}{\dnCh}{\dnCh}{\dnCh}{\dnCh}{\dnCh}{\dnCh}{\dnCh}{\dnCh}{\dnCh}{\dnCh}{\dnCh}{\dnCh}} \\
 \dnI{18} & \dnW: & {{\dnCh}{\dnCh}{\dnCh}{\dnCh}{\dnCh}{\dnCh}{\dnCh}{\dnCh}{\dnCh}{\dnCh}{\dnCh}{\dnCh}{\dnCh}{\dnCh}{\dnCh}{\dnCh}{\dnCh}{\dnCh}{\dnCh}{\dnCh}{\dnCh}{\dnCh}{\dnCh}{\dnCh}{\dnCh}{\dnCh}{\dnCh}{\dnCh}{\dnCh}{\dnCh}{\dnCh}{\dnCh}} & \dnW[E]{}: & {{\dnCh}{\dnCh}{\dnCh}{\dnCh}{\dnCh}{\dnCh}{\dnCh}{\dnCh}{\dnCh}{\dnCh}{\dnCh}{\dnCh}{\dnCh}{\dnCh}{\dnCh}{\dnCh}{\dnCh}{\dnCh}{\dnCh}{\dnCh}{\dnCh}{\dnCh}{\dnCh}{\dnCh}{\dnCh}{\dnCh}{\dnCh}{\dnCh}{\dnCh}{\dnCh}{\dnCh}{\dnCh}} & \dnW[W]{}: & {{\dnCh}{\dnCh}{\dnCh}{\dnCh}{\dnCh}{\dnCh}{\dnCh}{\dnCh}{\dnCh}{\dnCh}{\dnCh}{\dnCh}{\dnCh}{\dnCh}{\dnCh}{\dnCh}{\dnCh}{\dnCh}{\dnCh}{\dnCh}{\dnCh}{\dnCh}{\dnCh}{\dnCh}{\dnCh}{\dnCh}{\dnCh}{\dnCh}{\dnCh}{\dnCh}{\dnCh}{\dnCh}} \\
 \dnI{19} & \dnW: & {{\dnCh}{\dnCh}{\dnCh}{\dnCh}{\dnCh}{\dnCh}{\dnCh}{\dnCh}{\dnCh}{\dnCh}{\dnCh}{\dnCh}{\dnCh}{\dnCh}{\dnCh}{\dnCh}{\dnCh}{\dnCh}{\dnCh}{\dnCh}{\dnCh}{\dnCh}{\dnCh}{\dnCh}{\dnCh}{\dnCh}{\dnCh}{\dnCh}{\dnCh}{\dnCh}{\dnCh}{\dnCh}} & \dnW[E]{}: & {{\dnCh}{\dnCh}{\dnCh}{\dnCh}{\dnCh}{\dnCh}{\dnCh}{\dnCh}{\dnCh}{\dnCh}{\dnCh}{\dnCh}{\dnCh}{\dnCh}{\dnCh}{\dnCh}{\dnCh}{\dnCh}{\dnCh}{\dnCh}{\dnCh}{\dnCh}{\dnCh}{\dnCh}{\dnCh}{\dnCh}{\dnCh}{\dnCh}{\dnCh}{\dnCh}{\dnCh}{\dnCh}} & \dnW[W]{}: & {{\dnCh}{\dnCh}{\dnCh}{\dnCh}{\dnCh}{\dnCh}{\dnCh}{\dnCh}{\dnCh}{\dnCh}{\dnCh}{\dnCh}{\dnCh}{\dnCh}{\dnCh}{\dnCh}{\dnCh}{\dnCh}{\dnCh}{\dnCh}{\dnCh}{\dnCh}{\dnCh}{\dnCh}{\dnCh}{\dnCh}{\dnCh}{\dnCh}{\dnCh}{\dnCh}{\dnCh}{\dnCh}} \\
 \dnI{20} & \dnW: & {{\dnCh}{\dnCh}{\dnCh}{\dnCh}{\dnCh}{\dnCh}{\dnCh}{\dnCh}{\dnCh}{\dnCh}{\dnCh}{\dnCh}{\dnCh}{\dnCh}{\dnCh}{\dnCh}{\dnCh}{\dnCh}{\dnCh}{\dnCh}{\dnCh}{\dnCh}{\dnCh}{\dnCh}{\dnCh}{\dnCh}{\dnCh}{\dnCh}{\dnCh}{\dnCh}{\dnCh}{\dnCh}} & \dnW[E]{}: & {{\dnCh}{\dnCh}{\dnCh}{\dnCh}{\dnCh}{\dnCh}{\dnCh}{\dnCh}{\dnCh}{\dnCh}{\dnCh}{\dnCh}{\dnCh}{\dnCh}{\dnCh}{\dnCh}{\dnCh}{\dnCh}{\dnCh}{\dnCh}{\dnCh}{\dnCh}{\dnCh}{\dnCh}{\dnCh}{\dnCh}{\dnCh}{\dnCh}{\dnCh}{\dnCh}{\dnCh}{\dnCh}} & \dnW[W]{}: & {{\dnCh}{\dnCh}{\dnCh}{\dnCh}{\dnCh}{\dnCh}{\dnCh}{\dnCh}{\dnCh}{\dnCh}{\dnCh}{\dnCh}{\dnCh}{\dnCh}{\dnCh}{\dnCh}{\dnCh}{\dnCh}{\dnCh}{\dnCh}{\dnCh}{\dnCh}{\dnCh}{\dnCh}{\dnCh}{\dnCh}{\dnCh}{\dnCh}{\dnCh}{\dnCh}{\dnCh}{\dnCh}} \\
 \dnI{21} & \dnW: & {{\dnCh}{\dnCh}{\dnCh}{\dnCh}{\dnCh}{\dnCh}{\dnCh}{\dnCh}{\dnCh}{\dnCh}{\dnCh}{\dnCh}{\dnCh}{\dnCh}{\dnCh}{\dnCh}{\dnCh}{\dnCh}{\dnCh}{\dnCh}{\dnCh}{\dnCh}{\dnCh}{\dnCh}{\dnCh}{\dnCh}{\dnCh}{\dnCh}{\dnCh}{\dnCh}{\dnCh}{\dnCh}} & \dnW[E]{}: & {{\dnCh}{\dnCh}{\dnCh}{\dnCh}{\dnCh}{\dnCh}{\dnCh}{\dnCh}{\dnCh}{\dnCh}{\dnCh}{\dnCh}{\dnCh}{\dnCh}{\dnCh}{\dnCh}{\dnCh}{\dnCh}{\dnCh}{\dnCh}{\dnCh}{\dnCh}{\dnCh}{\dnCh}{\dnCh}{\dnCh}{\dnCh}{\dnCh}{\dnCh}{\dnCh}{\dnCh}{\dnCh}} & \dnW[W]{}: & {{\dnCh}{\dnCh}{\dnCh}{\dnCh}{\dnCh}{\dnCh}{\dnCh}{\dnCh}{\dnCh}{\dnCh}{\dnCh}{\dnCh}{\dnCh}{\dnCh}{\dnCh}{\dnCh}{\dnCh}{\dnCh}{\dnCh}{\dnCh}{\dnCh}{\dnCh}{\dnCh}{\dnCh}{\dnCh}{\dnCh}{\dnCh}{\dnCh}{\dnCh}{\dnCh}{\dnCh}{\dnCh}} \\
 \dnI{22} & \dnW: & {{\dnCh}{\dnCh}{\dnCh}{\dnCh}{\dnCh}{\dnCh}{\dnCh}{\dnCh}{\dnCh}{\dnCh}{\dnCh}{\dnCh}{\dnCh}{\dnCh}{\dnCh}{\dnCh}{\dnCh}{\dnCh}{\dnCh}{\dnCh}{\dnCh}{\dnCh}{\dnCh}{\dnCh}{\dnCh}{\dnCh}{\dnCh}{\dnCh}{\dnCh}{\dnCh}{\dnCh}{\dnCh}} & \dnW[E]{}: & {{\dnCh}{\dnCh}{\dnCh}{\dnCh}{\dnCh}{\dnCh}{\dnCh}{\dnCh}{\dnCh}{\dnCh}{\dnCh}{\dnCh}{\dnCh}{\dnCh}{\dnCh}{\dnCh}{\dnCh}{\dnCh}{\dnCh}{\dnCh}{\dnCh}{\dnCh}{\dnCh}{\dnCh}{\dnCh}{\dnCh}{\dnCh}{\dnCh}{\dnCh}{\dnCh}{\dnCh}{\dnCh}} & \dnW[W]{}: & {{\dnCh}{\dnCh}{\dnCh}{\dnCh}{\dnCh}{\dnCh}{\dnCh}{\dnCh}{\dnCh}{\dnCh}{\dnCh}{\dnCh}{\dnCh}{\dnCh}{\dnCh}{\dnCh}{\dnCh}{\dnCh}{\dnCh}{\dnCh}{\dnCh}{\dnCh}{\dnCh}{\dnCh}{\dnCh}{\dnCh}{\dnCh}{\dnCh}{\dnCh}{\dnCh}{\dnCh}{\dnCh}} \\
\hline
\end{tabular}
}
\caption{SHA-256 hash collision over 23 rounds}
\label{fig:tc23}
\end{center}
\end{figure}
}

\texttt{
\begin{figure}[!ht]
\begin{center}
\scalebox{.55}{
\begin{tabular}{|r|c|c|c|c|c|c|}
\hline
  $i$  &      &  A  &    & E &   & W \\
\hline
 \dnI{-4} & \dnW: & {{\dnCh}{\dnCh}{\dnCh}{\dnCh}{\dnCh}{\dnCh}{\dnCh}{\dnCh}{\dnCh}{\dnCh}{\dnCh}{\dnCh}{\dnCh}{\dnCh}{\dnCh}{\dnCh}{\dnCh}{\dnCh}{\dnCh}{\dnCh}{\dnCh}{\dnCh}{\dnCh}{\dnCh}{\dnCh}{\dnCh}{\dnCh}{\dnCh}{\dnCh}{\dnCh}{\dnCh}{\dnCh}} & \dnW[E]{}: & {{\dnCh}{\dnCh}{\dnCh}{\dnCh}{\dnCh}{\dnCh}{\dnCh}{\dnCh}{\dnCh}{\dnCh}{\dnCh}{\dnCh}{\dnCh}{\dnCh}{\dnCh}{\dnCh}{\dnCh}{\dnCh}{\dnCh}{\dnCh}{\dnCh}{\dnCh}{\dnCh}{\dnCh}{\dnCh}{\dnCh}{\dnCh}{\dnCh}{\dnCh}{\dnCh}{\dnCh}{\dnCh}} &  &     \\
 \dnI{-3} & \dnW: & {{\dnCh}{\dnCh}{\dnCh}{\dnCh}{\dnCh}{\dnCh}{\dnCh}{\dnCh}{\dnCh}{\dnCh}{\dnCh}{\dnCh}{\dnCh}{\dnCh}{\dnCh}{\dnCh}{\dnCh}{\dnCh}{\dnCh}{\dnCh}{\dnCh}{\dnCh}{\dnCh}{\dnCh}{\dnCh}{\dnCh}{\dnCh}{\dnCh}{\dnCh}{\dnCh}{\dnCh}{\dnCh}} & \dnW[E]{}: & {{\dnCh}{\dnCh}{\dnCh}{\dnCh}{\dnCh}{\dnCh}{\dnCh}{\dnCh}{\dnCh}{\dnCh}{\dnCh}{\dnCh}{\dnCh}{\dnCh}{\dnCh}{\dnCh}{\dnCh}{\dnCh}{\dnCh}{\dnCh}{\dnCh}{\dnCh}{\dnCh}{\dnCh}{\dnCh}{\dnCh}{\dnCh}{\dnCh}{\dnCh}{\dnCh}{\dnCh}{\dnCh}} &  &     \\
 \dnI{-2} & \dnW: & {{\dnCh}{\dnCh}{\dnCh}{\dnCh}{\dnCh}{\dnCh}{\dnCh}{\dnCh}{\dnCh}{\dnCh}{\dnCh}{\dnCh}{\dnCh}{\dnCh}{\dnCh}{\dnCh}{\dnCh}{\dnCh}{\dnCh}{\dnCh}{\dnCh}{\dnCh}{\dnCh}{\dnCh}{\dnCh}{\dnCh}{\dnCh}{\dnCh}{\dnCh}{\dnCh}{\dnCh}{\dnCh}} & \dnW[E]{}: & {{\dnCh}{\dnCh}{\dnCh}{\dnCh}{\dnCh}{\dnCh}{\dnCh}{\dnCh}{\dnCh}{\dnCh}{\dnCh}{\dnCh}{\dnCh}{\dnCh}{\dnCh}{\dnCh}{\dnCh}{\dnCh}{\dnCh}{\dnCh}{\dnCh}{\dnCh}{\dnCh}{\dnCh}{\dnCh}{\dnCh}{\dnCh}{\dnCh}{\dnCh}{\dnCh}{\dnCh}{\dnCh}} &  &     \\
 \dnI{-1} & \dnW: & {{\dnCh}{\dnCh}{\dnCh}{\dnCh}{\dnCh}{\dnCh}{\dnCh}{\dnCh}{\dnCh}{\dnCh}{\dnCh}{\dnCh}{\dnCh}{\dnCh}{\dnCh}{\dnCh}{\dnCh}{\dnCh}{\dnCh}{\dnCh}{\dnCh}{\dnCh}{\dnCh}{\dnCh}{\dnCh}{\dnCh}{\dnCh}{\dnCh}{\dnCh}{\dnCh}{\dnCh}{\dnCh}} & \dnW[E]{}: & {{\dnCh}{\dnCh}{\dnCh}{\dnCh}{\dnCh}{\dnCh}{\dnCh}{\dnCh}{\dnCh}{\dnCh}{\dnCh}{\dnCh}{\dnCh}{\dnCh}{\dnCh}{\dnCh}{\dnCh}{\dnCh}{\dnCh}{\dnCh}{\dnCh}{\dnCh}{\dnCh}{\dnCh}{\dnCh}{\dnCh}{\dnCh}{\dnCh}{\dnCh}{\dnCh}{\dnCh}{\dnCh}} &  &     \\
 \dnI{0} & \dnW: & {{\dnCh}{\dnCh}{\dnCh}{\dnCh}{\dnCh}{\dnCh}{\dnCh}{\dnCh}{\dnCh}{\dnCh}{\dnCh}{\dnCh}{\dnCh}{\dnCh}{\dnCh}{\dnCh}{\dnCh}{\dnCh}{\dnCh}{\dnCh}{\dnCh}{\dnCh}{\dnCh}{\dnCh}{\dnCh}{\dnCh}{\dnCh}{\dnCh}{\dnCh}{\dnCh}{\dnCh}{\dnCh}} & \dnW[E]{}: & {{\dnCh}{\dnCh}{\dnCh}{\dnCh}{\dnCh}{\dnCh}{\dnCh}{\dnCh}{\dnCh}{\dnCh}{\dnCh}{\dnCh}{\dnCh}{\dnCh}{\dnCh}{\dnCh}{\dnCh}{\dnCh}{\dnCh}{\dnCh}{\dnCh}{\dnCh}{\dnCh}{\dnCh}{\dnCh}{\dnCh}{\dnCh}{\dnCh}{\dnCh}{\dnCh}{\dnCh}{\dnCh}} & \dnW[W]{}: & {{\dnCh}{\dnCh}{\dnCh}{\dnCh}{\dnCh}{\dnCh}{\dnCh}{\dnCh}{\dnCh}{\dnCh}{\dnCh}{\dnCh}{\dnCh}{\dnCh}{\dnCh}{\dnCh}{\dnCh}{\dnCh}{\dnCh}{\dnCh}{\dnCh}{\dnCh}{\dnCh}{\dnCh}{\dnCh}{\dnCh}{\dnCh}{\dnCh}{\dnCh}{\dnCh}{\dnCh}{\dnCh}} \\
 \dnI{1} & \dnW: & {{\dnCh}{\dnCh}{\dnCh}{\dnCh}{\dnCh}{\dnCh}{\dnCh}{\dnCh}{\dnCh}{\dnCh}{\dnCh}{\dnCh}{\dnCh}{\dnCh}{\dnCh}{\dnCh}{\dnCh}{\dnCh}{\dnCh}{\dnCh}{\dnCh}{\dnCh}{\dnCh}{\dnCh}{\dnCh}{\dnCh}{\dnCh}{\dnCh}{\dnCh}{\dnCh}{\dnCh}{\dnCh}} & \dnW[E]{}: & {{\dnCh}{\dnCh}{\dnCh}{\dnCh}{\dnCh}{\dnCh}{\dnCh}{\dnCh}{\dnCh}{\dnCh}{\dnCh}{\dnCh}{\dnCh}{\dnCh}{\dnCh}{\dnCh}{\dnCh}{\dnCh}{\dnCh}{\dnCh}{\dnCh}{\dnCh}{\dnCh}{\dnCh}{\dnCh}{\dnCh}{\dnCh}{\dnCh}{\dnCh}{\dnCh}{\dnCh}{\dnCh}} & \dnW[W]{}: & {{\dnCh}{\dnCh}{\dnCh}{\dnCh}{\dnCh}{\dnCh}{\dnCh}{\dnCh}{\dnCh}{\dnCh}{\dnCh}{\dnCh}{\dnCh}{\dnCh}{\dnCh}{\dnCh}{\dnCh}{\dnCh}{\dnCh}{\dnCh}{\dnCh}{\dnCh}{\dnCh}{\dnCh}{\dnCh}{\dnCh}{\dnCh}{\dnCh}{\dnCh}{\dnCh}{\dnCh}{\dnCh}} \\
 \dnI{2} & \dnW: & {{\dnCh}{\dnCh}{\dnCh}{\dnCh}{\dnCh}{\dnCh}{\dnCh}{\dnCh}{\dnCh}{\dnCh}{\dnCh}{\dnCh}{\dnCh}{\dnCh}{\dnCh}{\dnCh}{\dnCh}{\dnCh}{\dnCh}{\dnCh}{\dnCh}{\dnCh}{\dnCh}{\dnCh}{\dnCh}{\dnCh}{\dnCh}{\dnCh}{\dnCh}{\dnCh}{\dnCh}{\dnCh}} & \dnW[E]{}: & {{\dnCh}{\dnCh}{\dnCh}{\dnCh}{\dnCh}{\dnCh}{\dnCh}{\dnCh}{\dnCh}{\dnCh}{\dnCh}{\dnCh}{\dnCh}{\dnCh}{\dnCh}{\dnCh}{\dnCh}{\dnCh}{\dnCh}{\dnCh}{\dnCh}{\dnCh}{\dnCh}{\dnCh}{\dnCh}{\dnCh}{\dnCh}{\dnCh}{\dnCh}{\dnCh}{\dnCh}{\dnCh}} & \dnW[W]{}: & {{\dnCh}{\dnCh}{\dnCh}{\dnCh}{\dnCh}{\dnCh}{\dnCh}{\dnCh}{\dnCh}{\dnCh}{\dnCh}{\dnCh}{\dnCh}{\dnCh}{\dnCh}{\dnCh}{\dnCh}{\dnCh}{\dnCh}{\dnCh}{\dnCh}{\dnCh}{\dnCh}{\dnCh}{\dnCh}{\dnCh}{\dnCh}{\dnCh}{\dnCh}{\dnCh}{\dnCh}{\dnCh}} \\
 \dnI{3} & \dnW: & {{\dnCh}{\dnCh}{\dnCh}{\dnCh}{\dnCh}{\dnCh}{\dnCh}{\dnCh}{\dnCh}{\dnCh}{\dnCh}{\dnCh}{\dnCh}{\dnCh}{\dnCh}{\dnCh}{\dnCh}{\dnCh}{\dnCh}{\dnCh}{\dnCh}{\dnCh}{\dnCh}{\dnCh}{\dnCh}{\dnCh}{\dnCh}{\dnCh}{\dnCh}{\dnCh}{\dnCh}{\dnCh}} & \dnW[E]{}: & {{\dnCh}{\dnCh}{\dnCh}{\dnCh}{\dnCh}{\dnCh}{\dnCh}{\dnCh}{\dnCh}{\dnCh}{\dnCh}{\dnCh}{\dnCh}{\dnCh}{\dnCh}{\dnCh}{\dnCh}{\dnCh}{\dnCh}{\dnCh}{\dnCh}{\dnCh}{\dnCh}{\dnCh}{\dnCh}{\dnCh}{\dnCh}{\dnCh}{\dnCh}{\dnCh}{\dnCh}{\dnCh}} & \dnW[W]{}: & {{\dnCh}{\dnCh}{\dnCh}{\dnCh}{\dnCh}{\dnCh}{\dnCh}{\dnCh}{\dnCh}{\dnCh}{\dnCh}{\dnCh}{\dnCh}{\dnCh}{\dnCh}{\dnCh}{\dnCh}{\dnCh}{\dnCh}{\dnCh}{\dnCh}{\dnCh}{\dnCh}{\dnCh}{\dnCh}{\dnCh}{\dnCh}{\dnCh}{\dnCh}{\dnCh}{\dnCh}{\dnCh}} \\
 \dnI{4} & \dnW: & {{\dnCh}{\dnCh}{\dnCh}{\dnCh}{\dnCh}{\dnCh}{\dnCh}{\dnCh}{\dnCh}{\dnCh}{\dnCh}{\dnCh}{\dnCh}{\dnCh}{\dnCh}{\dnCh}{\dnCh}{\dnCh}{\dnCh}{\dnCh}{\dnCh}{\dnCh}{\dnCh}{\dnCh}{\dnCh}{\dnCh}{\dnCh}{\dnCh}{\dnCh}{\dnCh}{\dnCh}{\dnCh}} & \dnW[E]{}: & {{\dnCh}{\dnCh}{\dnCh}{\dnCh}{\dnCh}{\dnCh}{\dnCh}{\dnCh}{\dnCh}{\dnCh}{\dnCh}{\dnCh}{\dnCh}{\dnCh}{\dnCh}{\dnCh}{\dnCh}{\dnCh}{\dnCh}{\dnCh}{\dnCh}{\dnCh}{\dnCh}{\dnCh}{\dnCh}{\dnCh}{\dnCh}{\dnCh}{\dnCh}{\dnCh}{\dnCh}{\dnCh}} & \dnW[W]{}: & {{\dnCh}{\dnCh}{\dnCh}{\dnCh}{\dnCh}{\dnCh}{\dnCh}{\dnCh}{\dnCh}{\dnCh}{\dnCh}{\dnCh}{\dnCh}{\dnCh}{\dnCh}{\dnCh}{\dnCh}{\dnCh}{\dnCh}{\dnCh}{\dnCh}{\dnCh}{\dnCh}{\dnCh}{\dnCh}{\dnCh}{\dnCh}{\dnCh}{\dnCh}{\dnCh}{\dnCh}{\dnCh}} \\
 \dnI{5} & \dnW: & {{\dnCh}{\dnCh}{\dnCh}{\dnCh}{\dnCh}{\dnCh}{\dnCh}{\dnCh}{\dnCh}{\dnCh}{\dnCh}{\dnCh}{\dnCh}{\dnCh}{\dnCh}{\dnCh}{\dnCh}{\dnCh}{\dnCh}{\dnCh}{\dnCh}{\dnCh}{\dnCh}{\dnCh}{\dnCh}{\dnCh}{\dnCh}{\dnCh}{\dnCh}{\dnCh}{\dnCh}{\dnCh}} & \dnW[E]{}: & {{\dnCh}{\dnCh}{\dnCh}{\dnCh}{\dnCh}{\dnCh}{\dnCh}{\dnCh}{\dnCh}{\dnCh}{\dnCh}{\dnCh}{\dnCh}{\dnCh}{\dnCh}{\dnCh}{\dnCh}{\dnCh}{\dnCh}{\dnCh}{\dnCh}{\dnCh}{\dnCh}{\dnCh}{\dnCh}{\dnCh}{\dnCh}{\dnCh}{\dnCh}{\dnCh}{\dnCh}{\dnCh}} & \dnW[W]{}: & {{\dnCh}{\dnCh}{\dnCh}{\dnCh}{\dnCh}{\dnCh}{\dnCh}{\dnCh}{\dnCh}{\dnCh}{\dnCh}{\dnCh}{\dnCh}{\dnCh}{\dnCh}{\dnCh}{\dnCh}{\dnCh}{\dnCh}{\dnCh}{\dnCh}{\dnCh}{\dnCh}{\dnCh}{\dnCh}{\dnCh}{\dnCh}{\dnCh}{\dnCh}{\dnCh}{\dnCh}{\dnCh}} \\
 \dnI{6} & \dnW: & {{\dnCh}{\dnCh}{\dnCh}{\dnCh}{\dnCh}{\dnCh}{\dnCh}{\dnCh}{\dnCh}{\dnCh}{\dnCh}{\dnCh}{\dnCh}{\dnCh}{\dnCh}{\dnCh}{\dnCh}{\dnCh}{\dnCh}{\dnCh}{\dnCh}{\dnCh}{\dnCh}{\dnCh}{\dnCh}{\dnCh}{\dnCh}{\dnCh}{\dnCh}{\dnCh}{\dnCh}{\dnCh}} & \dnW[E]{}: & {{\dnCh}{\dnCh}{\dnCh}{\dnCh}{\dnCh}{\dnCh}{\dnCh}{\dnCh}{\dnCh}{\dnCh}{\dnCh}{\dnCh}{\dnCh}{\dnCh}{\dnCh}{\dnCh}{\dnCh}{\dnCh}{\dnCh}{\dnCh}{\dnCh}{\dnCh}{\dnCh}{\dnCh}{\dnCh}{\dnCh}{\dnCh}{\dnCh}{\dnCh}{\dnCh}{\dnCh}{\dnCh}} & \dnW[W]{}: & {{\dnCh}{\dnCh}{\dnCh}{\dnCh}{\dnCh}{\dnCh}{\dnCh}{\dnCh}{\dnCh}{\dnCh}{\dnCh}{\dnCh}{\dnCh}{\dnCh}{\dnCh}{\dnCh}{\dnCh}{\dnCh}{\dnCh}{\dnCh}{\dnCh}{\dnCh}{\dnCh}{\dnCh}{\dnCh}{\dnCh}{\dnCh}{\dnCh}{\dnCh}{\dnCh}{\dnCh}{\dnCh}} \\
 \dnI{7} & \dnW: & {{\dnCx}{\dnCq}{\dnCq}{\dnCq}{\dnCq}{\dnCq}{\dnCq}{\dnCq}{\dnCq}{\dnCq}{\dnCq}{\dnCq}{\dnCq}{\dnCq}{\dnCq}{\dnCq}{\dnCq}{\dnCq}{\dnCq}{\dnCq}{\dnCq}{\dnCq}{\dnCq}{\dnCq}{\dnCq}{\dnCq}{\dnCq}{\dnCq}{\dnCq}{\dnCq}{\dnCq}{\dnCq}} & \dnW[E]{}: & {{\dnCq}{\dnCq}{\dnCq}{\dnCq}{\dnCq}{\dnCq}{\dnCq}{\dnCq}{\dnCq}{\dnCq}{\dnCq}{\dnCq}{\dnCq}{\dnCq}{\dnCq}{\dnCq}{\dnCq}{\dnCq}{\dnCq}{\dnCq}{\dnCq}{\dnCq}{\dnCq}{\dnCq}{\dnCq}{\dnCq}{\dnCq}{\dnCq}{\dnCq}{\dnCq}{\dnCq}{\dnCq}} & \dnW[W]{}: & {{\dnCq}{\dnCq}{\dnCq}{\dnCq}{\dnCq}{\dnCq}{\dnCq}{\dnCq}{\dnCq}{\dnCq}{\dnCq}{\dnCq}{\dnCq}{\dnCq}{\dnCq}{\dnCq}{\dnCq}{\dnCq}{\dnCq}{\dnCq}{\dnCq}{\dnCq}{\dnCq}{\dnCq}{\dnCq}{\dnCq}{\dnCq}{\dnCq}{\dnCq}{\dnCq}{\dnCq}{\dnCq}} \\
 \dnI{8} & \dnW: & {{\dnCh}{\dnCh}{\dnCh}{\dnCh}{\dnCh}{\dnCh}{\dnCh}{\dnCh}{\dnCh}{\dnCh}{\dnCh}{\dnCh}{\dnCh}{\dnCh}{\dnCh}{\dnCh}{\dnCh}{\dnCh}{\dnCh}{\dnCh}{\dnCh}{\dnCh}{\dnCh}{\dnCh}{\dnCh}{\dnCh}{\dnCh}{\dnCh}{\dnCh}{\dnCh}{\dnCh}{\dnCh}} & \dnW[E]{}: & {{\dnCq}{\dnCq}{\dnCq}{\dnCq}{\dnCq}{\dnCq}{\dnCq}{\dnCq}{\dnCq}{\dnCq}{\dnCq}{\dnCq}{\dnCq}{\dnCq}{\dnCq}{\dnCq}{\dnCq}{\dnCq}{\dnCq}{\dnCq}{\dnCq}{\dnCq}{\dnCq}{\dnCq}{\dnCq}{\dnCq}{\dnCq}{\dnCq}{\dnCq}{\dnCq}{\dnCq}{\dnCq}} & \dnW[W]{}: & {{\dnCq}{\dnCq}{\dnCq}{\dnCq}{\dnCq}{\dnCq}{\dnCq}{\dnCq}{\dnCq}{\dnCq}{\dnCq}{\dnCq}{\dnCq}{\dnCq}{\dnCq}{\dnCq}{\dnCq}{\dnCq}{\dnCq}{\dnCq}{\dnCq}{\dnCq}{\dnCq}{\dnCq}{\dnCq}{\dnCq}{\dnCq}{\dnCq}{\dnCq}{\dnCq}{\dnCq}{\dnCq}} \\
 \dnI{9} & \dnW: & {{\dnCh}{\dnCh}{\dnCh}{\dnCh}{\dnCh}{\dnCh}{\dnCh}{\dnCh}{\dnCh}{\dnCh}{\dnCh}{\dnCh}{\dnCh}{\dnCh}{\dnCh}{\dnCh}{\dnCh}{\dnCh}{\dnCh}{\dnCh}{\dnCh}{\dnCh}{\dnCh}{\dnCh}{\dnCh}{\dnCh}{\dnCh}{\dnCh}{\dnCh}{\dnCh}{\dnCh}{\dnCh}} & \dnW[E]{}: & {{\dnCq}{\dnCq}{\dnCq}{\dnCq}{\dnCq}{\dnCq}{\dnCq}{\dnCq}{\dnCq}{\dnCq}{\dnCq}{\dnCq}{\dnCq}{\dnCq}{\dnCq}{\dnCq}{\dnCq}{\dnCq}{\dnCq}{\dnCq}{\dnCq}{\dnCq}{\dnCq}{\dnCq}{\dnCq}{\dnCq}{\dnCq}{\dnCq}{\dnCq}{\dnCq}{\dnCq}{\dnCq}} & \dnW[W]{}: & {{\dnCh}{\dnCh}{\dnCh}{\dnCh}{\dnCh}{\dnCh}{\dnCh}{\dnCh}{\dnCh}{\dnCh}{\dnCh}{\dnCh}{\dnCh}{\dnCh}{\dnCh}{\dnCh}{\dnCh}{\dnCh}{\dnCh}{\dnCh}{\dnCh}{\dnCh}{\dnCh}{\dnCh}{\dnCh}{\dnCh}{\dnCh}{\dnCh}{\dnCh}{\dnCh}{\dnCh}{\dnCh}} \\
 \dnI{10} & \dnW: & {{\dnCh}{\dnCh}{\dnCh}{\dnCh}{\dnCh}{\dnCh}{\dnCh}{\dnCh}{\dnCh}{\dnCh}{\dnCh}{\dnCh}{\dnCh}{\dnCh}{\dnCh}{\dnCh}{\dnCh}{\dnCh}{\dnCh}{\dnCh}{\dnCh}{\dnCh}{\dnCh}{\dnCh}{\dnCh}{\dnCh}{\dnCh}{\dnCh}{\dnCh}{\dnCh}{\dnCh}{\dnCh}} & \dnW[E]{}: & {{\dnCq}{\dnCq}{\dnCq}{\dnCq}{\dnCq}{\dnCq}{\dnCq}{\dnCq}{\dnCq}{\dnCq}{\dnCq}{\dnCq}{\dnCq}{\dnCq}{\dnCq}{\dnCq}{\dnCq}{\dnCq}{\dnCq}{\dnCq}{\dnCq}{\dnCq}{\dnCq}{\dnCq}{\dnCq}{\dnCq}{\dnCq}{\dnCq}{\dnCq}{\dnCq}{\dnCq}{\dnCq}} & \dnW[W]{}: & {{\dnCq}{\dnCq}{\dnCq}{\dnCq}{\dnCq}{\dnCq}{\dnCq}{\dnCq}{\dnCq}{\dnCq}{\dnCq}{\dnCq}{\dnCq}{\dnCq}{\dnCq}{\dnCq}{\dnCq}{\dnCq}{\dnCq}{\dnCq}{\dnCq}{\dnCq}{\dnCq}{\dnCq}{\dnCq}{\dnCq}{\dnCq}{\dnCq}{\dnCq}{\dnCq}{\dnCq}{\dnCq}} \\
 \dnI{11} & \dnW: & {{\dnCh}{\dnCh}{\dnCh}{\dnCh}{\dnCh}{\dnCh}{\dnCh}{\dnCh}{\dnCh}{\dnCh}{\dnCh}{\dnCh}{\dnCh}{\dnCh}{\dnCh}{\dnCh}{\dnCh}{\dnCh}{\dnCh}{\dnCh}{\dnCh}{\dnCh}{\dnCh}{\dnCh}{\dnCh}{\dnCh}{\dnCh}{\dnCh}{\dnCh}{\dnCh}{\dnCh}{\dnCh}} & \dnW[E]{}: & {{\dnCq}{\dnCq}{\dnCq}{\dnCq}{\dnCq}{\dnCq}{\dnCq}{\dnCq}{\dnCq}{\dnCq}{\dnCq}{\dnCq}{\dnCq}{\dnCq}{\dnCq}{\dnCq}{\dnCq}{\dnCq}{\dnCq}{\dnCq}{\dnCq}{\dnCq}{\dnCq}{\dnCq}{\dnCq}{\dnCq}{\dnCq}{\dnCq}{\dnCq}{\dnCq}{\dnCq}{\dnCq}} & \dnW[W]{}: & {{\dnCh}{\dnCh}{\dnCh}{\dnCh}{\dnCh}{\dnCh}{\dnCh}{\dnCh}{\dnCh}{\dnCh}{\dnCh}{\dnCh}{\dnCh}{\dnCh}{\dnCh}{\dnCh}{\dnCh}{\dnCh}{\dnCh}{\dnCh}{\dnCh}{\dnCh}{\dnCh}{\dnCh}{\dnCh}{\dnCh}{\dnCh}{\dnCh}{\dnCh}{\dnCh}{\dnCh}{\dnCh}} \\
 \dnI{12} & \dnW: & {{\dnCh}{\dnCh}{\dnCh}{\dnCh}{\dnCh}{\dnCh}{\dnCh}{\dnCh}{\dnCh}{\dnCh}{\dnCh}{\dnCh}{\dnCh}{\dnCh}{\dnCh}{\dnCh}{\dnCh}{\dnCh}{\dnCh}{\dnCh}{\dnCh}{\dnCh}{\dnCh}{\dnCh}{\dnCh}{\dnCh}{\dnCh}{\dnCh}{\dnCh}{\dnCh}{\dnCh}{\dnCh}} & \dnW[E]{}: & {{\dnCh}{\dnCh}{\dnCh}{\dnCh}{\dnCh}{\dnCh}{\dnCh}{\dnCh}{\dnCh}{\dnCh}{\dnCh}{\dnCh}{\dnCh}{\dnCh}{\dnCh}{\dnCh}{\dnCh}{\dnCh}{\dnCh}{\dnCh}{\dnCh}{\dnCh}{\dnCh}{\dnCh}{\dnCh}{\dnCh}{\dnCh}{\dnCh}{\dnCh}{\dnCh}{\dnCh}{\dnCh}} & \dnW[W]{}: & {{\dnCh}{\dnCh}{\dnCh}{\dnCh}{\dnCh}{\dnCh}{\dnCh}{\dnCh}{\dnCh}{\dnCh}{\dnCh}{\dnCh}{\dnCh}{\dnCh}{\dnCh}{\dnCh}{\dnCh}{\dnCh}{\dnCh}{\dnCh}{\dnCh}{\dnCh}{\dnCh}{\dnCh}{\dnCh}{\dnCh}{\dnCh}{\dnCh}{\dnCh}{\dnCh}{\dnCh}{\dnCh}} \\
 \dnI{13} & \dnW: & {{\dnCh}{\dnCh}{\dnCh}{\dnCh}{\dnCh}{\dnCh}{\dnCh}{\dnCh}{\dnCh}{\dnCh}{\dnCh}{\dnCh}{\dnCh}{\dnCh}{\dnCh}{\dnCh}{\dnCh}{\dnCh}{\dnCh}{\dnCh}{\dnCh}{\dnCh}{\dnCh}{\dnCh}{\dnCh}{\dnCh}{\dnCh}{\dnCh}{\dnCh}{\dnCh}{\dnCh}{\dnCh}} & \dnW[E]{}: & {{\dnCh}{\dnCh}{\dnCh}{\dnCh}{\dnCh}{\dnCh}{\dnCh}{\dnCh}{\dnCh}{\dnCh}{\dnCh}{\dnCh}{\dnCh}{\dnCh}{\dnCh}{\dnCh}{\dnCh}{\dnCh}{\dnCh}{\dnCh}{\dnCh}{\dnCh}{\dnCh}{\dnCh}{\dnCh}{\dnCh}{\dnCh}{\dnCh}{\dnCh}{\dnCh}{\dnCh}{\dnCh}} & \dnW[W]{}: & {{\dnCh}{\dnCh}{\dnCh}{\dnCh}{\dnCh}{\dnCh}{\dnCh}{\dnCh}{\dnCh}{\dnCh}{\dnCh}{\dnCh}{\dnCh}{\dnCh}{\dnCh}{\dnCh}{\dnCh}{\dnCh}{\dnCh}{\dnCh}{\dnCh}{\dnCh}{\dnCh}{\dnCh}{\dnCh}{\dnCh}{\dnCh}{\dnCh}{\dnCh}{\dnCh}{\dnCh}{\dnCh}} \\
 \dnI{14} & \dnW: & {{\dnCh}{\dnCh}{\dnCh}{\dnCh}{\dnCh}{\dnCh}{\dnCh}{\dnCh}{\dnCh}{\dnCh}{\dnCh}{\dnCh}{\dnCh}{\dnCh}{\dnCh}{\dnCh}{\dnCh}{\dnCh}{\dnCh}{\dnCh}{\dnCh}{\dnCh}{\dnCh}{\dnCh}{\dnCh}{\dnCh}{\dnCh}{\dnCh}{\dnCh}{\dnCh}{\dnCh}{\dnCh}} & \dnW[E]{}: & {{\dnCh}{\dnCh}{\dnCh}{\dnCh}{\dnCh}{\dnCh}{\dnCh}{\dnCh}{\dnCh}{\dnCh}{\dnCh}{\dnCh}{\dnCh}{\dnCh}{\dnCh}{\dnCh}{\dnCh}{\dnCh}{\dnCh}{\dnCh}{\dnCh}{\dnCh}{\dnCh}{\dnCh}{\dnCh}{\dnCh}{\dnCh}{\dnCh}{\dnCh}{\dnCh}{\dnCh}{\dnCh}} & \dnW[W]{}: & {{\dnCh}{\dnCh}{\dnCh}{\dnCh}{\dnCh}{\dnCh}{\dnCh}{\dnCh}{\dnCh}{\dnCh}{\dnCh}{\dnCh}{\dnCh}{\dnCh}{\dnCh}{\dnCh}{\dnCh}{\dnCh}{\dnCh}{\dnCh}{\dnCh}{\dnCh}{\dnCh}{\dnCh}{\dnCh}{\dnCh}{\dnCh}{\dnCh}{\dnCh}{\dnCh}{\dnCh}{\dnCh}} \\
 \dnI{15} & \dnW: & {{\dnCh}{\dnCh}{\dnCh}{\dnCh}{\dnCh}{\dnCh}{\dnCh}{\dnCh}{\dnCh}{\dnCh}{\dnCh}{\dnCh}{\dnCh}{\dnCh}{\dnCh}{\dnCh}{\dnCh}{\dnCh}{\dnCh}{\dnCh}{\dnCh}{\dnCh}{\dnCh}{\dnCh}{\dnCh}{\dnCh}{\dnCh}{\dnCh}{\dnCh}{\dnCh}{\dnCh}{\dnCh}} & \dnW[E]{}: & {{\dnCh}{\dnCh}{\dnCh}{\dnCh}{\dnCh}{\dnCh}{\dnCh}{\dnCh}{\dnCh}{\dnCh}{\dnCh}{\dnCh}{\dnCh}{\dnCh}{\dnCh}{\dnCh}{\dnCh}{\dnCh}{\dnCh}{\dnCh}{\dnCh}{\dnCh}{\dnCh}{\dnCh}{\dnCh}{\dnCh}{\dnCh}{\dnCh}{\dnCh}{\dnCh}{\dnCh}{\dnCh}} & \dnW[W]{}: & {{\dnCq}{\dnCq}{\dnCq}{\dnCq}{\dnCq}{\dnCq}{\dnCq}{\dnCq}{\dnCq}{\dnCq}{\dnCq}{\dnCq}{\dnCq}{\dnCq}{\dnCq}{\dnCq}{\dnCq}{\dnCq}{\dnCq}{\dnCq}{\dnCq}{\dnCq}{\dnCq}{\dnCq}{\dnCq}{\dnCq}{\dnCq}{\dnCq}{\dnCq}{\dnCq}{\dnCq}{\dnCq}} \\
 \dnI{16} & \dnW: & {{\dnCh}{\dnCh}{\dnCh}{\dnCh}{\dnCh}{\dnCh}{\dnCh}{\dnCh}{\dnCh}{\dnCh}{\dnCh}{\dnCh}{\dnCh}{\dnCh}{\dnCh}{\dnCh}{\dnCh}{\dnCh}{\dnCh}{\dnCh}{\dnCh}{\dnCh}{\dnCh}{\dnCh}{\dnCh}{\dnCh}{\dnCh}{\dnCh}{\dnCh}{\dnCh}{\dnCh}{\dnCh}} & \dnW[E]{}: & {{\dnCh}{\dnCh}{\dnCh}{\dnCh}{\dnCh}{\dnCh}{\dnCh}{\dnCh}{\dnCh}{\dnCh}{\dnCh}{\dnCh}{\dnCh}{\dnCh}{\dnCh}{\dnCh}{\dnCh}{\dnCh}{\dnCh}{\dnCh}{\dnCh}{\dnCh}{\dnCh}{\dnCh}{\dnCh}{\dnCh}{\dnCh}{\dnCh}{\dnCh}{\dnCh}{\dnCh}{\dnCh}} & \dnW[W]{}: & {{\dnCh}{\dnCh}{\dnCh}{\dnCh}{\dnCh}{\dnCh}{\dnCh}{\dnCh}{\dnCh}{\dnCh}{\dnCh}{\dnCh}{\dnCh}{\dnCh}{\dnCh}{\dnCh}{\dnCh}{\dnCh}{\dnCh}{\dnCh}{\dnCh}{\dnCh}{\dnCh}{\dnCh}{\dnCh}{\dnCh}{\dnCh}{\dnCh}{\dnCh}{\dnCh}{\dnCh}{\dnCh}} \\
 \dnI{17} & \dnW: & {{\dnCh}{\dnCh}{\dnCh}{\dnCh}{\dnCh}{\dnCh}{\dnCh}{\dnCh}{\dnCh}{\dnCh}{\dnCh}{\dnCh}{\dnCh}{\dnCh}{\dnCh}{\dnCh}{\dnCh}{\dnCh}{\dnCh}{\dnCh}{\dnCh}{\dnCh}{\dnCh}{\dnCh}{\dnCh}{\dnCh}{\dnCh}{\dnCh}{\dnCh}{\dnCh}{\dnCh}{\dnCh}} & \dnW[E]{}: & {{\dnCh}{\dnCh}{\dnCh}{\dnCh}{\dnCh}{\dnCh}{\dnCh}{\dnCh}{\dnCh}{\dnCh}{\dnCh}{\dnCh}{\dnCh}{\dnCh}{\dnCh}{\dnCh}{\dnCh}{\dnCh}{\dnCh}{\dnCh}{\dnCh}{\dnCh}{\dnCh}{\dnCh}{\dnCh}{\dnCh}{\dnCh}{\dnCh}{\dnCh}{\dnCh}{\dnCh}{\dnCh}} & \dnW[W]{}: & {{\dnCh}{\dnCh}{\dnCh}{\dnCh}{\dnCh}{\dnCh}{\dnCh}{\dnCh}{\dnCh}{\dnCh}{\dnCh}{\dnCh}{\dnCh}{\dnCh}{\dnCh}{\dnCh}{\dnCh}{\dnCh}{\dnCh}{\dnCh}{\dnCh}{\dnCh}{\dnCh}{\dnCh}{\dnCh}{\dnCh}{\dnCh}{\dnCh}{\dnCh}{\dnCh}{\dnCh}{\dnCh}} \\
 \dnI{18} & \dnW: & {{\dnCh}{\dnCh}{\dnCh}{\dnCh}{\dnCh}{\dnCh}{\dnCh}{\dnCh}{\dnCh}{\dnCh}{\dnCh}{\dnCh}{\dnCh}{\dnCh}{\dnCh}{\dnCh}{\dnCh}{\dnCh}{\dnCh}{\dnCh}{\dnCh}{\dnCh}{\dnCh}{\dnCh}{\dnCh}{\dnCh}{\dnCh}{\dnCh}{\dnCh}{\dnCh}{\dnCh}{\dnCh}} & \dnW[E]{}: & {{\dnCh}{\dnCh}{\dnCh}{\dnCh}{\dnCh}{\dnCh}{\dnCh}{\dnCh}{\dnCh}{\dnCh}{\dnCh}{\dnCh}{\dnCh}{\dnCh}{\dnCh}{\dnCh}{\dnCh}{\dnCh}{\dnCh}{\dnCh}{\dnCh}{\dnCh}{\dnCh}{\dnCh}{\dnCh}{\dnCh}{\dnCh}{\dnCh}{\dnCh}{\dnCh}{\dnCh}{\dnCh}} & \dnW[W]{}: & {{\dnCh}{\dnCh}{\dnCh}{\dnCh}{\dnCh}{\dnCh}{\dnCh}{\dnCh}{\dnCh}{\dnCh}{\dnCh}{\dnCh}{\dnCh}{\dnCh}{\dnCh}{\dnCh}{\dnCh}{\dnCh}{\dnCh}{\dnCh}{\dnCh}{\dnCh}{\dnCh}{\dnCh}{\dnCh}{\dnCh}{\dnCh}{\dnCh}{\dnCh}{\dnCh}{\dnCh}{\dnCh}} \\
 \dnI{19} & \dnW: & {{\dnCh}{\dnCh}{\dnCh}{\dnCh}{\dnCh}{\dnCh}{\dnCh}{\dnCh}{\dnCh}{\dnCh}{\dnCh}{\dnCh}{\dnCh}{\dnCh}{\dnCh}{\dnCh}{\dnCh}{\dnCh}{\dnCh}{\dnCh}{\dnCh}{\dnCh}{\dnCh}{\dnCh}{\dnCh}{\dnCh}{\dnCh}{\dnCh}{\dnCh}{\dnCh}{\dnCh}{\dnCh}} & \dnW[E]{}: & {{\dnCh}{\dnCh}{\dnCh}{\dnCh}{\dnCh}{\dnCh}{\dnCh}{\dnCh}{\dnCh}{\dnCh}{\dnCh}{\dnCh}{\dnCh}{\dnCh}{\dnCh}{\dnCh}{\dnCh}{\dnCh}{\dnCh}{\dnCh}{\dnCh}{\dnCh}{\dnCh}{\dnCh}{\dnCh}{\dnCh}{\dnCh}{\dnCh}{\dnCh}{\dnCh}{\dnCh}{\dnCh}} & \dnW[W]{}: & {{\dnCh}{\dnCh}{\dnCh}{\dnCh}{\dnCh}{\dnCh}{\dnCh}{\dnCh}{\dnCh}{\dnCh}{\dnCh}{\dnCh}{\dnCh}{\dnCh}{\dnCh}{\dnCh}{\dnCh}{\dnCh}{\dnCh}{\dnCh}{\dnCh}{\dnCh}{\dnCh}{\dnCh}{\dnCh}{\dnCh}{\dnCh}{\dnCh}{\dnCh}{\dnCh}{\dnCh}{\dnCh}} \\
 \dnI{20} & \dnW: & {{\dnCh}{\dnCh}{\dnCh}{\dnCh}{\dnCh}{\dnCh}{\dnCh}{\dnCh}{\dnCh}{\dnCh}{\dnCh}{\dnCh}{\dnCh}{\dnCh}{\dnCh}{\dnCh}{\dnCh}{\dnCh}{\dnCh}{\dnCh}{\dnCh}{\dnCh}{\dnCh}{\dnCh}{\dnCh}{\dnCh}{\dnCh}{\dnCh}{\dnCh}{\dnCh}{\dnCh}{\dnCh}} & \dnW[E]{}: & {{\dnCh}{\dnCh}{\dnCh}{\dnCh}{\dnCh}{\dnCh}{\dnCh}{\dnCh}{\dnCh}{\dnCh}{\dnCh}{\dnCh}{\dnCh}{\dnCh}{\dnCh}{\dnCh}{\dnCh}{\dnCh}{\dnCh}{\dnCh}{\dnCh}{\dnCh}{\dnCh}{\dnCh}{\dnCh}{\dnCh}{\dnCh}{\dnCh}{\dnCh}{\dnCh}{\dnCh}{\dnCh}} & \dnW[W]{}: & {{\dnCh}{\dnCh}{\dnCh}{\dnCh}{\dnCh}{\dnCh}{\dnCh}{\dnCh}{\dnCh}{\dnCh}{\dnCh}{\dnCh}{\dnCh}{\dnCh}{\dnCh}{\dnCh}{\dnCh}{\dnCh}{\dnCh}{\dnCh}{\dnCh}{\dnCh}{\dnCh}{\dnCh}{\dnCh}{\dnCh}{\dnCh}{\dnCh}{\dnCh}{\dnCh}{\dnCh}{\dnCh}} \\
 \dnI{21} & \dnW: & {{\dnCh}{\dnCh}{\dnCh}{\dnCh}{\dnCh}{\dnCh}{\dnCh}{\dnCh}{\dnCh}{\dnCh}{\dnCh}{\dnCh}{\dnCh}{\dnCh}{\dnCh}{\dnCh}{\dnCh}{\dnCh}{\dnCh}{\dnCh}{\dnCh}{\dnCh}{\dnCh}{\dnCh}{\dnCh}{\dnCh}{\dnCh}{\dnCh}{\dnCh}{\dnCh}{\dnCh}{\dnCh}} & \dnW[E]{}: & {{\dnCh}{\dnCh}{\dnCh}{\dnCh}{\dnCh}{\dnCh}{\dnCh}{\dnCh}{\dnCh}{\dnCh}{\dnCh}{\dnCh}{\dnCh}{\dnCh}{\dnCh}{\dnCh}{\dnCh}{\dnCh}{\dnCh}{\dnCh}{\dnCh}{\dnCh}{\dnCh}{\dnCh}{\dnCh}{\dnCh}{\dnCh}{\dnCh}{\dnCh}{\dnCh}{\dnCh}{\dnCh}} & \dnW[W]{}: & {{\dnCh}{\dnCh}{\dnCh}{\dnCh}{\dnCh}{\dnCh}{\dnCh}{\dnCh}{\dnCh}{\dnCh}{\dnCh}{\dnCh}{\dnCh}{\dnCh}{\dnCh}{\dnCh}{\dnCh}{\dnCh}{\dnCh}{\dnCh}{\dnCh}{\dnCh}{\dnCh}{\dnCh}{\dnCh}{\dnCh}{\dnCh}{\dnCh}{\dnCh}{\dnCh}{\dnCh}{\dnCh}} \\
 \dnI{22} & \dnW: & {{\dnCh}{\dnCh}{\dnCh}{\dnCh}{\dnCh}{\dnCh}{\dnCh}{\dnCh}{\dnCh}{\dnCh}{\dnCh}{\dnCh}{\dnCh}{\dnCh}{\dnCh}{\dnCh}{\dnCh}{\dnCh}{\dnCh}{\dnCh}{\dnCh}{\dnCh}{\dnCh}{\dnCh}{\dnCh}{\dnCh}{\dnCh}{\dnCh}{\dnCh}{\dnCh}{\dnCh}{\dnCh}} & \dnW[E]{}: & {{\dnCh}{\dnCh}{\dnCh}{\dnCh}{\dnCh}{\dnCh}{\dnCh}{\dnCh}{\dnCh}{\dnCh}{\dnCh}{\dnCh}{\dnCh}{\dnCh}{\dnCh}{\dnCh}{\dnCh}{\dnCh}{\dnCh}{\dnCh}{\dnCh}{\dnCh}{\dnCh}{\dnCh}{\dnCh}{\dnCh}{\dnCh}{\dnCh}{\dnCh}{\dnCh}{\dnCh}{\dnCh}} & \dnW[W]{}: & {{\dnCh}{\dnCh}{\dnCh}{\dnCh}{\dnCh}{\dnCh}{\dnCh}{\dnCh}{\dnCh}{\dnCh}{\dnCh}{\dnCh}{\dnCh}{\dnCh}{\dnCh}{\dnCh}{\dnCh}{\dnCh}{\dnCh}{\dnCh}{\dnCh}{\dnCh}{\dnCh}{\dnCh}{\dnCh}{\dnCh}{\dnCh}{\dnCh}{\dnCh}{\dnCh}{\dnCh}{\dnCh}} \\
 \dnI{23} & \dnW: & {{\dnCh}{\dnCh}{\dnCh}{\dnCh}{\dnCh}{\dnCh}{\dnCh}{\dnCh}{\dnCh}{\dnCh}{\dnCh}{\dnCh}{\dnCh}{\dnCh}{\dnCh}{\dnCh}{\dnCh}{\dnCh}{\dnCh}{\dnCh}{\dnCh}{\dnCh}{\dnCh}{\dnCh}{\dnCh}{\dnCh}{\dnCh}{\dnCh}{\dnCh}{\dnCh}{\dnCh}{\dnCh}} & \dnW[E]{}: & {{\dnCh}{\dnCh}{\dnCh}{\dnCh}{\dnCh}{\dnCh}{\dnCh}{\dnCh}{\dnCh}{\dnCh}{\dnCh}{\dnCh}{\dnCh}{\dnCh}{\dnCh}{\dnCh}{\dnCh}{\dnCh}{\dnCh}{\dnCh}{\dnCh}{\dnCh}{\dnCh}{\dnCh}{\dnCh}{\dnCh}{\dnCh}{\dnCh}{\dnCh}{\dnCh}{\dnCh}{\dnCh}} & \dnW[W]{}: & {{\dnCh}{\dnCh}{\dnCh}{\dnCh}{\dnCh}{\dnCh}{\dnCh}{\dnCh}{\dnCh}{\dnCh}{\dnCh}{\dnCh}{\dnCh}{\dnCh}{\dnCh}{\dnCh}{\dnCh}{\dnCh}{\dnCh}{\dnCh}{\dnCh}{\dnCh}{\dnCh}{\dnCh}{\dnCh}{\dnCh}{\dnCh}{\dnCh}{\dnCh}{\dnCh}{\dnCh}{\dnCh}} \\
\hline
\end{tabular}
}
\caption{SHA-256 hash collision over 24 rounds}
\label{fig:tc24}
\end{center}
\end{figure}
}

\chapter{Runtimes retrieved}
\label{app:runtimes}

TODO
\begin{verbatim}
FAILED means unexpected exit code or missing results.
TIMEOUT mean did not terminate within 1 day.
SUCCESS means SAT within 1 day.

Results cover 50 CNFs
               cmsat5   83 occurences
              minisat  169 occurences
     lingeling-ats1o4  267 occurences
         minisat-core   85 occurences
        glucose-syrup   84 occurences
     lingeling-ats1o1  1145 occurences
   treengeling-ats1o1   84 occurences
     lingeling-ats1o2  250 occurences
    plingeling-ats1o1   85 occurences
               cmsat4   84 occurences
              glucose   84 occurences
FAILED=303  TIMEOUT=722  SUCCESS=1395


\end{verbatim}
\tiny{
\begin{table}[p]
  \begin{center}
    \begin{tabular}{l|cccccccc}
        \textbf{SAT solver} & \textbf{CNF file} & \textbf{simplified} & \textbf{phase} & \textbf{vars/clauses} & \textbf{propagations} & \textbf{decisions} & \textbf{restarts} & \textbf{Runtime (sec)} \\
      \hline
  cmsat4                         & algotocnf\_1.cnf               & none       & no    & 48704 / 253984 &           &           &            & timeout \\ % a8be7416eae9919a0caa3bd0b8637ab442032c9233f342549b59c8ec
  cmsat4                         & algotocnf\_1.cnf               & cmsat      & no    & 24503 / 111931 &           &           &            & timeout \\ % 5e3667f271cbc169b7c334faee19773cdc5ab554f94b2c932728fe5e
  cmsat4                         & algotocnf\_1.cnf               & minisat    & no    & 20895 / 118236 &           &           &            & timeout \\ % d9436aad4a57f47d9f1f6b85b43ded925063e94691ab3fabc1e3aa81
  cmsat4                         & algotocnf\_1.cnf               & satelite   & no    & 27495 / 153262 &           &           &            & timeout \\ % c10546657067d7c1e10aa679300aee7d4189e050dd45ba2cca779294
  cmsat4                         & algotocnf\_18-t9.cnf           & none       & no    & 107839 / 590953 & 140000    & 2441566   & 539        & 26 \\ % 9e8032a9087d0f7ea8ff0f33aa1019d4a5a3d57af8e02496db9703c0
  cmsat4                         & algotocnf\_18-t9.cnf           & cmsat      & no    & 61789 / 342139 & 0         & 2514784   & 226        & 17 \\ % 9d327b5c25933b581b2922634d1f4ad364c3b2d456e9d792eb308c1d
  cmsat4                         & algotocnf\_18-t9.cnf           & lingeling  & no    & 107839 / 344544 & 90000     & 1565224   & 519        & 28 \\ % 903a1c78259e4fe79d3a25f2b1d9a35f3c527af03bb11c2d79abcd90
  cmsat4                         & algotocnf\_18-t9.cnf           & minisat    & no    & 61845 / 391134 & 120000    & 2046800   & 781        & 43 \\ % b63c9799924b593dc452b6aebd154773b217dba69a9f2f1ce1ae3bc6
  cmsat4                         & algotocnf\_18-t9.cnf           & satelite   & no    & 69670 / 457972 & 0         & 1330116   & 195        & 13 \\ % f965bd7199b99982ef287a17dc73d8f525ceffd84a80afd672d84d0c
  cmsat4                         & algotocnf\_18-t9\_diff-desc.cnf & none       & no    & 107839 / 846487 & 140000    & 3639049   & 784        & 43 \\ % 8532a45aaa3b41f900152f51b0f08b580e9df3c938cb43f292b46a5e
  cmsat4                         & algotocnf\_18-t9\_diff-desc.cnf & cmsat      & no    & 57114 / 303489 & 120000    & 2048109   & 364        & 30 \\ % 15981f2f7dc3e7cb6ac06a3e76a85dd80066ed98b848268d04b10715
  cmsat4                         & algotocnf\_18-t9\_diff-desc.cnf & lingeling  & no    & 107839 / 404666 & 0         & 784078    & 242        & 13 \\ % dfe922f7fe2265c9d117051a0e76bf6472834ec84c12e043eac9f87f
  cmsat4                         & algotocnf\_18-t9\_diff-desc.cnf & minisat    & no    & 78055 / 511431 & 0         & 1027604   & 248        & 19 \\ % 78e45ada2268f0b1da5f34aab88e8c2ea2b6515bf6e9de728d90aca5
  cmsat4                         & algotocnf\_18-t9\_diff-desc.cnf & satelite   & no    & 85386 / 692408 & 140000    & 2964218   & 480        & 60 \\ % faf52268c99017ba15b1b450f1fdbc56bdcef399160c9ec08a9dcc10
  cmsat4                         & algotocnf\_1\_diff-desc.cnf    & none       & no    & 48704 / 373920 & 860000    & 31850316  & 43299      & 3750 \\ % 9c0b9456086f69edaf19cd3b62859c44abe858fb70116a238793cd45
  cmsat4                         & algotocnf\_1\_diff-desc.cnf    & cmsat      & no    & 15402 / 73142 & 520000    & 11675857  & 12466      & 900 \\ % 515f5cc8c57c9ae7042dd49b09b5b4ada106529aa275a793e2aa046b
  cmsat4                         & algotocnf\_1\_diff-desc.cnf    & lingeling  & no    & 48704 / 101182 & 530000    & 21126661  & 24487      & 1724 \\ % f89a032141ac530d98b2846c8c4d4fe2e20cb51ea72ced8e21adcf24
  cmsat4                         & algotocnf\_1\_diff-desc.cnf    & minisat    & no    & 25107 / 158393 & 800000    & 22305937  & 28680      & 1852 \\ % 28d3ff9f2925de3916f6611da35536f63c69c377a73b5f85aeb45387
  cmsat4                         & algotocnf\_1\_diff-desc.cnf    & satelite   & no    & 30950 / 226602 & 780000    & 20569654  & 25267      & 1419 \\ % 1206dbc0d6fba52f8731a1d60fd30b1f380af779d38002f480026c27
  cmsat4                         & algotocnf\_2.cnf               & none       & no    & 48704 / 254210 &           &           &            & timeout \\ % d5c046144656768ff3920fd7d7eb7fde35747aae74bc3d4e57844e29
  cmsat4                         & algotocnf\_2.cnf               & cmsat      & no    & 24323 / 111661 &           &           &            & timeout \\ % f9625d474f6a1c3eeee723c68b23d14625e3a5217ded4f5434afb8af
  cmsat4                         & algotocnf\_2.cnf               & minisat    & no    & 20895 / 118688 &           &           &            & timeout \\ % 82558a81dc773e79b92b0f299081dbf18c00ef0cf305e346d1df814e
  cmsat4                         & algotocnf\_2.cnf               & satelite   & no    & 27775 / 154276 &           &           &            & timeout \\ % 2cfb9f8bab37104ee5f57212d31c8e89f08ce48e422a03cfbf0607b2
  cmsat4                         & algotocnf\_21-t9.cnf           & none       & no    & 116800 / 636838 &           &           &            & timeout \\ % df87b2104a3e017b53cbf9bc3aac16081769eb72aadf08ad56263ff2
  cmsat4                         & algotocnf\_21-t9.cnf           & cmsat      & no    & 63555 / 358308 &           &           &            & timeout \\ % 8b62e387073c080e5ea8ae7a0cd4ba1592c5e398ba57a2491a2fbc12
  cmsat4                         & algotocnf\_21-t9.cnf           & minisat    & no    & 65572 / 414470 &           &           &            & timeout \\ % 430bc665d503881ee35024be9d37a21442f2fa6adebb77ebcd618a86
  cmsat4                         & algotocnf\_21-t9.cnf           & satelite   & no    & 74562 / 486714 &           &           &            & timeout \\ % 04dfc6c531c8d4c749e3c8994e3800812bfe5d8d0bb6db69e02caf83
  cmsat4                         & algotocnf\_21-t9\_diff-desc.cnf & none       & no    & 116800 / 911629 & 1540000   & 79320703  & 168537     & 29424 \\ % e6abc2db576b2ca33b351c34bb1c6141034331a6f32a91485dfab8df
  cmsat4                         & algotocnf\_21-t9\_diff-desc.cnf & cmsat      & no    & 59528 / 318049 & 1510000   & 64453001  & 115182     & 27359 \\ % 5bb2d101d0ecae3871ddca2608d7485c93a2aad44267de17a7463998
  cmsat4                         & algotocnf\_21-t9\_diff-desc.cnf & lingeling  & no    & 116800 / 404126 & 1290000   & 69278226  & 120139     & 25272 \\ % 6877f853aab53cb3d471b3c686bf24f73602df9b84f73d9fd30e5877
  cmsat4                         & algotocnf\_21-t9\_diff-desc.cnf & minisat    & no    & 81650 / 534965 & 1560000   & 82066672  & 168220     & 47023 \\ % b2cdbfd5ae43e67363eb6cfa0a9ce056a65058b4726c5a12d8403b25
  cmsat4                         & algotocnf\_21-t9\_diff-desc.cnf & satelite   & no    & 90236 / 720755 & 1540000   & 64208938  & 132149     & 20296 \\ % 0f04e343332a00aa2085011e9dd4d94af0d9f5915dae36791ce55192
  cmsat4                         & algotocnf\_23-t9.cnf           & none       & no    & 122774 / 667438 &           &           &            & timeout \\ % 41485df32b1320a8c25696f28321af5a2d448a4330fe8dcd746c3c43
  cmsat4                         & algotocnf\_23-t9.cnf           & cmsat      & no    & 64880 / 369775 &           &           &            & timeout \\ % fa4d0a346bedd144b7771b7fca7cf86362d8a6aa96c778b005c5b99a
  cmsat4                         & algotocnf\_23-t9.cnf           & minisat    & no    & 67979 / 429734 &           &           &            & timeout \\ % 07b2d6187c7a39ba8124561ca31857b30f3ff300fc5bfc53d7ef7787
  cmsat4                         & algotocnf\_23-t9.cnf           & satelite   & no    & 77835 / 506217 &           &           &            & timeout \\ % 5806d39a93f3eded34f750daf517feac4f07b5229bfbdd8f74b4d320
    \end{tabular}
  \end{center}
\end{table}
}
%Text such that \LaTeX{} won't screw up.
\newpage

\begin{table}[p]
  \begin{center}
    \begin{tabular}{l|cccccccc}
        \textbf{SAT solver} & \textbf{CNF file} & \textbf{simplified} & \textbf{phase} & \textbf{vars/clauses} & \textbf{propagations} & \textbf{decisions} & \textbf{restarts} & \textbf{Runtime (sec)} \\
      \hline
  cmsat4                         & algotocnf\_23-t9\_diff-desc.cnf & none       & no    & 122774 / 955067 &           &           &            & timeout \\ % 7d2fcc16402bd3b49f8c92a592ad9372dc96cec8f981d427589cf7e7
  cmsat4                         & algotocnf\_23-t9\_diff-desc.cnf & cmsat      & no    & 57510 / 307118 &           &           &            & timeout \\ % f2824cb3f655e0a8ac4fbea52dcc5ac41073cc15d6d5c42366ce1dcd
  cmsat4                         & algotocnf\_23-t9\_diff-desc.cnf & minisat    & no    & 84080 / 550727 &           &           &            & timeout \\ % 42906a7c58f02b308a31d68f6316d4642d74bd7193139e6fa3c166e2
  cmsat4                         & algotocnf\_23-t9\_diff-desc.cnf & satelite   & no    & 93504 / 740545 &           &           &            & timeout \\ % dbe9bb10ede97e2a70637b1e5133b3b335474f01976652245f203a7a
  cmsat4                         & algotocnf\_24-t9.cnf           & none       & no    & 125761 / 682722 &           &           &            & timeout \\ % 34a0ddd515657c7219863dbc75015eae950548664723fb6d10795bd6
  cmsat4                         & algotocnf\_24-t9.cnf           & cmsat      & no    & 65579 / 377228 &           &           &            & timeout \\ % 457035ee7907153fcc12fa87e5d9d46d12b4b2203c23a370001f5f8b
  cmsat4                         & algotocnf\_24-t9.cnf           & minisat    & no    & 69023 / 436292 &           &           &            & timeout \\ % 2994b853be23a1549b1d3e1b0bdb580349f4e95a1b01e87c18a30a3c
  cmsat4                         & algotocnf\_24-t9.cnf           & satelite   & no    & 79486 / 515977 &           &           &            & timeout \\ % 4081a43018db35b73febdd919d8de782084e420749895a177567a209
  cmsat4                         & algotocnf\_24-t9\_diff-desc.cnf & none       & no    & 125761 / 976770 &           &           &            & timeout \\ % df927e6f18a778df1146b3d90c6c32962f83e86a6038a31566411d71
  cmsat4                         & algotocnf\_24-t9\_diff-desc.cnf & cmsat      & no    & 56447 / 302014 &           &           &            & timeout \\ % 718ef1e2d3c23037edb4d934cf5556498ccb78d47dfd39a5ce41ea2f
  cmsat4                         & algotocnf\_24-t9\_diff-desc.cnf & minisat    & no    & 85312 / 558731 &           &           &            & timeout \\ % 69bc4d7b015a44c59902cf0499cd7bd5b0e5510a52e268370d2975f6
  cmsat4                         & algotocnf\_24-t9\_diff-desc.cnf & satelite   & no    & 95154 / 750557 &           &           &            & timeout \\ % 743da2b7329484f146fe8ea62d6a1134e2c212ba644cac27ab87fdf0
  cmsat4                         & algotocnf\_27-t10-diff.cnf     & none       & no    & 134722 / 728349 & 0         & 0         & 1          & 1 \\ % eac14926339500a1acd8b6998e437692cc40a66191ad0a6df842e943
  cmsat4                         & algotocnf\_27-t10-diff.cnf     & cmsat      & no    & 0 / 0      & 0         & 0         & 1          & 0 \\ % b5a80ace889d63c78487bf2b2608c84c167c2cf6f8df9c3b58f94186
  cmsat4                         & algotocnf\_27-t10-diff.cnf     & lingeling  & no    & 134722 / 0 & 0         & 0         & 1          & 0 \\ % 3bec9a95d0ab3c2f003a65980d557404519da59cb905bc28a896fad8
  cmsat4                         & algotocnf\_27-t10-diff.cnf     & minisat    & no    & 0 / 0      & 0         & 0         & 1          & 0 \\ % 55068ef37ca69ef3e19dc6814904a773cf8fce9b995c472dae5597f2
  cmsat4                         & algotocnf\_27-t10-diff.cnf     & satelite   & no    & 0 / 0      & 0         & 0         & 1          & 0 \\ % 91e6fa002dc4e65075d1d98d07bea5aa3b687858a59c9383960f32d6
  cmsat4                         & algotocnf\_27-t10-diff\_diff-desc.cnf & cmsat      & no    & 0 / 0      & 0         & 0         & 1          & 0 \\ % 119b97307b6cdf5d416387f360f3083dc51157df6ed5410d0633bd49
  cmsat4                         & algotocnf\_27-t10-diff\_diff-desc.cnf & lingeling  & no    & 134722 / 0 & 0         & 0         & 1          & 0 \\ % 81119eb4d074fdac2df1cfcee056ebf8adf04bced4c8f390f9809212
  cmsat4                         & algotocnf\_27-t10-diff\_diff-desc.cnf & minisat    & no    & 0 / 0      & 0         & 0         & 1          & 0 \\ % 7c19d9a2c4e4b397ac725f94a8f1ec9ff57c1649745df9eaba144016
  cmsat4                         & algotocnf\_27-t10-diff\_diff-desc.cnf & satelite   & no    & 0 / 0      & 0         & 0         & 1          & 0 \\ % 95c235675e947721a91bc49e7639e727e5d178709a313111bf0949aa
  cmsat4                         & algotocnf\_27-t10-simplified.cnf & none       & no    & 134722 / 728542 &           &           &            & timeout \\ % 139009542b37a0b7179e0eaf69c5a1884c3c7450d95c219687ae97df
  cmsat4                         & algotocnf\_27-t10-simplified.cnf & cmsat      & no    & 67175 / 392680 &           &           &            & timeout \\ % 95b0d3147a6196360d20e8943cac799bc2491e2f6f24a493fc132e82
  cmsat4                         & algotocnf\_27-t10-simplified.cnf & minisat    & no    & 72909 / 460404 &           &           &            & timeout \\ % b8854c5b5e98ffea68d044a756bf11f33ff53a48410197361f0c59ed
  cmsat4                         & algotocnf\_27-t10-simplified.cnf & satelite   & no    & 84442 / 544998 &           &           &            & timeout \\ % 1955fe23db5432840d6bc2028732337b122c61729916248e2c356b75
  cmsat4                         & algotocnf\_27-t10-simplified\_diff-desc.cnf & none       & no    & 134722 / 1041847 &           &           &            & timeout \\ % aca4ab10e8b82e05620bf281fe849fb3a39fbe7bf6cb337d8f3bb266
  cmsat4                         & algotocnf\_27-t10-simplified\_diff-desc.cnf & cmsat      & no    & 56223 / 302443 &           &           &            & timeout \\ % 2ef3e8c7a60b5a2eef0d1e253545948c01b9841c6989f95a1fac62a4
  cmsat4                         & algotocnf\_27-t10-simplified\_diff-desc.cnf & minisat    & no    & 89667 / 586309 &           &           &            & timeout \\ % a04b2dd6843b18635dc75893ead62986ed9359478e11324b54c078ce
  cmsat4                         & algotocnf\_27-t10-simplified\_diff-desc.cnf & satelite   & no    & 100735 / 789589 &           &           &            & timeout \\ % 87457ff7e3c3d63ed235cd9e4795213e2fc61ee52308747df55a4bc4
  cmsat4                         & algotocnf\_27-t10.cnf          & none       & no    & 134722 / 728030 &           &           &            & timeout \\ % 5722db87957078b771188c427c6b3eda4b8e61898d382fbc802c2014
  cmsat4                         & algotocnf\_27-t10.cnf          & cmsat      & no    & 72044 / 411276 &           &           &            & timeout \\ % 79938a97621cc3643c4327b50bb615263bcf3453778a193182223718
  cmsat4                         & algotocnf\_27-t10.cnf          & minisat    & no    & 73897 / 465248 &           &           &            & timeout \\ % 185ef9793d5c4973c0652744c287f2470a25c52db7b27ecc294f5a72
  cmsat4                         & algotocnf\_27-t10.cnf          & satelite   & no    & 85274 / 549734 &           &           &            & timeout \\ % b631f9afe7ef6bde7aa908a525f9613dcef0ab3dcfdb3b1390a93580
  cmsat4                         & algotocnf\_27-t10\_diff-desc.cnf & none       & no    & 134722 / 1041335 &           &           &            & timeout \\ % 2ea4caaac14f0f8d63d8f79036d6e3e07bf2071d84725f4d182cbb2e
  cmsat4                         & algotocnf\_27-t10\_diff-desc.cnf & cmsat      & no    & 57956 / 311219 &           &           &            & timeout \\ % c0be2abb3e4120c553142429324144992c081b8616f0eb0214a9fb46
    \end{tabular}
  \end{center}
\end{table}

%Text such that \LaTeX{} won't screw up.
\newpage

\begin{table}[p]
  \begin{center}
    \begin{tabular}{l|cccccccc}
        \textbf{SAT solver} & \textbf{CNF file} & \textbf{simplified} & \textbf{phase} & \textbf{vars/clauses} & \textbf{propagations} & \textbf{decisions} & \textbf{restarts} & \textbf{Runtime (sec)} \\
      \hline
  cmsat4                         & algotocnf\_27-t10\_diff-desc.cnf & minisat    & no    & 94850 / 619480 &           &           &            & timeout \\ % b2194617604be63c175321ac8009f80034d7bb001894bfe06fa0762a
  cmsat4                         & algotocnf\_27-t10\_diff-desc.cnf & satelite   & no    & 105733 / 857675 &           &           &            & timeout \\ % fb2b8f02f16fc04a7515d1949745fb2085824c2a3f2cb6376bd656ef
  cmsat4                         & algotocnf\_2\_diff-desc.cnf    & none       & no    & 48704 / 374146 & 950000    & 37014401  & 53532      & 5499 \\ % 964753c16bd56272b9d5b480d7e131098271b0bba7987d39247bb110
  cmsat4                         & algotocnf\_2\_diff-desc.cnf    & cmsat      & no    & 15342 / 72870 & 550000    & 12729243  & 13482      & 982 \\ % 3cedbf1076dbe38494f32b137e967f499d50e9830f30c8318d39c043
  cmsat4                         & algotocnf\_2\_diff-desc.cnf    & lingeling  & no    & 48704 / 106546 & 490000    & 15443185  & 16887      & 1363 \\ % 181004203ecd952fcfc47d39524b54c686948517e738950e6a4e5409
  cmsat4                         & algotocnf\_2\_diff-desc.cnf    & minisat    & no    & 23531 / 149436 & 720000    & 15364610  & 17914      & 1467 \\ % 9b4f6049f601b92485fdd210841f32c416bbb2aa04aa6fc1ccb10a0d
  cmsat4                         & algotocnf\_2\_diff-desc.cnf    & satelite   & no    & 29428 / 204577 & 700000    & 12819538  & 14227      & 741 \\ % 0b65341a9f94567dc504c8bc7c01d9a3e37fc192029e229da00c05f1
  cmsat4                         & algotocnf\_3.cnf               & none       & no    & 48704 / 254656 & 70000     & 2080042   & 626        & 24 \\ % 0318b5d869ff885f721734718fb97b1aeff614f97c4731815f05af8b
  cmsat4                         & algotocnf\_3.cnf               & cmsat      & no    & 23893 / 111403 & 70000     & 1140160   & 444        & 18 \\ % 955fbe66e167db3abd9157b6b5c2771b4765487d56daecb377db9a9a
  cmsat4                         & algotocnf\_3.cnf               & lingeling  & no    & 48704 / 100941 & 0         & 631534    & 188        & 5 \\ % 38029959cfd94704a4740fb06f00d7ee2e503127c1f25d573ab9ce92
  cmsat4                         & algotocnf\_3.cnf               & minisat    & no    & 20895 / 119580 & 60000     & 730694    & 334        & 12 \\ % 1d635cd8fb24a930138bbb885b848296ab11b52e89c50439d46935ba
  cmsat4                         & algotocnf\_3.cnf               & satelite   & no    & 28647 / 156910 & 60000     & 3084641   & 652        & 32 \\ % 6b45546558be97ee6f6dbc455ed0cdfafe47583d539946ee0525434d
  cmsat4                         & algotocnf\_3\_diff-desc.cnf    & none       & no    & 48704 / 374592 & 0         & 1597496   & 447        & 18 \\ % a6d034725c9d02c9b24913d8459bffc0e754ab8819143376c20ec930
  cmsat5                         & algotocnf\_1.cnf               & none       & no    & 48704 / 253984 &           & 18893645  & 23185      & 1064 \\ % dd65dd3547d2446b98586186a2e1c1582c3cd423287e6e09fad2176a
  cmsat5                         & algotocnf\_1.cnf               & cmsat      & no    & 24503 / 111931 &           & 16820944  & 16686      & 973 \\ % 1d3752326201b44e49b52a67a06f78742295d281b8b64d5cb6da02c9
  cmsat5                         & algotocnf\_1.cnf               & minisat    & no    & 20895 / 118236 &           & 30415171  & 33140      & 4470 \\ % a13b96d03246256b5a3ab6d21cc97ba4494b02a989b0ef6338e1b6e6
  cmsat5                         & algotocnf\_1.cnf               & satelite   & no    & 27495 / 153262 &           & 32077946  & 34269      & 3920 \\ % e11556c27aecaaa0e87b94b7b965f0f2242c486d82cef5e8e3076e46
  cmsat5                         & algotocnf\_18-t9.cnf           & none       & no    & 107839 / 590953 &           & 2428824   & 637        & 37 \\ % b37ae276e36bf564b19c0c0da987196cf01555caa7593d72d84b4eee
  cmsat5                         & algotocnf\_18-t9.cnf           & cmsat      & no    & 61789 / 342139 &           & 1682348   & 499        & 34 \\ % 7e599057f0477da364afdce4e74553c90e99ec27296176d66fd8c68b
  cmsat5                         & algotocnf\_18-t9.cnf           & lingeling  & no    & 107839 / 344544 &           & 1481090   & 582        & 27 \\ % 50713132da31e22b53075788c5a8331f30ce0921e9dc5d383ecfa0c3
  cmsat5                         & algotocnf\_18-t9.cnf           & minisat    & no    & 61845 / 391134 &           & 1514709   & 433        & 32 \\ % c45dd05b61ea4662b1f8baa37afb3541c399205d17f2334b5b30428e
  cmsat5                         & algotocnf\_18-t9.cnf           & satelite   & no    & 69670 / 457972 &           & 969684    & 185        & 10 \\ % c2f7ad0df06de34be2dbac129753e6ca1d34cc0bcfe6b11b73572022
  cmsat5                         & algotocnf\_18-t9\_diff-desc.cnf & none       & no    & 107839 / 846487 &           & 2935002   & 627        & 37 \\ % e393650144aaf468987ac02d0fa0940821eb1ffa21c9aec810081857
  cmsat5                         & algotocnf\_18-t9\_diff-desc.cnf & cmsat      & no    & 57114 / 303489 &           & 2169645   & 393        & 32 \\ % 1eb899aafbccd650b1de1391d4c6191cffcca9cd324acc6f4b020071
  cmsat5                         & algotocnf\_18-t9\_diff-desc.cnf & lingeling  & no    & 107839 / 404666 &           & 1331765   & 347        & 39 \\ % 67636cf4cbbb0b74bb53bff80340bd6f66b256cfd6dcc82ca08373eb
  cmsat5                         & algotocnf\_18-t9\_diff-desc.cnf & satelite   & no    & 85386 / 692408 &           & 2265178   & 445        & 36 \\ % dd3da7e1877abdf2a671be508add55d58765c523a2170b69149d403c
  cmsat5                         & algotocnf\_1\_diff-desc.cnf    & none       & no    & 48704 / 373920 &           & 6183574   & 5599       & 231 \\ % aa3b78767b8e278c02ad0f4d64bc587a30c2459e647aa460e6422c19
  cmsat5                         & algotocnf\_1\_diff-desc.cnf    & cmsat      & no    & 15402 / 73142 &           & 3548865   & 2969       & 155 \\ % 7b73f0009bce267fc3657df3d110c350c7d926a1503d4f67bbac2191
  cmsat5                         & algotocnf\_1\_diff-desc.cnf    & lingeling  & no    & 48704 / 101182 &           & 2899864   & 2153       & 107 \\ % 7e9101830ff0fa57ce7a03c656d979215226e24943b460fc65cdc772
  cmsat5                         & algotocnf\_1\_diff-desc.cnf    & minisat    & no    & 25107 / 158393 &           & 4174263   & 3465       & 125 \\ % 3e166db528d5ee694fd263a1a51663b67853e3ba0dfc534e6191f634
  cmsat5                         & algotocnf\_1\_diff-desc.cnf    & satelite   & no    & 30950 / 226602 &           & 6698827   & 5794       & 233 \\ % 4e66e1d3f1385240789982b42d736f68fb59f8154e8ee767c9ab54de
  cmsat5                         & algotocnf\_2.cnf               & none       & no    & 48704 / 254210 &           & 13968754  & 15787      & 571 \\ % d35410b792d69daae7910c543ba9d5d72bd536d7558a5a50dfa02664
  cmsat5                         & algotocnf\_2.cnf               & cmsat      & no    & 24323 / 111661 &           & 16420056  & 16911      & 1079 \\ % 89b72f50b32b5f5ebf2f41dd6e3cdb124223d3bbfdb5a78de3000648
  cmsat5                         & algotocnf\_2.cnf               & minisat    & no    & 20895 / 118688 &           & 31064564  & 33957      & 7280 \\ % 731684bed773b479fd02e269f4d045016b1d5c4572a996fe39737af8
  cmsat5                         & algotocnf\_2.cnf               & satelite   & no    & 27775 / 154276 &           & 20428518  & 21263      & 1668 \\ % 2f956d4a0785186a60c0ba9e67544d12c948bd71414eb6a3764ab9f0
    \end{tabular}
  \end{center}
\end{table}

%Text such that \LaTeX{} won't screw up.
\newpage

\begin{table}[p]
  \begin{center}
    \begin{tabular}{l|cccccccc}
        \textbf{SAT solver} & \textbf{CNF file} & \textbf{simplified} & \textbf{phase} & \textbf{vars/clauses} & \textbf{propagations} & \textbf{decisions} & \textbf{restarts} & \textbf{Runtime (sec)} \\
      \hline
  cmsat5                         & algotocnf\_21-t9.cnf           & none       & no    & 116800 / 636838 &           &           & 697        & timeout \\ % d9fa53911829875b6b15a6f29bdae11bfba04be71228fba336dd72ed
  cmsat5                         & algotocnf\_21-t9.cnf           & cmsat      & no    & 63555 / 358308 &           &           &            & timeout \\ % 78ee8e6c5d4e428160d53c12a51da977f108f5889249c2847aa96d4d
  cmsat5                         & algotocnf\_21-t9.cnf           & minisat    & no    & 65572 / 414470 &           &           & 629        & timeout \\ % f8485bb3695f53fbdab5073ae370d3368607b62639a7e3cd7cfdaca4
  cmsat5                         & algotocnf\_21-t9.cnf           & satelite   & no    & 74562 / 486714 &           &           & 597        & timeout \\ % d59026f7a07d4b3681ec6a1bcfa381d8964b272b1a0217e1ef9a3784
  cmsat5                         & algotocnf\_21-t9\_diff-desc.cnf & none       & no    & 116800 / 911629 &           & 23821572  & 37787      & 7855 \\ % 1b2e8ac95dbda1aa8a9ef8c2ddaa6b44d0c051552265f805fb6e0c99
  cmsat5                         & algotocnf\_21-t9\_diff-desc.cnf & cmsat      & no    & 59528 / 318049 &           & 21099143  & 25196      & 4457 \\ % e61d6b143ef2cca90a198c9e9744f43a5b96303b7833f25607ecb163
  cmsat5                         & algotocnf\_21-t9\_diff-desc.cnf & lingeling  & no    & 116800 / 404126 &           & 21868261  & 29926      & 7257 \\ % c13f33f580f50e342b589fa6d13949220bbf4975e14a0b18c549f172
  cmsat5                         & algotocnf\_21-t9\_diff-desc.cnf & minisat    & no    & 81650 / 534965 &           & 16256636  & 27000      & 3327 \\ % dcabafdbabe50047d42047d78d905cbe70b3dd3bc48344691dcf537f
  cmsat5                         & algotocnf\_21-t9\_diff-desc.cnf & satelite   & no    & 90236 / 720755 &           & 8231920   & 10501      & 909 \\ % b2c58e7a46a24b48445856b7336793895df34576c7a53f8df118e9d5
  cmsat5                         & algotocnf\_23-t9.cnf           & none       & no    & 122774 / 667438 &           &           & 555        & timeout \\ % 3fe0bef6a6dc09f1322fae5e9d257f2deabd229db01b84440287cf93
  cmsat5                         & algotocnf\_23-t9.cnf           & cmsat      & no    & 64880 / 369775 &           &           &            & timeout \\ % d5e26478b91a3cb00617553a211ed0072d50d16bdc9a12cef1925794
  cmsat5                         & algotocnf\_23-t9.cnf           & minisat    & no    & 67979 / 429734 &           &           & 625        & timeout \\ % 13572ccd41bf8efc5bb49cb25a3bbc3d548a3d062dab3867cbb435fa
  cmsat5                         & algotocnf\_23-t9.cnf           & satelite   & no    & 77835 / 506217 &           &           & 729        & timeout \\ % 97aa1eb1a1c3a80f29f53e09ef81766a74e5593f4065ea9e2e0c2fbf
  cmsat5                         & algotocnf\_23-t9\_diff-desc.cnf & none       & no    & 122774 / 955067 &           & 38544343  & 63944      & 26212 \\ % 18963c101a0e26abc7afc8d743578455ff71ef6bb088f13c4667ecac
  cmsat5                         & algotocnf\_23-t9\_diff-desc.cnf & cmsat      & no    & 57510 / 307118 &           & 51071860  & 75446      & 45576 \\ % d9844db15414670b3e66d993001683279c208039c94cc4da3904fd27
  cmsat5                         & algotocnf\_23-t9\_diff-desc.cnf & minisat    & no    & 84080 / 550727 &           & 40612782  & 68749      & 21254 \\ % 63c81268ba4d2e7cb4fe6395245f5b6be001c0eb5af014a760c294cc
  cmsat5                         & algotocnf\_23-t9\_diff-desc.cnf & satelite   & no    & 93504 / 740545 &           & 21805012  & 36348      & 5503 \\ % 6f36c46f819e81a34f11b00559862d887287ca245fbfe0b1d5f737f7
  cmsat5                         & algotocnf\_24-t9.cnf           & none       & no    & 125761 / 682722 &           &           & 636        & timeout \\ % 3b9e199c59f16d0fca5cfd0c99455902bb28e0d4625524242215ff80
  cmsat5                         & algotocnf\_24-t9.cnf           & cmsat      & no    & 65579 / 377228 &           &           &            & timeout \\ % efd19a15f0e46a6a994b1f2bc5966c79178cc23c85da270dc8be730f
  cmsat5                         & algotocnf\_24-t9.cnf           & minisat    & no    & 69023 / 436292 &           &           & 576        & timeout \\ % c09f80263d2d8e9966a8885cf81b8013687e23f42d6306fa38a0dd1a
  cmsat5                         & algotocnf\_24-t9.cnf           & satelite   & no    & 79486 / 515977 &           &           & 627        & timeout \\ % f9544e9089111947746b12b353a8956b32216050841c0c564443f43a
  cmsat5                         & algotocnf\_24-t9\_diff-desc.cnf & none       & no    & 125761 / 976770 &           & 55539005  & 98123      & 37194 \\ % 1328a16d7f06b4be5e1faa1101be76ca6e70190237e46a71a47da1a2
  cmsat5                         & algotocnf\_24-t9\_diff-desc.cnf & cmsat      & no    & 56447 / 302014 &           & 72716973  & 118102     & 72908 \\ % 3a3d05f98895d77a274093d205f305010a03432b6d2835de078c0c1f
  cmsat5                         & algotocnf\_24-t9\_diff-desc.cnf & minisat    & no    & 85312 / 558731 &           & 45377443  & 76719      & 26074 \\ % 4bcc393de85da42c37b4bd0673461f2a6f323fac4a9c7d987427ab03
  cmsat5                         & algotocnf\_24-t9\_diff-desc.cnf & satelite   & no    & 95154 / 750557 &           & 35455174  & 59757      & 23267 \\ % 68e21c1027bd2ae5ec9b5db11415141d1df155fb36029b59834c8de1
  cmsat5                         & algotocnf\_27-t10-diff.cnf     & none       & no    & 134722 / 728349 &           & 0         & 1          & 1 \\ % 9362fb209e13f9184e899c1951e775788f370d50d8492f887cba50fb
  cmsat5                         & algotocnf\_27-t10-diff.cnf     & cmsat      & no    & 0 / 0      &           & 0         & 1          & 0 \\ % 99410c103376dd023f5e56569c6ad85835da863c5ac016717bdf6c44
  cmsat5                         & algotocnf\_27-t10-diff.cnf     & lingeling  & no    & 134722 / 0 &           & 0         & 1          & 0 \\ % dc5ade5ba4ef907f3c6da0c6b2d66c122a66876f365174ba02668c19
  cmsat5                         & algotocnf\_27-t10-diff.cnf     & minisat    & no    & 0 / 0      &           & 0         & 1          & 0 \\ % b8cfde1c2532a2971fce2b35551d3d5e208d8aaed4ea119c57fb0248
  cmsat5                         & algotocnf\_27-t10-diff.cnf     & satelite   & no    & 0 / 0      &           & 0         & 1          & 0 \\ % 3164775d5b42ceb512cb45ec78e29e094e5f78330d24a86cd8605542
  cmsat5                         & algotocnf\_27-t10-diff\_diff-desc.cnf & cmsat      & no    & 0 / 0      &           & 0         & 1          & 0 \\ % c68a13dbf5250cd4f0ebf6548d0045b4c901def30aa41758c7278a4b
  cmsat5                         & algotocnf\_27-t10-diff\_diff-desc.cnf & lingeling  & no    & 134722 / 0 &           & 0         & 1          & 0 \\ % 347b04358c408698db4eaab30e6d58210cfc7b5eba5575b9febe686e
  cmsat5                         & algotocnf\_27-t10-diff\_diff-desc.cnf & minisat    & no    & 0 / 0      &           & 0         & 1          & 0 \\ % d3560958b55da1b3254a7a4fa00adf9174fe300405cf150623037cb8
  cmsat5                         & algotocnf\_27-t10-diff\_diff-desc.cnf & satelite   & no    & 0 / 0      &           & 0         & 1          & 0 \\ % f9682e3c73ecbdc2cd886b4922c7140a4225eeb51f41fc3574d9af1c
  cmsat5                         & algotocnf\_27-t10-simplified.cnf & none       & no    & 134722 / 728542 &           &           & 562        & timeout \\ % fc15fc2d9c1f976d1cf14735d8b9ae25d423ecb6677e0969180a547a
    \end{tabular}
  \end{center}
\end{table}

%Text such that \LaTeX{} won't screw up.
\newpage

\begin{table}[p]
  \begin{center}
    \begin{tabular}{l|cccccccc}
        \textbf{SAT solver} & \textbf{CNF file} & \textbf{simplified} & \textbf{phase} & \textbf{vars/clauses} & \textbf{propagations} & \textbf{decisions} & \textbf{restarts} & \textbf{Runtime (sec)} \\
      \hline
  cmsat5                         & algotocnf\_27-t10-simplified.cnf & cmsat      & no    & 67175 / 392680 &           &           &            & timeout \\ % a4205728d62a96e849ded36ccbb3a12db7187f9034ff59f8c69643b1
  cmsat5                         & algotocnf\_27-t10-simplified.cnf & minisat    & no    & 72909 / 460404 &           &           & 970        & timeout \\ % a4a96fedff782a512ce0a83c805b146889cf44636112102089951865
  cmsat5                         & algotocnf\_27-t10-simplified.cnf & satelite   & no    & 84442 / 544998 &           &           & 544        & timeout \\ % 12651e453d6ecb625dc2ec25e404aa184ea7b90046cc952fa7db3983
  cmsat5                         & algotocnf\_27-t10-simplified\_diff-desc.cnf & none       & no    & 134722 / 1041847 &           &           & 723        & timeout \\ % 170eaa58c51a84f843380f0c08564f2ed2cc8f01deaaa251c8de44cc
  cmsat5                         & algotocnf\_27-t10-simplified\_diff-desc.cnf & cmsat      & no    & 56223 / 302443 &           &           &            & timeout \\ % 68ab9def9699e8ae0a1508bf9d24db493c6d2076477e0b167a1ec37d
  cmsat5                         & algotocnf\_27-t10-simplified\_diff-desc.cnf & minisat    & no    & 89667 / 586309 &           &           & 535        & timeout \\ % 6b359868e47a463435644dcc654ac9d20c193260e16ac273602f49b7
  cmsat5                         & algotocnf\_27-t10-simplified\_diff-desc.cnf & satelite   & no    & 100735 / 789589 &           &           & 619        & timeout \\ % 70886d3c0860f70ce9e29a8a8e0e179c6178caf8055075c4b2a2ef6a
  cmsat5                         & algotocnf\_27-t10.cnf          & none       & no    & 134722 / 728030 &           &           & 1086       & timeout \\ % 79d7ea25b894a37be6d85e336b86fb7f1febb8cdcfa3a72f37c47d5b
  cmsat5                         & algotocnf\_27-t10.cnf          & cmsat      & no    & 72044 / 411276 &           &           &            & timeout \\ % 81352263671734274320949e824dfb66cc08548dfee25ffb456b577a
  cmsat5                         & algotocnf\_27-t10.cnf          & minisat    & no    & 73897 / 465248 &           &           & 831        & timeout \\ % cb9a08eb20884b38febccabbbb97fa9340d3a41a574f07a563f578a4
  cmsat5                         & algotocnf\_27-t10.cnf          & satelite   & no    & 85274 / 549734 &           &           & 984        & timeout \\ % e9baeb1f77526fa5b120d17cee596e478a060c38ff0e3a215f1355dd
  cmsat5                         & algotocnf\_27-t10\_diff-desc.cnf & none       & no    & 134722 / 1041335 &           &           & 680        & timeout \\ % 30541a84ab1a640f5a0536b24eb193f10252cdc2d1f6e0a031d2c744
  cmsat5                         & algotocnf\_27-t10\_diff-desc.cnf & cmsat      & no    & 57956 / 311219 &           &           &            & timeout \\ % fa2cc3ad9d9dddb96007b7083827304f96624636a721280f39127fe1
  cmsat5                         & algotocnf\_27-t10\_diff-desc.cnf & minisat    & no    & 94850 / 619480 &           &           & 679        & timeout \\ % fee81fc5a317da26a7ca9f28535c4e99de96d57eca9a37f9050a8373
  cmsat5                         & algotocnf\_27-t10\_diff-desc.cnf & satelite   & no    & 105733 / 857675 &           &           & 678        & timeout \\ % 90f28b4f35bbef4c32cc0efd48a7b974844ee1e3a460769fea476115
  cmsat5                         & algotocnf\_2\_diff-desc.cnf    & none       & no    & 48704 / 374146 &           & 7410370   & 7959       & 395 \\ % 9f47466f8654be1f0703b3b18de95295a3af4de768715b8ea3199965
  cmsat5                         & algotocnf\_2\_diff-desc.cnf    & cmsat      & no    & 15342 / 72870 &           & 4281712   & 3843       & 201 \\ % c66b0d263663137b71aaa5702ff1b5942d80a1b38a18c293cc4ce833
  cmsat5                         & algotocnf\_2\_diff-desc.cnf    & lingeling  & no    & 48704 / 106546 &           & 6364462   & 5611       & 413 \\ % e039b331df0239a9e273192ad65146cda57cd5a63c3a3e4c3241da28
  cmsat5                         & algotocnf\_2\_diff-desc.cnf    & minisat    & no    & 23531 / 149436 &           & 2452843   & 1744       & 80 \\ % 8cd5e8e8a4254ace82a09028831b8051150f16dfebc232a282e463a5
  cmsat5                         & algotocnf\_2\_diff-desc.cnf    & satelite   & no    & 29428 / 204577 &           & 4885340   & 3755       & 141 \\ % 71374a1e335ae97413515a811e56d6c48205aa0ad6821af694032afe
  cmsat5                         & algotocnf\_3.cnf               & none       & no    & 48704 / 254656 &           & 2178965   & 598        & 29 \\ % 2ee02e002a280f91f745a10c1268843a1822787c7e319373fb96559f
  cmsat5                         & algotocnf\_3.cnf               & cmsat      & no    & 23893 / 111403 &           & 1156826   & 398        & 16 \\ % a2b47f34b45ffaf96893d573b3ceac9ad88950a74c41d022b288410e
  cmsat5                         & algotocnf\_3.cnf               & lingeling  & no    & 48704 / 100941 &           & 410433    & 154        & 4 \\ % 90d2c889cfa79199789b1a021bbbf21c687389889b34068005b09c84
  cmsat5                         & algotocnf\_3.cnf               & minisat    & no    & 20895 / 119580 &           & 881366    & 326        & 10 \\ % 62e49ad3722e8d9c8729b39efa3bf58e0ad61f40655967fe4be2c0d5
  cmsat5                         & algotocnf\_3.cnf               & satelite   & no    & 28647 / 156910 &           & 1450065   & 437        & 18 \\ % d248953af26fcb424a8f921a060924ed52ed37cd54ccced3fdd3aeaf
  cmsat5                         & algotocnf\_3\_diff-desc.cnf    & none       & no    & 48704 / 374592 &           & 1611652   & 425        & 20 \\ % f7d34ab20c2c01e3829f2806f8ae25fb56ee4e5689beb51df23f9a21
  glucose                        & algotocnf\_1.cnf               & none       & no    & 48704 / 225124 &           &           &            & 86400 \\ % bf0397e93984a54f67f2fcfa7fc191facc0a819b73c34c40271c0388
  glucose                        & algotocnf\_1.cnf               & cmsat      & no    & 24503 / 111931 &           &           &            & 86400 \\ % 9acf0be9e0f16c20d589a47295a19036e5292442249be1e93c7cb070
  glucose                        & algotocnf\_1.cnf               & minisat    & no    & 20895 / 118236 &           &           &            & 86400 \\ % b467520d30bd674bc0a95e1b4d7139d5e356697e367197fda9dea4eb
  glucose                        & algotocnf\_1.cnf               & satelite   & no    & 27495 / 153262 &           &           &            & 86400 \\ % 580a9b4ce979cd7fe61d05a238e3f1985ec2d72fa7e73a050b402068
  glucose                        & algotocnf\_18-t9.cnf           & none       & no    & 107839 / 582936 & 14727839  & 990491    & 272        & 8 \\ % e45eabde0367796264a30445d57fbb7528be7434ac8a61d11c105c79
  glucose                        & algotocnf\_18-t9.cnf           & cmsat      & no    & 61789 / 342139 & 26259195  & 1773186   & 497        & 16 \\ % b6c1cee06ac586eaaa618867a3fe37e7e6fab4bdbea5194828dc6134
  glucose                        & algotocnf\_18-t9.cnf           & lingeling  & no    & 102616 / 344544 & 14568101  & 1212089   & 395        & 9 \\ % e046351ca4cca378b9bc9c277f945c99a03c6c74a7f81610ae1346e2
  glucose                        & algotocnf\_18-t9.cnf           & minisat    & no    & 61845 / 391134 & 38366719  & 2318012   & 686        & 30 \\ % 3f122c1ea802d7cefa9d5714dc213895875a1ca98680c6ff7d0293f4
  glucose                        & algotocnf\_18-t9.cnf           & satelite   & no    & 69670 / 457972 & 36089889  & 2123915   & 633        & 21 \\ % 0b39121673773035a983e5f36aaa72b673ea49b8ed69703094d0d39c
    \end{tabular}
  \end{center}
\end{table}

%Text such that \LaTeX{} won't screw up.
\newpage

\begin{table}[p]
  \begin{center}
    \begin{tabular}{l|cccccccc}
        \textbf{SAT solver} & \textbf{CNF file} & \textbf{simplified} & \textbf{phase} & \textbf{vars/clauses} & \textbf{propagations} & \textbf{decisions} & \textbf{restarts} & \textbf{Runtime (sec)} \\
      \hline
  glucose                        & algotocnf\_18-t9\_diff-desc.cnf & none       & no    & 107839 / 831263 & 10471476  & 644404    & 74         & 4 \\ % f3ab387c716d008a7cc709e504704551fd2f203c4dabb2b9a18c3989
  glucose                        & algotocnf\_18-t9\_diff-desc.cnf & cmsat      & no    & 57114 / 303489 & 20143829  & 1246909   & 243        & 11 \\ % 1347449bbb34c2b4a5126e753da189dd6fbcdc956718a9d36b2aeeed
  glucose                        & algotocnf\_18-t9\_diff-desc.cnf & lingeling  & no    & 99452 / 404666 & 7239871   & 764796    & 166        & 7 \\ % 53d3e55cfa87e834385b6da6c2a10ffb78c7e11ca6c13d94a5a35374
  glucose                        & algotocnf\_18-t9\_diff-desc.cnf & minisat    & no    & 78055 / 511431 & 15403560  & 934062    & 99         & 12 \\ % 5d3b3ee1f28feb0ef758d886ef77131db3345fe610342b0e957436ba
  glucose                        & algotocnf\_18-t9\_diff-desc.cnf & satelite   & no    & 85386 / 692408 & 30131933  & 1539016   & 302        & 26 \\ % fc35451ed1dc9838e5e91c60dcc1c3d7369704ada40bbb4c967f50d2
  glucose                        & algotocnf\_1\_diff-desc.cnf    & none       & no    & 48704 / 330692 & 1231386505 & 47264123  & 56508      & 3050 \\ % 8fe7b531373fb67eec823dfe1827d86ffbae6197a2ce822d10767468
  glucose                        & algotocnf\_1\_diff-desc.cnf    & cmsat      & no    & 15402 / 73142 & 1038552689 & 32670843  & 33629      & 2317 \\ % 8c0077bf9cb1b7f9b5919fe5a4b5b4087c08f31a6c33cf236d43e76a
  glucose                        & algotocnf\_1\_diff-desc.cnf    & lingeling  & no    & 40554 / 101182 & 1251948097 & 59378788  & 62092      & 4858 \\ % 814732dfa0aff44d2fcad0460bc447854d95bc1308dfc8bf156ae67e
  glucose                        & algotocnf\_1\_diff-desc.cnf    & minisat    & no    & 25107 / 158393 & 785694216 & 29540624  & 30952      & 1199 \\ % 77a0de6eea8ccdd1532d7825f4827243ee49f7772154fc71ae365883
  glucose                        & algotocnf\_1\_diff-desc.cnf    & satelite   & no    & 30950 / 226602 & 751454677 & 28310860  & 29558      & 1126 \\ % 0d57f0f8115d6f165457a97b7bb3cfe549c2ae44baf0934063d51ff8
  glucose                        & algotocnf\_2.cnf               & none       & no    & 48704 / 224220 &           &           &            & 86400 \\ % f869d4379d6f8a0073e13855c6b15d9fdc4774a69dce920e7d9aaad8
  glucose                        & algotocnf\_2.cnf               & cmsat      & no    & 24323 / 111661 &           &           &            & 86400 \\ % 6509e9ebb0e18adefdc8fa285e7c49ef6329efaea4f43864c981862f
  glucose                        & algotocnf\_2.cnf               & minisat    & no    & 20895 / 118688 &           &           &            & 86400 \\ % b879ffb8e5572a7111c92350d97b6038f4c75d767779b221ca82bb05
  glucose                        & algotocnf\_2.cnf               & satelite   & no    & 27775 / 154276 &           &           &            & 86400 \\ % 90086563d47664dc016a241d6054cb7b4c91380168cd4dcfcf8a90b5
  glucose                        & algotocnf\_21-t9.cnf           & none       & no    & 116800 / 627112 &           &           &            & 86400 \\ % fa572c98f7d9723df7724dbf96b15f0c8d55ad2bb930e94ff7ffc003
  glucose                        & algotocnf\_21-t9.cnf           & cmsat      & no    & 63555 / 358308 &           &           &            & 86400 \\ % 622b5ae4005c2e75d2ba8995a3a0e1f65e0dcab112d6f8cf6ae6baca
  glucose                        & algotocnf\_21-t9.cnf           & minisat    & no    & 65572 / 414470 &           &           &            & 86400 \\ % c07c485817de0880fba7f452774c3dc9000b62bbf6ad95282608ee46
  glucose                        & algotocnf\_21-t9.cnf           & satelite   & no    & 74562 / 486714 &           &           &            & 86400 \\ % 40516aeb09dfee80fe2ffaf049295601e350dab4e5b72d0e3f705ed1
  glucose                        & algotocnf\_21-t9\_diff-desc.cnf & none       & no    & 116800 / 894072 &           &           &            & 86400 \\ % 46f8890c94b917200dcd3adf8656d67b083263c18a306c50c7425fbc
  glucose                        & algotocnf\_21-t9\_diff-desc.cnf & cmsat      & no    & 59528 / 318049 & 2790904343 & 82667240  & 130502     & 17967 \\ % e447b4e76af564d87a0b2c293c2904b59431684b31cd17e0ad37e47d
  glucose                        & algotocnf\_21-t9\_diff-desc.cnf & lingeling  & no    & 104860 / 404126 & 1838242760 & 74906191  & 112992     & 10388 \\ % c7ac7ee2ca2796adc0255b1c7d2fd6f786cf8c3a38edfc938e3de8e9
  glucose                        & algotocnf\_21-t9\_diff-desc.cnf & minisat    & no    & 81650 / 534965 & 3062139546 & 87991203  & 136762     & 19210 \\ % 872a3505db2626c545265e2598fa4aa540e2e5f5cebf32b852f86f2d
  glucose                        & algotocnf\_21-t9\_diff-desc.cnf & satelite   & no    & 90236 / 720755 & 3468599865 & 98765464  & 161923     & 16278 \\ % 3aa02c8fc195268e64a0b038c2a596c24fadc90acdb04be83c7d4943
  glucose                        & algotocnf\_23-t9.cnf           & none       & no    & 122774 / 656628 &           &           &            & 86400 \\ % dfc0e162e204d05caff5ba8f70489136643d8cdd906c31ca223c4c52
  glucose                        & algotocnf\_23-t9.cnf           & cmsat      & no    & 64880 / 369775 &           &           &            & 86400 \\ % f71cbd8592ca2fea3cfd002c329f7fd1e266d6383b3fc4f52448123f
  glucose                        & algotocnf\_23-t9.cnf           & minisat    & no    & 67979 / 429734 &           &           &            & 86400 \\ % 76485ce3493d5cf37273cd9d46fa46b0a5532b0f01645f39c32d2c54
  glucose                        & algotocnf\_23-t9.cnf           & satelite   & no    & 77835 / 506217 &           &           &            & 86400 \\ % 77b3ab5b52146f81e52cc9fbd6a86c6297a1125b553c4d9127bf4f5b
  glucose                        & algotocnf\_23-t9\_diff-desc.cnf & none       & no    & 122774 / 936270 &           &           &            & 86400 \\ % 49f59ca12a4a430a8859210ac57ff7b7bed852f7dde9b271ee8d3429
  glucose                        & algotocnf\_23-t9\_diff-desc.cnf & cmsat      & no    & 57510 / 307118 &           &           &            & 86400 \\ % 09b4547d1bf87872dd4fdbab08153f96549187379053276c42427de5
  glucose                        & algotocnf\_23-t9\_diff-desc.cnf & minisat    & no    & 84080 / 550727 &           &           &            & 86400 \\ % 93e17442b058f16e10b838d1e51e645f844a507b370a195c2c8b4727
  glucose                        & algotocnf\_23-t9\_diff-desc.cnf & satelite   & no    & 93504 / 740545 &           &           &            & 86400 \\ % 35f4ebad02d56b9494535124944096cc566aad5faeabb5361ee8a11d
  glucose                        & algotocnf\_24-t9.cnf           & none       & no    & 125761 / 671432 &           &           &            & 86400 \\ % 1716ee0070cfa1691ef8db1fb32a8140bd85fb73928d1c8620cb93e1
  glucose                        & algotocnf\_24-t9.cnf           & cmsat      & no    & 65579 / 377228 &           &           &            & 86400 \\ % 0236dfb335d8188d190903d268f7e16b2b94e5afcee5d08c1e87d8f1
  glucose                        & algotocnf\_24-t9.cnf           & minisat    & no    & 69023 / 436292 &           &           &            & 86400 \\ % 6f8dd85f3f1ac12908587d8fc7413a4aeb40464768f85c4bf85bad9b
  glucose                        & algotocnf\_24-t9.cnf           & satelite   & no    & 79486 / 515977 &           &           &            & 86400 \\ % 036e72609090961d5f33cda92e95bfd93dcf34d44a78f3cc7b57757e
    \end{tabular}
  \end{center}
\end{table}

%Text such that \LaTeX{} won't screw up.
\newpage

\begin{table}[p]
  \begin{center}
    \begin{tabular}{l|cccccccc}
        \textbf{SAT solver} & \textbf{CNF file} & \textbf{simplified} & \textbf{phase} & \textbf{vars/clauses} & \textbf{propagations} & \textbf{decisions} & \textbf{restarts} & \textbf{Runtime (sec)} \\
      \hline
  glucose                        & algotocnf\_24-t9\_diff-desc.cnf & none       & no    & 125761 / 957493 &           &           &            & 86400 \\ % c372897126c37a6367aac4a55410b99d588c2741bc7f3c7a97953e31
  glucose                        & algotocnf\_24-t9\_diff-desc.cnf & cmsat      & no    & 56447 / 302014 &           &           &            & 86400 \\ % 92a741dea81d43928d5d4414f828da87f73b6dd1a0ab984b7d38574d
  glucose                        & algotocnf\_24-t9\_diff-desc.cnf & minisat    & no    & 85312 / 558731 &           &           &            & 86400 \\ % ea3ac64a8c0f3e29ca1c4a513c55b822db1604c0c6cc01bad10ae086
  glucose                        & algotocnf\_24-t9\_diff-desc.cnf & satelite   & no    & 95154 / 750557 &           &           &            & 86400 \\ % 708ec1dee5bc41dfabd7fe4053cf85cf340b41e330cb2ecbf7c2c88b
  glucose                        & algotocnf\_27-t10-diff.cnf     & none       & no    & 134722 / 8320 & 134722    & 1         & 1          & 0 \\ % 323cbe30b40e32e35d8c3fa9c6913d4bab2813868ebba90115dbe4dd
  glucose                        & algotocnf\_27-t10-diff.cnf     & cmsat      & no    & 0 / 0      & 0         & 1         & 1          & 0 \\ % 75804ad77fa1c20e35832e286890e594b835655f54392a8f4538f31b
  glucose                        & algotocnf\_27-t10-diff.cnf     & lingeling  & no    & 0 / 0      & 0         & 1         & 1          & 0 \\ % 12da98c367905a79a90d60ba93fc8f0df5e2b6708bf59617f75a64c1
  glucose                        & algotocnf\_27-t10-diff.cnf     & minisat    & no    & 0 / 0      & 0         & 1         & 1          & 0 \\ % aa11676fa9162bb47fb8505a5e3abc6ccfedab8003504fc602ca0ca7
  glucose                        & algotocnf\_27-t10-diff.cnf     & satelite   & no    & 0 / 0      & 0         & 1         & 1          & 0 \\ % cdfba0fde46b22d420d79e14ad77f666ff7a60eacdafb2a8a96216ac
  glucose                        & algotocnf\_27-t10-diff\_diff-desc.cnf & cmsat      & no    & 0 / 0      & 0         & 1         & 1          & 0 \\ % d92d913edcd302e34f06c231e31e703a9a892df081b871e90d0b1a0a
  glucose                        & algotocnf\_27-t10-diff\_diff-desc.cnf & lingeling  & no    & 0 / 0      & 0         & 1         & 1          & 0 \\ % 29f44867e2e7ae7be43a6086a478375a0de9459d7006c78052e46c0e
  glucose                        & algotocnf\_27-t10-diff\_diff-desc.cnf & minisat    & no    & 0 / 0      & 0         & 1         & 1          & 0 \\ % ea1d054f9bda170a39dd2aae6695ba9e463641bb7f4c6fccf8352bab
  glucose                        & algotocnf\_27-t10-diff\_diff-desc.cnf & satelite   & no    & 0 / 0      & 0         & 1         & 1          & 0 \\ % a6cd9214fa7692b9af2df500127e022a1ea926455c70fc3e66da3ce9
  glucose                        & algotocnf\_27-t10-simplified.cnf & none       & no    & 134722 / 715830 &           &           &            & 86400 \\ % d56475b88b3113060e6368efb525258459aa277bad590ea8a191e58d
  glucose                        & algotocnf\_27-t10-simplified.cnf & cmsat      & no    & 67175 / 392680 &           &           &            & 86400 \\ % ede38989d8ed2541e58c7a43fd369a209281e7c398f450426684492b
  glucose                        & algotocnf\_27-t10-simplified.cnf & minisat    & no    & 72909 / 460404 &           &           &            & 86400 \\ % 58030b8ade89ab4f449127776a52e28dc0af2152f2dd3ae0a80af555
  glucose                        & algotocnf\_27-t10-simplified.cnf & satelite   & no    & 84442 / 544998 &           &           &            & 86400 \\ % e91c3abe73687b3b8995f0da122ebc851dfc8a8abb55c907a4fd2ee8
  glucose                        & algotocnf\_27-t10-simplified\_diff-desc.cnf & none       & no    & 134722 / 1020826 &           &           &            & 86400 \\ % d8a3cb39d1ceffdf5b9aa0ee2da236ce4f8e4c20b30a13a2135fac97
  glucose                        & algotocnf\_27-t10-simplified\_diff-desc.cnf & cmsat      & no    & 56223 / 302443 &           &           &            & 86400 \\ % 22a6b6c961516005fe9a81f2b8fc9990f1bdc988eb9993b6c10c24fd
  glucose                        & algotocnf\_27-t10-simplified\_diff-desc.cnf & minisat    & no    & 89667 / 586309 &           &           &            & 86400 \\ % d1dccbbd370b2eb70de54951d652845750556c045b2a972b70846db2
  glucose                        & algotocnf\_27-t10-simplified\_diff-desc.cnf & satelite   & no    & 100735 / 789589 &           &           &            & 86400 \\ % 8ef8981a472d74e023ff6ca795728161f022eeca29a5efda84995bbe
  glucose                        & algotocnf\_27-t10.cnf          & none       & no    & 134722 / 717878 &           &           &            & 86400 \\ % 43392bd796208436fda2a52645cb96552f8566a9b2bb0df13f600fdf
  glucose                        & algotocnf\_27-t10.cnf          & cmsat      & no    & 72044 / 411276 &           &           &            & 86400 \\ % cce9d5ee7cf147bb965a46594c8809151313d44b114e72dd00db9406
  glucose                        & algotocnf\_27-t10.cnf          & minisat    & no    & 73897 / 465248 &           &           &            & 86400 \\ % 56debaf966090b32d1568f4272e3adc98bf404207fef0d86d73a2d1b
  glucose                        & algotocnf\_27-t10.cnf          & satelite   & no    & 85274 / 549734 &           &           &            & 86400 \\ % 399db484d6d4670338387082121b499f0d31d130cb415f5c60aada78
  glucose                        & algotocnf\_27-t10\_diff-desc.cnf & none       & no    & 134722 / 1022874 &           &           &            & 86400 \\ % 1f365529b8bbebf12f006e0e5134af2bdbb701c009eee824bad87e72
  glucose                        & algotocnf\_27-t10\_diff-desc.cnf & cmsat      & no    & 57956 / 311219 &           &           &            & 86400 \\ % eaa2a1dae65c7a57137caf2ef37d3acf89fa87bf237d49e416650082
  glucose                        & algotocnf\_27-t10\_diff-desc.cnf & minisat    & no    & 94850 / 619480 &           &           &            & 86400 \\ % 32f7f111f5ccc3ff9f3f7692a5182b11c65a93bb86435606963c0825
  glucose                        & algotocnf\_27-t10\_diff-desc.cnf & satelite   & no    & 105733 / 857675 &           &           &            & 86400 \\ % 30e3bdf59a414376c58f8a0d37ad4fa9ed0ec0287a07bdddf4e0e168
  glucose                        & algotocnf\_2\_diff-desc.cnf    & none       & no    & 48704 / 329788 & 1124234867 & 42935596  & 51702      & 2753 \\ % 89608acf0088ddffe0b6ee0e9872bcd95d5bb333be18d1e7577bf82f
  glucose                        & algotocnf\_2\_diff-desc.cnf    & cmsat      & no    & 15342 / 72870 & 939102575 & 29596011  & 30194      & 1340 \\ % 6dd1c5d2067b7eff9ae1ff3aa8de3f8b8e4c5d789540c6eab1242df3
  glucose                        & algotocnf\_2\_diff-desc.cnf    & lingeling  & no    & 40767 / 106546 & 993525934 & 48331362  & 48734      & 3420 \\ % 4181a9ebe490fed20b62b0a1d2f82d7b9057ae9bc3fbb906a0d0d59e
  glucose                        & algotocnf\_2\_diff-desc.cnf    & minisat    & no    & 23531 / 149436 & 714911601 & 26684384  & 26335      & 1529 \\ % 691dc514cdf86ab05d91dfac7ea88656099613422b4b3658345f28e2
  glucose                        & algotocnf\_2\_diff-desc.cnf    & satelite   & no    & 29428 / 204577 & 599836102 & 22217082  & 21451      & 570 \\ % c3fc351a3ca23dae6ffc4ba434c5402db6875da2555d062990297bad
  glucose                        & algotocnf\_3.cnf               & none       & no    & 48704 / 222436 & 19131489  & 1199216   & 402        & 10 \\ % e376d005f0fd66db3e3a64eb502ac763d2b098d77873951fe7c9e9f5
    \end{tabular}
  \end{center}
\end{table}

%Text such that \LaTeX{} won't screw up.
\newpage

\begin{table}[p]
  \begin{center}
    \begin{tabular}{l|cccccccc}
        \textbf{SAT solver} & \textbf{CNF file} & \textbf{simplified} & \textbf{phase} & \textbf{vars/clauses} & \textbf{propagations} & \textbf{decisions} & \textbf{restarts} & \textbf{Runtime (sec)} \\
      \hline
  glucose                        & algotocnf\_3.cnf               & cmsat      & no    & 23893 / 111403 & 23157475  & 1275904   & 377        & 11 \\ % 42a7d7a91da7285a89d9325e6ce13e2bd47311946071a9e69171675a
  glucose                        & algotocnf\_3.cnf               & lingeling  & no    & 44000 / 100941 & 9187396   & 731846    & 160        & 4 \\ % 8dfb200e3b9fa5e7df16f3ac2b74649ea7cf3bff6350729fe1be17ed
  glucose                        & algotocnf\_3.cnf               & minisat    & no    & 20895 / 119580 & 16616478  & 1045134   & 293        & 7 \\ % 150f1174309b2525be9bbe5f7222332c9f48edc8423f955b4b76f81c
  glucose                        & algotocnf\_3.cnf               & satelite   & no    & 28647 / 156910 & 15878157  & 1009336   & 290        & 9 \\ % d7c2a8e8fab8e026db26eeea0594355d0164263ce5a2f983d721ad5e
  glucose                        & algotocnf\_3\_diff-desc.cnf    & none       & no    & 48704 / 328004 & 9899755   & 531225    & 114        & 8 \\ % 37b391606bfeee0763e4614ee05d8c0148c2e1a2ccba55bd003b8d97
  glucose-syrup                  & algotocnf\_1.cnf               & none       & no    & 48704 / 225124 &           & 173944439 & 259907     & timeout \\ % 50ebc956525027e93e38812de4dc24c2265e903e8283d5cfc86e9bb2
  glucose-syrup                  & algotocnf\_1.cnf               & cmsat      & no    & 24503 / 111931 &           & 128905334 & 162433     & timeout \\ % e29bb153920756e6a0b38ed84916c735a5f099120b99f72aa611cce8
  glucose-syrup                  & algotocnf\_1.cnf               & minisat    & no    & 20895 / 118236 &           & 130691654 & 166049     & timeout \\ % 99236bb22220e906e752b609c2ddf16ea043a71b2273f1fddbae5916
  glucose-syrup                  & algotocnf\_1.cnf               & satelite   & no    & 27495 / 153262 &           & 130323045 & 168717     & timeout \\ % 851f2c67fd8e2ec29a47f5ffd766cdf72cd6217152e63e36d78490b2
  glucose-syrup                  & algotocnf\_18-t9.cnf           & none       & no    & 107839 / 582936 &           & 629363    & 201        & 14 \\ % b1b437b793ab842b6c5a231ee05a32b2dce4ac6ba6da2d31d270dd53
  glucose-syrup                  & algotocnf\_18-t9.cnf           & cmsat      & no    & 61789 / 342139 &           & 643117    & 142        & 23 \\ % e98153882b9674ed6bdc833287543218365988a7ce92feb4443eda82
  glucose-syrup                  & algotocnf\_18-t9.cnf           & lingeling  & no    & 102616 / 344544 &           & 497597    & 159        & 15 \\ % 5e6acfd81c66aea9e2300eb18784892900c55663d6d1684ab69f969f
  glucose-syrup                  & algotocnf\_18-t9.cnf           & minisat    & no    & 61845 / 391134 &           & 722943    & 209        & 37 \\ % 9ec15e804f4865b309ec201739d1b39ba584759ee1caf12c5b9a9e87
  glucose-syrup                  & algotocnf\_18-t9.cnf           & satelite   & no    & 69670 / 457972 &           & 686160    & 187        & 29 \\ % 063ad40a769d6f965c6e41656bc86c79bb6f94feb2f7eed5c23306d6
  glucose-syrup                  & algotocnf\_18-t9\_diff-desc.cnf & none       & no    & 107839 / 831263 &           & 744555    & 132        & 38 \\ % 7796b5105762c8084bd35da38a819d374a4c6d3c304283dc5a2c4c48
  glucose-syrup                  & algotocnf\_18-t9\_diff-desc.cnf & cmsat      & no    & 57114 / 303489 &           & 2012701   & 457        & 78 \\ % e24b1216f4970703cfd88a96e68e5d7c826fadef6b93c29c1bcfe764
  glucose-syrup                  & algotocnf\_18-t9\_diff-desc.cnf & lingeling  & no    & 99452 / 404666 &           & 481297    & 110        & 21 \\ % ea2dd164f27ccd6e69f894c633617b8ccbe1618c970c38d5acc6f333
  glucose-syrup                  & algotocnf\_18-t9\_diff-desc.cnf & minisat    & no    & 78055 / 511431 &           & 534307    & 90         & 33 \\ % d895bc79c2b382130d7258123c32de53dd3cb9f5d4f770a645225035
  glucose-syrup                  & algotocnf\_18-t9\_diff-desc.cnf & satelite   & no    & 85386 / 692408 &           & 597733    & 83         & 37 \\ % 087003769dc14912c1b6a7cc74a95fc4cac49a581c2be610a084b4c9
  glucose-syrup                  & algotocnf\_1\_diff-desc.cnf    & none       & no    & 48704 / 330692 &           & 28145105  & 35868      & 10325 \\ % 0fccef88fe15b12a91d0ebcff1799dfd894a201841322b0a789457a0
  glucose-syrup                  & algotocnf\_1\_diff-desc.cnf    & cmsat      & no    & 15402 / 73142 &           & 8568127   & 7891       & 2029 \\ % fbe5688c26ddc13d440fa57e9224fdf264525d66aa13638de6b781d0
  glucose-syrup                  & algotocnf\_1\_diff-desc.cnf    & lingeling  & no    & 40554 / 101182 &           & 12094215  & 11272      & 2661 \\ % 8da486614ba52905041667e69d78da230b5c9386dd339df8752569d6
  glucose-syrup                  & algotocnf\_1\_diff-desc.cnf    & minisat    & no    & 25107 / 158393 &           & 7647590   & 7036       & 1394 \\ % 96b9e20c43d8bca80edc71d8830fe4b10332e6fb1242a65570b5bf23
  glucose-syrup                  & algotocnf\_1\_diff-desc.cnf    & satelite   & no    & 30950 / 226602 &           & 13372279  & 14049      & 3417 \\ % babeee6cae79aa1fb455a7b696f65eff9789330c6aa9f514fcbe8cdf
  glucose-syrup                  & algotocnf\_2.cnf               & none       & no    & 48704 / 224220 &           & 174375109 & 256807     & timeout \\ % 1cd2ed89d07f26ce548416824973f323069d14e1761b9f829745ac81
  glucose-syrup                  & algotocnf\_2.cnf               & cmsat      & no    & 24323 / 111661 &           & 149723742 & 195305     & timeout \\ % ed48bfea22628c2738493c6b2316704dbd4b2b56f477ced2fa769720
  glucose-syrup                  & algotocnf\_2.cnf               & minisat    & no    & 20895 / 118688 &           & 141489532 & 182961     & 67164 \\ % da874c8a40671a1ce9eb619c716c79e4fc5f57d9ef6beed7a3eb9683
  glucose-syrup                  & algotocnf\_2.cnf               & satelite   & no    & 27775 / 154276 &           & 171722347 & 229222     & timeout \\ % bf3bd4546912abab012fa4865790327f7e017057f86954ec1b983232
  glucose-syrup                  & algotocnf\_21-t9.cnf           & none       & no    & 116800 / 627112 &           & 82220645  & 133403     & timeout \\ % 1c5b1691c2bb4e24b227fc8b8a267798f906428e978c16d6aa5b7c3b
  glucose-syrup                  & algotocnf\_21-t9.cnf           & cmsat      & no    & 63555 / 358308 &           & 87272393  & 138790     & timeout \\ % cb089fda095a1e875cdf0138a901cbd9b7c42965f6c569b98a502783
  glucose-syrup                  & algotocnf\_21-t9.cnf           & minisat    & no    & 65572 / 414470 &           & 78852843  & 110030     & timeout \\ % 03e42f68a73e01ed18e86669c2273a4f2ceca2be0c76be8d81c2d7f6
  glucose-syrup                  & algotocnf\_21-t9.cnf           & satelite   & no    & 74562 / 486714 &           & 80656433  & 115935     & timeout \\ % dbe34d4d78282eac0c8309ab3e7600672cef85fb24700af6d6a4491f
  glucose-syrup                  & algotocnf\_21-t9\_diff-desc.cnf & none       & no    & 116800 / 894072 &           & 48531865  & 79100      & 25909 \\ % a497113af67119f0fb58ceddf3aea672880c007e7a8b8c794369d72c
  glucose-syrup                  & algotocnf\_21-t9\_diff-desc.cnf & cmsat      & no    & 59528 / 318049 &           & 24615086  & 32966      & 13281 \\ % 69fa773833c90b0ca2d93b03794a29f6d51eb70f90a0490ae5e64785
  glucose-syrup                  & algotocnf\_21-t9\_diff-desc.cnf & lingeling  & no    & 104860 / 404126 &           & 30587590  & 43486      & 12123 \\ % 55666895d7c74c39ec2d2dcd1fb5ee38d133e39313236cd2ae88c49c
    \end{tabular}
  \end{center}
\end{table}

%Text such that \LaTeX{} won't screw up.
\newpage

\begin{table}[p]
  \begin{center}
    \begin{tabular}{l|cccccccc}
        \textbf{SAT solver} & \textbf{CNF file} & \textbf{simplified} & \textbf{phase} & \textbf{vars/clauses} & \textbf{propagations} & \textbf{decisions} & \textbf{restarts} & \textbf{Runtime (sec)} \\
      \hline
  glucose-syrup                  & algotocnf\_21-t9\_diff-desc.cnf & minisat    & no    & 81650 / 534965 &           & 34360445  & 50526      & 15501 \\ % d58ab772dd259e03a05029546b9e3e972451155d28c0c78b9a7a4977
  glucose-syrup                  & algotocnf\_21-t9\_diff-desc.cnf & satelite   & no    & 90236 / 720755 &           & 48591283  & 74843      & 27215 \\ % f845599ebf251285e8ada0bd82fbdda741b6d406f81041615b3c181c
  glucose-syrup                  & algotocnf\_23-t9.cnf           & none       & no    & 122774 / 656628 &           & 82295157  & 132051     & timeout \\ % 297ac25cd9180dca9b11cd170cf3847b5735e0a1914cbbb1a8c54006
  glucose-syrup                  & algotocnf\_23-t9.cnf           & cmsat      & no    & 64880 / 369775 &           & 73164991  & 108673     & timeout \\ % 9c4b86d4178d02372f54e09bd926951193c6f1c37390e926b30306d4
  glucose-syrup                  & algotocnf\_23-t9.cnf           & minisat    & no    & 67979 / 429734 &           & 81337922  & 118510     & timeout \\ % 2ffd9915699cd51205485834558447db0c5ebec4d7ea6f110736c34b
  glucose-syrup                  & algotocnf\_23-t9.cnf           & satelite   & no    & 77835 / 506217 &           & 78390241  & 113533     & timeout \\ % 2c302d44918aa0e7762bc7a7690ae8798ad399f4cf2225c670f02ed9
  glucose-syrup                  & algotocnf\_23-t9\_diff-desc.cnf & none       & no    & 122774 / 936270 &           & 77708491  & 135578     & timeout \\ % 6c826fdeb3986bf2bce6210a762cea5461df040f4b3ea687c3459286
  glucose-syrup                  & algotocnf\_23-t9\_diff-desc.cnf & cmsat      & no    & 57510 / 307118 &           & 78826943  & 126056     & timeout \\ % ce21807995df5790503bcf06356c87a9f064c1debf0e1d6da6745c3b
  glucose-syrup                  & algotocnf\_23-t9\_diff-desc.cnf & minisat    & no    & 84080 / 550727 &           & 80435602  & 131805     & timeout \\ % d0f51f44e43f144e3eff7dd5b59563a1f3b181057ba30624112192ae
  glucose-syrup                  & algotocnf\_23-t9\_diff-desc.cnf & satelite   & no    & 93504 / 740545 &           & 78665802  & 129109     & timeout \\ % eba2a2a1151996e43044f76e8c0f3be4c36c8cff7b27aa2b76b262fd
  glucose-syrup                  & algotocnf\_24-t9.cnf           & none       & no    & 125761 / 671432 &           & 84613298  & 129311     & timeout \\ % 23cdb7de37fdc6cf3a94466ca36accb30b9d0d5b5436effe23ac19d3
  glucose-syrup                  & algotocnf\_24-t9.cnf           & cmsat      & no    & 65579 / 377228 &           & 100827233 & 147009     & timeout \\ % c19b2572b2cd634ff3853f3f96544f3ed914f12b72a8748c518de42a
  glucose-syrup                  & algotocnf\_24-t9.cnf           & minisat    & no    & 69023 / 436292 &           & 85458489  & 118118     & timeout \\ % 220203195c78f6722a9a34641eaa4184802e682a2679bc4fc8abae16
  glucose-syrup                  & algotocnf\_24-t9.cnf           & satelite   & no    & 79486 / 515977 &           & 96089748  & 140970     & timeout \\ % 4ee8ef0bcbbb6c246c05b461701ce0e087098195ba03b4744cf2c188
  glucose-syrup                  & algotocnf\_24-t9\_diff-desc.cnf & none       & no    & 125761 / 957493 &           & 100515468 & 178519     & timeout \\ % eedeb59424902e168adb28bf51d5aba96b41ea3b5fc5252faf5e69d4
  glucose-syrup                  & algotocnf\_24-t9\_diff-desc.cnf & cmsat      & no    & 56447 / 302014 &           & 78765879  & 132913     & timeout \\ % 7c22d79062a7fe42df4bcb25e4789ef0a8d6ee207e68d4bfe8a0bb82
  glucose-syrup                  & algotocnf\_24-t9\_diff-desc.cnf & minisat    & no    & 85312 / 558731 &           & 98178126  & 164169     & timeout \\ % 0db195bc30298c17460b38f211098753bb350d9fd4207ec2cb7a02e5
  glucose-syrup                  & algotocnf\_24-t9\_diff-desc.cnf & satelite   & no    & 95154 / 750557 &           & 106350807 & 179879     & timeout \\ % 0844182ead77ef11bf2829bb06efae7819bec5c2dd70978876ec26d1
  glucose-syrup                  & algotocnf\_27-t10-diff.cnf     & none       & no    & 134722 / 8320 &           & 0         & 0          & 0 \\ % 36d158c0d354171b4c85ebc1e4fc87f56eb32982522e2dad8aaaff48
  glucose-syrup                  & algotocnf\_27-t10-diff.cnf     & cmsat      & no    & 0 / 0      &           & 0         & 0          & 0 \\ % f10a084c87b1236f40bf157b92908c43b35cd4ee0211285431f9f3ad
  glucose-syrup                  & algotocnf\_27-t10-diff.cnf     & lingeling  & no    & 0 / 0      &           & 0         & 0          & 0 \\ % e232ba66e6948fb763de05831f7b0ebcacca4f59e8188a2027ca32cb
  glucose-syrup                  & algotocnf\_27-t10-diff.cnf     & minisat    & no    & 0 / 0      &           & 0         & 0          & 0 \\ % 0ceb44e48526c1abd99d37b7f5293d1521471f2755ea6e8274440c6a
  glucose-syrup                  & algotocnf\_27-t10-diff.cnf     & satelite   & no    & 0 / 0      &           & 0         & 0          & 0 \\ % 63dace2ee7539a29f576f27b162ef2ee285bdc3373f4a99118bbdf58
  glucose-syrup                  & algotocnf\_27-t10-diff\_diff-desc.cnf & cmsat      & no    & 0 / 0      &           & 0         & 0          & 0 \\ % fa0d7a0898848e0b16847dcae07fc5a69cd4fb16a89366ff0ab46bf1
  glucose-syrup                  & algotocnf\_27-t10-diff\_diff-desc.cnf & lingeling  & no    & 0 / 0      &           & 0         & 0          & 0 \\ % a07d500804c92aec6294bcf55c2349ab3dec6512ae8cd82f999fb73e
  glucose-syrup                  & algotocnf\_27-t10-diff\_diff-desc.cnf & minisat    & no    & 0 / 0      &           & 0         & 0          & 0 \\ % 8bb9d84e7986f25636be520f21ce69074a37f894ed1f88ee18612070
  glucose-syrup                  & algotocnf\_27-t10-diff\_diff-desc.cnf & satelite   & no    & 0 / 0      &           & 0         & 0          & 0 \\ % 4304e3b243d312eedf50876789bbb1b729fd9b4715d61a974ba87546
  glucose-syrup                  & algotocnf\_27-t10-simplified.cnf & none       & no    & 134722 / 715830 &           & 119526558 & 141070     & timeout \\ % 2b4489150f366bf4ad953a38d0d54676a7fd9c06ea989723e82f942d
  glucose-syrup                  & algotocnf\_27-t10-simplified.cnf & cmsat      & no    & 67175 / 392680 &           & 91037046  & 147343     & timeout \\ % 62e0c9c57160144e999d47126c59480cda5da6e1e6de8b6a64c64346
  glucose-syrup                  & algotocnf\_27-t10-simplified.cnf & minisat    & no    & 72909 / 460404 &           & 96223551  & 94552      & timeout \\ % 57f6bd83eb2fe7172cf2bdf90b21ce2a5e7b7524a28e88a962ac5a41
  glucose-syrup                  & algotocnf\_27-t10-simplified.cnf & satelite   & no    & 84442 / 544998 &           & 97896761  & 92353      & timeout \\ % 000fbb49f53577c4890a5bf784f94d2dd764ec5fd054e3c0386b50a9
  glucose-syrup                  & algotocnf\_27-t10-simplified\_diff-desc.cnf & none       & no    & 134722 / 1020826 &           & 89825381  & 145503     & timeout \\ % 50066a82cf1829be894225f7e4ba3728c5bd8c91679cc8ed61901c85
  glucose-syrup                  & algotocnf\_27-t10-simplified\_diff-desc.cnf & cmsat      & no    & 56223 / 302443 &           & 67175956  & 109663     & timeout \\ % 0fe6ada24d49c160d43acdf4d018cc26271d8969320633cfc5de602d
  glucose-syrup                  & algotocnf\_27-t10-simplified\_diff-desc.cnf & minisat    & no    & 89667 / 586309 &           & 99815898  & 166069     & timeout \\ % 4ba3b23b3df351e75e1c3e72c85f566a24f7fe60646d39a11c21361d
  glucose-syrup                  & algotocnf\_27-t10-simplified\_diff-desc.cnf & satelite   & no    & 100735 / 789589 &           & 99608876  & 166684     & timeout \\ % f8a754ba4f35a9571d518b4bfce597f839494321b0abd0dd44100a75
    \end{tabular}
  \end{center}
\end{table}

%Text such that \LaTeX{} won't screw up.
\newpage

\begin{table}[p]
  \begin{center}
    \begin{tabular}{l|cccccccc}
        \textbf{SAT solver} & \textbf{CNF file} & \textbf{simplified} & \textbf{phase} & \textbf{vars/clauses} & \textbf{propagations} & \textbf{decisions} & \textbf{restarts} & \textbf{Runtime (sec)} \\
      \hline
  glucose-syrup                  & algotocnf\_27-t10.cnf          & none       & no    & 134722 / 717878 &           & 77657134  & 105575     & timeout \\ % 2916ec214b6368ae9cb3de0e1fe1aa925ef5d9ce00dd8afc865b37e9
  glucose-syrup                  & algotocnf\_27-t10.cnf          & cmsat      & no    & 72044 / 411276 &           & 104111969 & 126466     & timeout \\ % 0913ca8fe8740d51317b54be0a88224be74e292161f4caf86c1559da
  glucose-syrup                  & algotocnf\_27-t10.cnf          & minisat    & no    & 73897 / 465248 &           & 118165803 & 127158     & timeout \\ % ebf32948d138f0d25f91a758d7adbdc89206784be9abfe280f156be8
  glucose-syrup                  & algotocnf\_27-t10.cnf          & satelite   & no    & 85274 / 549734 &           & 77044829  & 93029      & timeout \\ % f30aea1c1915b98c5b24897e5a5f1a6dc150b3f1254701f0e68cd056
  glucose-syrup                  & algotocnf\_27-t10\_diff-desc.cnf & none       & no    & 134722 / 1022874 &           & 80346128  & 139519     & timeout \\ % 47e4dc01e430bef898899887ecfab3dceb967f0ddf5797eb188d0690
  glucose-syrup                  & algotocnf\_27-t10\_diff-desc.cnf & cmsat      & no    & 57956 / 311219 &           & 81214152  & 132456     & timeout \\ % d828d5e6781d50f0859c46ae795dbf36f110b693e0533a023c7cbe5a
  glucose-syrup                  & algotocnf\_27-t10\_diff-desc.cnf & minisat    & no    & 94850 / 619480 &           & 103396479 & 185354     & timeout \\ % c7a19434dcd8f96f7bb3ed3c3027c0554080b79e2564f3dd061923d4
  glucose-syrup                  & algotocnf\_27-t10\_diff-desc.cnf & satelite   & no    & 105733 / 857675 &           & 91342072  & 167639     & timeout \\ % 3d4943663659e6948d8bbabb2106b64470dfe406647d63de90c37d2f
  glucose-syrup                  & algotocnf\_2\_diff-desc.cnf    & none       & no    & 48704 / 329788 &           & 22348276  & 27402      & 6985 \\ % b0fada78e55bb235985ea57f900b4fadaa4261be44753f144dcca591
  glucose-syrup                  & algotocnf\_2\_diff-desc.cnf    & cmsat      & no    & 15342 / 72870 &           & 12180094  & 11970      & 2306 \\ % f5a4897f85e10c1aab90eb3a2414050ff08713bba93e5e50c1b7c370
  glucose-syrup                  & algotocnf\_2\_diff-desc.cnf    & lingeling  & no    & 40767 / 106546 &           & 19654829  & 20212      & 5726 \\ % 1c39476c873100191bedad636ce9041562ee4b7cbc90daa9e9539ee4
  glucose-syrup                  & algotocnf\_2\_diff-desc.cnf    & minisat    & no    & 23531 / 149436 &           & 11475483  & 11520      & 1903 \\ % 62cde986af3f1a17fea44fd33c2b8b9633a1696bf43ab04623f1442c
  glucose-syrup                  & algotocnf\_2\_diff-desc.cnf    & satelite   & no    & 29428 / 204577 &           & 9428928   & 9082       & 1808 \\ % a95053bc808a86881cbbd9f6562bdf753b31aacfc085c9d473a0a9bd
  glucose-syrup                  & algotocnf\_3.cnf               & none       & no    & 48704 / 222436 &           & 611213    & 153        & 31 \\ % 015a99b6ceceecdba0ad58f4594e7ad98af71ba4deb58376da06c0a5
  glucose-syrup                  & algotocnf\_3.cnf               & cmsat      & no    & 23893 / 111403 &           & 619646    & 133        & 28 \\ % 9486ecdf1047ba162701da7a260a5d1dff68f79f3e47074e4ef99fc0
  glucose-syrup                  & algotocnf\_3.cnf               & lingeling  & no    & 44000 / 100941 &           & 319727    & 63         & 10 \\ % 56db9c841f45039daa593a61ed0dc3d8c8869cf65cb9ad24d993e2a3
  glucose-syrup                  & algotocnf\_3.cnf               & minisat    & no    & 20895 / 119580 &           & 586941    & 155        & 25 \\ % 5d12f7358e2343a460fabe30a239cd1be22705fe43a937a373ab83c3
  glucose-syrup                  & algotocnf\_3.cnf               & satelite   & no    & 28647 / 156910 &           & 724170    & 179        & 29 \\ % 4da4e8c031927554da6e9f483866c6ccd91d4a65eca696bbea7e0c40
  glucose-syrup                  & algotocnf\_3\_diff-desc.cnf    & none       & no    & 48704 / 328004 &           & 321422    & 71         & 21 \\ % 74014981b3d03530d84f80362bd742b555ea7ad871af91fc9ad226e4
  lingeling-ats1o1               & algotocnf\_1.cnf               & cmsat      & no    & 24503 / 111931 & 830112341 & 17162591  &            & 906 \\ % 0888ccd7ebd4a609d6090ac4e92a755e7675e0f4c6297efa3e4e90a8
  lingeling-ats1o1               & algotocnf\_1.cnf               & cmsat      & yes   & 24503 / 111931 & 820517590 & 17498492  &            & 867 \\ % d9e99da0069e0ebde53d9a68f3f530ec26dacf5fd9442b0db4750a4a
  lingeling-ats1o1               & algotocnf\_1.cnf               & minisat    & no    & 20895 / 118236 & 575783290 & 16139460  &            & 860 \\ % daf0b5758cc6a1ee2bceeb7d9a7ea687ddcf940eab91d2a9c67ed69f
  lingeling-ats1o1               & algotocnf\_1.cnf               & minisat    & yes   & 20895 / 118236 & 333100908 & 11587868  &            & 551 \\ % 65d1affa7c5d180dc4eec4cd16c0ddec0afd8faaf0684b4cf63202cb
  lingeling-ats1o1               & algotocnf\_1.cnf               & satelite   & no    & 27495 / 153262 & 855959855 & 20330973  &            & 1297 \\ % ffafd0811ba627c748bc2656fcdfc4444d0edc5fa8146a3b09363189
  lingeling-ats1o1               & algotocnf\_1.cnf               & satelite   & yes   & 27495 / 153262 & 455268862 & 14056621  &            & 538 \\ % ef3c644cd4f61718e611d367ee47ddb99fbfdef58f9acd6d30158a62
  lingeling-ats1o1               & algotocnf\_18-t9.cnf           & none       & no    & 107839 / 590953 & 6579999   & 519222    &            & 29 \\ % a86efc5e422d702b40dfacfc7a45a767d51075f7581f2c6ceeb63d12
  lingeling-ats1o1               & algotocnf\_18-t9.cnf           & none       & yes   & 107839 / 590953 & 6637382   & 541483    &            & 29 \\ % b4b0c3411402de03d84a6e5e55f8a318702a554408bfe9b854f3d573
  lingeling-ats1o1               & algotocnf\_18-t9.cnf           & cmsat      & no    & 61789 / 342139 & 19838315  & 1377715   &            & 54 \\ % bfd14fb4519d673a4259d4934f1584d1bf965a174f7c5664cd1a3728
  lingeling-ats1o1               & algotocnf\_18-t9.cnf           & cmsat      & yes   & 61789 / 342139 & 20128556  & 895906    &            & 57 \\ % b6ab967691e571392b53325c674f67ec0b2da030b8c721105ec67fc3
  lingeling-ats1o1               & algotocnf\_18-t9.cnf           & lingeling  & no    & 102616 / 344544 & 4891154   & 415677    &            & 19 \\ % 8fe2b61ffd2bd260405f5d63e73468cae8ef721513986e1d4f4a3602
  lingeling-ats1o1               & algotocnf\_18-t9.cnf           & lingeling  & yes   & 102616 / 344544 & 3026807   & 187113    &            & 17 \\ % 20c2fa39221a4110e23fedbe50cec7448ba1bc8660f15ff7cfead7e5
  lingeling-ats1o1               & algotocnf\_18-t9.cnf           & minisat    & no    & 61845 / 391134 & 6037622   & 425290    &            & 21 \\ % 55072758ede63fc35731cb76b81fb2fc4b750831b1d806d34ce6065b
  lingeling-ats1o1               & algotocnf\_18-t9.cnf           & minisat    & yes   & 61845 / 391134 & 8386925   & 494325    &            & 25 \\ % de52e2792d76d4612d4f6860643db6502a6aec1c7b479856af5e87c9
  lingeling-ats1o1               & algotocnf\_18-t9.cnf           & satelite   & no    & 69670 / 457972 & 4521332   & 334680    &            & 29 \\ % dcfed67f629ce4b508f7d7796bbc3c856af8d9e3baeecf21a2813deb
  lingeling-ats1o1               & algotocnf\_18-t9.cnf           & satelite   & yes   & 69670 / 457972 & 6584687   & 568455    &            & 24 \\ % 4c7c4adb74daf68cd45ae170977bef1225e451ba8761d470bfb72ccc
    \end{tabular}
  \end{center}
\end{table}

%Text such that \LaTeX{} won't screw up.
\newpage

\begin{table}[p]
  \begin{center}
    \begin{tabular}{l|cccccccc}
        \textbf{SAT solver} & \textbf{CNF file} & \textbf{simplified} & \textbf{phase} & \textbf{vars/clauses} & \textbf{propagations} & \textbf{decisions} & \textbf{restarts} & \textbf{Runtime (sec)} \\
      \hline
  lingeling-ats1o1               & algotocnf\_18-t9\_diff-desc-sc\_ocnf.cnf & none       & no    & 196473 / 1197423 & 66966666158 & 10840122346 &            & 124876 \\ % aaa659f5d7fd6d341a0fc0cc48be9bd996b3ca08c446bdc28863be5b
  lingeling-ats1o1               & algotocnf\_18-t9\_diff-desc-sc\_ocnf.cnf & none       & yes   & 196473 / 1197423 & 65508201410 & 10232669290 &            & 124859 \\ % 5023729ab8d9398fbd3d945d5b84b700627584927ea9f3d360a1cef2
  lingeling-ats1o1               & algotocnf\_18-t9\_diff-desc.cnf & none       & yes   & 107839 / 846487 & 72334119  & 509989    &            & 99 \\ % c1cdc6c4fdc37464138ee23270ed2ee4c3d5b3a7814eaba670fd7f6a
  lingeling-ats1o1               & algotocnf\_18-t9\_diff-desc.cnf & cmsat      & no    & 57114 / 303489 & 49338586  & 2387529   &            & 99 \\ % b06ca5c711ecc6ff686a149e3a751e1e62b924eae75956966404dcda
  lingeling-ats1o1               & algotocnf\_18-t9\_diff-desc.cnf & cmsat      & yes   & 57114 / 303489 & 301812566 & 3965383   &            & 78 \\ % e9983c3e65e9a9e5697d3924d7a8e653ed68feff502866c9b7534a67
  lingeling-ats1o1               & algotocnf\_18-t9\_diff-desc.cnf & lingeling  & no    & 99452 / 404666 & 40453782  & 1995270   &            & 62 \\ % 192faed8f0861b468959bcf7c7ef4906acc9c3d1ee1175059c9c2858
  lingeling-ats1o1               & algotocnf\_18-t9\_diff-desc.cnf & lingeling  & yes   & 99452 / 404666 & 24115557  & 181714    &            & 96 \\ % 198b0e3a39c44603ac47a2876b514739b19915fc5b88aff9edb415bb
  lingeling-ats1o1               & algotocnf\_18-t9\_diff-desc.cnf & minisat    & no    & 78055 / 511431 & 9156673   & 774601    &            & 183 \\ % 10fa620fbd7c658fd630e3c992f5ef06d3dc0d5076363c3ada73534c
  lingeling-ats1o1               & algotocnf\_18-t9\_diff-desc.cnf & minisat    & yes   & 78055 / 511431 & 276293998 & 1576114   &            & 885 \\ % ae097abd1093074bb40da95a62486137d62b8191b634278db0717255
  lingeling-ats1o1               & algotocnf\_18-t9\_diff-desc.cnf & satelite   & no    & 85386 / 692408 & 44249412  & 1148783   &            & 50 \\ % 8d70abcb50e2d7d2e7986abaae89bdf5c9bc7f947ee6b8edba9fffd2
  lingeling-ats1o1               & algotocnf\_18-t9\_diff-desc.cnf & satelite   & yes   & 85386 / 692408 & 126162246 & 981514    &            & 150 \\ % 4bfefb318c85fdf5af2d9217965792d32836108b3006c94f9c472e12
  lingeling-ats1o1               & algotocnf\_18-t9\_diff-desc\_ocnf.cnf & none       & no    & 107839 / 846487 & 66891741743 & 10848404602 &            & 124869 \\ % 4a86090b0f9196cb9f3cf10fd08ff9258c267993aac9b03c3caa0697
  lingeling-ats1o1               & algotocnf\_18-t9\_diff-desc\_ocnf.cnf & none       & yes   & 107839 / 846487 & 68433381304 & 10719447822 &            & 124822 \\ % 9adcb088e76a73f1683fbbad034121683cd7f9a52f5f4f766a23c5c6
  lingeling-ats1o1               & algotocnf\_18-t9\_ocnf.cnf     & none       & no    & 107839 / 590953 & 61726388410 & 53023507367 &            & 89845 \\ % cb4e4331984bde02a5a06d2690f511362a64930d7742227957fa2e90
  lingeling-ats1o1               & algotocnf\_18-t9\_ocnf.cnf     & none       & yes   & 107839 / 590953 & 54217462922 & 46160327639 &            & 89818 \\ % 5bb9c34434105fd38bda60eafa7338404bf2523f842f3f700195d083
  lingeling-ats1o1               & algotocnf\_1\_diff-desc-sc\_ocnf.cnf & none       & no    & 85760 / 522144 & 70194750  & 5542349   &            & 65 \\ % dda14aeed6bbd4a610eee25a4828bec97baf2b91cff0520d544a1735
  lingeling-ats1o1               & algotocnf\_1\_diff-desc-sc\_ocnf.cnf & none       & yes   & 85760 / 522144 & 61858501  & 5611067   &            & 58 \\ % 3de89f920430fd545092913e0b6bd4f9d29ccc55e8684658a15a0aed
  lingeling-ats1o1               & algotocnf\_1\_diff-desc.cnf    & none       & no    & 48704 / 373920 & 29525944  & 1719106   &            & 97 \\ % cc3c0cf1416eac75af117d6b8e64ca9eb98afc4adb97199e3f9226c7
  lingeling-ats1o1               & algotocnf\_1\_diff-desc.cnf    & none       & yes   & 48704 / 373920 & 63328779  & 2536475   &            & 82 \\ % 8736e30df76a18b51ec2ce7a65613d0806301aff17fa004ca9d94d8a
  lingeling-ats1o1               & algotocnf\_1\_diff-desc.cnf    & cmsat      & no    & 15402 / 73142 & 60741384  & 2474533   &            & 62 \\ % 0afd0602412894b8bf21f4a5b250eae0c2faecdbf58408fc1fc1c2ae
  lingeling-ats1o1               & algotocnf\_1\_diff-desc.cnf    & cmsat      & yes   & 15402 / 73142 & 83911694  & 2838140   &            & 70 \\ % fa1c7008753e25f5c33b8b0774c5217885aa5883a6db03fbfdfc9f10
  lingeling-ats1o1               & algotocnf\_1\_diff-desc.cnf    & lingeling  & no    & 40554 / 101182 & 56989140  & 3075627   &            & 73 \\ % 753622078bada8d7c43682c57d981aa10e79d1e999d9986dc3bd4159
  lingeling-ats1o1               & algotocnf\_1\_diff-desc.cnf    & lingeling  & yes   & 40554 / 101182 & 29324158  & 1976906   &            & 45 \\ % 2ea9f68710038bb1d1a153de95354f5eeb8a0c51577783eb90a99581
  lingeling-ats1o1               & algotocnf\_1\_diff-desc.cnf    & minisat    & no    & 25107 / 158393 & 72206425  & 2441239   &            & 74 \\ % a11892b4594332d5bc643c63830da312822b507b3c039fba84882f3b
  lingeling-ats1o1               & algotocnf\_1\_diff-desc.cnf    & minisat    & yes   & 25107 / 158393 & 59630317  & 1972292   &            & 63 \\ % 8b5488a0567d3b97ce594609387f7517c555299cdb98a0a9c43e0d9f
  lingeling-ats1o1               & algotocnf\_1\_diff-desc.cnf    & satelite   & no    & 30950 / 226602 & 44335947  & 1920992   &            & 69 \\ % c7fa2b13c2f96b3c7f3edfa34ca2df2299eac8dc38c78268030bc8d0
  lingeling-ats1o1               & algotocnf\_1\_diff-desc.cnf    & satelite   & yes   & 30950 / 226602 & 60556070  & 2674102   &            & 62 \\ % cae20c476154f6ca2b52d8b2683619581be4a8085a9f85a8598d88ce
  lingeling-ats1o1               & algotocnf\_1\_diff-desc\_ocnf.cnf & none       & no    & 48704 / 373920 & 95072840  & 7417561   &            & 163 \\ % 353352726560d082abad93a9a8a0a284d1a3b93cf866f92bc6c5ab58
  lingeling-ats1o1               & algotocnf\_1\_diff-desc\_ocnf.cnf & none       & yes   & 48704 / 373920 & 223255639 & 11718590  &            & 146 \\ % 6767e4ae15016362f8300fd47ee62f37e6986b3aeab5c1cbc67f828f
  lingeling-ats1o1               & algotocnf\_1\_ocnf.cnf         & none       & no    & 48704 / 253984 & 242646021 & 9046602   &            & 854 \\ % cf96c5b1eb996259470e2a0e2fff5101a288a7ed313dd6ee30ebd4f6
  lingeling-ats1o1               & algotocnf\_1\_ocnf.cnf         & none       & yes   & 48704 / 253984 & 156311236 & 5780694   &            & 525 \\ % e042242ac9c6adf203dddeb7d8c6f4c8f41e21cc8828e40f5385701a
  lingeling-ats1o1               & algotocnf\_2.cnf               & none       & no    & 48704 / 254210 & 96220197  & 4112922   &            & 260 \\ % 42392cb867c51d5791e1b08e3815bb445ba6e91b2eb7ec3ae55f2605
  lingeling-ats1o1               & algotocnf\_2.cnf               & none       & yes   & 48704 / 254210 & 100536825 & 4630586   &            & 820 \\ % bc630e7cab148b09a702af397acf40ee28fd66691a85f6fe184acee9
  lingeling-ats1o1               & algotocnf\_2.cnf               & cmsat      & no    & 24323 / 111661 & 199273327 & 5678337   &            & 822 \\ % 4db34ad5d36416db654ac12ccfb8ee4a3b3cb010963e289b29cb3ef6
  lingeling-ats1o1               & algotocnf\_2.cnf               & cmsat      & yes   & 24323 / 111661 & 266257133 & 7894754   &            & 325 \\ % 3b23a69c53108413910160f4a63cac6f4dbf9d1c1031670b5b1a6975
    \end{tabular}
  \end{center}
\end{table}

%Text such that \LaTeX{} won't screw up.
\newpage

\begin{table}[p]
  \begin{center}
    \begin{tabular}{l|cccccccc}
        \textbf{SAT solver} & \textbf{CNF file} & \textbf{simplified} & \textbf{phase} & \textbf{vars/clauses} & \textbf{propagations} & \textbf{decisions} & \textbf{restarts} & \textbf{Runtime (sec)} \\
      \hline
  lingeling-ats1o1               & algotocnf\_2.cnf               & minisat    & no    & 20895 / 118688 & 254117613 & 9856171   &            & 508 \\ % 920dfbfffb6d3cea01d7ebd5f160143b7b10e08c1f70d13843bdb9f6
  lingeling-ats1o1               & algotocnf\_2.cnf               & minisat    & yes   & 20895 / 118688 & 137661864 & 5159552   &            & 232 \\ % 7023a5a443236cc1ff986dafe550da7ac90cba6c5607659a294215aa
  lingeling-ats1o1               & algotocnf\_2.cnf               & satelite   & no    & 27775 / 154276 & 369167237 & 10750678  &            & 792 \\ % a32072cf6064225e11ddb2e4e1ec5a81403b5b056903c83981d904b9
  lingeling-ats1o1               & algotocnf\_2.cnf               & satelite   & yes   & 27775 / 154276 & 582590815 & 16895000  &            & 679 \\ % b92293c3d2d65a2482bf7afaf3c94f2ae5aa42fec084423bb843af55
  lingeling-ats1o1               & algotocnf\_21-t9.cnf           & none       & no    & 116800 / 636838 & 6567493736 & 154119871 &            & 27599 \\ % e796768a1c8d357d2044c53ae416f5331a35eecf9bf836045e8c6658
  lingeling-ats1o1               & algotocnf\_21-t9.cnf           & none       & yes   & 116800 / 636838 & 1224816609 & 19965195  &            & 19247 \\ % b26c6d4e425014ea77126ac64800f3fe27d2453093b999686af7f4a9
  lingeling-ats1o1               & algotocnf\_21-t9.cnf           & cmsat      & no    & 63555 / 358308 & 8081065262 & 128186673 &            & 32099 \\ % 4dabd450c6f55331103e30477dfae53379d95467e5d13f455f96380c
  lingeling-ats1o1               & algotocnf\_21-t9.cnf           & cmsat      & yes   & 63555 / 358308 & 5169630907 & 84207170  &            & 30244 \\ % d1a711e30db7fc8ff08bd27392529479782da7b2fa7d15da7e796775
  lingeling-ats1o1               & algotocnf\_21-t9.cnf           & minisat    & no    & 65572 / 414470 & 7555279670 & 155614117 &            & 30205 \\ % 0a86026ce5abf6422e87b14bf2289f33cbf0fbeafb6ba54ad2a8bf6d
  lingeling-ats1o1               & algotocnf\_21-t9.cnf           & minisat    & yes   & 65572 / 414470 & 3766091891 & 64175495  &            & 20193 \\ % 8549f6d8c9c16f9d87ed9fad3615d293bccd1206e6d2b7b422aace3c
  lingeling-ats1o1               & algotocnf\_21-t9.cnf           & satelite   & no    & 74562 / 486714 & 9417657748 & 212696810 &            & 44494 \\ % b497c575757bb4302e901d3b67436f4210cadb9281e17048f727135a
  lingeling-ats1o1               & algotocnf\_21-t9.cnf           & satelite   & yes   & 74562 / 486714 & 5033023924 & 125409709 &            & 29234 \\ % e8444b717e51cb4e82af4d078df829368f4889919685ddb5e09b355d
  lingeling-ats1o1               & algotocnf\_21-t9\_diff-desc.cnf & none       & no    & 116800 / 911629 & 3799761108 & 15913739  &            & 2282 \\ % 4b786a2cab486f843ddd27df32dc6c625272fe4687a1311576760aef
  lingeling-ats1o1               & algotocnf\_21-t9\_diff-desc.cnf & none       & yes   & 116800 / 911629 & 897799847 & 3355217   &            & 1172 \\ % b3922d512053180ecb9569cf154cc3d9ab4c77120a110096be7f0303
  lingeling-ats1o1               & algotocnf\_21-t9\_diff-desc.cnf & cmsat      & no    & 59528 / 318049 & 193901610 & 4750865   &            & 892 \\ % bcdb7cd03449f01f858e952d3ff5413fd846ab0a8a288fb26a959489
  lingeling-ats1o1               & algotocnf\_21-t9\_diff-desc.cnf & cmsat      & yes   & 59528 / 318049 & 711520249 & 12255201  &            & 537 \\ % 8ac0d5f36ffe86bd00bad86ffc89f1a6e65070c35aa35141c8442a2d
  lingeling-ats1o1               & algotocnf\_21-t9\_diff-desc.cnf & minisat    & no    & 81650 / 534965 & 1719807367 & 5323332   &            & 596 \\ % 1c3a0dda839ec43841ae5cee3ad3b82a66f98f843ccdc249c0c7e042
  lingeling-ats1o1               & algotocnf\_21-t9\_diff-desc.cnf & minisat    & yes   & 81650 / 534965 & 1086561165 & 10408544  &            & 1240 \\ % b9b9e312708bf59671339c0e83424f43ccfc45cb872ece6cafb0d174
  lingeling-ats1o1               & algotocnf\_21-t9\_diff-desc.cnf & satelite   & no    & 90236 / 720755 & 819924550 & 6702460   &            & 791 \\ % 80335e49df47cfdc8dc58e875db2a82f8d8e833f35ddc2e0ad952662
  lingeling-ats1o1               & algotocnf\_21-t9\_diff-desc.cnf & satelite   & yes   & 90236 / 720755 & 463392096 & 5426586   &            & 1124 \\ % 71869ae8cdca77f22bdf558dc582c5f2e8d3dc70a156e3e8383dd207
  lingeling-ats1o1               & algotocnf\_21-t9\_diff-desc\_ocnf.cnf & none       & no    & 116800 / 911629 & 47871721121 & 6760244481 &            & 89839 \\ % 564ddd74a5670b8ac98a4b5e6efce5f521e64391347736112b1fad88
  lingeling-ats1o1               & algotocnf\_21-t9\_diff-desc\_ocnf.cnf & none       & yes   & 116800 / 911629 & 50571983496 & 6640481141 &            & 89869 \\ % 8b670970392d8ecc524920880c2ac8ba2ff8bd673f20012e612fc66a
  lingeling-ats1o1               & algotocnf\_23-t9.cnf           & none       & no    & 122774 / 667438 & 11406748416 & 217188108 &            & 59312 \\ % 1791e15f3e7fc2b98fec8c7faeea38db4e010a70309b4432bcbe8c1d
  lingeling-ats1o1               & algotocnf\_23-t9.cnf           & none       & yes   & 122774 / 667438 & 9255197381 & 185869252 &            & 57045 \\ % 6b50082bc243f077a4577b9d01dd88ce23314a9f1960e129f563424f
  lingeling-ats1o1               & algotocnf\_23-t9.cnf           & cmsat      & no    & 64880 / 369775 & 9492444885 & 145903933 &            & 53893 \\ % 6f5e8dac8f27654a7ebc35296c89397dbb20c6c9f8ded59ec85380f8
  lingeling-ats1o1               & algotocnf\_23-t9.cnf           & cmsat      & yes   & 64880 / 369775 & 14551473440 & 224468851 &            & 83137 \\ % afb7ae6ecb6ced448470c3d7d884d1e4c3288dda901722d0b6617e06
  lingeling-ats1o1               & algotocnf\_23-t9.cnf           & minisat    & no    & 67979 / 429734 & 7755711599 & 138305657 &            & 55343 \\ % 37c469cd4296d8e4956d7dcd73f91369f0b38d2d969734c0bd452cb0
  lingeling-ats1o1               & algotocnf\_23-t9.cnf           & satelite   & no    & 77835 / 506217 & 9669970149 & 182406674 &            & 81836 \\ % e6798c79de331d49d599b20decf0c0123bda7d5fd6d0e906b3e7cae0
  lingeling-ats1o1               & algotocnf\_23-t9.cnf           & satelite   & yes   & 77835 / 506217 & 10029572768 & 175753619 &            & 62471 \\ % 515c79611bcf9d0a1f8459f85736fdbc754978a9b684d7fe3e83115b
  lingeling-ats1o1               & algotocnf\_23-t9\_diff-desc-sc\_ocnf.cnf & none       & no    & 222613 / 1350343 & 35575573063 & 4725372248 &            & 89822 \\ % b5dbe947be45283ca4ae30b6ebce11360ef6f31391d81d3840afb308
  lingeling-ats1o1               & algotocnf\_23-t9\_diff-desc-sc\_ocnf.cnf & none       & yes   & 222613 / 1350343 & 41256543353 & 5230217945 &            & 89844 \\ % 69b3aaa963d41984535b0fb01c4e55bff572f08b2a27f2d9110eceab
  lingeling-ats1o1               & algotocnf\_23-t9\_diff-desc.cnf & none       & no    & 122774 / 955067 & 945949965 & 17154521  &            & 1314 \\ % 8b8a4ef58e202963ae86435143a34f04594cb0e211516c2e00be6213
  lingeling-ats1o1               & algotocnf\_23-t9\_diff-desc.cnf & none       & yes   & 122774 / 955067 & 695528751 & 10321022  &            & 1506 \\ % 56ebbb031d2d881a4165eaf5ce1cf17dbf2ee32ef3b0e27566df94f0
  lingeling-ats1o1               & algotocnf\_23-t9\_diff-desc.cnf & cmsat      & no    & 57510 / 307118 & 706776958 & 13809628  &            & 1005 \\ % aef50dde52024e5bd402430e3e49e3aa7da082978055c2fa89d06f23
  lingeling-ats1o1               & algotocnf\_23-t9\_diff-desc.cnf & cmsat      & yes   & 57510 / 307118 & 1221950463 & 17665069  &            & 1546 \\ % 3044a79acecf1f1423e83e6eb351fa123552c624b492b56f0940f9ac
    \end{tabular}
  \end{center}
\end{table}

%Text such that \LaTeX{} won't screw up.
\newpage

\begin{table}[p]
  \begin{center}
    \begin{tabular}{l|cccccccc}
        \textbf{SAT solver} & \textbf{CNF file} & \textbf{simplified} & \textbf{phase} & \textbf{vars/clauses} & \textbf{propagations} & \textbf{decisions} & \textbf{restarts} & \textbf{Runtime (sec)} \\
      \hline
  lingeling-ats1o1               & algotocnf\_23-t9\_diff-desc.cnf & minisat    & no    & 84080 / 550727 & 499208360 & 6985772   &            & 1238 \\ % def3f8b273d79fb407030dca3beca03356d791a8c8b6906e4598bb33
  lingeling-ats1o1               & algotocnf\_23-t9\_diff-desc.cnf & minisat    & yes   & 84080 / 550727 & 1358098860 & 17095121  &            & 1624 \\ % ffb11ef061b518dc8c56bf1ddefa0dce57885f8b69394a53da112fc2
  lingeling-ats1o1               & algotocnf\_23-t9\_diff-desc.cnf & satelite   & no    & 93504 / 740545 & 2975485032 & 20909550  &            & 3184 \\ % d0bda72ce0e51157b74f08d315e445a8f9ccca6c7aae0e2ed549f262
  lingeling-ats1o1               & algotocnf\_23-t9\_diff-desc.cnf & satelite   & yes   & 93504 / 740545 & 1125859276 & 5811303   &            & 4692 \\ % 1e8516eb8998c34d94a7f17f2662c80b32bbf06d32df946835301ffc
  lingeling-ats1o1               & algotocnf\_23-t9\_diff-desc\_ocnf.cnf & none       & no    & 122774 / 955067 & 39379759048 & 5606250709 &            & 89761 \\ % df2fe6348dacf65078246f054f75fa72e645a31a43eec311b26764a5
  lingeling-ats1o1               & algotocnf\_23-t9\_diff-desc\_ocnf.cnf & none       & yes   & 122774 / 955067 & 40325598456 & 5781976401 &            & 89810 \\ % 83b0aa2a490a96165f961deda7467871eddfa468ee42bfa6d669f076
  lingeling-ats1o1               & algotocnf\_23-t9\_ocnf.cnf     & none       & no    & 122774 / 667438 & 51774535244 & 43295595444 &            & 89857 \\ % 514b6da7520a66874d1c5961d57aa222f5eacc75ebf21f21808ffd8e
  lingeling-ats1o1               & algotocnf\_23-t9\_ocnf.cnf     & none       & yes   & 122774 / 667438 & 50026224361 & 41228755262 &            & 89850 \\ % 2ffd53a6c58cfccec6b6a29e867993adbc4a09ac53b1858cd932f473
  lingeling-ats1o1               & algotocnf\_24-t9.cnf           & none       & no    & 125761 / 682722 & 8197106091 & 147635256 &            & 66052 \\ % b026751ca55a291dd5c3c9ac1f89793517ad2d2061821063b08107f0
  lingeling-ats1o1               & algotocnf\_24-t9.cnf           & none       & yes   & 125761 / 682722 & 11176742237 & 212766331 &            & 81935 \\ % d21d15e272294ebff91380949da4f82cb7b675fafd03a168d2b468f3
  lingeling-ats1o1               & algotocnf\_24-t9.cnf           & cmsat      & no    & 65579 / 377228 & 15021825301 & 237337855 &            & 89566 \\ % 037c7f3013bc0fb78e1b2969f174154f7161310a9fb38e5de4ee55f4
  lingeling-ats1o1               & algotocnf\_24-t9.cnf           & cmsat      & yes   & 65579 / 377228 & 13856965209 & 199420781 &            & 87614 \\ % 24816fb7f4b87406083898bee1e3b9eae7b31a6edeb9933f70dd170e
  lingeling-ats1o1               & algotocnf\_24-t9.cnf           & minisat    & no    & 69023 / 436292 & 10756062377 & 194157926 &            & 71580 \\ % c17f0ec2f05878c5637d9ab633e5b1b4c2746a28077a32636dfea75a
  lingeling-ats1o1               & algotocnf\_24-t9.cnf           & minisat    & yes   & 69023 / 436292 & 10170680100 & 181353444 &            & 68580 \\ % e4e9a66f2ada502a2b70344a3ffccde671dfb4193f18d6ea5d949cff
  lingeling-ats1o1               & algotocnf\_24-t9.cnf           & satelite   & no    & 79486 / 515977 & 14801677241 & 269786968 &            & 97381 \\ % ddf41d64a28d2410c569e02b8460592077f580f5d7fcd9b20ad5996d
  lingeling-ats1o1               & algotocnf\_24-t9.cnf           & satelite   & yes   & 79486 / 515977 & 8262557163 & 139215714 &            & 93099 \\ % 578f9b80f1c8274a01636e7472b1f083cdf6ca6d07413f97be28576e
  lingeling-ats1o1               & algotocnf\_24-t9\_diff-desc-sc\_ocnf.cnf & none       & no    & 227841 / 1380914 & 34420266744 & 4135271712 &            & 89783 \\ % ce5062fcc69c2143576628d3bbbc22443586c8d2a460eb78448c6c5d
  lingeling-ats1o1               & algotocnf\_24-t9\_diff-desc-sc\_ocnf.cnf & none       & yes   & 227841 / 1380914 & 39557677456 & 5551974857 &            & 89816 \\ % bbd5799c1218eaf635b2d15712300a8dd99683ef8fa16f5a38f804f2
  lingeling-ats1o1               & algotocnf\_24-t9\_diff-desc.cnf & none       & no    & 125761 / 976770 & 4481773501 & 16901463  &            & 2632 \\ % 3ffacd892460941b20d0142b7a3b7df5d1e3df56a37ce9b8a228c795
  lingeling-ats1o1               & algotocnf\_24-t9\_diff-desc.cnf & none       & yes   & 125761 / 976770 & 781065944 & 11902522  &            & 1484 \\ % 2054fe924ad60fac82921ef734d3db0c28a2b66a8f5835ca8d9317fd
  lingeling-ats1o1               & algotocnf\_24-t9\_diff-desc.cnf & cmsat      & no    & 56447 / 302014 & 484584260 & 11579124  &            & 1105 \\ % c04a95e788a3356f02d944e11d662146b146a43d567ebd23c9d76db7
  lingeling-ats1o1               & algotocnf\_24-t9\_diff-desc.cnf & cmsat      & yes   & 56447 / 302014 & 892139481 & 11410685  &            & 1632 \\ % b9a8b3f018f1f8229c884ccc381939a9d73f8ff99d8d3520172a51a7
  lingeling-ats1o1               & algotocnf\_24-t9\_diff-desc.cnf & minisat    & no    & 85312 / 558731 & 1230585614 & 13221628  &            & 2961 \\ % a6a4e1fe982837d4e9cd6ec690bf6432900c510afd4f6c3c73d83528
  lingeling-ats1o1               & algotocnf\_24-t9\_diff-desc.cnf & minisat    & yes   & 85312 / 558731 & 995261008 & 7578211   &            & 1807 \\ % 85ef4323fc5ddbd7055edea0db06f5595a5b034e588028ff9471ae20
  lingeling-ats1o1               & algotocnf\_24-t9\_diff-desc.cnf & satelite   & no    & 95154 / 750557 & 1061878602 & 14146291  &            & 1291 \\ % 31c127ee50c4130c610e8f89784b79c94515f204cd577c8b650c6eca
  lingeling-ats1o1               & algotocnf\_24-t9\_diff-desc.cnf & satelite   & yes   & 95154 / 750557 & 362102146 & 4738718   &            & 1627 \\ % 2655682b16e2b2128efd936d0ebdfa1c27586ae22fc761bc004bc9c6
  lingeling-ats1o1               & algotocnf\_24-t9\_diff-desc\_ocnf.cnf & none       & no    & 125761 / 976770 & 39115183280 & 5095117316 &            & 89845 \\ % c3763e25f9c56c7b65f4b6210e66bd1c0bf452aee7d4b2bb87be9532
  lingeling-ats1o1               & algotocnf\_24-t9\_diff-desc\_ocnf.cnf & none       & yes   & 125761 / 976770 & 34060175655 & 3810110610 &            & 89835 \\ % 2843fc8486904acd7cbaa459d041297d7a36ba9766b39e58e3a808fd
  lingeling-ats1o1               & algotocnf\_27-t10-diff.cnf     & none       & no    & 134722 / 728349 & 134722    & 0         &            & 0 \\ % b4879ee42684caf08ad435f05752eafae5da98d342fbbf86ccbfc9cf
  lingeling-ats1o1               & algotocnf\_27-t10-diff.cnf     & none       & yes   & 134722 / 728349 & 134722    & 0         &            & 0 \\ % c5e2a623c54de43dec780795153d3f1f14050b4b9689096b2c688cf9
  lingeling-ats1o1               & algotocnf\_27-t10-diff.cnf     & cmsat      & no    & 0 / 0      & 0         & 0         &            & 0 \\ % e1b98133973922b2d3a2379a204a8253e28481915f1430afb822f18a
  lingeling-ats1o1               & algotocnf\_27-t10-diff.cnf     & cmsat      & yes   & 0 / 0      & 0         & 0         &            & 0 \\ % ced4a1cc1734caf273d750765d63cdcb9ba5695b9cae25ae6d0adeb8
  lingeling-ats1o1               & algotocnf\_27-t10-diff.cnf     & lingeling  & no    & 0 / 0      & 0         & 0         &            & 0 \\ % 025b081918846434f3f1224418c9df06f3d00fd9f0ee0d9abf5ecfbc
  lingeling-ats1o1               & algotocnf\_27-t10-diff.cnf     & lingeling  & yes   & 0 / 0      & 0         & 0         &            & 0 \\ % b1d13c3a2c8eb2d4d2eb715da965d99fa94b13094461d549bf0ab4a5
  lingeling-ats1o1               & algotocnf\_27-t10-diff.cnf     & minisat    & no    & 0 / 0      & 0         & 0         &            & 0 \\ % 0fe29861861f38385ce174ace0ca38cf37e8fb7482ce7e44c20c42d6
    \end{tabular}
  \end{center}
\end{table}

%Text such that \LaTeX{} won't screw up.
\newpage

\begin{table}[p]
  \begin{center}
    \begin{tabular}{l|cccccccc}
        \textbf{SAT solver} & \textbf{CNF file} & \textbf{simplified} & \textbf{phase} & \textbf{vars/clauses} & \textbf{propagations} & \textbf{decisions} & \textbf{restarts} & \textbf{Runtime (sec)} \\
      \hline
  lingeling-ats1o1               & algotocnf\_27-t10-diff.cnf     & minisat    & yes   & 0 / 0      & 0         & 0         &            & 0 \\ % 53391643a8f9d73e734513979735989b8060e255bf46179ca7d0b7d0
  lingeling-ats1o1               & algotocnf\_27-t10-diff.cnf     & satelite   & no    & 0 / 0      & 0         & 0         &            & 0 \\ % dbbeccd3a4e0c2c514cb1d75c3a6c299cd0008ab7a5c6b373e9b8d76
  lingeling-ats1o1               & algotocnf\_27-t10-diff.cnf     & satelite   & yes   & 0 / 0      & 0         & 0         &            & 0 \\ % 7fd74627c888483ef751d1097522d78a6bfe7594959f1f6fbd75339b
  lingeling-ats1o1               & algotocnf\_27-t10-diff\_diff-desc.cnf & none       & no    & 134722 / 1041654 & 134722    & 0         &            & 0 \\ % 4e3605c373104d23bdb2cb06d9119873088bee845a71410c0665b45d
  lingeling-ats1o1               & algotocnf\_27-t10-diff\_diff-desc.cnf & none       & yes   & 134722 / 1041654 & 134722    & 0         &            & 0 \\ % 11927e9d70e9086b790e5a4000ea9034ad57f221e7783053a91ebcc6
  lingeling-ats1o1               & algotocnf\_27-t10-diff\_diff-desc.cnf & cmsat      & no    & 0 / 0      & 0         & 0         &            & 0 \\ % fd7aee88f394863c41e4c4488cbd56c585bbc6ac09f0e743450565c4
  lingeling-ats1o1               & algotocnf\_27-t10-diff\_diff-desc.cnf & cmsat      & yes   & 0 / 0      & 0         & 0         &            & 0 \\ % 04e8175135155718e5a0271dea12c70f9fb4ba064a84aaba0e3bc3df
  lingeling-ats1o1               & algotocnf\_27-t10-diff\_diff-desc.cnf & lingeling  & no    & 0 / 0      & 0         & 0         &            & 0 \\ % a6c061135b91c0918c40a5b406d6a38562e792913b8266deee4f6aea
  lingeling-ats1o1               & algotocnf\_27-t10-diff\_diff-desc.cnf & lingeling  & yes   & 0 / 0      & 0         & 0         &            & 0 \\ % eca50401a08db695ac2cd319038d364c3bffe45d178bf5748ee340de
  lingeling-ats1o1               & algotocnf\_27-t10-diff\_diff-desc.cnf & minisat    & no    & 0 / 0      & 0         & 0         &            & 0 \\ % 8d80c276b5abf66c73da5cac21bebaa6b59bce0022c6e83099f07d01
  lingeling-ats1o1               & algotocnf\_27-t10-diff\_diff-desc.cnf & minisat    & yes   & 0 / 0      & 0         & 0         &            & 0 \\ % 8f0e05407b5975dbe9814eb6f427f640dee7c56b4e448f8b0810248d
  lingeling-ats1o1               & algotocnf\_27-t10-diff\_diff-desc.cnf & satelite   & no    & 0 / 0      & 0         & 0         &            & 0 \\ % 5eb6195a507f9b03d4eb6d9e58f85fc700e36cbaa49d033fcda88075
  lingeling-ats1o1               & algotocnf\_27-t10-diff\_diff-desc.cnf & satelite   & yes   & 0 / 0      & 0         & 0         &            & 0 \\ % e2773b319a92e699f4f0b85538906d0db6115f82fee439cab127c0c4
  lingeling-ats1o1               & algotocnf\_27-t10-diff\_diff-desc\_ocnf.cnf & none       & no    & 134722 / 1041654 & 134722    & 0         &            & 0 \\ % b8a8581bdd9f9892d9eefb4eba341e2b9c12f290a9b8af59aeb32991
  lingeling-ats1o1               & algotocnf\_27-t10-diff\_diff-desc\_ocnf.cnf & none       & yes   & 134722 / 1041654 & 134722    & 0         &            & 0 \\ % 70f25f51a356e52cc51a39fe5d21fbeb01eb55aa0b473072412a52b6
  lingeling-ats1o1               & algotocnf\_27-t10-diff\_ocnf.cnf & none       & no    & 134722 / 728349 & 134722    & 0         &            & 0 \\ % 612d2be661a6aafa76c5f8db7d019770462fe7f189198425c31a7652
  lingeling-ats1o1               & algotocnf\_27-t10-diff\_ocnf.cnf & none       & yes   & 134722 / 728349 & 134722    & 0         &            & 0 \\ % f00fa51898d4efd7515e2a485ca1fffde25e13ae4b308aea86cac649
  lingeling-ats1o1               & algotocnf\_27-t10-simplified.cnf & none       & no    & 134722 / 728542 & 16823476404 & 184806557 &            & 89822 \\ % b722f2a9b443cb3e30cb054aace9ab54cd1ae8d59c780da222a7231d
  lingeling-ats1o1               & algotocnf\_27-t10-simplified.cnf & none       & yes   & 134722 / 728542 & 15834304998 & 188185237 &            & 89856 \\ % 1c663ac2e39a6a19591f4263a5b2252e4b04d02a690889fcbf522b97
  lingeling-ats1o1               & algotocnf\_27-t10-simplified.cnf & cmsat      & no    & 67175 / 392680 & 19313586484 & 183131032 &            & 89839 \\ % 0a1379d21e4ca5dc195f62df2637eda33904e921e66b4932ecbac61b
  lingeling-ats1o1               & algotocnf\_27-t10-simplified.cnf & cmsat      & yes   & 67175 / 392680 & 19132485681 & 186800227 &            & 89844 \\ % 3c4a420d6c4603e6c6da9a39d4e4c67b3b7ae4d612cfc43c856c775e
  lingeling-ats1o1               & algotocnf\_27-t10-simplified.cnf & minisat    & no    & 72909 / 460404 & 14222697283 & 183992472 &            & 89815 \\ % c924ba83b150c26dff8feb6dadcb58146f07375a18244ff5d75d5e02
  lingeling-ats1o1               & algotocnf\_27-t10-simplified.cnf & minisat    & yes   & 72909 / 460404 & 17237171315 & 189984455 &            & 89816 \\ % 61a32a3e9f04226cc5ec2745224fdd692552cd109f6e44cf1789b799
  lingeling-ats1o1               & algotocnf\_27-t10-simplified.cnf & satelite   & no    & 84442 / 544998 & 16088380688 & 180695961 &            & 89856 \\ % 5682e03785fbe2ca5ae777dbcb56b13d7342e274865fa0d276dcdaf2
  lingeling-ats1o1               & algotocnf\_27-t10-simplified.cnf & satelite   & yes   & 84442 / 544998 & 16591195824 & 186326350 &            & 89850 \\ % 5344f270c1ea66ea30d45f0389fa13eaa185c5f2bba41e94f769ac3a
  lingeling-ats1o1               & algotocnf\_27-t10-simplified\_diff-desc-sc\_ocnf.cnf & none       & no    & 243525 / 1472595 & 21979301660 & 872505185 &            & 89766 \\ % 8a28c380b6fe04dfea65925a762815623b4d98441a3a46872668061d
  lingeling-ats1o1               & algotocnf\_27-t10-simplified\_diff-desc-sc\_ocnf.cnf & none       & yes   & 243525 / 1472595 & 23865760536 & 810280857 &            & 89784 \\ % 3299e0bf8265af75e348d2943b1594159fb64a42063d9c91def5388f
  lingeling-ats1o1               & algotocnf\_27-t10-simplified\_diff-desc.cnf & none       & no    & 134722 / 1041847 & 17319765194 & 131235725 &            & 89803 \\ % 73b889346e6e0fcc7308cfff931d6c7574523ed828c8cd3cc68fccd1
  lingeling-ats1o1               & algotocnf\_27-t10-simplified\_diff-desc.cnf & none       & yes   & 134722 / 1041847 & 16028749104 & 138208540 &            & 89762 \\ % b1283c80d09caccd10699bea0a672b085e76f735541c643c682be5be
  lingeling-ats1o1               & algotocnf\_27-t10-simplified\_diff-desc.cnf & satelite   & no    & 100735 / 789589 & 18132759801 & 147300411 &            & 89863 \\ % d848eed3f01747e0b8a00ce00af2e431fc3d25746dd79347c18c2f48
  lingeling-ats1o1               & algotocnf\_27-t10-simplified\_diff-desc.cnf & satelite   & yes   & 100735 / 789589 & 16824683913 & 130788952 &            & 89829 \\ % 959d57a3a8d32cf2c6b9ea7dd91af7339925fae26e878ef194b10c1d
  lingeling-ats1o1               & algotocnf\_27-t10-simplified\_diff-desc\_ocnf.cnf & none       & no    & 134722 / 1041847 & 33968799275 & 2584200794 &            & 124722 \\ % 18234b98bda957422a78e05ff754fcb698874f86065f2198b720b8a2
  lingeling-ats1o1               & algotocnf\_27-t10-simplified\_diff-desc\_ocnf.cnf & none       & yes   & 134722 / 1041847 & 42088801397 & 5091590370 &            & 124723 \\ % ee6c6388921b6ffd3d58dcb9dac3dc96255d1f8a909206a2cb61e4f5
  lingeling-ats1o1               & algotocnf\_27-t10-simplified\_ocnf.cnf & none       & no    & 134722 / 728542 & 58684314796 & 47290646982 &            & 89799 \\ % 037e93a18c8d22fce930528d9b356b8d36b1f48938637019b10532f3
  lingeling-ats1o1               & algotocnf\_27-t10-simplified\_ocnf.cnf & none       & yes   & 134722 / 728542 & 54586733606 & 46857594500 &            & 89852 \\ % f98e5ff7fba1d67b546240c4e711fd427cd183c3eb45e8732fa69c70
    \end{tabular}
  \end{center}
\end{table}

%Text such that \LaTeX{} won't screw up.
\newpage

\begin{table}[p]
  \begin{center}
    \begin{tabular}{l|cccccccc}
        \textbf{SAT solver} & \textbf{CNF file} & \textbf{simplified} & \textbf{phase} & \textbf{vars/clauses} & \textbf{propagations} & \textbf{decisions} & \textbf{restarts} & \textbf{Runtime (sec)} \\
      \hline
  lingeling-ats1o1               & algotocnf\_27-t10.cnf          & none       & no    & 134722 / 728030 & 19316277585 & 177166537 &            & 89851 \\ % 84c16114ae2d1eeedcc03c75dcb99abc7f122b92ce8e201955ed5cd5
  lingeling-ats1o1               & algotocnf\_27-t10.cnf          & none       & yes   & 134722 / 728030 & 23955661635 & 161913123 &            & 89774 \\ % a620f00bcb33c4c8224eac17e39103aeeb35cd74a721abb6ef6d9490
  lingeling-ats1o1               & algotocnf\_27-t10.cnf          & cmsat      & no    & 72044 / 411276 & 28870600750 & 126791459 &            & 89821 \\ % 157828d9b23a5ab08266162a08a141f0673bda9f1eead8267cd2f2d6
  lingeling-ats1o1               & algotocnf\_27-t10.cnf          & cmsat      & yes   & 72044 / 411276 & 19814684216 & 183355756 &            & 89814 \\ % fc989285fcbd7a2277fcc0d038fec9443262b47526ef82f29820c0ae
  lingeling-ats1o1               & algotocnf\_27-t10.cnf          & minisat    & no    & 73897 / 465248 & 22961923695 & 123420917 &            & 89847 \\ % e5ac426823d16c0c07fe8bfa59ddfb33858f697bb4cb14a11a93e23b
  lingeling-ats1o1               & algotocnf\_27-t10.cnf          & minisat    & yes   & 73897 / 465248 & 14935831688 & 171869387 &            & 89819 \\ % e3c56cd19c8a30d4bbecd40bee3d7805c5e8a2e27574c13b080ef365
  lingeling-ats1o1               & algotocnf\_27-t10\_diff-desc-sc\_ocnf.cnf & none       & no    & 243525 / 1472083 & 28578985339 & 923919864 &            & 124721 \\ % f9356b714e9eecf4c2ee7e9cecbc769f0b5e9ec27ccdf853b62d7f91
  lingeling-ats1o1               & algotocnf\_27-t10\_diff-desc-sc\_ocnf.cnf & none       & yes   & 243525 / 1472083 & 29827989620 & 629722196 &            & 124741 \\ % 3efdedb0a43cc35a445425efa9b35b5eeaf5d5d8e5e83c650d3285ea
  lingeling-ats1o1               & algotocnf\_27-t10\_diff-desc.cnf & none       & no    & 134722 / 1041335 & 16515872973 & 140698747 &            & 89759 \\ % fe500612f91dd73be6cb0a21c24d5de07dab8ece6169671b35b8c5e9
  lingeling-ats1o1               & algotocnf\_27-t10\_diff-desc.cnf & none       & yes   & 134722 / 1041335 & 15963674533 & 132056886 &            & 89799 \\ % c6b5178a2ad3bed0c149acaf3aff21c0aafb05fa28001c157684c5f0
  lingeling-ats1o1               & algotocnf\_27-t10\_diff-desc.cnf & cmsat      & no    & 57956 / 311219 & 17952660151 & 125912589 &            & 89863 \\ % 91c457bf08672235b43651a14b8ef1ef9aa1ec3eef55a18d0c90790f
  lingeling-ats1o1               & algotocnf\_27-t10\_diff-desc.cnf & cmsat      & yes   & 57956 / 311219 & 19174829873 & 134061138 &            & 89856 \\ % 7edf45237255a1822cad396ea6722495e09c829f9b79f6178c0eebeb
  lingeling-ats1o1               & algotocnf\_27-t10\_diff-desc.cnf & minisat    & no    & 94850 / 619480 & 17827952175 & 156209895 &            & 89821 \\ % d068223e4bdab5c70f4985290de13af82a8c82b05c96c502a0f1c713
  lingeling-ats1o1               & algotocnf\_27-t10\_diff-desc.cnf & satelite   & no    & 105733 / 857675 & 15990560517 & 133885162 &            & 89846 \\ % a031963a438ea11d4ba4d2187ca2c0a3026409789e42e8a478c4ca43
  lingeling-ats1o1               & algotocnf\_27-t10\_diff-desc.cnf & satelite   & yes   & 105733 / 857675 & 15042527283 & 135707358 &            & 89829 \\ % bbb17e93f272e7ae20f94cc0c2ec06a74fa5693f30f25eb014c7efec
  lingeling-ats1o1               & algotocnf\_27-t10\_diff-desc\_ocnf.cnf & none       & no    & 134722 / 1041335 & 40679378533 & 4208531765 &            & 124778 \\ % 8caeaa8a32e315a009c6b583ba1d8405eaebe3ed4617eab6c942d18d
  lingeling-ats1o1               & algotocnf\_27-t10\_diff-desc\_ocnf.cnf & none       & yes   & 134722 / 1041335 & 34132366227 & 2449401414 &            & 124776 \\ % 63dcd5cbf9901265a4bf7785b178e069465732e093e003c1be402e2a
  lingeling-ats1o1               & algotocnf\_27-t10\_ocnf.cnf    & none       & no    & 134722 / 728030 & 43845075238 & 36403442138 &            & 89817 \\ % d399697d63738ee1071a65746cd1d1849086b1db62c834fbe1a78261
  lingeling-ats1o1               & algotocnf\_27-t10\_ocnf.cnf    & none       & yes   & 134722 / 728030 & 45133105805 & 32513726277 &            & 89841 \\ % ca25667d414801a1d6dd197e2ac5e8e23ba5d784cedc75d0d44bd5f4
  lingeling-ats1o1               & algotocnf\_2\_diff-desc-sc\_ocnf.cnf & none       & no    & 85760 / 522370 & 40294556  & 4535019   &            & 52 \\ % e5554a3c24afe50e8343d4c34ea69da4e8d82efbbaab5b7e668ca3ca
  lingeling-ats1o1               & algotocnf\_2\_diff-desc-sc\_ocnf.cnf & none       & yes   & 85760 / 522370 & 28905439  & 3807250   &            & 47 \\ % 70987cd6ba5167dffd8493b859f9e2359f2dfadde2457a34bdabbd9f
  lingeling-ats1o1               & algotocnf\_2\_diff-desc.cnf    & none       & no    & 48704 / 374146 & 23547825  & 1386644   &            & 82 \\ % ee255979807d8301fb3f6d564028249e4a800078b2ca2d351c2a0716
  lingeling-ats1o1               & algotocnf\_2\_diff-desc.cnf    & none       & yes   & 48704 / 374146 & 77176485  & 2287767   &            & 121 \\ % 37a4bca2f1d9f8b725f31479fcefe4be76c2c37baa751b498d442866
  lingeling-ats1o1               & algotocnf\_2\_diff-desc.cnf    & cmsat      & no    & 15342 / 72870 & 86049890  & 2520691   &            & 101 \\ % 691a3608b99f0d5804a71ce6f5f62c38b39ac14892b596cd9f4cefc1
  lingeling-ats1o1               & algotocnf\_2\_diff-desc.cnf    & cmsat      & yes   & 15342 / 72870 & 160131421 & 3787880   &            & 84 \\ % cb177dda1a9d4bb87571ccdffa075b929d8555c4b62c8a15f7160bdf
  lingeling-ats1o1               & algotocnf\_2\_diff-desc.cnf    & lingeling  & no    & 40767 / 106546 & 53660381  & 2385632   &            & 99 \\ % d1ea40f476619f37236d5ea3c30ee83e911337bf8c2c88ad8ace4d2d
  lingeling-ats1o1               & algotocnf\_2\_diff-desc.cnf    & lingeling  & yes   & 40767 / 106546 & 98751058  & 3428241   &            & 101 \\ % a513e083357fe7a83bee18c8b7f274c7df0b75452de7deda023ff09f
  lingeling-ats1o1               & algotocnf\_2\_diff-desc.cnf    & minisat    & no    & 23531 / 149436 & 16434516  & 974640    &            & 48 \\ % 6414fcdc86d12a46843ff7122f7b0c7d70718cd93ce255612c2d9884
  lingeling-ats1o1               & algotocnf\_2\_diff-desc.cnf    & minisat    & yes   & 23531 / 149436 & 53164776  & 1855578   &            & 71 \\ % d272826c82069f64a316fb7650af6b659f430ac6422de67b13936f44
  lingeling-ats1o1               & algotocnf\_2\_diff-desc.cnf    & satelite   & no    & 29428 / 204577 & 80790609  & 2273744   &            & 83 \\ % 45f2b60348c4bf695e957f0c46d089ae4ee73bae8f9fa0079fb7d8d2
  lingeling-ats1o1               & algotocnf\_2\_diff-desc.cnf    & satelite   & yes   & 29428 / 204577 & 176766056 & 3019784   &            & 95 \\ % 32b5664cac46746fd12a571f0863a320d16e5a66878d37aedc7a43db
  lingeling-ats1o1               & algotocnf\_2\_diff-desc\_ocnf.cnf & none       & no    & 48704 / 374146 & 135861289 & 9542091   &            & 129 \\ % 0bdb762155af09441de73894d465cdd3b084efd05c38ca3c7174f222
  lingeling-ats1o1               & algotocnf\_2\_diff-desc\_ocnf.cnf & none       & yes   & 48704 / 374146 & 232012337 & 14044684  &            & 136 \\ % 25e6f33f97b515296f3cad2993e3437e60592207f2b281bcc8306067
  lingeling-ats1o1               & algotocnf\_2\_ocnf.cnf         & none       & no    & 48704 / 254210 & 266721567 & 8861947   &            & 505 \\ % e6158b44a1a8d7f14c81c3a54bd8caf345c08da3b4e57fce6d89d749
  lingeling-ats1o1               & algotocnf\_2\_ocnf.cnf         & none       & yes   & 48704 / 254210 & 118917460 & 5517576   &            & 588 \\ % 3569bc5aacde5ee63c3b08c00b0fd93e5a34b7184ec04acf288fb9a6
    \end{tabular}
  \end{center}
\end{table}

%Text such that \LaTeX{} won't screw up.
\newpage

\begin{table}[p]
  \begin{center}
    \begin{tabular}{l|cccccccc}
        \textbf{SAT solver} & \textbf{CNF file} & \textbf{simplified} & \textbf{phase} & \textbf{vars/clauses} & \textbf{propagations} & \textbf{decisions} & \textbf{restarts} & \textbf{Runtime (sec)} \\
      \hline
  lingeling-ats1o1               & algotocnf\_3.cnf               & none       & no    & 48704 / 254656 & 13847781  & 1126960   &            & 20 \\ % f7aab7c2cb2431c31174b240a1691b0d4103979d10f4f69fe9395438
  lingeling-ats1o1               & algotocnf\_3.cnf               & none       & yes   & 48704 / 254656 & 13005974  & 1020287   &            & 22 \\ % 0170d525d0e8a8a6841274daf4e249f486aeb1beb348b954c3c61801
  lingeling-ats1o1               & algotocnf\_3.cnf               & cmsat      & no    & 23893 / 111403 & 25817726  & 1503571   &            & 33 \\ % f74207519eb84c21a47705d580593784e3aeb8e8899bc12a8c813457
  lingeling-ats1o1               & algotocnf\_3.cnf               & cmsat      & yes   & 23893 / 111403 & 27298128  & 1693132   &            & 32 \\ % cc5a14a1bc1497dd1f11fb9550634dfad88c97206df37aa91a661852
  lingeling-ats1o1               & algotocnf\_3.cnf               & lingeling  & no    & 44000 / 100941 & 10750713  & 942145    &            & 9 \\ % cc37dace3fd720c454127a03e34100ea8c74e5afab2fc9a97e231a64
  lingeling-ats1o1               & algotocnf\_3.cnf               & lingeling  & yes   & 44000 / 100941 & 9981891   & 906171    &            & 10 \\ % 317da1adc678aec213f0651041249ae1d886621145544139f63a654c
  lingeling-ats1o1               & algotocnf\_3.cnf               & minisat    & no    & 20895 / 119580 & 11745732  & 985232    &            & 23 \\ % 0f47adcc64ea3b2108fd2de1ed2ccb2071795095403ba4a8277dc92f
  lingeling-ats1o1               & algotocnf\_3.cnf               & minisat    & yes   & 20895 / 119580 & 10246533  & 903649    &            & 20 \\ % d330aa06d09afea93676997f4efd6555e5d72b99d297b16d97681778
  lingeling-ats1o1               & algotocnf\_3.cnf               & satelite   & no    & 28647 / 156910 & 17912734  & 1389672   &            & 21 \\ % 2c48fd442f5ba83d0efcd67b18b3138e12694494914c65c3e5fae19f
  lingeling-ats1o1               & algotocnf\_3.cnf               & satelite   & yes   & 28647 / 156910 & 15774948  & 1301568   &            & 21 \\ % 9065fa3a8997f91c392992f21e17cac8122b6d3779c9dc329213025a
  lingeling-ats1o1               & algotocnf\_3\_diff-desc-sc\_ocnf.cnf & none       & no    & 85760 / 522816 & 3147770   & 235518    &            & 8 \\ % 4686c87895830e73fbd1fc4054b5d7ca96f0c8208d0fd905db81396a
  lingeling-ats1o1               & algotocnf\_3\_diff-desc-sc\_ocnf.cnf & none       & yes   & 85760 / 522816 & 3776333   & 316002    &            & 8 \\ % ff0fcf4ab21a29e6cf488a6b44297e3bda057bf1eea25b48c0415a81
  lingeling-ats1o1               & algotocnf\_3\_diff-desc.cnf    & none       & no    & 48704 / 374592 & 5210253   & 416999    &            & 7 \\ % 60659781b8af6f4a1c04215614d5f8745789c82fd30b075c7734ef9a
  lingeling-ats1o1               & algotocnf\_3\_diff-desc.cnf    & none       & yes   & 48704 / 374592 & 3443782   & 267614    &            & 8 \\ % 1415cc2b8435aa9f0850bb49c47d03f1af13db4b10c8208cd4f76b20
  lingeling-ats1o1               & algotocnf\_3\_diff-desc\_ocnf.cnf & none       & no    & 48704 / 374592 & 6295677   & 467773    &            & 12 \\ % 90a06d62666a251679655dc5096d3303a71548c7da07431ac665480a
  lingeling-ats1o1               & algotocnf\_3\_diff-desc\_ocnf.cnf & none       & yes   & 48704 / 374592 & 5561367   & 417632    &            & 10 \\ % 631054bf34ee5943a3a5c5095537fa051cb357ac4fcfbaa52cb13349
  lingeling-ats1o1               & algotocnf\_3\_ocnf.cnf         & none       & no    & 48704 / 254656 & 13788432  & 1131169   &            & 24 \\ % 86cd1a4c8878dd8944682aa977e587197776c3db8bceb079bf5ae0da
  lingeling-ats1o1               & algotocnf\_3\_ocnf.cnf         & none       & yes   & 48704 / 254656 & 12247565  & 975269    &            & 22 \\ % df55460b22362f14465357599568b2d93fdaf9ac0d0dc1e893eff6bf
  lingeling-ats1o2               & algotocnf\_18-t9\_diff-desc\_ocnf.cnf & none       & no    & 107839 / 846487 & 291996044 & 39949235  &            & 95830 \\ % f305760b6ff6820f902403208d0a33bcce7d5013a02fcfc15749cc03
  lingeling-ats1o2               & algotocnf\_1\_diff-desc\_ocnf.cnf & none       & no    & 48704 / 373920 & 1385390130 & 167057864 &            & 95278 \\ % 67b030e53a75dde0dc89e047bc7e9b46673bebea1ca9f120aaa1a629
  lingeling-ats1o2               & algotocnf\_1\_ocnf.cnf         & none       & no    & 48704 / 253984 & 223311823 & 99914039  &            & 95768 \\ % ab122ca3973326a61681602f18564eaf54f19f654dd6c671768f5bcf
  lingeling-ats1o2               & algotocnf\_21-t9.cnf           & none       & no    & 116800 / 636838 & 6567493736 & 154119871 &            & 35305 \\ % bc3d37d19130e2b1b91c3b56edf2496b4b45d237bf280842cad1631c
  lingeling-ats1o2               & algotocnf\_21-t9.cnf           & satelite   & no    & 74562 / 486714 & 9417657748 & 212696810 &            & 69722 \\ % 06dab86ac126635a5c3f89eb5e5c37da722abfbbe00cc35a1d18c487
  lingeling-ats1o2               & algotocnf\_21-t9\_diff-desc-sc\_ocnf.cnf & none       & no    & 212157 / 1289169 & 13886232120 & 98969798  &            & 95850 \\ % 0dbe02eb1459e4e8f77aa18d8c7887513cb27c22bce831bee2b5c272
  lingeling-ats1o2               & algotocnf\_21-t9\_diff-desc.cnf & lingeling  & no    & 104860 / 404126 & 2064556252 & 6943123   &            & 3323 \\ % 8ef70e6180df75d1e1e85c017b71e263df531842588a6bd1ce04273e
  lingeling-ats1o2               & algotocnf\_21-t9\_diff-desc.cnf & satelite   & no    & 90236 / 720755 & 2367713566 & 15624031  &            & 3678 \\ % eb592cc2a879478efb77f7af33259b27440c11435c5d002e338c5764
  lingeling-ats1o2               & algotocnf\_21-t9\_diff-desc\_ocnf.cnf & none       & no    & 116800 / 911629 & 3524212437 & 568511577 &            & 159894 \\ % b056c180c71e9079f4c8a4a1427572ca26fbd94619e26b46e8ba21dd
  lingeling-ats1o2               & algotocnf\_21-t9\_ocnf.cnf     & none       & no    & 116800 / 636838 & 126885002 & 40686508  &            & 95835 \\ % ac141b2964ee583dc577c3ec8a9ae8f75894d55f93d93769bf529043
  lingeling-ats1o2               & algotocnf\_23-t9.cnf           & satelite   & no    & 77835 / 506217 & 11185442968 & 235463767 &            & 95855 \\ % f45d1f7c875b1fedac096e49cda86ee39678dd88aff330606e923f85
  lingeling-ats1o2               & algotocnf\_23-t9\_diff-desc-sc\_ocnf.cnf & none       & no    & 222613 / 1350343 & 13598362021 & 106046634 &            & 95838 \\ % ae20b3813fbbd254a3c9539d48dff5093b8583d845b60351c695be54
  lingeling-ats1o2               & algotocnf\_23-t9\_diff-desc.cnf & satelite   & no    & 93504 / 740545 & 2975485032 & 20909550  &            & 4274 \\ % 60913c5e664e0831e6f109ff1ccf94bf88537684df15f187900565f8
  lingeling-ats1o2               & algotocnf\_24-t9.cnf           & none       & no    & 125761 / 682722 & 8197106091 & 147635256 &            & 75817 \\ % 671ff9fb3f95e5dddbe903e5ee7bc234b6c67c10f1871830678128cc
  lingeling-ats1o2               & algotocnf\_24-t9.cnf           & cmsat      & no    & 65579 / 377228 & 13872340130 & 218557892 &            & 94010 \\ % 7cde08f2686ca49798dda7351a29ebaf9273b43f08edeb247c391331
  lingeling-ats1o2               & algotocnf\_24-t9.cnf           & minisat    & no    & 69023 / 436292 & 10756062377 & 194157926 &            & 80231 \\ % c5daba2748dce1f0bd5b0da59d22a6fcc46dddb4f0f234d965d1736b
  lingeling-ats1o2               & algotocnf\_24-t9\_diff-desc-sc\_ocnf.cnf & none       & no    & 227841 / 1380914 & 15199857218 & 85385887  &            & 95795 \\ % de964373b18dfa37c07e30abe6a8da43d04e14b34ad9808a01b04426
    \end{tabular}
  \end{center}
\end{table}

%Text such that \LaTeX{} won't screw up.
\newpage

\begin{table}[p]
  \begin{center}
    \begin{tabular}{l|cccccccc}
        \textbf{SAT solver} & \textbf{CNF file} & \textbf{simplified} & \textbf{phase} & \textbf{vars/clauses} & \textbf{propagations} & \textbf{decisions} & \textbf{restarts} & \textbf{Runtime (sec)} \\
      \hline
  lingeling-ats1o2               & algotocnf\_24-t9\_diff-desc.cnf & cmsat      & no    & 56447 / 302014 & 1358176160 & 16585410  &            & 2419 \\ % 0ca2f6d17bfe51f40be2b88941f7fb9cfd67c7740123b1193aae3feb
  lingeling-ats1o2               & algotocnf\_24-t9\_diff-desc\_ocnf.cnf & none       & no    & 125761 / 976770 & 1182449031 & 200345630 &            & 95834 \\ % be51c97f87f1597c710fe8a84d51b72483ebc4bf876e84c00073b7cf
  lingeling-ats1o2               & algotocnf\_24-t9\_ocnf.cnf     & none       & no    & 125761 / 682722 & 165196086 & 69746929  &            & 160096 \\ % c607c8ed90a8eb6a1aac16159005030d6db666debf40b78994289a72
  lingeling-ats1o2               & algotocnf\_27-t10-diff.cnf     & none       & no    & 134722 / 728349 & 134722    & 0         &            & 0 \\ % 2051648f4701efe60b268cf83ce9d0510c0e45b9b3989b59b34dbc4c
  lingeling-ats1o2               & algotocnf\_27-t10-diff.cnf     & cmsat      & no    & 0 / 0      & 0         & 0         &            & 0 \\ % b85ac3a6b7f7f485877808379c2bdb9b09db4014f217db086820c4d3
  lingeling-ats1o2               & algotocnf\_27-t10-diff.cnf     & lingeling  & no    & 0 / 0      & 0         & 0         &            & 0 \\ % f9dd00b18b2944688c836dfbd701cd7de0771d71a5343f925602fe56
  lingeling-ats1o2               & algotocnf\_27-t10-diff.cnf     & satelite   & no    & 0 / 0      & 0         & 0         &            & 0 \\ % ffb225ee6e8fb50722cf492c176068046432110cb689f2dd934ab0c6
  lingeling-ats1o2               & algotocnf\_27-t10-diff\_diff-desc.cnf & none       & no    & 134722 / 1041654 & 134722    & 0         &            & 0 \\ % 14ae706ff857e624b6e96d93e06f0d7e6ca06146cfe33915da6eeda9
  lingeling-ats1o2               & algotocnf\_27-t10-diff\_diff-desc.cnf & cmsat      & no    & 0 / 0      & 0         & 0         &            & 0 \\ % f2f41514b9139fec8c9ce1c826657a1b3ad667d2ee1f7082c3202c86
  lingeling-ats1o2               & algotocnf\_27-t10-diff\_diff-desc.cnf & lingeling  & no    & 0 / 0      & 0         & 0         &            & 0 \\ % 7309c65f42c330f189d9064c89ec00b0d24be1919be3f1dcd115bbb0
  lingeling-ats1o2               & algotocnf\_27-t10-diff\_diff-desc.cnf & minisat    & no    & 0 / 0      & 0         & 0         &            & 0 \\ % 5575a12e7c6406fbaf8c647d793b53bb0efc89d55cdc78cf9f0a5630
  lingeling-ats1o2               & algotocnf\_27-t10-diff\_diff-desc.cnf & satelite   & no    & 0 / 0      & 0         & 0         &            & 0 \\ % fee950ed53297f0ca78e657c068ccb5acbcd4971c572d1157fd6fdee
  lingeling-ats1o2               & algotocnf\_27-t10-diff\_diff-desc\_ocnf.cnf & none       & no    & 134722 / 1041654 & 134722    & 0         &            & 0 \\ % ee25c13b6e8bf35945f1835f9164db39305c1e511c5445b97ee77d7c
  lingeling-ats1o2               & algotocnf\_27-t10-diff\_ocnf.cnf & none       & no    & 134722 / 728349 & 134722    & 0         &            & 0 \\ % 7612cf55883954a2565e505a2c466668d8ee41a74842e29e2249f48c
  lingeling-ats1o2               & algotocnf\_27-t10-simplified.cnf & none       & no    & 134722 / 728542 & 15631488002 & 178920666 &            & 133720 \\ % b4672580e30a3782c6e332c590aca54b3402a6e5733057e59a0c6b53
  lingeling-ats1o2               & algotocnf\_27-t10-simplified.cnf & cmsat      & no    & 67175 / 392680 & 17298246061 & 166455368 &            & 95842 \\ % 8931615973bda31234dbff26a1c2b2b52f5c0eae68b2b828e513d04c
  lingeling-ats1o2               & algotocnf\_27-t10-simplified.cnf & minisat    & no    & 72909 / 460404 & 14700588884 & 189460193 &            & 95730 \\ % 9e20fe23a2a78db5c73452a884ce247cf8ff585c8a0323ecd5ea5d3f
  lingeling-ats1o2               & algotocnf\_27-t10-simplified.cnf & satelite   & no    & 84442 / 544998 & 18770086979 & 208376737 &            & 160111 \\ % bcfbc05e027189f0809af743f55e17f4ab5204d1f574341d28bdf2ca
  lingeling-ats1o2               & algotocnf\_27-t10-simplified\_diff-desc-sc\_ocnf.cnf & none       & no    & 243525 / 1472595 & 16720604413 & 179707465 &            & 95635 \\ % 291859059e7328a2a71b83bf38d6d0db002e637ddc145fc758c262c8
  lingeling-ats1o2               & algotocnf\_27-t10-simplified\_diff-desc.cnf & cmsat      & no    & 56223 / 302443 & 17670976251 & 122737674 &            & 95849 \\ % df22c72d15a4fd4330e0488e5bb7b44e6663efd561e676d067cc5d30
  lingeling-ats1o2               & algotocnf\_27-t10-simplified\_diff-desc.cnf & satelite   & no    & 100735 / 789589 & 17634625706 & 143035794 &            & 95632 \\ % 998feaba7da152b8c57636765ac461cdcdaa23fe438f9fe21347742f
  lingeling-ats1o2               & algotocnf\_27-t10-simplified\_diff-desc\_ocnf.cnf & none       & no    & 134722 / 1041847 & 1542644428 & 258943393 &            & 95837 \\ % 05993b1c06f1d2ab8e8cbad264499d2656f4eb8b5dd2b69cc456a155
  lingeling-ats1o2               & algotocnf\_27-t10-simplified\_ocnf.cnf & none       & no    & 134722 / 728542 & 90264066  & 31405671  &            & 95828 \\ % 5f18a4ac667782f0f88ab9c416572eaffa6881283497aac19b53116a
  lingeling-ats1o2               & algotocnf\_27-t10.cnf          & none       & no    & 134722 / 728030 & 19410271191 & 178330077 &            & 95855 \\ % 49801bb9855de1c0a1146b201753bc38ffcc72d282a9c7e5328240e1
  lingeling-ats1o2               & algotocnf\_27-t10.cnf          & cmsat      & no    & 72044 / 411276 & 27820295377 & 123319542 &            & 95719 \\ % 164f095a94625b3e535358d355a1632d05da735f6c302fbe5c4309a1
  lingeling-ats1o2               & algotocnf\_27-t10\_diff-desc-sc\_ocnf.cnf & none       & no    & 243525 / 1472083 & 22329136399 & 118835758 &            & 160081 \\ % 55b2b06684400857a7c1caa4a5512075b10d461e5b1b794f2067eeda
  lingeling-ats1o2               & algotocnf\_27-t10\_diff-desc.cnf & none       & no    & 134722 / 1041335 & 22655826502 & 182697495 &            & 160066 \\ % 718b1126fd6fb2c6c293bfdea5dc50fdfd104a8af746c28bdf60bf20
  lingeling-ats1o2               & algotocnf\_27-t10\_diff-desc.cnf & cmsat      & no    & 57956 / 311219 & 20741562150 & 150381459 &            & 159965 \\ % ce1cff9db38ce3bd9c96ab65766a170e9704805d49f4dbf85eaadc1d
  lingeling-ats1o2               & algotocnf\_27-t10\_diff-desc.cnf & satelite   & no    & 105733 / 857675 & 12294423275 & 103385615 &            & 95796 \\ % 1c2e960f558c0d4507e4831c9e543969705dc6f8872a14d4dc378b8b
  lingeling-ats1o2               & algotocnf\_27-t10\_diff-desc\_ocnf.cnf & none       & no    & 134722 / 1041335 & 1852807780 & 287706832 &            & 95815 \\ % 3d18808b2be742831763cf5253fd8e01eee906646edb824d9b7a4450
  lingeling-ats1o2               & algotocnf\_27-t10\_ocnf.cnf    & none       & no    & 134722 / 728030 & 457892629 & 260401799 &            & 95763 \\ % c70aa3721cb87c1857a67a4e91bcd735bb78cc2f549094559422d996
  lingeling-ats1o2               & algotocnf\_2\_diff-desc\_ocnf.cnf & none       & no    & 48704 / 374146 & 4135680082 & 309203481 &            & 160067 \\ % 648d8b4b8d0e3e30ca44a868c1cf65160c65c0ba8755a8b7f5c2548e
  lingeling-ats1o2               & algotocnf\_2\_ocnf.cnf         & none       & no    & 48704 / 254210 & 500067111 & 190079058 &            & 159187 \\ % a584ea1b393aafc03907bfe29a72f629a090623d4a1ea520fc6869b8
  lingeling-ats1o2               & algotocnf\_3\_ocnf.cnf         & none       & no    & 48704 / 254656 & 609087535 & 105726046 &            & 154129 \\ % cb25097966daca729b51bded2437a3b07faf8f69f89cbc251e765009
  lingeling-ats1o4               & algotocnf\_18-t9\_diff-desc-sc\_ocnf.cnf & none       & no    & 196473 / 1197423 & 12277603197 & 680249279 &            & 95849 \\ % 18d67c55aab1c93b6a1beee2526ad125663ff66b2ef5fac3e22ad5df
    \end{tabular}
  \end{center}
\end{table}

%Text such that \LaTeX{} won't screw up.
\newpage

\begin{table}[p]
  \begin{center}
    \begin{tabular}{l|cccccccc}
        \textbf{SAT solver} & \textbf{CNF file} & \textbf{simplified} & \textbf{phase} & \textbf{vars/clauses} & \textbf{propagations} & \textbf{decisions} & \textbf{restarts} & \textbf{Runtime (sec)} \\
      \hline
  lingeling-ats1o4               & algotocnf\_18-t9\_ocnf.cnf     & none       & no    & 107839 / 590953 & 498021629 & 227303444 &            & 160120 \\ % f062f37c192f658930e44d9ea3e8c3dce8d7a0f6f8542d9a6ce43aab
  lingeling-ats1o4               & algotocnf\_1\_diff-desc\_ocnf.cnf & none       & no    & 48704 / 373920 & 1195216194 & 147205149 &            & 95836 \\ % d621970960f4baa856a3a044ac60e1b139807e7adc2ea11778fdc531
  lingeling-ats1o4               & algotocnf\_1\_ocnf.cnf         & none       & no    & 48704 / 253984 & 203583051 & 92099610  &            & 95691 \\ % a783bbdb641a5fe0738fd4f67f377aae6fce410ef3e73f1a78f16f51
  lingeling-ats1o4               & algotocnf\_21-t9.cnf           & cmsat      & no    & 63555 / 358308 & 8081065262 & 128186673 &            & 45547 \\ % 67bc14e5c33de78e4fc6e6ce1ff42480ac3f87e186fbfbe3c9e63230
  lingeling-ats1o4               & algotocnf\_21-t9.cnf           & satelite   & no    & 74562 / 486714 & 9417657748 & 212696810 &            & 55627 \\ % c22aa78695b39f64a9b285982bdcc6a0a3376552443aec36e088d135
  lingeling-ats1o4               & algotocnf\_21-t9\_diff-desc-sc\_ocnf.cnf & none       & no    & 212157 / 1289169 & 15965407992 & 100295965 &            & 95842 \\ % 8af617c121e11ff03f25fd1d347576b3bfab54c2ec42cef6d4034e0c
  lingeling-ats1o4               & algotocnf\_21-t9\_diff-desc.cnf & none       & no    & 116800 / 911629 & 3799761108 & 15913739  &            & 3110 \\ % 748d4400d7ae068bdde31adf8df40ed60ffb8f337304e7d01c99ab05
  lingeling-ats1o4               & algotocnf\_21-t9\_diff-desc.cnf & cmsat      & no    & 59528 / 318049 & 768486697 & 12516663  &            & 1533 \\ % 677698045e2bac418e4feb0a6a20941de21b73d2df23e06a6e6249f6
  lingeling-ats1o4               & algotocnf\_21-t9\_diff-desc.cnf & lingeling  & no    & 104860 / 404126 & 2064556252 & 6943123   &            & 3631 \\ % 91729a3f342b8d0176e27bcabb008a2d234863103a3f217450953e93
  lingeling-ats1o4               & algotocnf\_21-t9\_diff-desc\_ocnf.cnf & none       & no    & 116800 / 911629 & 2491633965 & 420970693 &            & 159883 \\ % d60e67951a3a6f88fb5a1bec1bb0d0cd341fe9872513d18934d52427
  lingeling-ats1o4               & algotocnf\_21-t9\_ocnf.cnf     & none       & no    & 116800 / 636838 & 143086967 & 48311472  &            & 95857 \\ % 68d651215da81090431179087ff9c3ad73bdf8df080c4701c4dafaf0
  lingeling-ats1o4               & algotocnf\_23-t9.cnf           & satelite   & no    & 77835 / 506217 & 10476896074 & 202355468 &            & 95840 \\ % dafdaa7412b36a2feb95b7858c2470ad54b722bd246e1225553ae8ba
  lingeling-ats1o4               & algotocnf\_23-t9\_diff-desc-sc\_ocnf.cnf & none       & no    & 222613 / 1350343 & 9636415989 & 96343356  &            & 95820 \\ % 35205e84dfcad09e8cfcb5961e987c969724ebed5622db2d3d51dc5f
  lingeling-ats1o4               & algotocnf\_23-t9\_diff-desc.cnf & satelite   & no    & 93504 / 740545 & 2543789170 & 24088216  &            & 5010 \\ % 50542e02358b6539181cc84b4b127946f125a9c397753973f19bea6b
  lingeling-ats1o4               & algotocnf\_23-t9\_diff-desc\_ocnf.cnf & none       & no    & 122774 / 955067 & 2515872107 & 415182711 &            & 95680 \\ % 52ff1ae40ad7a170ac7aaeca181e05d612100961c639a2ff98675f6d
  lingeling-ats1o4               & algotocnf\_23-t9\_ocnf.cnf     & none       & no    & 122774 / 667438 & 89997096  & 28131959  &            & 95676 \\ % 99507b43dfa2c25668f1304446c6a595fe918bb2d839b6e4980ca6aa
  lingeling-ats1o4               & algotocnf\_24-t9.cnf           & none       & no    & 125761 / 682722 & 8197106091 & 147635256 &            & 66665 \\ % 734ff1fed79e9eabafdcfdcac9d76c1ae8c2d5501e2c039490b61bea
  lingeling-ats1o4               & algotocnf\_24-t9.cnf           & cmsat      & no    & 65579 / 377228 & 15021825301 & 237337855 &            & 93952 \\ % 4df6930a52f1ad9aa846c297caa8df02bda203d391f6660b68464d6d
  lingeling-ats1o4               & algotocnf\_24-t9.cnf           & minisat    & no    & 69023 / 436292 & 10756062377 & 194157926 &            & 97562 \\ % 4a04d6ee33f036c0c43a0354f5fdc657df1ea2f908bec682e3be54b4
  lingeling-ats1o4               & algotocnf\_24-t9\_diff-desc-sc\_ocnf.cnf & none       & no    & 227841 / 1380914 & 11689416476 & 84053166  &            & 95831 \\ % 9c588f5e181a9d5876884a3bddbc69f54bbe7ccabf7520a63d747e03
  lingeling-ats1o4               & algotocnf\_24-t9\_diff-desc.cnf & cmsat      & no    & 56447 / 302014 & 1358176160 & 16585410  &            & 2497 \\ % 55f419f802b0b4bf8ed76ec950ebfc930ab83682ae96b3483ea3d8ce
  lingeling-ats1o4               & algotocnf\_24-t9\_diff-desc\_ocnf.cnf & none       & no    & 125761 / 976770 & 3442060193 & 558496632 &            & 159881 \\ % 75c17a5eb14bbc8254f2366736715bf340a28cff50512c275fc39f30
  lingeling-ats1o4               & algotocnf\_24-t9\_ocnf.cnf     & none       & no    & 125761 / 682722 & 364438349 & 143978605 &            & 95825 \\ % 19bd1e333503ad1a4f70444c19cbbaf10f0ea2b5e3d83feef9d7a7c8
  lingeling-ats1o4               & algotocnf\_27-t10-diff.cnf     & none       & no    & 134722 / 728349 & 134722    & 0         &            & 0 \\ % ddb24161db9b7d183c35a9a1544e37fbe9e02090d18d968352966bc6
  lingeling-ats1o4               & algotocnf\_27-t10-diff.cnf     & cmsat      & no    & 0 / 0      & 0         & 0         &            & 0 \\ % 6b992cd2e79581aeb82de311e0f48e6a1861dc422fab71519d8c9f3d
  lingeling-ats1o4               & algotocnf\_27-t10-diff.cnf     & lingeling  & no    & 0 / 0      & 0         & 0         &            & 0 \\ % fc19bb728fcb5a992f928283e44111e56964cfccc4246b0b8b1b830d
  lingeling-ats1o4               & algotocnf\_27-t10-diff.cnf     & satelite   & no    & 0 / 0      & 0         & 0         &            & 0 \\ % a1fa870e089f16eff41b3eac951c5a96e155c7e3e8abf667f87ec0be
  lingeling-ats1o4               & algotocnf\_27-t10-diff\_diff-desc-sc\_ocnf.cnf & none       & no    & 243525 / 1472402 & 243525    & 0         &            & 0 \\ % 90895e6d588831415b91e82bebe662cf5489748f1d919bfefea0913f
  lingeling-ats1o4               & algotocnf\_27-t10-diff\_diff-desc.cnf & none       & no    & 134722 / 1041654 & 134722    & 0         &            & 0 \\ % 5d053d805e0c817456913e9b3f0cdedcbce07cf4725e915e9f0df34d
  lingeling-ats1o4               & algotocnf\_27-t10-diff\_diff-desc.cnf & cmsat      & no    & 0 / 0      & 0         & 0         &            & 0 \\ % 5ddf6031ca47b7575cc21ea927df53980c0a72fc0f125d19503c43ac
  lingeling-ats1o4               & algotocnf\_27-t10-diff\_diff-desc.cnf & lingeling  & no    & 0 / 0      & 0         & 0         &            & 0 \\ % 61d92044243a8468e1aec2711ba0cb54e723ec63e22fc36ba5ea177f
  lingeling-ats1o4               & algotocnf\_27-t10-diff\_diff-desc.cnf & minisat    & no    & 0 / 0      & 0         & 0         &            & 0 \\ % be39d1c3361283e156ff05d953a2913db10e28da893b1ddd576aeefd
  lingeling-ats1o4               & algotocnf\_27-t10-diff\_diff-desc.cnf & satelite   & no    & 0 / 0      & 0         & 0         &            & 0 \\ % e89a8139eee959ce046e9fc84741ce150e0991778d61cb3fb2b3a26e
  lingeling-ats1o4               & algotocnf\_27-t10-diff\_diff-desc\_ocnf.cnf & none       & no    & 134722 / 1041654 & 134722    & 0         &            & 0 \\ % 075b5df765e231dc8db552eca444243940b7b54289bf145bdcb984d9
  lingeling-ats1o4               & algotocnf\_27-t10-diff\_ocnf.cnf & none       & no    & 134722 / 728349 & 134722    & 0         &            & 0 \\ % ab8cecb87df7373ff492d18c50dd4a3f37c4ac5684eb1809fbc67b32
    \end{tabular}
  \end{center}
\end{table}

%Text such that \LaTeX{} won't screw up.
\newpage

\begin{table}[p]
  \begin{center}
    \begin{tabular}{l|cccccccc}
        \textbf{SAT solver} & \textbf{CNF file} & \textbf{simplified} & \textbf{phase} & \textbf{vars/clauses} & \textbf{propagations} & \textbf{decisions} & \textbf{restarts} & \textbf{Runtime (sec)} \\
      \hline
  lingeling-ats1o4               & algotocnf\_27-t10-simplified.cnf & none       & no    & 134722 / 728542 & 14735940174 & 171726678 &            & 95833 \\ % 79ae0470e38bcc67aeac8ea96a751556692712b3d42b650bc7d9c2c7
  lingeling-ats1o4               & algotocnf\_27-t10-simplified.cnf & minisat    & no    & 72909 / 460404 & 10327638460 & 137249195 &            & 95447 \\ % 3c209df63d87050c21f12a08be6349c910395df744c3e20361f52fd0
  lingeling-ats1o4               & algotocnf\_27-t10-simplified.cnf & satelite   & no    & 84442 / 544998 & 20755182678 & 230484571 &            & 160006 \\ % af0c111e82bfa91cb60cfe811b994245780288f36f7967ed8caf8684
  lingeling-ats1o4               & algotocnf\_27-t10-simplified\_diff-desc-sc\_ocnf.cnf & none       & no    & 243525 / 1472595 & 13159666083 & 139061603 &            & 95862 \\ % bd2455a20d40ecf947cacedb8df0fd66192b29d8b3ac27ecc1b07f72
  lingeling-ats1o4               & algotocnf\_27-t10-simplified\_diff-desc.cnf & none       & no    & 134722 / 1041847 & 20840528886 & 162794536 &            & 160035 \\ % a031144abebfe28ea62cce124a1f70e68aa98b93576ecbe67ee96cec
  lingeling-ats1o4               & algotocnf\_27-t10-simplified\_diff-desc.cnf & satelite   & no    & 100735 / 789589 & 16754672686 & 136890301 &            & 95620 \\ % 0053b998bd0d40a0b866cedf4b001d8c57d9e0955bac08c5d21acac2
  lingeling-ats1o4               & algotocnf\_27-t10-simplified\_diff-desc\_ocnf.cnf & none       & no    & 134722 / 1041847 & 1554191630 & 261007191 &            & 95793 \\ % 5df1d76509e73bfad60f2549a78abc0d4f0b7361f9ae82e981cb1ba7
  lingeling-ats1o4               & algotocnf\_27-t10-simplified\_ocnf.cnf & none       & no    & 134722 / 728542 & 90784648  & 31646766  &            & 124460 \\ % 1d6a8c2457fcff84048d2777baba114b7c67730e212babd36fffa7d7
  lingeling-ats1o4               & algotocnf\_27-t10.cnf          & none       & no    & 134722 / 728030 & 21570758246 & 205922861 &            & 95730 \\ % a83759891dd7f7b59b1232e92391689118b5767d5b2fd97c264242a7
  lingeling-ats1o4               & algotocnf\_27-t10.cnf          & cmsat      & no    & 72044 / 411276 & 23322994944 & 104889093 &            & 95801 \\ % e3d50af143fb93b02a43b5ec917c999e07f4d429c317d2f2f385a00c
  lingeling-ats1o4               & algotocnf\_27-t10.cnf          & minisat    & no    & 73897 / 465248 & 25224394009 & 132548995 &            & 160008 \\ % a7c697a286b4e8ef67d03004f834f695e2f04e9fb00a6926a8f637e3
  lingeling-ats1o4               & algotocnf\_27-t10.cnf          & satelite   & no    & 85274 / 549734 & 16980330264 & 121022090 &            & 95829 \\ % baec08d0f3a2d63a04b7fe7a08b0e665e5b16a20b066040e53b9a5fa
  lingeling-ats1o4               & algotocnf\_27-t10\_diff-desc-sc\_ocnf.cnf & none       & no    & 243525 / 1472083 & 21470599930 & 117190875 &            & 159964 \\ % bbc569f58d16a152b351cfb3cb454b44e795d3dee0a1c78e38ce9281
  lingeling-ats1o4               & algotocnf\_27-t10\_diff-desc.cnf & cmsat      & no    & 57956 / 311219 & 23040861935 & 152804664 &            & 95806 \\ % befcaf26209d151f2fef6017b3a9826b17561748b29702571954dfe9
  lingeling-ats1o4               & algotocnf\_27-t10\_diff-desc.cnf & minisat    & no    & 94850 / 619480 & 13120004235 & 119040417 &            & 95844 \\ % fb0b45091c24ce4221d871f7488dbe05cd0cd7335c9eb9382c10e897
  lingeling-ats1o4               & algotocnf\_27-t10\_diff-desc.cnf & satelite   & no    & 105733 / 857675 & 21379439073 & 175877300 &            & 160035 \\ % 058fde864cab827615b01f68f97ce2c04a3d0411ff1955d722019cf6
  lingeling-ats1o4               & algotocnf\_27-t10\_diff-desc\_ocnf.cnf & none       & no    & 134722 / 1041335 & 3505169535 & 554888001 &            & 144013 \\ % ddeefd95643d884b598b9291c1c7495d73ac67da65c3825f98722a3b
  lingeling-ats1o4               & algotocnf\_27-t10\_ocnf.cnf    & none       & no    & 134722 / 728030 & 282468976 & 153223726 &            & 95639 \\ % 721004ecdc64dec809257abf6f23d10cd1263ff96addbacb22416596
  lingeling-ats1o4               & algotocnf\_3\_ocnf.cnf         & none       & no    & 48704 / 254656 & 956630637 & 90542534  &            & 127054 \\ % c3729f11d6e8efc3d4d69c82cc8dd29d218459b13b6489de0a59c728
  minisat                        & algotocnf\_1.cnf               & none       & no    & 48704 / 225124 & 1976338916 & 80207179  &            & 4519 \\ % b111d1f0e4fe8e84aa99a72b3688acebb950034e6519ffea6e4c2fa9
  minisat                        & algotocnf\_1.cnf               & none       & no    &            &           &           &            & timeout \\ % b90d9ea16b34c060029c3a152ad126a257a5a9bb3b7d75fa0117f017
  minisat                        & algotocnf\_1.cnf               & cmsat      & no    &            &           &           &            & timeout \\ % 1f46743f5042b6cb3fe13f25326848d5dbf3f18064d9b5d330e20676
  minisat                        & algotocnf\_1.cnf               & cmsat      & no    & 24503 / 111931 & 5412022717 & 174687681 &            & 7649 \\ % 86020c3420b0f034c669b316e88a864c3879c54c984bc19715ea0a43
  minisat                        & algotocnf\_1.cnf               & minisat    & no    & 20895 / 118236 & 1001195109 & 41539126  &            & 1476 \\ % 683e335675d12fc859a14cdd82bbff503909f86c8e48fcad920431c3
  minisat                        & algotocnf\_1.cnf               & minisat    & no    &            &           &           &            & timeout \\ % 9c12c78b3c20acea33fbfd278e5b8a61231aed68997d52e11700ef23
  minisat                        & algotocnf\_1.cnf               & satelite   & no    &            &           &           &            & timeout \\ % 6057ca956c327d61a8ef431f4e360fc2e227b19595b105462d794e5b
  minisat                        & algotocnf\_1.cnf               & satelite   & no    & 27495 / 153262 & 961077704 & 41599245  &            & 1293 \\ % 67ad2ccd4a360fc9d78682f3563afc205d1c00ad7a472f287941d810
  minisat                        & algotocnf\_18-t9.cnf           & none       & no    & 107839 / 582936 & 3813726   & 250759    &            & 3 \\ % 89e4260963e613a876af6ed95b560e435c003f4f3b429017a426c3b1
  minisat                        & algotocnf\_18-t9.cnf           & none       & no    &            &           &           &            & timeout \\ % 8a17ad4c211c2d11a8577c7e2e59b512be5fd3f5756eb8c0ab21292a
  minisat                        & algotocnf\_18-t9.cnf           & cmsat      & no    & 61789 / 342139 & 7216828   & 573598    &            & 6 \\ % 578786d50d531039730b6608eeb2286299e96d878469cc43da0a86a0
  minisat                        & algotocnf\_18-t9.cnf           & cmsat      & no    &            &           &           &            & timeout \\ % d1e1b322dc397fcec0346b697ff824080ae583ffcbf25a5cd9b2b965
  minisat                        & algotocnf\_18-t9.cnf           & lingeling  & no    &            &           &           &            & timeout \\ % 2dce86c3214ce3ebda7bd598901f68c9447b6df5b0febe311f1eee57
  minisat                        & algotocnf\_18-t9.cnf           & lingeling  & no    & 102616 / 344544 & 16428911  & 1003041   &            & 10 \\ % 3e1c23106d6c9ad8222fbc8260226114fae259d01c08da0859ccdb47
  minisat                        & algotocnf\_18-t9.cnf           & minisat    & no    & 61845 / 391134 & 11474979  & 666416    &            & 6 \\ % 76a2d3d9e7f3ab67ba31e7e0c1a181b8e6d3eb8522e27cad59db1830
  minisat                        & algotocnf\_18-t9.cnf           & minisat    & no    &            &           &           &            & timeout \\ % d3fcac2405180fb649d6a60b82fb2327c94fb8604bf4246d37698028
    \end{tabular}
  \end{center}
\end{table}

%Text such that \LaTeX{} won't screw up.
\newpage

\begin{table}[p]
  \begin{center}
    \begin{tabular}{l|cccccccc}
        \textbf{SAT solver} & \textbf{CNF file} & \textbf{simplified} & \textbf{phase} & \textbf{vars/clauses} & \textbf{propagations} & \textbf{decisions} & \textbf{restarts} & \textbf{Runtime (sec)} \\
      \hline
  minisat                        & algotocnf\_18-t9.cnf           & satelite   & no    & 69670 / 457972 & 17335955  & 893213    &            & 12 \\ % 786df8708ce15c252b84d89382ea461195f983b8f849fa84d6eadc37
  minisat                        & algotocnf\_18-t9.cnf           & satelite   & no    &            &           &           &            & timeout \\ % e99f569e3e44d3c3f4cac9f666eb159d1df0ca7f585b19cf7163f4d9
  minisat                        & algotocnf\_18-t9\_diff-desc.cnf & none       & no    &            &           &           &            & timeout \\ % 1a11d82482270afb3a2db0499bd7170f88d83a0f383ead6cca18bfd7
  minisat                        & algotocnf\_18-t9\_diff-desc.cnf & none       & no    & 107839 / 831263 & 18671534  & 857309    &            & 14 \\ % 93129fc6f90c919e5ae24ddf626abdada313c53c7bba89b9dca4b414
  minisat                        & algotocnf\_18-t9\_diff-desc.cnf & cmsat      & no    &            &           &           &            & timeout \\ % 253c11ef7f8808d5b84eb8a800f83fe90acb37009768b6cac7729b14
  minisat                        & algotocnf\_18-t9\_diff-desc.cnf & cmsat      & no    & 57114 / 303489 & 21547194  & 1001883   &            & 10 \\ % ccf0576d78e0dfea91174074920e217791ef705f57f1d90e2ce43a9f
  minisat                        & algotocnf\_18-t9\_diff-desc.cnf & lingeling  & no    & 99452 / 404666 & 7408845   & 609210    &            & 5 \\ % 0e8ce06c044cfda37ca0f79f3e7101153e3b3af9f092634ce815b971
  minisat                        & algotocnf\_18-t9\_diff-desc.cnf & lingeling  & no    &            &           &           &            & timeout \\ % 9d66d00b43fa5a460a3c6339877f84043106ad88e8f332d578894bbc
  minisat                        & algotocnf\_18-t9\_diff-desc.cnf & minisat    & no    & 78055 / 511431 & 6527792   & 383605    &            & 4 \\ % 1fe28ccef631613b02ce8c4e449707a21e3cc237cfae4b07b67cc0ca
  minisat                        & algotocnf\_18-t9\_diff-desc.cnf & minisat    & no    &            &           &           &            & timeout \\ % 40fc205af86e8cfb0bb311d0b9f82ea4f8e39710d4ae592cc0fb7e47
  minisat                        & algotocnf\_18-t9\_diff-desc.cnf & satelite   & no    &            &           &           &            & timeout \\ % 1cc49cdd366557b39b94f002b0c55f917859b4da045be42422138a31
  minisat                        & algotocnf\_18-t9\_diff-desc.cnf & satelite   & no    & 85386 / 692408 & 7396646   & 401013    &            & 6 \\ % c3511e9cdf39df9a85a7ebd38dfc442bd96fc04b135c258a7f7487eb
  minisat                        & algotocnf\_1\_diff-desc.cnf    & none       & no    &            &           &           &            & timeout \\ % 00c70735846675e14b9f1ff9af192107d696ed74159a7d08de191478
  minisat                        & algotocnf\_1\_diff-desc.cnf    & none       & no    & 48704 / 330692 & 145509993 & 4711702   &            & 206 \\ % 1e8692295c0b3bcb0c2bc34ff4974df900c9d54c2a20bbfb28c55f14
  minisat                        & algotocnf\_1\_diff-desc.cnf    & cmsat      & no    & 15402 / 73142 & 221480273 & 6830663   &            & 145 \\ % 2382dd49213836776a7f8f3b1984491c05c73274d74f559295bd3a4c
  minisat                        & algotocnf\_1\_diff-desc.cnf    & cmsat      & no    &            &           &           &            & timeout \\ % 320fc6a681ea0351c03f87158f8bd9095b92c69971b292bba5fad431
  minisat                        & algotocnf\_1\_diff-desc.cnf    & lingeling  & no    & 40554 / 101182 & 407709713 & 18029831  &            & 681 \\ % 268880ee0011c49c773d281fce9e1ecbb46dad8b87cec3855ac7ef6b
  minisat                        & algotocnf\_1\_diff-desc.cnf    & lingeling  & no    &            &           &           &            & timeout \\ % c35388b95ef610a233a0bfc3daf092affc4dd8845a2be1a3e2be5e83
  minisat                        & algotocnf\_1\_diff-desc.cnf    & minisat    & no    &            &           &           &            & timeout \\ % 4784cfcd9e7a602b72b8227cbb7d4a8a41527e42c254aa15640b62bf
  minisat                        & algotocnf\_1\_diff-desc.cnf    & minisat    & no    & 25107 / 158393 & 166787763 & 5399857   &            & 202 \\ % 6bbe3e9ea2e26f5faba3e0ba1de91e6907ea5e6b64bb85d7e6810d99
  minisat                        & algotocnf\_1\_diff-desc.cnf    & satelite   & no    & 30950 / 226602 & 267008243 & 8016584   &            & 449 \\ % 0fd97c7e693a667bc83ae5e66a4181b5ed434344a0e0e5b746c2faae
  minisat                        & algotocnf\_1\_diff-desc.cnf    & satelite   & no    &            &           &           &            & timeout \\ % 8ef0d753b6275b6e2a4da164323922d0b235b7e0b4578d976d8b5af5
  minisat                        & algotocnf\_2.cnf               & none       & no    &            &           &           &            & timeout \\ % 08ff7d31ddc848908386cab8ce349a0a9e46f48a03b6c88f03561f38
  minisat                        & algotocnf\_2.cnf               & none       & no    & 48704 / 224220 & 201587749 & 9459774   &            & 312 \\ % c26ee7b5fac811ead00bcb7ee154b09cc0f2861d04ca230b6c9495e0
  minisat                        & algotocnf\_2.cnf               & cmsat      & no    & 24323 / 111661 & 2750894175 & 90426293  &            & 5166 \\ % 0fcc820c0aac8b66bd3182cecf7af81725514325bcf4eafa39d91f19
  minisat                        & algotocnf\_2.cnf               & cmsat      & no    &            &           &           &            & timeout \\ % 5baf9d4eaf95782bda05cfcc6d375afdc4af208dfccb5e488b68555a
  minisat                        & algotocnf\_2.cnf               & minisat    & no    &            &           &           &            & timeout \\ % 60d6a006f8d5a280ab49788fcaa7af04116b2de7c7e666b5edb344d3
  minisat                        & algotocnf\_2.cnf               & minisat    & no    & 20895 / 118688 & 491695491 & 21320818  &            & 826 \\ % a8cef0396d6671cbf2ffc0e13f0f77a76fad653b67401509ab0c632f
  minisat                        & algotocnf\_2.cnf               & satelite   & no    & 27775 / 154276 & 1522804488 & 65886001  &            & 2541 \\ % 6b723b6f822e5440858d2e57e02696a090cf6976d1b012b84ce64917
  minisat                        & algotocnf\_2.cnf               & satelite   & no    &            &           &           &            & timeout \\ % 7f2d7654381a1725172cd77c292edf445a48734360f1d63490f93270
  minisat                        & algotocnf\_21-t9.cnf           & none       & no    &            &           &           &            & 86400 \\ % 312079e41ef302e652e0cb37c702a56ac6dcf53eeb700859e33e153c
  minisat                        & algotocnf\_21-t9.cnf           & none       & no    &            &           &           &            & timeout \\ % f4e0235347bde5c294485e5b8ba1a0d320a1613c5d85ce38d6e0c699
  minisat                        & algotocnf\_21-t9.cnf           & cmsat      & no    &            &           &           &            & timeout \\ % 61f5ee320a53cc38139f08975d5315ecabc1da17e8e4996ecfba5c13
  minisat                        & algotocnf\_21-t9.cnf           & cmsat      & no    &            &           &           &            & 86400 \\ % db8d1485bdbe6f79240d7405540038302204561b79cf6f8f9f88fac7
  minisat                        & algotocnf\_21-t9.cnf           & minisat    & no    &            &           &           &            & 86400 \\ % 6ac8ce3208d69a1942de42e158bdb93d16ad0793aa98e0c10e0db245
    \end{tabular}
  \end{center}
\end{table}

%Text such that \LaTeX{} won't screw up.
\newpage

\begin{table}[p]
  \begin{center}
    \begin{tabular}{l|cccccccc}
        \textbf{SAT solver} & \textbf{CNF file} & \textbf{simplified} & \textbf{phase} & \textbf{vars/clauses} & \textbf{propagations} & \textbf{decisions} & \textbf{restarts} & \textbf{Runtime (sec)} \\
      \hline
  minisat                        & algotocnf\_21-t9.cnf           & minisat    & no    &            &           &           &            & timeout \\ % b808f4cbdcf4c3d133409a9f5cab5c529b1078bdbd29d5eb6542e5da
  minisat                        & algotocnf\_21-t9.cnf           & satelite   & no    &            &           &           &            & 86400 \\ % 1a391c92f1997ffb91f4980af8cda2345de16c8f4d1c5bac2ce436e8
  minisat                        & algotocnf\_21-t9.cnf           & satelite   & no    &            &           &           &            & timeout \\ % 38432566a07052f645cecef28c8ca92666da8d04ad6bc4e96dea651c
  minisat                        & algotocnf\_21-t9\_diff-desc.cnf & none       & no    &            &           &           &            & timeout \\ % 7410d133ede3043f57ac4665a667d230d2d75908341ba5ed05eede16
  minisat                        & algotocnf\_21-t9\_diff-desc.cnf & none       & no    & 116800 / 894072 & 9147981007 & 108081289 &            & 62327 \\ % b7a3e6c28cb3a5efd64656d20a4a6f510c8ec0936f63ff1a1e98b776
  minisat                        & algotocnf\_21-t9\_diff-desc.cnf & cmsat      & no    & 59528 / 318049 & 577921212 & 12829996  &            & 2288 \\ % 9d9902a9b6f4a6bf1bf6304ff6d3933c8417a319a35387229a0638ce
  minisat                        & algotocnf\_21-t9\_diff-desc.cnf & cmsat      & no    &            &           &           &            & timeout \\ % e696e394f9aae94fc617e5f6eac5beb10cec2e412282556fd51d0305
  minisat                        & algotocnf\_21-t9\_diff-desc.cnf & lingeling  & no    & 104860 / 404126 & 1058635766 & 26665555  &            & 7096 \\ % cb3cf8e152717c38aaf8235c950d4198336088ae2fc0553a4ee5ff3c
  minisat                        & algotocnf\_21-t9\_diff-desc.cnf & lingeling  & no    &            &           &           &            & timeout \\ % dfcbeb8d8248e792fb2f41553ccfded32b5bc9ff9f3b57e0a24c35b8
  minisat                        & algotocnf\_21-t9\_diff-desc.cnf & minisat    & no    & 81650 / 534965 & 3082155154 & 42479717  &            & 17990 \\ % 57ac475183a99b616cc7102174d54bd6d8648d9fd05c12e316a2c1c4
  minisat                        & algotocnf\_21-t9\_diff-desc.cnf & minisat    & no    &            &           &           &            & timeout \\ % e07989b1f2dd6e4c17b5dda0ec48040c56b0acaa168fce9f15f6a5fb
  minisat                        & algotocnf\_21-t9\_diff-desc.cnf & satelite   & no    & 90236 / 720755 & 4261434766 & 51061640  &            & 23399 \\ % 400de9abefee9353a34ce5d0fe60ea6983b968797cc0122c512a588b
  minisat                        & algotocnf\_21-t9\_diff-desc.cnf & satelite   & no    &            &           &           &            & timeout \\ % 779b919f330b1e68b21c0be21626c05940f2a46401686958c5835b62
  minisat                        & algotocnf\_23-t9.cnf           & none       & no    &            &           &           &            & 86400 \\ % 1ff89a0aab68c8f82ded463175c239c1fe2e5f014071180bae29db96
  minisat                        & algotocnf\_23-t9.cnf           & none       & no    &            &           &           &            & timeout \\ % f3a8f3af1af97cdb9f8eee0c8c3ba27bbbf235c4ada5a4e0eaffeaf8
  minisat                        & algotocnf\_23-t9.cnf           & cmsat      & no    &            &           &           &            & 86400 \\ % 944b2e5e5eb013d240fd2880503d8a1c3b240c96e1f1c1d0b2a9da12
  minisat                        & algotocnf\_23-t9.cnf           & cmsat      & no    &            &           &           &            & timeout \\ % e3b7e3248ed2287abebb6c69211dad66de49dd8ccc2a6671d58c2c75
  minisat                        & algotocnf\_23-t9.cnf           & minisat    & no    &            &           &           &            & 86400 \\ % e2f0338c5340b49fea38967bf71b2eac388620fa2ea27978620e46c3
  minisat                        & algotocnf\_23-t9.cnf           & minisat    & no    &            &           &           &            & timeout \\ % ff169e69d3bb0911edb308fb0c53b731a60fdfa8ddb49d306cb96792
  minisat                        & algotocnf\_23-t9.cnf           & satelite   & no    &            &           &           &            & 86400 \\ % 9e3d0dea36a79b582be831796727e1f6d8edf16e682818fbb654a41e
  minisat                        & algotocnf\_23-t9.cnf           & satelite   & no    &            &           &           &            & timeout \\ % cec342019ffb6e03eb00e6ee34e0a2bcfafd7785af1de96aa7f508e7
  minisat                        & algotocnf\_23-t9\_diff-desc.cnf & none       & no    &            &           &           &            & 86400 \\ % d1a23253662cc4fa72cdfb4b0f7f7fd294f14695993f424437430163
  minisat                        & algotocnf\_23-t9\_diff-desc.cnf & none       & no    &            &           &           &            & timeout \\ % d2cbb9b0a16a78ed970d83795783f6c1352b5af27df731e83867e075
  minisat                        & algotocnf\_23-t9\_diff-desc.cnf & cmsat      & no    & 57510 / 307118 & 3433286956 & 54032360  &            & 16806 \\ % 1dc648903f8eba9544d829766e6c892e2ed242fbf1ca49a42d2cb534
  minisat                        & algotocnf\_23-t9\_diff-desc.cnf & cmsat      & no    &            &           &           &            & timeout \\ % 5f0b3d6e3fc814c022d60cc72aa5c2b00d38e1004969987526700f93
  minisat                        & algotocnf\_23-t9\_diff-desc.cnf & minisat    & no    &            &           &           &            & timeout \\ % a209f7d807d6e3f4e9a94533d6f6327a24a785a5a80923878563a72a
  minisat                        & algotocnf\_23-t9\_diff-desc.cnf & minisat    & no    &            &           &           &            & 86400 \\ % d3eefd42016324fd73afb2b5d72d9f5bb1f2b5d75d4e096e13144159
  minisat                        & algotocnf\_23-t9\_diff-desc.cnf & satelite   & no    &            &           &           &            & 86400 \\ % 6f612d03db28b7f88853b95db942e057526a4ba53cf72c1e954e4906
  minisat                        & algotocnf\_23-t9\_diff-desc.cnf & satelite   & no    &            &           &           &            & timeout \\ % b06a6977c463aed2d89c7e0fb0845b09f5d8ac84d68c0cb36e6ec6a1
  minisat                        & algotocnf\_24-t9.cnf           & none       & no    &            &           &           &            & 86400 \\ % a1991e684ce067a49365f3051c3b64319d03fc8a6c1b6029b38a8e75
  minisat                        & algotocnf\_24-t9.cnf           & none       & no    &            &           &           &            & timeout \\ % a2c5022a9f2a0db5c0d2135e60a781623100bc8e70ee420a53a78a60
  minisat                        & algotocnf\_24-t9.cnf           & cmsat      & no    &            &           &           &            & timeout \\ % 147506400bce2359fb63b3bb9bd4fd8a35ef9a1904160b0c6544029d
  minisat                        & algotocnf\_24-t9.cnf           & cmsat      & no    &            &           &           &            & 86400 \\ % 56d6915b0be4a537e71dffde86ffe5f81bdf9294813586f27d2eb420
  minisat                        & algotocnf\_24-t9.cnf           & minisat    & no    &            &           &           &            & 86400 \\ % 17b5dc6ab8866ba59dbc915374ffe45a17f22968b54a6b3a6eee0fd8
  minisat                        & algotocnf\_24-t9.cnf           & minisat    & no    &            &           &           &            & timeout \\ % 3e3c388c18841fc087b4d58a4f484a8d540cd7ba9160e37bfc8ff49e
    \end{tabular}
  \end{center}
\end{table}

%Text such that \LaTeX{} won't screw up.
\newpage

\begin{table}[p]
  \begin{center}
    \begin{tabular}{l|cccccccc}
        \textbf{SAT solver} & \textbf{CNF file} & \textbf{simplified} & \textbf{phase} & \textbf{vars/clauses} & \textbf{propagations} & \textbf{decisions} & \textbf{restarts} & \textbf{Runtime (sec)} \\
      \hline
  minisat                        & algotocnf\_24-t9.cnf           & satelite   & no    &            &           &           &            & timeout \\ % 397559546582b95b512e72f2077386e1db71fb9965784ef3b7f55f61
  minisat                        & algotocnf\_24-t9.cnf           & satelite   & no    &            &           &           &            & 86400 \\ % 7b202d0c8b8e8a4cf36003d27bea416fcdc04b337416772e7c75cef7
  minisat                        & algotocnf\_24-t9\_diff-desc.cnf & none       & no    &            &           &           &            & timeout \\ % 28cc689a02f8b4418c1762a9dffc9bdbab13e4073368bc35f09cfb61
  minisat                        & algotocnf\_24-t9\_diff-desc.cnf & none       & no    &            &           &           &            & 86400 \\ % 338442506f9639b6842c7c424e541acbd6dca07624c31c7552de4d5d
  minisat                        & algotocnf\_24-t9\_diff-desc.cnf & cmsat      & no    & 56447 / 302014 & 3471709121 & 54649404  &            & 11414 \\ % 00faa18eb503bcbabaf277f57da033005d466441fd5ec3885cd6c885
  minisat                        & algotocnf\_24-t9\_diff-desc.cnf & cmsat      & no    &            &           &           &            & timeout \\ % a418625d35d6f39120ddbeb3c5d7f2ec96ae553d98301dfe5fe893b8
  minisat                        & algotocnf\_24-t9\_diff-desc.cnf & minisat    & no    &            &           &           &            & timeout \\ % 17fa35ffce9bee471c6cf20aa86e4605e1b3b861afbff949941f109e
  minisat                        & algotocnf\_24-t9\_diff-desc.cnf & minisat    & no    &            &           &           &            & 86400 \\ % 215e56fdbad78cfca0ff4983fd02682685479adb64dc2d72be7efa23
  minisat                        & algotocnf\_24-t9\_diff-desc.cnf & satelite   & no    &            &           &           &            & 86400 \\ % 6b8a033dbbd57fc66a180deb93d9bae15bc04194db45c096de078999
  minisat                        & algotocnf\_24-t9\_diff-desc.cnf & satelite   & no    &            &           &           &            & timeout \\ % 989f5557e59104a71612622be2280de00c4c38b4f9552599c0697536
  minisat                        & algotocnf\_27-t10-diff.cnf     & none       & no    & 134722 / 8320 & 134722    & 1         &            & 0 \\ % 15d01747cbe3ce9609c0ff566a7819bedbefaa45d5490410e9a49997
  minisat                        & algotocnf\_27-t10-diff.cnf     & none       & no    &            &           &           &            & timeout \\ % cd779c9ba44fa42e6adbc5ef433c735f0d43135fa13819f569928bf4
  minisat                        & algotocnf\_27-t10-diff.cnf     & cmsat      & no    &            &           &           &            & timeout \\ % b08f30654966050490459668b15f30a93cb940d59a8e8b6a9f2b1329
  minisat                        & algotocnf\_27-t10-diff.cnf     & cmsat      & no    & 0 / 0      & 0         & 1         &            & 0 \\ % e1b4a8df7c08054221b3b05fbfb04096803b0e0d341727fa0df148bc
  minisat                        & algotocnf\_27-t10-diff.cnf     & lingeling  & no    & 0 / 0      & 0         & 1         &            & 0 \\ % 100fdcba5b8495598d1c5b84c8a35b70a39eb246e7a3021b97a12a3e
  minisat                        & algotocnf\_27-t10-diff.cnf     & lingeling  & no    &            &           &           &            & timeout \\ % 992b7d2356908cf21d0921cc22b44c96ffc9d3d184cec051cf67975c
  minisat                        & algotocnf\_27-t10-diff.cnf     & minisat    & no    &            &           &           &            & timeout \\ % 13f460c867e17cc33ea4313258e273bb6ff818f8f8cec980b5f20fcb
  minisat                        & algotocnf\_27-t10-diff.cnf     & minisat    & no    & 0 / 0      & 0         & 1         &            & 0 \\ % 2995f9ab402f627152ddfe212a87b142031f5223b133c106256a6500
  minisat                        & algotocnf\_27-t10-diff.cnf     & satelite   & no    &            &           &           &            & timeout \\ % 202553b821f5d92ea4a50cf655b59c7703cda81495a40a29a8c43775
  minisat                        & algotocnf\_27-t10-diff.cnf     & satelite   & no    & 0 / 0      & 0         & 1         &            & 0 \\ % 3bcf6b0aa70a00397c7fc19931344ea0e65df939e19572a35c0f78ba
  minisat                        & algotocnf\_27-t10-diff\_diff-desc.cnf & none       & no    & 134722 / 8320 & 134722    & 1         &            & 0 \\ % 19a77cb76559bc957054a4cfba0c71cfefca89784ac050bf762d799b
  minisat                        & algotocnf\_27-t10-diff\_diff-desc.cnf & cmsat      & no    &            &           &           &            & timeout \\ % 0e4cd334aa0e117f1f51f84c8c15b1b0ea734e7a003c45370e3ff59e
  minisat                        & algotocnf\_27-t10-diff\_diff-desc.cnf & cmsat      & no    & 0 / 0      & 0         & 1         &            & 0 \\ % 1736e359bc64a7967235205586e31b6f374a3a297350ad5811f8f8f5
  minisat                        & algotocnf\_27-t10-diff\_diff-desc.cnf & lingeling  & no    &            &           &           &            & timeout \\ % 214d534d6f0d11bdd45f58112e4a71b37eb69c8fe5959ac15e45a599
  minisat                        & algotocnf\_27-t10-diff\_diff-desc.cnf & lingeling  & no    & 0 / 0      & 0         & 1         &            & 0 \\ % 9546014a14d3d65673492d9f1e8915ca07c725f6c66c021136e53a63
  minisat                        & algotocnf\_27-t10-diff\_diff-desc.cnf & minisat    & no    & 0 / 0      & 0         & 1         &            & 0 \\ % b28208e66e973cb02ae280c9d63378d6e532a466c926e0ec35e26ed5
  minisat                        & algotocnf\_27-t10-diff\_diff-desc.cnf & minisat    & no    &            &           &           &            & timeout \\ % eb9832fe06f69a6431022c7e9523ab3f100aab2a7fe7426a85cfb73a
  minisat                        & algotocnf\_27-t10-diff\_diff-desc.cnf & satelite   & no    &            &           &           &            & timeout \\ % 3d7ee8abf5cc8f1f423fb9fc54109bcfcf192e2ffe6ab5501f15c9e6
  minisat                        & algotocnf\_27-t10-diff\_diff-desc.cnf & satelite   & no    & 0 / 0      & 0         & 1         &            & 0 \\ % 44a757dc95ee9a9f1a9a962bab44a2e3f4c6dd140fd9246d54d345da
  minisat                        & algotocnf\_27-t10-simplified.cnf & none       & no    &            &           &           &            & 86400 \\ % 02b34d88023be9d2482260c177c1e710620d4b57ff4f0c55ae6ad156
  minisat                        & algotocnf\_27-t10-simplified.cnf & none       & no    &            &           &           &            & timeout \\ % 950f56ae5d483165aa87d7a7b19d33919a9952b2ff36fa06c6186dce
  minisat                        & algotocnf\_27-t10-simplified.cnf & cmsat      & no    &            &           &           &            & timeout \\ % 4275a0ccf8618dd95293dd579cb97a1aa81622c974551a53c95b99cd
  minisat                        & algotocnf\_27-t10-simplified.cnf & cmsat      & no    &            &           &           &            & 86400 \\ % 568fa17c275e65c887f6ef8d29b1f0ef2a1133b983b06043dd87a5fa
  minisat                        & algotocnf\_27-t10-simplified.cnf & minisat    & no    &            &           &           &            & timeout \\ % b2d96c83c3aeca433c4d9b8d04b458f76f46fda79b619e9c005617c8
  minisat                        & algotocnf\_27-t10-simplified.cnf & minisat    & no    &            &           &           &            & 86400 \\ % befe9da1c286d8f429a4bd3401325780633ce87e22be82fd75a61e7e
    \end{tabular}
  \end{center}
\end{table}

%Text such that \LaTeX{} won't screw up.
\newpage

\begin{table}[p]
  \begin{center}
    \begin{tabular}{l|cccccccc}
        \textbf{SAT solver} & \textbf{CNF file} & \textbf{simplified} & \textbf{phase} & \textbf{vars/clauses} & \textbf{propagations} & \textbf{decisions} & \textbf{restarts} & \textbf{Runtime (sec)} \\
      \hline
  minisat                        & algotocnf\_27-t10-simplified.cnf & satelite   & no    &            &           &           &            & timeout \\ % 04836f1e4ffdc8db71d82ab890d75759f6a007ea213a724e56478052
  minisat                        & algotocnf\_27-t10-simplified.cnf & satelite   & no    &            &           &           &            & 86400 \\ % 5494c3712635f4022432432e62eef56b8f7d1ae077a88f4f699394c2
  minisat                        & algotocnf\_27-t10-simplified\_diff-desc.cnf & none       & no    &            &           &           &            & 86400 \\ % 9f85b15fd24183a2313db3af42644d7517d09e830915afa6a6b02e23
  minisat                        & algotocnf\_27-t10-simplified\_diff-desc.cnf & none       & no    &            &           &           &            & timeout \\ % e69bb8473818bea23ca23c74a8e9e47a2c87cc89dbe7fcf34ae8fb7d
  minisat                        & algotocnf\_27-t10-simplified\_diff-desc.cnf & cmsat      & no    &            &           &           &            & 86400 \\ % 0626505794799bfbe79fec7717ad7e613d6398a482ea8920390e1639
  minisat                        & algotocnf\_27-t10-simplified\_diff-desc.cnf & cmsat      & no    &            &           &           &            & timeout \\ % b9d53ab3c744ce85053052e6cbe3503a661ebfcbd885d5f286811def
  minisat                        & algotocnf\_27-t10-simplified\_diff-desc.cnf & minisat    & no    &            &           &           &            & 86400 \\ % 8b4a3e799baece7d39e0590a26c7facb4ffaa958554befd181891257
  minisat                        & algotocnf\_27-t10-simplified\_diff-desc.cnf & minisat    & no    &            &           &           &            & timeout \\ % 9d9c16ff68394399291c81390a0e4163bc882bd1bc1756866983f803
  minisat                        & algotocnf\_27-t10-simplified\_diff-desc.cnf & satelite   & no    &            &           &           &            & timeout \\ % 7ca32247958e7c58f2fcfc3cdd635c4fa61e44fdce6e64474c5a53ae
  minisat                        & algotocnf\_27-t10-simplified\_diff-desc.cnf & satelite   & no    &            &           &           &            & 86400 \\ % ece721cabbc53ae3ab3d938caec4c0c18dfa9b04497e75d9c6b107bc
  minisat                        & algotocnf\_27-t10.cnf          & none       & no    &            &           &           &            & timeout \\ % 08912a4d8319ca2ac5312f679b54dd7cd66e8d76ce32fce8165e6f26
  minisat                        & algotocnf\_27-t10.cnf          & none       & no    &            &           &           &            & 86400 \\ % eb4f0f2a5bb257a3f2e83809842743fa289d32a909ed6609f784eb52
  minisat                        & algotocnf\_27-t10.cnf          & cmsat      & no    &            &           &           &            & timeout \\ % 5c30b5a1bef79f2198edbf575246818898bccec6b37894b536f5f594
  minisat                        & algotocnf\_27-t10.cnf          & cmsat      & no    &            &           &           &            & 86400 \\ % f8842f20b45240a53800dfb6f0f2fb4339be64eeea88aacbaec7f6fc
  minisat                        & algotocnf\_27-t10.cnf          & minisat    & no    &            &           &           &            & timeout \\ % 6b3d44c3b4b6e2568e1916d37424625749763e95829780210e516d71
  minisat                        & algotocnf\_27-t10.cnf          & minisat    & no    &            &           &           &            & 86400 \\ % e3cc6ddb7ce25183b8421a814c76bed8324568b57863da277e871e8e
  minisat                        & algotocnf\_27-t10.cnf          & satelite   & no    &            &           &           &            & 86400 \\ % 11fc018dc06f4d4149147706500995b94ddef94c7a9defca7536c561
  minisat                        & algotocnf\_27-t10.cnf          & satelite   & no    &            &           &           &            & timeout \\ % 780f10770e0d2230f9642d83c60a20917c24fcffa3c747f8473bf4fb
  minisat                        & algotocnf\_27-t10\_diff-desc.cnf & none       & no    &            &           &           &            & 86400 \\ % 268c287bddb09c75dd9b39053dcb2c9da3fa3cc983a4690b42116e6c
  minisat                        & algotocnf\_27-t10\_diff-desc.cnf & none       & no    &            &           &           &            & timeout \\ % 9c7c8e18de1ed0a66510fa171c2381454dc4d4c22e804511b238f290
  minisat                        & algotocnf\_27-t10\_diff-desc.cnf & cmsat      & no    &            &           &           &            & timeout \\ % 0e60cdfe4e845c1dfea77988990e6fb7c0416cfa01e7e219a4e14c04
  minisat                        & algotocnf\_27-t10\_diff-desc.cnf & cmsat      & no    &            &           &           &            & 86400 \\ % 5168fcaf607f1cceac01ebba20c2e220f9ed22ca7c357a8f2979228d
  minisat                        & algotocnf\_27-t10\_diff-desc.cnf & minisat    & no    &            &           &           &            & timeout \\ % a6c1b9bd595414ea14fb0c8f4805b9cc421488ff2d635b0387677b29
  minisat                        & algotocnf\_27-t10\_diff-desc.cnf & minisat    & no    &            &           &           &            & 86400 \\ % eea711d26d63f2cb270d1c29558713a2fa432269baf8baa04a6fff42
  minisat                        & algotocnf\_27-t10\_diff-desc.cnf & satelite   & no    &            &           &           &            & timeout \\ % 195485ff9e65156d0e2f004e247a7f61baa7654a13f9eba9b6f6d0c7
  minisat                        & algotocnf\_27-t10\_diff-desc.cnf & satelite   & no    &            &           &           &            & 86400 \\ % 85534f8844160d69e952729c913d8a8f1109481d520e12443d753474
  minisat                        & algotocnf\_2\_diff-desc.cnf    & none       & no    & 48704 / 329788 & 44141579  & 1901159   &            & 38 \\ % 081b07503b2b521428371ad6d5681225d084ec3a6c2fe83eb27d3180
  minisat                        & algotocnf\_2\_diff-desc.cnf    & none       & no    &            &           &           &            & timeout \\ % 86b49bd2f2d62e5e3bbac3cd5c30c4dc269060612ec0f18a9b5b3499
  minisat                        & algotocnf\_2\_diff-desc.cnf    & cmsat      & no    &            &           &           &            & timeout \\ % 49825d861fe86f1a061b671cf6274ee023d07439cb7ee526570fb5e9
  minisat                        & algotocnf\_2\_diff-desc.cnf    & cmsat      & no    & 15342 / 72870 & 99844401  & 3416907   &            & 84 \\ % c0f7e47c6ec663d74831c8070e3a2fede25479aee93b6820340ba619
  minisat                        & algotocnf\_2\_diff-desc.cnf    & lingeling  & no    &            &           &           &            & timeout \\ % 704e9ce20a9d09455fee46780af083998e3849f40cef1ffa36d780c6
  minisat                        & algotocnf\_2\_diff-desc.cnf    & lingeling  & no    & 40767 / 106546 & 118762201 & 5942147   &            & 122 \\ % 815c4f09df4a016650b8525979b701883dd95432a4526eb53a853b6b
  minisat                        & algotocnf\_2\_diff-desc.cnf    & minisat    & no    & 23531 / 149436 & 169599488 & 5161470   &            & 180 \\ % 0999aa6d6e9681774d50e8214f4bf9d62f752998f98780eab1d42cc7
  minisat                        & algotocnf\_2\_diff-desc.cnf    & minisat    & no    &            &           &           &            & timeout \\ % 8b9735220552123799f3d19786f55c901cdd0580a8c8e0801f601ace
  minisat                        & algotocnf\_2\_diff-desc.cnf    & satelite   & no    & 29428 / 204577 & 118408491 & 4024150   &            & 153 \\ % b1c0e81dd3f73e923587abbc977c5cede33ff1c24f85a70ab51b80a0
    \end{tabular}
  \end{center}
\end{table}

%Text such that \LaTeX{} won't screw up.
\newpage

\begin{table}[p]
  \begin{center}
    \begin{tabular}{l|cccccccc}
        \textbf{SAT solver} & \textbf{CNF file} & \textbf{simplified} & \textbf{phase} & \textbf{vars/clauses} & \textbf{propagations} & \textbf{decisions} & \textbf{restarts} & \textbf{Runtime (sec)} \\
      \hline
  minisat                        & algotocnf\_2\_diff-desc.cnf    & satelite   & no    &            &           &           &            & timeout \\ % f4e224753b8c60176307ef6f160b7fd46c2158a669e10db9f85aced8
  minisat                        & algotocnf\_3.cnf               & none       & no    & 48704 / 222436 & 13991247  & 853580    &            & 9 \\ % 8f51d7012499c8b13ee6e7d5e7db080d454284fe9e67066a5595064c
  minisat                        & algotocnf\_3.cnf               & none       & no    &            &           &           &            & timeout \\ % e8e4dd33d90584c706abe064df3faae738796def9c24b8b2baa23852
  minisat                        & algotocnf\_3.cnf               & cmsat      & no    & 23893 / 111403 & 15330537  & 894641    &            & 10 \\ % 0b31e71827bff181e0d457fe714a1558034bed0ee72669eff14a3989
  minisat                        & algotocnf\_3.cnf               & cmsat      & no    &            &           &           &            & timeout \\ % bd2b1bf6a533804f426587c80994a40cd05e5799f03a776a124b49fa
  minisat                        & algotocnf\_3.cnf               & lingeling  & no    &            &           &           &            & timeout \\ % 4cac1c7b205196dc264efdeb845c5b749e9762c1007298145d1a3f60
  minisat                        & algotocnf\_3.cnf               & lingeling  & no    & 44000 / 100941 & 3922854   & 328431    &            & 2 \\ % dbeeb6d8b7727a0f1165ca86deb993ec8c24183a23a2c59848f4fca3
  minisat                        & algotocnf\_3.cnf               & minisat    & no    & 20895 / 119580 & 11055439  & 705917    &            & 7 \\ % a0707a0fe074b58d3430dfaf0fdb901d4cc4033daf2247def7e36c22
  minisat                        & algotocnf\_3.cnf               & minisat    & no    &            &           &           &            & timeout \\ % af9b8f2b0bcbee727397102136f78340c438bf02e39e9c7ba7054250
  minisat                        & algotocnf\_3.cnf               & satelite   & no    &            &           &           &            & timeout \\ % 7672561f6a3d076f45d68353864f563425c5c55df1960483c840ae5c
  minisat                        & algotocnf\_3.cnf               & satelite   & no    & 28647 / 156910 & 15840699  & 1019521   &            & 8 \\ % ff4c575d84fc7e3bbc79ac00d963f3c1355f0eb1af81f3924130b7a0
  minisat                        & algotocnf\_3\_diff-desc.cnf    & none       & no    &            &           &           &            & timeout \\ % 1ac8b11834c082c1585471bc70c4526f52eeeb7ac209d7de2ecde580
  minisat                        & algotocnf\_3\_diff-desc.cnf    & none       & no    & 48704 / 328004 & 8355604   & 451770    &            & 4 \\ % f7bf2cfc7733580441fd33fb93199fffb56c53c52d11d54953ad2e89
  minisat-core                   & algotocnf\_1.cnf               & none       & no    & 48704 / 225124 & 12249636981 & 194879682 &            & 19683 \\ % 92843528701a645eebdc4b0c4aeafce4557162e15212609470860bb8
  minisat-core                   & algotocnf\_1.cnf               & cmsat      & no    & 24503 / 111931 & 1345904991 & 45407643  &            & 2206 \\ % d71dd0554732b4d7233d4ef8a074274dcebfd8857d247a8caca92d7d
  minisat-core                   & algotocnf\_1.cnf               & minisat    & no    & 20895 / 118236 & 1001195109 & 41539126  &            & 1579 \\ % b9f8d4fb4182fa6723cc39ca9f1eec9ac4181734eb2b947db8daa19b
  minisat-core                   & algotocnf\_1.cnf               & satelite   & no    & 27495 / 153262 & 2852530112 & 91754777  &            & 4628 \\ % 50401b6d7b7298858916dde823eee8d1098e09835fb4542bd5ae8640
  minisat-core                   & algotocnf\_18-t9.cnf           & none       & no    & 107839 / 582936 & 197561807 & 5771226   &            & 79 \\ % d0e8168dbe7c6afdaa93c478731210760af73449bf0061c8da84fed8
  minisat-core                   & algotocnf\_18-t9.cnf           & cmsat      & no    & 61789 / 342139 & 3714952   & 336242    &            & 2 \\ % 4400efcf9248e92906351215efa6f61727797070b130f3c6d27c69f5
  minisat-core                   & algotocnf\_18-t9.cnf           & lingeling  & no    & 102616 / 344544 & 27126157  & 3729263   &            & 19 \\ % 8de591b056292544c46e9a1069246e3633cc99339b2be380de1af71d
  minisat-core                   & algotocnf\_18-t9.cnf           & minisat    & no    & 61845 / 391134 & 11474979  & 666416    &            & 7 \\ % ae69c7e4dda77eb47f13d79c3251c48e0cb053a97af5706ed35253c4
  minisat-core                   & algotocnf\_18-t9.cnf           & satelite   & no    & 69670 / 457972 & 48139459  & 2806650   &            & 18 \\ % a11e0d071f29994bfc04dce494f11373561005f536b134760810eb2b
  minisat-core                   & algotocnf\_18-t9\_diff-desc.cnf & none       & no    & 107839 / 831263 & 134492463 & 5454225   &            & 46 \\ % 63217a08851f7553f0f5cf50b1074ab59402af8d2c4c45c9a4aaa8fe
  minisat-core                   & algotocnf\_18-t9\_diff-desc.cnf & cmsat      & no    & 57114 / 303489 & 22660129  & 1033163   &            & 15 \\ % cee340d9d44c4efed37f6ed7a2de08ad81af9e84882c8a25641e0d5d
  minisat-core                   & algotocnf\_18-t9\_diff-desc.cnf & lingeling  & no    & 99452 / 404666 & 4783352   & 1249213   &            & 3 \\ % c3e2a5ff267d515b2f89731c187b80211d1346f49d3b4ed2a5a0c229
  minisat-core                   & algotocnf\_18-t9\_diff-desc.cnf & minisat    & no    & 78055 / 511431 & 6527792   & 383605    &            & 4 \\ % 320e0b632a00324efc078496178c8412007b99891f0bf002d6b5c86d
  minisat-core                   & algotocnf\_18-t9\_diff-desc.cnf & satelite   & no    & 85386 / 692408 & 68801966  & 3555378   &            & 47 \\ % d31725c58ab860500c00d23cf18ff94a7af676aeeb902b8a48d7a31c
  minisat-core                   & algotocnf\_1\_diff-desc.cnf    & none       & no    & 48704 / 330692 & 556973568 & 10855280  &            & 834 \\ % c3c28b67df90396846fef5060e174db77138235ea8d089e70bf88007
  minisat-core                   & algotocnf\_1\_diff-desc.cnf    & cmsat      & no    & 15402 / 73142 & 565587287 & 15477212  &            & 652 \\ % 6acd4c4bc11573a51b7f9c26f5e25da392f95af7682504ad1263b5f3
  minisat-core                   & algotocnf\_1\_diff-desc.cnf    & lingeling  & no    & 40554 / 101182 & 130669677 & 6453337   &            & 187 \\ % 93133ce850e1d9240c0ef6b4b7ba8747341a633fca289626e27dadce
  minisat-core                   & algotocnf\_1\_diff-desc.cnf    & minisat    & no    & 25107 / 158393 & 166787763 & 5399857   &            & 185 \\ % d0ccd3adb4a63d45d20b0441f3b82d6be402278cd83e6b1dd789c706
  minisat-core                   & algotocnf\_1\_diff-desc.cnf    & satelite   & no    & 30950 / 226602 & 218192599 & 7694052   &            & 241 \\ % c5d51bcab919c2e22128858dfcac69a4d9474adb8c14d0b8544ed36c
  minisat-core                   & algotocnf\_2.cnf               & none       & no    & 48704 / 224220 & 4474027487 & 81273454  &            & 7817 \\ % a1ea8b59043da04b0fd89dd2a2cceca4d2943af2667187a61f905834
  minisat-core                   & algotocnf\_2.cnf               & cmsat      & no    & 24323 / 111661 & 3271551070 & 109512958 &            & 4722 \\ % 7656af6ac1604d858880aa93258ce3ec954589087548015cf2137c5e
  minisat-core                   & algotocnf\_2.cnf               & minisat    & no    & 20895 / 118688 & 491695491 & 21320818  &            & 622 \\ % 331d25010ae307f15fd645f6e1e2757803cfefd2c556be5347848f0b
    \end{tabular}
  \end{center}
\end{table}

%Text such that \LaTeX{} won't screw up.
\newpage

\begin{table}[p]
  \begin{center}
    \begin{tabular}{l|cccccccc}
        \textbf{SAT solver} & \textbf{CNF file} & \textbf{simplified} & \textbf{phase} & \textbf{vars/clauses} & \textbf{propagations} & \textbf{decisions} & \textbf{restarts} & \textbf{Runtime (sec)} \\
      \hline
  minisat-core                   & algotocnf\_2.cnf               & satelite   & no    & 27775 / 154276 & 1356174519 & 45617994  &            & 1965 \\ % 00ee90aaecba4e4f3d843823cab191e36558dd8020fb6cce0f66329a
  minisat-core                   & algotocnf\_21-t9.cnf           & none       & no    &            &           &           &            & 86400 \\ % 89aa557127c903bc59f286136a8d7aade0657930afd196d8875fde9b
  minisat-core                   & algotocnf\_21-t9.cnf           & cmsat      & no    & 63555 / 358308 & 9720846001 & 182867764 &            & 60458 \\ % 6a07ecba7c6a08d70fd27c7c22f66caf3710f4a7a1f136c6496e4c7f
  minisat-core                   & algotocnf\_21-t9.cnf           & minisat    & no    &            &           &           &            & 86400 \\ % cb31aaa83dc8b9c55fe4440d922d5b8fd58758abbc0267e9d477fec2
  minisat-core                   & algotocnf\_21-t9.cnf           & satelite   & no    &            &           &           &            & 86400 \\ % 7ea7db8f0bbe9117366de49164c5a8665bda9a9b04cc18cbb8a860ba
  minisat-core                   & algotocnf\_21-t9\_diff-desc.cnf & none       & no    & 116800 / 894072 & 6184915141 & 69982430  &            & 38376 \\ % ec97e37d9050d5b198da19be071d691b80f7aea54ca7c49e00cd25f9
  minisat-core                   & algotocnf\_21-t9\_diff-desc.cnf & cmsat      & no    & 59528 / 318049 & 3033335128 & 52731840  &            & 10776 \\ % 0d96e02ead9c1ec6f012d3833dce7bc129457f69c0e5e88a97c458da
  minisat-core                   & algotocnf\_21-t9\_diff-desc.cnf & lingeling  & no    & 104860 / 404126 & 2309115640 & 56884034  &            & 14440 \\ % 02387f33c8ede6ac833a554ae23ad54dfbf8678d30a2154d0988e903
  minisat-core                   & algotocnf\_21-t9\_diff-desc.cnf & minisat    & no    & 81650 / 534965 & 3082155154 & 42479717  &            & 18914 \\ % 28da3413de75e42f3bba2e403c6b7b996446f78ad28e6996d51aa943
  minisat-core                   & algotocnf\_21-t9\_diff-desc.cnf & satelite   & no    & 90236 / 720755 & 2028870107 & 33135732  &            & 9626 \\ % e4b9c9f1cbbfe8ad3a2df21f3ca0c69c3ba2ca16038a65e0c1de572b
  minisat-core                   & algotocnf\_23-t9.cnf           & none       & no    &            &           &           &            & 86400 \\ % e809794e0e5088b2ecad07e3d35b5c80a8942090a7f38db370c824a8
  minisat-core                   & algotocnf\_23-t9.cnf           & cmsat      & no    &            &           &           &            & 86400 \\ % 282ebdacfb51fad93528823e42bf510650ee4d5f528a4d3f65a24242
  minisat-core                   & algotocnf\_23-t9.cnf           & minisat    & no    &            &           &           &            & 86400 \\ % 02047b664fc8f08dd4405a1efdc051fd59134be5769176e91aad6c8c
  minisat-core                   & algotocnf\_23-t9.cnf           & satelite   & no    &            &           &           &            & 86400 \\ % 6ae7daf284c6f6860623440881a328b6d5cd9274593593f1089846e2
  minisat-core                   & algotocnf\_23-t9\_diff-desc.cnf & none       & no    &            &           &           &            & 86400 \\ % 6a7a93c091e46bed22a48d098a77a481acc808b13b11d89ade708ffb
  minisat-core                   & algotocnf\_23-t9\_diff-desc.cnf & cmsat      & no    & 57510 / 307118 & 6686701995 & 106401427 &            & 38818 \\ % b54bd208dd37bf964e304e2450ebbb44aae366da035ce8d90cc818a2
  minisat-core                   & algotocnf\_23-t9\_diff-desc.cnf & minisat    & no    &            &           &           &            & 86400 \\ % 866e9daa649140b984d679b02e557cf79e3e936f0a8001ab0ecf8b8d
  minisat-core                   & algotocnf\_23-t9\_diff-desc.cnf & satelite   & no    &            &           &           &            & 86400 \\ % 97588a46b7556f1c4dca8acd81f458533fab199a850a442c742ac12d
  minisat-core                   & algotocnf\_24-t9.cnf           & none       & no    &            &           &           &            & 86400 \\ % 4cef4d92bab7b713f99cd37a919d7eb57f7ba1116d0ff33f7ba316d6
  minisat-core                   & algotocnf\_24-t9.cnf           & cmsat      & no    &            &           &           &            & 86400 \\ % 5cf14bae003c76838575285fcf1be7f401078f2ad81da133395b0fae
  minisat-core                   & algotocnf\_24-t9.cnf           & minisat    & no    &            &           &           &            & 86400 \\ % 35f95337fa37e2f5556a2c0c463c5d210282f188354753ad3d92b3d0
  minisat-core                   & algotocnf\_24-t9.cnf           & satelite   & no    &            &           &           &            & 86400 \\ % 5939f105b298cc4ab280164b397ab9f74fbf5b3d209dbd2076caa74c
  minisat-core                   & algotocnf\_24-t9\_diff-desc.cnf & none       & no    &            &           &           &            & 86400 \\ % 9b9b5868241cae54be43460510bc547189263bee7b7ee134c43d8752
  minisat-core                   & algotocnf\_24-t9\_diff-desc.cnf & cmsat      & no    & 56447 / 302014 & 3308202991 & 53182192  &            & 11410 \\ % ac4eec9a63283ccd5b9da20bd5e7b11512bacf40604eb76b08d3905c
  minisat-core                   & algotocnf\_24-t9\_diff-desc.cnf & minisat    & no    &            &           &           &            & 86400 \\ % 65904dcede8aff9171f828a7afef17d5140d1e23518cf45b47830707
  minisat-core                   & algotocnf\_24-t9\_diff-desc.cnf & satelite   & no    &            &           &           &            & 86400 \\ % e68790c4413eca1622f4717114ba5947bb7baf78b5d945046ee73f13
  minisat-core                   & algotocnf\_27-t10-diff.cnf     & none       & no    & 134722 / 8320 & 134722    & 1         &            & 0 \\ % e6bc5f89bc76ec8548ef8e0d8e60a651c87f3cf6c90d75c8283287ed
  minisat-core                   & algotocnf\_27-t10-diff.cnf     & cmsat      & no    & 0 / 0      & 0         & 1         &            & 0 \\ % 9fcfee46ed67b137ea476970787e4654a4c1688e65b7a261cf2ec546
  minisat-core                   & algotocnf\_27-t10-diff.cnf     & lingeling  & no    & 0 / 0      & 0         & 1         &            & 0 \\ % 126fe605077329ea9c5f15cc3684599a7c8d5bc4dbf657989a06243a
  minisat-core                   & algotocnf\_27-t10-diff.cnf     & minisat    & no    & 0 / 0      & 0         & 1         &            & 0 \\ % 8844586e120ca2abb893f9f408804aff5d51ad6d7b16114b6a1270bd
  minisat-core                   & algotocnf\_27-t10-diff.cnf     & satelite   & no    & 0 / 0      & 0         & 1         &            & 0 \\ % 700c7aefbce005830ca0ef0bb9b3ebd208dcff26356369e8b0ea37e7
  minisat-core                   & algotocnf\_27-t10-diff\_diff-desc.cnf & none       & no    & 134722 / 8320 & 134722    & 1         &            & 0 \\ % 1e50d07e9488483e0acdd0e5702828b735df09d6d763b7d864e5bd47
  minisat-core                   & algotocnf\_27-t10-diff\_diff-desc.cnf & cmsat      & no    & 0 / 0      & 0         & 1         &            & 0 \\ % 4540e386591c2bb570d0068b9fdfbe5720cda17d689a144f7191d153
  minisat-core                   & algotocnf\_27-t10-diff\_diff-desc.cnf & lingeling  & no    & 0 / 0      & 0         & 1         &            & 0 \\ % c13d8c4a1d7b40f0871ab1a4ccbf8502637b74ce7da7b5a24817dd06
  minisat-core                   & algotocnf\_27-t10-diff\_diff-desc.cnf & minisat    & no    & 0 / 0      & 0         & 1         &            & 0 \\ % 4043cbcd69d95fb63da409e54a9f3158fac8419f42b6b1a9ee719318
    \end{tabular}
  \end{center}
\end{table}

%Text such that \LaTeX{} won't screw up.
\newpage

\begin{table}[p]
  \begin{center}
    \begin{tabular}{l|cccccccc}
        \textbf{SAT solver} & \textbf{CNF file} & \textbf{simplified} & \textbf{phase} & \textbf{vars/clauses} & \textbf{propagations} & \textbf{decisions} & \textbf{restarts} & \textbf{Runtime (sec)} \\
      \hline
  minisat-core                   & algotocnf\_27-t10-diff\_diff-desc.cnf & satelite   & no    & 0 / 0      & 0         & 1         &            & 0 \\ % e1cf64b7e7ada39ec8af4c8c2fae84f53dbbff23b35797165c9acf58
  minisat-core                   & algotocnf\_27-t10-simplified.cnf & none       & no    &            &           &           &            & 86400 \\ % 3b6b0d2d9f6f8a6213f2d9eea3f43be278e629bb540b195adf778bf2
  minisat-core                   & algotocnf\_27-t10-simplified.cnf & cmsat      & no    &            &           &           &            & 86400 \\ % 0fb5581e96188d4609a4a6e19d30725a87e3d5390230a3539d0c9b7f
  minisat-core                   & algotocnf\_27-t10-simplified.cnf & minisat    & no    &            &           &           &            & 86400 \\ % a7db5f7ba4710cd07ca7c3663a92e485c7e97358d16078e63d6a3401
  minisat-core                   & algotocnf\_27-t10-simplified.cnf & satelite   & no    &            &           &           &            & 86400 \\ % 61ca0665db016ff49f5cba0c9256efd9537d8ae1932bb5a40b30f97e
  minisat-core                   & algotocnf\_27-t10-simplified\_diff-desc.cnf & none       & no    &            &           &           &            & 86400 \\ % f8b4db7a94dedba746c3dfe50b0fb139589db2a612b6ba2d3b293afb
  minisat-core                   & algotocnf\_27-t10-simplified\_diff-desc.cnf & cmsat      & no    &            &           &           &            & 86400 \\ % eafaaa6ec6c86c729ac717220253f3e310d7beb10614d103807f66ff
  minisat-core                   & algotocnf\_27-t10-simplified\_diff-desc.cnf & minisat    & no    &            &           &           &            & 86400 \\ % 6f4cbc8aac945ac965973c2d855221c82ab0d2d560037e4c91511362
  minisat-core                   & algotocnf\_27-t10-simplified\_diff-desc.cnf & satelite   & no    &            &           &           &            & 86400 \\ % 40a5540eb06d3fcd7b318cd987fb644e73e52b67956443d165ae1e1f
  minisat-core                   & algotocnf\_27-t10.cnf          & none       & no    &            &           &           &            & 86400 \\ % 411694111851fb1f234bc515e9861de991b2d16b837337318d938536
  minisat-core                   & algotocnf\_27-t10.cnf          & cmsat      & no    &            &           &           &            & 86400 \\ % 295df816203921763e0c0c89409cf104e696a9a03b03e9d03b1ec37a
  minisat-core                   & algotocnf\_27-t10.cnf          & minisat    & no    &            &           &           &            & 86400 \\ % 4d28111eaf2762db8ead4d729b594f346eb4c243d33c761af876cc16
  minisat-core                   & algotocnf\_27-t10.cnf          & satelite   & no    &            &           &           &            & 86400 \\ % 99bd9867642dc7f6e09ce8214f2e997ba2e64a6da965b3661533ff15
  minisat-core                   & algotocnf\_27-t10\_diff-desc.cnf & none       & no    &            &           &           &            & 86400 \\ % 06902a1c4c2b99e562d8acdde90a6ef6ebf75cb09451efb8fc05ee29
  minisat-core                   & algotocnf\_27-t10\_diff-desc.cnf & cmsat      & no    &            &           &           &            & 86400 \\ % 92a6377b25e1a912da42277d62f433836018f13ad36e775139bfed0e
  minisat-core                   & algotocnf\_27-t10\_diff-desc.cnf & minisat    & no    &            &           &           &            & 86400 \\ % 4371cf155fafb01628ef7ff25eb90a18d0bfca3eec4a77ff1f7696be
  minisat-core                   & algotocnf\_27-t10\_diff-desc.cnf & satelite   & no    &            &           &           &            & 86400 \\ % 5ec4ca292114bc55cc475645b68d1a5a77261ccdfa764f89e5d9148c
  minisat-core                   & algotocnf\_2\_diff-desc.cnf    & none       & no    & 48704 / 329788 & 642600542 & 12472624  &            & 946 \\ % cc70aa7ef474c852c7a3329bc923bdfda35fa1ee61af4c896ee9e6b0
  minisat-core                   & algotocnf\_2\_diff-desc.cnf    & cmsat      & no    & 15342 / 72870 & 100117349 & 3345445   &            & 83 \\ % eb43cf2621d19b1757d66263a5a22accefb933492fc51208382da94d
  minisat-core                   & algotocnf\_2\_diff-desc.cnf    & lingeling  & no    & 40767 / 106546 & 128399770 & 6982842   &            & 190 \\ % 8d9e3076c7997ea60b09ff8f5063971134c0a50470cbd6307e0fe5db
  minisat-core                   & algotocnf\_2\_diff-desc.cnf    & minisat    & no    & 23531 / 149436 & 169599488 & 5161470   &            & 251 \\ % 776fd7d035cf5ddd986bc26d0c4636688f6f9a98f06cbe2643c73965
  minisat-core                   & algotocnf\_2\_diff-desc.cnf    & satelite   & no    & 29428 / 204577 & 214332265 & 7443515   &            & 173 \\ % 22324189bb398c78222bfbe0e425c13cf38a5122571662640de9b5a3
  minisat-core                   & algotocnf\_3.cnf               & none       & no    & 48704 / 222436 & 134609706 & 4257400   &            & 65 \\ % 94410dec89a249cc381a9fcb11f21a9ea05cc89a53d779f8cfbdc9a5
  minisat-core                   & algotocnf\_3.cnf               & cmsat      & no    & 23893 / 111403 & 16248891  & 886675    &            & 7 \\ % 1872a77b967d63c0fa334e637bdafba6b56e201f2d9729daa29e43fe
  minisat-core                   & algotocnf\_3.cnf               & lingeling  & no    & 44000 / 100941 & 4973985   & 734132    &            & 2 \\ % 17e277f641c860c348a0cca8aecc3cf5e772511ca0909999dfbf1a9a
  minisat-core                   & algotocnf\_3.cnf               & minisat    & no    & 20895 / 119580 & 11055439  & 705917    &            & 5 \\ % 3daed7ca29f094ac0315aeb6a73085d68432edb143be3b7362c41528
  minisat-core                   & algotocnf\_3.cnf               & satelite   & no    & 28647 / 156910 & 81453944  & 3957973   &            & 40 \\ % bf2bfae417fc60de1c8c777770e89526b10f21b81f6fe2c89f102131
  minisat-core                   & algotocnf\_3\_diff-desc.cnf    & none       & no    & 48704 / 328004 & 82077674  & 2936309   &            & 53 \\ % 80f89129f5265878e0db5f3f57874b103ae474922cb02f09d83a4d8d
  plingeling-ats1o1              & algotocnf\_1.cnf               & cmsat      & no    & 24503 / 111931 & 8548008070 & 37237312  &            & 1232 \\ % ca21656e4756468dfb224a5cb9e8d87ffc4bef5d081425f3c13dfe83
  plingeling-ats1o1              & algotocnf\_1.cnf               & minisat    & no    & 20895 / 118236 & 8304337810 & 44390611  &            & 1384 \\ % 3f2c3884f9300cb1a4b37e0c07b5ccf1f9a7ec7344933ec71f032765
  plingeling-ats1o1              & algotocnf\_1.cnf               & satelite   & no    & 27495 / 153262 & 12438835790 & 58208495  &            & 2030 \\ % 05952731c478437bd27eb8a30c6699b15acb96625ab48f769002f7f8
  plingeling-ats1o1              & algotocnf\_18-t9.cnf           & none       & no    & 107839 / 590953 & 589105230 & 3489049   &            & 176 \\ % 683a62e1cb4c9d2217968a77fc9e72cd2a5157e93fe7cef2144ea36a
  plingeling-ats1o1              & algotocnf\_18-t9.cnf           & cmsat      & no    & 61789 / 342139 & 357501470 & 2878644   &            & 101 \\ % e4168ea727963fa16680f65bccb2d0b7eed37ce8572a5573f2efc25f
  plingeling-ats1o1              & algotocnf\_18-t9.cnf           & lingeling  & no    & 107839 / 344544 & 199229360 & 2083451   &            & 21 \\ % 253bed3059fc7ae4fe73361b67d672fb253c2927ac6dd180e446d905
  plingeling-ats1o1              & algotocnf\_18-t9.cnf           & minisat    & no    & 61845 / 391134 & 30363010  & 231634    &            & 5 \\ % 3774881f9de4ef6b58f99854eef2cca05f975072ec631ed787ec98aa
    \end{tabular}
  \end{center}
\end{table}

%Text such that \LaTeX{} won't screw up.
\newpage

\begin{table}[p]
  \begin{center}
    \begin{tabular}{l|cccccccc}
        \textbf{SAT solver} & \textbf{CNF file} & \textbf{simplified} & \textbf{phase} & \textbf{vars/clauses} & \textbf{propagations} & \textbf{decisions} & \textbf{restarts} & \textbf{Runtime (sec)} \\
      \hline
  plingeling-ats1o1              & algotocnf\_18-t9.cnf           & satelite   & no    & 69670 / 457972 & 396313800 & 3125943   &            & 77 \\ % ac60865df834cbdb2f20b04b4f4f1b7bc3ec97d7a5081bdd1ddeedfc
  plingeling-ats1o1              & algotocnf\_18-t9\_diff-desc.cnf & cmsat      & no    & 57114 / 303489 & 949615550 & 7232823   &            & 126 \\ % f12b8b4571385309d762a5a8b97fd9e40e815ddee54b82a2bbb0de99
  plingeling-ats1o1              & algotocnf\_18-t9\_diff-desc.cnf & lingeling  & no    & 107839 / 404666 & 745621300 & 5465546   &            & 65 \\ % 1246145b1a6d4f8445161cd667980749231b25d1038b0a9b59356804
  plingeling-ats1o1              & algotocnf\_18-t9\_diff-desc.cnf & minisat    & no    & 78055 / 511431 & 52703620  & 310182    &            & 7 \\ % 4f66e00647527096ead024242114e824b00f0a12c93c887bfc5bd0d6
  plingeling-ats1o1              & algotocnf\_18-t9\_diff-desc.cnf & satelite   & no    & 85386 / 692408 & 375280530 & 1941381   &            & 109 \\ % d7f6876e2b034eaa09085dfe828a7e166a03075550699aae0958e2cd
  plingeling-ats1o1              & algotocnf\_1\_diff-desc.cnf    & none       & no    & 48704 / 373920 & 4289131630 & 18654755  &            & 534 \\ % b8997e17e93fde3550ee79a4cbbf399bb2054e50a5c686718a5835d2
  plingeling-ats1o1              & algotocnf\_1\_diff-desc.cnf    & cmsat      & no    & 15402 / 73142 & 3548147180 & 15149776  &            & 376 \\ % fc7441ff4bd14751eb66db74ca8f789d8d7681965820e8a9a0de0885
  plingeling-ats1o1              & algotocnf\_1\_diff-desc.cnf    & lingeling  & no    & 48704 / 101182 & 3180845780 & 21922776  &            & 376 \\ % 614a7b48b8fa0f5a83b4f1b992064cb33c3abc1db446807a8125f85c
  plingeling-ats1o1              & algotocnf\_1\_diff-desc.cnf    & minisat    & no    & 25107 / 158393 & 1965238720 & 10013850  &            & 284 \\ % 5b8a8d86f8b3bb9f1cf648419fec93926e2d920b8433068518926c56
  plingeling-ats1o1              & algotocnf\_1\_diff-desc.cnf    & satelite   & no    & 30950 / 226602 & 2491219360 & 12716739  &            & 324 \\ % 6260ce50d2e579540ee6b5138d6e56f89923a6398bb32cb598eeec24
  plingeling-ats1o1              & algotocnf\_2.cnf               & none       & no    & 48704 / 254210 & 11278846250 & 47300319  &            & 1860 \\ % c6b986d9fcba1f5ac1914b6e6b7521a093db84c19f37d45f81bb3d9a
  plingeling-ats1o1              & algotocnf\_2.cnf               & cmsat      & no    & 24323 / 111661 & 8145031310 & 31795532  &            & 1100 \\ % 62e0169f37e6cfd8f993fb4d12748ef0b1ed9d9f003152f6796758a8
  plingeling-ats1o1              & algotocnf\_2.cnf               & minisat    & no    & 20895 / 118688 & 3124819420 & 16746404  &            & 568 \\ % fc802ece3fd75265706e1f0fa3e47bc10b3f140892f132573f88fa70
  plingeling-ats1o1              & algotocnf\_2.cnf               & satelite   & no    & 27775 / 154276 & 6888905530 & 33989411  &            & 1152 \\ % 1a8943bf0eeb75cf061dee2ce30939ab9a28114d2f7555d65a92c100
  plingeling-ats1o1              & algotocnf\_21-t9.cnf           & none       & no    & 116800 / 636838 & 130188640360 & 274034016 &            & 42025 \\ % d9d90bf0968e2fa364b8a24c4f6bc13ed972b0b9860ccda7e7a7dd92
  plingeling-ats1o1              & algotocnf\_21-t9.cnf           & cmsat      & no    & 63555 / 358308 & 17255856970 & 31570089  &            & 3423 \\ % e9d7a37ba1980c3624a6c10437ff1203a5a399180afac8d513a1254f
  plingeling-ats1o1              & algotocnf\_21-t9.cnf           & minisat    & no    & 65572 / 414470 & 191274957430 & 384124771 &            & 67220 \\ % 8c07210c095b179fbd7cfebf8a05b07da7da8c5d1c82dfbd6191c3ba
  plingeling-ats1o1              & algotocnf\_21-t9.cnf           & satelite   & no    & 74562 / 486714 & 132702819760 & 256625513 &            & 47064 \\ % b9964f89b32c64018ab0e1e970518fc4e2b7b665e92d627b581d18af
  plingeling-ats1o1              & algotocnf\_21-t9\_diff-desc.cnf & none       & no    & 116800 / 911629 & 18514838920 & 20859643  &            & 2545 \\ % 9d5d0e63a34024cd87b8fc3930e85c96c558a9c1785e8672be97ba70
  plingeling-ats1o1              & algotocnf\_21-t9\_diff-desc.cnf & cmsat      & no    & 59528 / 318049 & 14450240650 & 25298696  &            & 1928 \\ % 8ef439e5ae5b70cb2693c44410f4dbc0b19c4f56260138d1c125ad80
  plingeling-ats1o1              & algotocnf\_21-t9\_diff-desc.cnf & minisat    & no    & 81650 / 534965 & 21718771620 & 27595849  &            & 2553 \\ % 814e9c11b56d70872970ef5396817f8814d68213d1caaf82fdec63a1
  plingeling-ats1o1              & algotocnf\_21-t9\_diff-desc.cnf & satelite   & no    & 90236 / 720755 & 23804708950 & 36221337  &            & 3204 \\ % 9f238e8b7c3a6776c2c9770d267d67f658d9daec8f85bdbfb13552ce
  plingeling-ats1o1              & algotocnf\_23-t9.cnf           & none       & no    & 122774 / 667438 & 231925212560 & 406491721 &            & 90019 \\ % 355e080b70ed653c48a8d637a080cf8ed633f049ac27bd853bda3ccd
  plingeling-ats1o1              & algotocnf\_23-t9.cnf           & cmsat      & no    & 64880 / 369775 & 251983202340 & 366928770 &            & 90007 \\ % 4e13fcc1ea3660e22a5fc018ef09f146c07da9175c9ffa76ac428f4c
  plingeling-ats1o1              & algotocnf\_23-t9.cnf           & minisat    & no    & 67979 / 429734 & 261991505790 & 473285487 &            & 125383 \\ % a297e60b7c460f55c376898f3eb34c43161cba48c18480661a121220
  plingeling-ats1o1              & algotocnf\_23-t9.cnf           & satelite   & no    & 77835 / 506217 & 228831325710 & 406736241 &            & 90002 \\ % 1d1c5f1c7792c5b315c34ba9458ccff7e11c039f3d596574157dcda0
  plingeling-ats1o1              & algotocnf\_23-t9\_diff-desc.cnf & none       & no    & 122774 / 955067 & 55888853190 & 53857251  &            & 8175 \\ % 29df17185515f4218e76f5b31187d2b1a05ec73b95ae4a78e362abf9
  plingeling-ats1o1              & algotocnf\_23-t9\_diff-desc.cnf & cmsat      & no    & 57510 / 307118 & 25488237880 & 38692192  &            & 3695 \\ % 74bdce987e723e6fde7917ef907d42243bc3a8f1405536da3ca98b5b
  plingeling-ats1o1              & algotocnf\_23-t9\_diff-desc.cnf & minisat    & no    & 84080 / 550727 & 39712137130 & 63433218  &            & 6475 \\ % 4925ad5378f6301a0acb5cbb7f6a08e52e33bbd756a43e81d92c0751
  plingeling-ats1o1              & algotocnf\_23-t9\_diff-desc.cnf & satelite   & no    & 93504 / 740545 & 53177836930 & 81353095  &            & 9536 \\ % 5b331e7c0b483fc9503ac4198bdd108198d5aae7a84005c4175022e5
  plingeling-ats1o1              & algotocnf\_24-t9.cnf           & none       & no    & 125761 / 682722 & 768191210920 & 1353464490 &            & 378393 \\ % fd47fd352512e06a3127ded13230502ba0c5b762a8eb9764f3982cb5
  plingeling-ats1o1              & algotocnf\_24-t9.cnf           & cmsat      & no    & 65579 / 377228 & 274278093010 & 404493616 &            & 90000 \\ % ee5275a88aa159693b30ba624ae67d9c57e81eb0c72ad5f2640c1fe5
  plingeling-ats1o1              & algotocnf\_24-t9.cnf           & minisat    & no    & 69023 / 436292 & 251062096120 & 415427110 &            & 90007 \\ % b470c5d9e50fda1f3928f2dc1fffa8e1694fb2bbd867828c4ff21c99
  plingeling-ats1o1              & algotocnf\_24-t9.cnf           & satelite   & no    & 79486 / 515977 & 295147128350 & 486067461 &            & 124488 \\ % 1567c0d98c332e6a49dceb970038af0ece6098bf21f14de856eb960d
  plingeling-ats1o1              & algotocnf\_24-t9\_diff-desc.cnf & none       & no    & 125761 / 976770 & 43878084770 & 65364303  &            & 8578 \\ % 9896fcb690560c07bfb43a42f3db4109f18274844312cacec87bfece
    \end{tabular}
  \end{center}
\end{table}

%Text such that \LaTeX{} won't screw up.
\newpage

\begin{table}[p]
  \begin{center}
    \begin{tabular}{l|cccccccc}
        \textbf{SAT solver} & \textbf{CNF file} & \textbf{simplified} & \textbf{phase} & \textbf{vars/clauses} & \textbf{propagations} & \textbf{decisions} & \textbf{restarts} & \textbf{Runtime (sec)} \\
      \hline
  plingeling-ats1o1              & algotocnf\_24-t9\_diff-desc.cnf & cmsat      & no    & 56447 / 302014 & 31835601210 & 43860424  &            & 5601 \\ % 72caec0c1b48d2ce13a2691ddfd789fbdf81b501600685f22a3342c9
  plingeling-ats1o1              & algotocnf\_24-t9\_diff-desc.cnf & minisat    & no    & 85312 / 558731 & 33053039330 & 46301510  &            & 4601 \\ % ea8a7cf0b46e30ffffda4f71e7a1aa073c023bb6038dbf294e6a0314
  plingeling-ats1o1              & algotocnf\_24-t9\_diff-desc.cnf & satelite   & no    & 95154 / 750557 & 29320347790 & 39722986  &            & 4212 \\ % 6b35d7212e351d66a11964feac822bd9f439796552b3805a2a874697
  plingeling-ats1o1              & algotocnf\_27-t10-diff.cnf     & none       & no    & 134722 / 728349 & 10777760  & 0         &            & 1 \\ % 3036a09a15cf3d0ca7b4c590f673e72bc0619e330af77307d2299aed
  plingeling-ats1o1              & algotocnf\_27-t10-diff.cnf     & cmsat      & no    & 0 / 0      & 0         & 0         &            & 0 \\ % 1b6f8e21836d47dc417ebb881bab61f2de7c8add07b0b49ba0e90090
  plingeling-ats1o1              & algotocnf\_27-t10-diff.cnf     & lingeling  & no    & 134722 / 0 & 0         & 0         &            & 0 \\ % 76fd2d8cbe93df5d40f9de289cf851997af260b085f72d9c511966c8
  plingeling-ats1o1              & algotocnf\_27-t10-diff.cnf     & minisat    & no    & 0 / 0      & 0         & 0         &            & 0 \\ % 6ba7b39b423078ab98fbd7741edce87fdd831f334c26531e3cc66582
  plingeling-ats1o1              & algotocnf\_27-t10-diff.cnf     & satelite   & no    & 0 / 0      & 0         & 0         &            & 0 \\ % 2cd77d0a53ffed6779f52f072cb5dea9f54d6ed95877988ec0bf5a11
  plingeling-ats1o1              & algotocnf\_27-t10-diff\_diff-desc.cnf & none       & no    & 134722 / 1041654 & 16166640  & 0         &            & 1 \\ % 20aa707d10b08bacb707c796dc22a0d31ba4574299cbf9b72a18d8c1
  plingeling-ats1o1              & algotocnf\_27-t10-diff\_diff-desc.cnf & cmsat      & no    & 0 / 0      & 0         & 0         &            & 0 \\ % eb085525e1b9eff731d697e3ebd6123c21a688565ff1a06d650de07f
  plingeling-ats1o1              & algotocnf\_27-t10-diff\_diff-desc.cnf & lingeling  & no    & 134722 / 0 & 0         & 0         &            & 0 \\ % fbcf66ff272a522fc345e6cb0f7824a5d97e967110c6d46cd17076b1
  plingeling-ats1o1              & algotocnf\_27-t10-diff\_diff-desc.cnf & minisat    & no    & 0 / 0      & 0         & 0         &            & 0 \\ % e71d4119e10a67704205f224bf97084b5916f65fd7a4d48d35fdace2
  plingeling-ats1o1              & algotocnf\_27-t10-diff\_diff-desc.cnf & satelite   & no    & 0 / 0      & 0         & 0         &            & 0 \\ % 506bcb6499dd9368191cf658971c426e2d31e8dcc5634f558b19414e
  plingeling-ats1o1              & algotocnf\_27-t10-simplified.cnf & none       & no    & 134722 / 728542 & 274272903670 & 297757325 &            & 90013 \\ % 3d8e1950326774b04f2d3a7ba96ab9768470d8cfc7dc3b6d3eb4753c
  plingeling-ats1o1              & algotocnf\_27-t10-simplified.cnf & cmsat      & no    & 67175 / 392680 & 325437288820 & 373129456 &            & 90008 \\ % a8e1e752f467316949dc44837e59e2c651c503afaf8d679f71a73c07
  plingeling-ats1o1              & algotocnf\_27-t10-simplified.cnf & minisat    & no    & 72909 / 460404 & 255937911220 & 334262173 &            & 90024 \\ % 1536680b7000188c4757ec0c4af75af50df5e5e7c3fc74ee4f61374b
  plingeling-ats1o1              & algotocnf\_27-t10-simplified.cnf & satelite   & no    & 84442 / 544998 & 327012759110 & 413380539 &            & 90025 \\ % 9606da4645c6e5b50df2a0440f2576728406cdea60d97cd008eb5414
  plingeling-ats1o1              & algotocnf\_27-t10-simplified\_diff-desc.cnf & none       & no    & 134722 / 1041847 & 292002205560 & 238183761 &            & 90005 \\ % 91d213d806ee4a2055c2656abb0cde7b7ea55666cc470b1bdf4a9d0b
  plingeling-ats1o1              & algotocnf\_27-t10-simplified\_diff-desc.cnf & satelite   & no    & 100735 / 789589 & 298711804850 & 276668293 &            & 90013 \\ % 197382133bb94bfde87bb2f632bf147b84d64d414153af338c108908
  plingeling-ats1o1              & algotocnf\_27-t10.cnf          & none       & no    & 134722 / 728030 & 293304590350 & 358115584 &            & 90024 \\ % 0a58e78ee9e6b1947395c1ad3610e5026cf86d6300304aa40192df28
  plingeling-ats1o1              & algotocnf\_27-t10.cnf          & cmsat      & no    & 72044 / 411276 & 292264318530 & 317366792 &            & 90018 \\ % 4aaac8e744f50c3f5154c8fb0e584843cac7cc5ccf2f5a4c07fdc685
  plingeling-ats1o1              & algotocnf\_27-t10.cnf          & minisat    & no    & 73897 / 465248 & 238305702090 & 312928566 &            & 90000 \\ % 6dece8673f052c1bb318d3d1c223f8015c22c18a29b9445b15091571
  plingeling-ats1o1              & algotocnf\_27-t10\_diff-desc.cnf & none       & no    & 134722 / 1041335 & 285995947130 & 237982215 &            & 90001 \\ % 150428085fa1cc8f2a6e75d7242868706d3e9161de3ddc86ef7c61be
  plingeling-ats1o1              & algotocnf\_27-t10\_diff-desc.cnf & cmsat      & no    & 57956 / 311219 & 311794167570 & 240205381 &            & 90014 \\ % 6343c89e26b20062a5d5d1adad3f4cf80720a290e3efeac139af6f04
  plingeling-ats1o1              & algotocnf\_27-t10\_diff-desc.cnf & minisat    & no    & 94850 / 619480 & 325772399620 & 308016744 &            & 90007 \\ % 98cc95c840999eb3b8254cbc501f4ed63eac0c4e85b3520a48c5731f
  plingeling-ats1o1              & algotocnf\_27-t10\_diff-desc.cnf & satelite   & no    & 105733 / 857675 & 284666279240 & 268853902 &            & 90000 \\ % bd4a242c025467a63f99002eb0038ab9a114b138d2bb3a12957ce16b
  plingeling-ats1o1              & algotocnf\_2\_diff-desc.cnf    & none       & no    & 48704 / 374146 & 4747715750 & 16718540  &            & 498 \\ % 18af33baa4b158667de3041a3db943a47cbab5e3627c620e943672b9
  plingeling-ats1o1              & algotocnf\_2\_diff-desc.cnf    & cmsat      & no    & 15342 / 72870 & 2035703930 & 9699276   &            & 210 \\ % 9bea09cc9569d44d5c40e7fdec547ec21b75e159424f00c9b1019c2d
  plingeling-ats1o1              & algotocnf\_2\_diff-desc.cnf    & lingeling  & no    & 48704 / 106546 & 2146163460 & 14572254  &            & 311 \\ % f5217e4dd24ea91becff34e38a840a3197272e584d11c50a403b3ace
  plingeling-ats1o1              & algotocnf\_2\_diff-desc.cnf    & minisat    & no    & 23531 / 149436 & 1077232830 & 6883591   &            & 157 \\ % 3de8d82d15be9460b769a2891ecf124738a115f26285af01b433d7e2
  plingeling-ats1o1              & algotocnf\_2\_diff-desc.cnf    & satelite   & no    & 29428 / 204577 & 1665434460 & 9949133   &            & 200 \\ % 35ca7396652a7d4c9c9d4e5b4eb80d5ce52f62abde965ebf96224bcb
  plingeling-ats1o1              & algotocnf\_3.cnf               & none       & no    & 48704 / 254656 & 452630440 & 3275498   &            & 88 \\ % 406db7a97fa084420a1ce7d8b997eb222233ae223e9f4c5f65e5054b
  plingeling-ats1o1              & algotocnf\_3.cnf               & cmsat      & no    & 23893 / 111403 & 1267595100 & 6911456   &            & 219 \\ % 88993abddab37585aef5aeed78ddb74567f721bb3866301753c6f3ff
  plingeling-ats1o1              & algotocnf\_3.cnf               & lingeling  & no    & 48704 / 100941 & 531571420 & 5142102   &            & 60 \\ % 68ae6af20473be995a5e934e1eff8d98106d8ebbc512be65311cf194
  plingeling-ats1o1              & algotocnf\_3.cnf               & minisat    & no    & 20895 / 119580 & 472452560 & 3571352   &            & 79 \\ % 9af567c9a3ee3a5708ff0dd934681b09af0fec08e9781d4089dfc946
    \end{tabular}
  \end{center}
\end{table}

%Text such that \LaTeX{} won't screw up.
\newpage

\begin{table}[p]
  \begin{center}
    \begin{tabular}{l|cccccccc}
        \textbf{SAT solver} & \textbf{CNF file} & \textbf{simplified} & \textbf{phase} & \textbf{vars/clauses} & \textbf{propagations} & \textbf{decisions} & \textbf{restarts} & \textbf{Runtime (sec)} \\
      \hline
  plingeling-ats1o1              & algotocnf\_3.cnf               & satelite   & no    & 28647 / 156910 & 637183710 & 5053181   &            & 108 \\ % 2a4cd60c8d9c84e13ab8311cb595714eca61417c031b4b6fb443b6dc
  plingeling-ats1o1              & algotocnf\_3\_diff-desc.cnf    & none       & no    & 48704 / 374592 & 220439160 & 1579205   &            & 36 \\ % 3a98c08d64f3bdd55426c50aa0bc34229fbda87d2a0456cd9a9f31d9
  treengeling-ats1o1             & algotocnf\_1.cnf               & cmsat      & no    & 24503.0 / 111931.0 & 1758628978 & 104577617 &            & 13401 \\ % 5babfec1db673d690fab718495086508071ac84a3daf890c91b3d8fc
  treengeling-ats1o1             & algotocnf\_1.cnf               & minisat    & no    & 20895.0 / 118236.0 & 1173989436 & 89581806  &            & 13790 \\ % 80e28d8110628cccfe09c1573b21ee3b0bf28c0ceeb3bb2e1c6f4ecb
  treengeling-ats1o1             & algotocnf\_1.cnf               & satelite   & no    & 27495.0 / 153262.0 & 1250013660 & 81367172  &            & 10840 \\ % 8940641ecfa852d5b7c26cf1a3ec91687c25ded7a853dccf7ce4bc60
  treengeling-ats1o1             & algotocnf\_18-t9.cnf           & none       & no    & 107839.0 / 590953.0 & 8099141   & 593628    &            & 56 \\ % f1392915db836b318673ccbbb2c0a1f6ab70e1d5bcd0cf1dbae1ca9b
  treengeling-ats1o1             & algotocnf\_18-t9.cnf           & cmsat      & no    & 61789.0 / 342139.0 & 40223002  & 2709485   &            & 207 \\ % 949dac4ee25957a3769f8d45b612fd76a3fcd02e1f47c7ac9c91ad7d
  treengeling-ats1o1             & algotocnf\_18-t9.cnf           & lingeling  & no    & 107839.0 / 344544.0 & 6113944   & 456336    &            & 38 \\ % 656711b415c91ef2558cf890f0f92a38a0946fe532e5c02df696f2b4
  treengeling-ats1o1             & algotocnf\_18-t9.cnf           & minisat    & no    & 61845.0 / 391134.0 & 6056879   & 440497    &            & 36 \\ % 5ba3c9569d5b43bf98b89f487d5a77a26e425ef6ac43536e4e6a4479
  treengeling-ats1o1             & algotocnf\_18-t9.cnf           & satelite   & no    & 69670.0 / 457972.0 & 12679339  & 898982    &            & 65 \\ % 33d24b5ee86b92c010be9c434ad97528006e88a2e188b866f13e13f4
  treengeling-ats1o1             & algotocnf\_18-t9\_diff-desc.cnf & cmsat      & no    & 57114.0 / 303489.0 & 198613370 & 9599953   &            & 516 \\ % 595776b3d6955c0f8c8e073989e70e51ab4d572b5779cedbf874ae4a
  treengeling-ats1o1             & algotocnf\_18-t9\_diff-desc.cnf & lingeling  & no    & 107839.0 / 404666.0 & 9802722   & 1073545   &            & 17 \\ % a8edef24b364497436e6f145a3d5ad7711f5627fbd6e08fa658a3ab3
  treengeling-ats1o1             & algotocnf\_18-t9\_diff-desc.cnf & minisat    & no    & 78055.0 / 511431.0 & 274181625 & 8748578   &            & 846 \\ % 8a553c1d2c975b45b14b7c40562b55bebe05d919dc61b6115c0736a4
  treengeling-ats1o1             & algotocnf\_1\_diff-desc.cnf    & none       & no    & 48704.0 / 373920.0 & 357794060 & 22299892  &            & 1913 \\ % 98ea228ecf2dc2e8a1c5ba5c6049c9b1ec9ec7ecd144ec8f9eebd94c
  treengeling-ats1o1             & algotocnf\_1\_diff-desc.cnf    & cmsat      & no    & 15402.0 / 73142.0 & 584854327 & 31344752  &            & 2492 \\ % 5ca1ef06add54e1acd7b0a5806f00967132fc4132c3959b56c9bcb4f
  treengeling-ats1o1             & algotocnf\_1\_diff-desc.cnf    & minisat    & no    & 25107.0 / 158393.0 & 183160734 & 8783302   &            & 784 \\ % 616ebba7aabeb7987bf442d9876d509ecf2974b00b742dc146eefa97
  treengeling-ats1o1             & algotocnf\_1\_diff-desc.cnf    & satelite   & no    & 30950.0 / 226602.0 & 568539546 & 38187947  &            & 3075 \\ % 355a0a4ff7d611492b92e77465eca4589c72fd35de9b56bd33c06937
  treengeling-ats1o1             & algotocnf\_2.cnf               & none       & no    & 48704.0 / 254210.0 & 1381556910 & 116981692 &            & 12574 \\ % 3da5d773defc2c994530265959396262260a27cb16aa9f6ca1f4b088
  treengeling-ats1o1             & algotocnf\_2.cnf               & cmsat      & no    & 24323.0 / 111661.0 & 215197068 & 10471628  &            & 1821 \\ % 0570253de19ac599f84eb320c8edf1872da6e5ee23a1a931dc1c7330
  treengeling-ats1o1             & algotocnf\_2.cnf               & minisat    & no    & 20895.0 / 118688.0 & 60162535  & 5307169   &            & 288 \\ % dd652f62bad8b725d3e96f9a74a168d0b2f7ac8ee10961c78d2406a9
  treengeling-ats1o1             & algotocnf\_2.cnf               & satelite   & no    & 27775.0 / 154276.0 & 1365878152 & 110332766 &            & 14926 \\ % a3a21644245d076e3b3888386fb87903bf340713912c19e5525a3160
  treengeling-ats1o1             & algotocnf\_21-t9.cnf           & satelite   & no    & 74562.0 / 486714.0 & 48952027607 & 930986783 &            & 105290 \\ % 781b6bbf7b56d21a2893d52fd2b29d23a81b9a93b43fd58dd4566db2
  treengeling-ats1o1             & algotocnf\_21-t9\_diff-desc.cnf & none       & no    & 116800.0 / 911629.0 & 733199001 & 14011210  &            & 6419 \\ % fb17c3feebce5d713cd90beb8a4c43494657f53d131b0e52ee24e9e0
  treengeling-ats1o1             & algotocnf\_21-t9\_diff-desc.cnf & cmsat      & no    & 59528.0 / 318049.0 & 422501062 & 10144692  &            & 5118 \\ % d3d4f7e84cf2dc35361fc2e04557710eef7037ca713d5b8d55d9de64
  treengeling-ats1o1             & algotocnf\_21-t9\_diff-desc.cnf & minisat    & no    & 81650.0 / 534965.0 & 21638451159 & 210865064 &            & 27567 \\ % 25f977eb1d03b95b17b50898233effabb3f7cb3ee179f7eeb829d35e
  treengeling-ats1o1             & algotocnf\_23-t9.cnf           & none       & no    & 122774.0 / 667438.0 & 43707331147 & 844241021 &            & 90010 \\ % 77611c14cf4a915761ba1bbafd04e2d226708c30ddb666252a01e0d8
  treengeling-ats1o1             & algotocnf\_23-t9.cnf           & cmsat      & no    & 64880.0 / 369775.0 & 42615823775 & 742277791 &            & 90007 \\ % e31167f61c6090a834d8ae28306f0db8c886f969331764ad0d560424
  treengeling-ats1o1             & algotocnf\_23-t9.cnf           & minisat    & no    & 67979.0 / 429734.0 & 43608664278 & 785402879 &            & 125381 \\ % 5838c62564ced3b8c7f2d6201e794851ace0ceb38bb97dd1c295ee7c
  treengeling-ats1o1             & algotocnf\_23-t9.cnf           & satelite   & no    & 77835.0 / 506217.0 & 34896230441 & 627260693 &            & 90009 \\ % da9c9ca5a4857626bd9e32750596bb2405fc18e521dc983457658e53
  treengeling-ats1o1             & algotocnf\_23-t9\_diff-desc.cnf & none       & no    & 122774.0 / 955067.0 & 20645511861 & 317969600 &            & 34900 \\ % e01700aea1441f69b86c7bad952fafac6754ca37c733d8ce6398f307
  treengeling-ats1o1             & algotocnf\_23-t9\_diff-desc.cnf & cmsat      & no    & 57510.0 / 307118.0 & 6220331064 & 120006196 &            & 21248 \\ % 5af7706ca825a9e6884e4705b3a36a6f05662de691076d4c15783b73
  treengeling-ats1o1             & algotocnf\_23-t9\_diff-desc.cnf & minisat    & no    & 84080.0 / 550727.0 & 11954245211 & 164439336 &            & 25018 \\ % af9ab7e6dfc5f55fbbe7de361b413cfe8064a381dfd4d8b3f98b9c39
  treengeling-ats1o1             & algotocnf\_24-t9.cnf           & cmsat      & no    & 65579.0 / 377228.0 & 39372304774 & 708381692 &            & 90028 \\ % dc4cde5bf5c80969b465ea2b237b423b4799de048f17f882b1cfa3ea
  treengeling-ats1o1             & algotocnf\_24-t9.cnf           & minisat    & no    & 69023.0 / 436292.0 & 42090192000 & 833793163 &            & 90000 \\ % d0fa3b9bd06e3a7199243eeaad39ba934dff6bec853bacb14b6cb477
  treengeling-ats1o1             & algotocnf\_24-t9.cnf           & satelite   & no    & 79486.0 / 515977.0 & 50060946197 & 933285412 &            & 124486 \\ % 348c078356abf55c006c2248306b3041af9f9c1b9683e425ca1f264f
    \end{tabular}
  \end{center}
\end{table}

%Text such that \LaTeX{} won't screw up.
\newpage

\begin{table}[p]
  \begin{center}
    \begin{tabular}{l|cccccccc}
        \textbf{SAT solver} & \textbf{CNF file} & \textbf{simplified} & \textbf{phase} & \textbf{vars/clauses} & \textbf{propagations} & \textbf{decisions} & \textbf{restarts} & \textbf{Runtime (sec)} \\
      \hline
  treengeling-ats1o1             & algotocnf\_24-t9\_diff-desc.cnf & none       & no    & 125761.0 / 976770.0 & 931422937 & 21687044  &            & 9647 \\ % 7fd2035787068a3b23318b384cfe1c36ea94eb67e628a64272548890
  treengeling-ats1o1             & algotocnf\_24-t9\_diff-desc.cnf & cmsat      & no    & 56447.0 / 302014.0 & 11998764516 & 200412370 &            & 29166 \\ % a55605a665da3ea5faf0fcaead5b6a7e47805b4dc0a9087493f0cc3e
  treengeling-ats1o1             & algotocnf\_24-t9\_diff-desc.cnf & minisat    & no    & 85312.0 / 558731.0 & 480343946 & 11873378  &            & 4269 \\ % 213409faec1b95787d8ea79d6d5648d1d367cf2d29dc51d2950a6265
  treengeling-ats1o1             & algotocnf\_24-t9\_diff-desc.cnf & satelite   & no    & 95154.0 / 750557.0 & 19978663137 & 319972352 &            & 35809 \\ % d19a35f5dd521934cf0d43b2be57566041d4741f5ad592a1a00e8300
  treengeling-ats1o1             & algotocnf\_27-t10-diff.cnf     & none       & no    & 134722.0 / 728349.0 & 269444    & 0         &            & 1 \\ % 9c49c8fa0a2dc5ff76d9754558f3369b7ad3489c64bb660383196668
  treengeling-ats1o1             & algotocnf\_27-t10-diff.cnf     & cmsat      & no    & 0.0 / 0.0  & 0         & 0         &            & 0 \\ % 7ca0423eb7d9b494c7e980f0971640743caefa535acdf21c9cfe46d2
  treengeling-ats1o1             & algotocnf\_27-t10-diff.cnf     & lingeling  & no    & 134722.0 / 0.0 & 0         & 0         &            & 0 \\ % dc28c340a57c3ea5439d33107ad3259e1ba71e9b1f699a9d9c0e1853
  treengeling-ats1o1             & algotocnf\_27-t10-diff.cnf     & minisat    & no    & 0.0 / 0.0  & 0         & 0         &            & 0 \\ % fc40c586737df67183bf0ef8cdfe6a3e23c392c4651674bb2a832ee8
  treengeling-ats1o1             & algotocnf\_27-t10-diff.cnf     & satelite   & no    & 0.0 / 0.0  & 0         & 0         &            & 0 \\ % 3e281ddfa13c9560a0ffb1667d4b0e134b21556ad22f7d17abe485cd
  treengeling-ats1o1             & algotocnf\_27-t10-diff\_diff-desc.cnf & cmsat      & no    & 0.0 / 0.0  & 0         & 0         &            & 0 \\ % 7fc7f01a3378cbdd1001f55e91de2e3fd9badeb3739f91b363cb18c7
  treengeling-ats1o1             & algotocnf\_27-t10-diff\_diff-desc.cnf & lingeling  & no    & 134722.0 / 0.0 & 0         & 0         &            & 0 \\ % 3b6415a8ce9803a577673e9f9d1ce61873543edbc3b0f19a3f64024c
  treengeling-ats1o1             & algotocnf\_27-t10-diff\_diff-desc.cnf & minisat    & no    & 0.0 / 0.0  & 0         & 0         &            & 0 \\ % 31d862825af72a3a506af4b595d2ee09536a874578490773cc65327e
  treengeling-ats1o1             & algotocnf\_27-t10-diff\_diff-desc.cnf & satelite   & no    & 0.0 / 0.0  & 0         & 0         &            & 0 \\ % 30b2cb4e68565e26917898585855f38d771fedcc05b81d3643b02963
  treengeling-ats1o1             & algotocnf\_27-t10-simplified.cnf & cmsat      & no    & 67175.0 / 392680.0 & 47220737968 & 902071733 &            & 90006 \\ % fdaa10201f4efe8878b62ad7f26dc1f0bdf97469a185373d7376fd35
  treengeling-ats1o1             & algotocnf\_27-t10-simplified.cnf & minisat    & no    & 72909.0 / 460404.0 & 54499993435 & 911713561 &            & 90009 \\ % f7a23e4842605e7ec448197c5dd7ef578146023338feae422087c318
  treengeling-ats1o1             & algotocnf\_27-t10-simplified.cnf & satelite   & no    & 84442.0 / 544998.0 & 53317549880 & 861479437 &            & 90000 \\ % 34639ef0c0092d5828824a98c5d6db1952c55f4868b2fafa0aba6652
  treengeling-ats1o1             & algotocnf\_27-t10-simplified\_diff-desc.cnf & satelite   & no    & 100735.0 / 789589.0 & 41070418116 & 452346314 &            & 90009 \\ % 8b4c42aa41df103c5e88aa8ba720b66f6898dbebbe822e4380cac429
  treengeling-ats1o1             & algotocnf\_27-t10.cnf          & none       & no    & 134722.0 / 728030.0 & 37211671569 & 544734908 &            & 90009 \\ % 3a2dfa3a55a5977e9e149263ebc4fb47a9178df459b344b2cda1cb35
  treengeling-ats1o1             & algotocnf\_27-t10.cnf          & cmsat      & no    & 72044.0 / 411276.0 & 47142955948 & 928709284 &            & 90015 \\ % c1f28fd123fbf32ed43571e22935bfa0443a6296f49434aef3b9efbe
  treengeling-ats1o1             & algotocnf\_27-t10.cnf          & minisat    & no    & 73897.0 / 465248.0 & 46569438485 & 826117288 &            & 90021 \\ % ba3facbd0dcfb130cb402558db0d6c5350627cf0aa7932162e57f29c
  treengeling-ats1o1             & algotocnf\_27-t10\_diff-desc.cnf & none       & no    & 134722.0 / 1041335.0 & 51310449087 & 582878435 &            & 89999 \\ % fbe2e0cbb87532031c57e4aeaf2900c1a2e156afa28e85f0182f9328
  treengeling-ats1o1             & algotocnf\_27-t10\_diff-desc.cnf & cmsat      & no    & 57956.0 / 311219.0 & 57581929090 & 596857407 &            & 90011 \\ % 9465be454f79f18889de09fa3ea2d5ddbadd04336a124061f3e74a41
  treengeling-ats1o1             & algotocnf\_27-t10\_diff-desc.cnf & minisat    & no    & 94850.0 / 619480.0 & 62067865752 & 695858870 &            & 90005 \\ % 609061feb1e5cb392f6c48a67f3fed4f4f73d4155327a8b473defc2b
  treengeling-ats1o1             & algotocnf\_2\_diff-desc.cnf    & cmsat      & no    & 15342.0 / 72870.0 & 170385323 & 7533047   &            & 725 \\ % 7b3d255f829552da539b6f60e8bd8e45f7d4fa37cc265f8076123820
  treengeling-ats1o1             & algotocnf\_2\_diff-desc.cnf    & lingeling  & no    & 48704.0 / 106546.0 & 64924746  & 4705886   &            & 265 \\ % d007ae2a1482255f21621338b0657df2726d94316f39af7fab18c970
  treengeling-ats1o1             & algotocnf\_2\_diff-desc.cnf    & minisat    & no    & 23531.0 / 149436.0 & 122498098 & 7212854   &            & 480 \\ % f5238ff6075a0ceaee42a387da54e8fe4cbb1fc58f03fa719a21b7b0
  treengeling-ats1o1             & algotocnf\_2\_diff-desc.cnf    & satelite   & no    & 29428.0 / 204577.0 & 134449086 & 7392471   &            & 741 \\ % b6151599d6cbad9557a983e001573eb1e85b4eb4ea0ae740f882ee39
  treengeling-ats1o1             & algotocnf\_3.cnf               & none       & no    & 48704.0 / 254656.0 & 18629811  & 1640289   &            & 64 \\ % 93f8a5d32f4fc0464e29cc1d7894624b9fb4105e4519d7729c2295a9
  treengeling-ats1o1             & algotocnf\_3.cnf               & cmsat      & no    & 23893.0 / 111403.0 & 47477822  & 2980122   &            & 160 \\ % 069fa622a2b9683f3e267d9188784d513afaf1611c737b322a5be203
  treengeling-ats1o1             & algotocnf\_3.cnf               & lingeling  & no    & 48704.0 / 100941.0 & 11161253  & 1139418   &            & 17 \\ % 92d00810658deedbb159ad75760a4e386ee98c22776471f1aa9987b3
  treengeling-ats1o1             & algotocnf\_3.cnf               & minisat    & no    & 20895.0 / 119580.0 & 11430482  & 1109257   &            & 42 \\ % f2e2a7933dcc022d6512bca4f6c845379ee2b78ab0ad0e8a2b816a25
  treengeling-ats1o1             & algotocnf\_3.cnf               & satelite   & no    & 28647.0 / 156910.0 & 32476484  & 2719280   &            & 122 \\ % 36283e8063b55a519dd40de006de684e3af539f0ac90e28c4373f95d

    \end{tabular}
  \end{center}
\end{table}


\end{appendices}


\backmatter
\printbibliography
%\listoffigures
%\listoftables
%\printindex

\end{document}
